\documentclass[
  dvipdfmx,
  a4paper,
  11pt,
  fleqn,
  openany
]{jsbook}

% =========================================================
% 数学・作図
% =========================================================

\usepackage{amsmath,amsthm,amssymb,amsfonts}
\usepackage{nameref}
\usepackage{tikz}


% =========================================================
% レイアウト
% =========================================================

\usepackage[
  margin=20truemm
]{geometry}

\usepackage{multicol}
\usepackage{enumitem}
\usepackage{url}


% =========================================================
% 見出し・章デザイン
% =========================================================

\usepackage{quotchap}


% =========================================================
% 色・ボックス
% =========================================================

\usepackage[table]{xcolor}
\usepackage{tcolorbox}

\tcbuselibrary{breakable}

\definecolor{shadecolor}{gray}{0.80}


% =========================================================
% 全体の組版
% =========================================================

% 段の高さを無理にそろえず、上から順に流す
\raggedcolumns

% 2段組の段間
\setlength{\columnsep}{12pt}

% 2段組の中央線
\setlength{\columnseprule}{0.4pt}

% 段落間隔
\setlength{\parskip}{0pt}

% 行送りの不自然な拡張を必要な箇所だけ許容
\setlength{\emergencystretch}{1em}


% =========================================================
% 数式周辺の余白
% =========================================================

\setlength{\abovedisplayskip}{10pt plus 2pt minus 2pt}
\setlength{\belowdisplayskip}{10pt plus 2pt minus 2pt}

\setlength{\abovedisplayshortskip}{8pt plus 2pt minus 2pt}
\setlength{\belowdisplayshortskip}{8pt plus 2pt minus 2pt}


% =========================================================
% 箇条書き
% =========================================================

% 例題・方針・解答・解法などで使用している
% enumerate のインデントは元の設計を維持する。
%
% leftmargin を指定しないことで、jsbook / enumitem の
% \leftmargini を基準とした元来の狭い幅を保持する。
\setlist[enumerate,1]{
  label=(\arabic*),
  leftmargin=\leftmargini,
  parsep=3pt
}


% =========================================================
% 引用
% =========================================================

% enumerate のインデント体系と整合するよう、
% \leftmargini を基準に smallquote の幅を決める。
\newlength{\smallquotemargin}

\setlength{\smallquotemargin}{0.5\leftmargini}

\newenvironment{smallquote}
  {%
    \begin{list}{}{%
      \setlength{\leftmargin}{\smallquotemargin}%
      \setlength{\rightmargin}{\smallquotemargin}%
      \setlength{\topsep}{0.35\baselineskip}%
      \setlength{\partopsep}{0pt}%
      \setlength{\parsep}{0pt}%
      \setlength{\listparindent}{0pt}%
    }%
    \item[]%
  }
  {\end{list}}


% =========================================================
% 箇条書き内で式番号を付ける
% =========================================================

\newcommand{\itemeq}[2][]{%
  \refstepcounter{equation}%
  \item
  \makebox[\linewidth]{%
    $\displaystyle #2$
    \hfill
    (\theequation)%
  }%
  \ifx\relax#1\relax
  \else
    \label{#1}%
  \fi
}


% =========================================================
% 証明終了記号
% =========================================================

\renewcommand{\qed}{%
  \unskip\nobreak\quad\qedsymbol
}


% =========================================================
% 目次・見出し番号
% =========================================================

\setcounter{tocdepth}{3}
\setcounter{secnumdepth}{4}

\newcommand{\chapterref}[1]{第\ref{#1}章~\nameref{#1}}
\newcommand{\sectionref}[1]{\ref{#1}~\nameref{#1}}

\title{漸化式の超体系的解説}

\begin{document}

\date{\today}
\author{MAS256}
\maketitle

\tableofcontents
\pagebreak
% chapter: はじめに

\chapter{はじめに}\label{章：はじめに}
\begin{bookcolumns}
  \section{本書について}\label{節：はじめに}
    本書は、\bookterm{recurrence}{漸化式}の解法と、その背後にある考え方をたどる本である。
    差を取る。倍率をそろえる。複数の項を組み合わせる。
    一つの式変形から、別の解法や、グラフ・\bookterm{linear}{線形}代数へのつながりが見えてくる。
    その道筋を、式と短い文章で描いていく。

    各型では、代表的な解法から始める。
    例題の方針には着眼点を、解答には結果を、解法には計算の過程を記した。
    例題を解いた後で、同じ操作を別の見方から捉え直す。
    計算の理由と、複数の方法の使い分けも扱う。

    紙とペンで式を確かめると、文章だけでは見えにくい変化も見えてくる。
    演習は、その理解を別の式へ広げるためにある。
    本書でいう「基本」は、後の議論を支える考え方を指す。
    計算の道具は\chapterref{章：前提知識と計算の基本}へ、
    用語と発展的な計算は\chapterref{章：用語と計算の参照}へ戻って確認できる。

    \bookterm{general}{一般項}を求めることから、\bookterm{limit}{極限}や整数の性質を調べることへ。
    それぞれの解法を結び付けながら、\bookterm{recurrence}{漸化式}を見渡していく。

  \section{章立て}\label{節：章立て}
    基本の計算から、解法、その選び方、数学的なつながりへ進む。
    巻末には、言葉と計算の参照先を置いた。
    \begin{itemize}
      \item \chapterref{章：前提知識と計算の基本}は参照用である.必要な知識が出てきたときに,対応する節へ戻ればよい.
      \item \chapterref{章：漸化式の解説}では,個々の\bookterm{recurrence}{漸化式}に対してのアプローチを学ぶ.
      \item \chapterref{章：実際の解き方}では,\chapterref{章：漸化式の解説}で学んだ内容をまとめ,実際に\bookterm{recurrence}{漸化式}が与えられた時にどうするかを学ぶ.
      \item \chapterref{章：付録}では、本編の操作を\bookterm{matrix}{行列}・\bookterm{difference-operator}{差分}・\bookterm{generating}{母関数}から見直す。
      \item \chapterref{章：用語と計算の参照}には、用語の意味と、\bookterm{logarithm}{対数}・\bookterm{trigonometric}{三角関数}・\bookterm{limit}{極限}・複素数の計算をまとめた。
    \end{itemize}
    それぞれの関係は
    \BookFigRoadmap
    のようになっている.

  \section{対象読者と読む経路}\label{節：対象読者}
    対象読者は,\bookterm{recurrence}{漸化式}を初めて学ぶ人から,既習の解法の理由や数学的なつながりを知りたい人まで対応している.
    特にこの本にマッチするのは,数学が好きな高校生である.
    関心や既習の内容に合わせて、次の場所から読める。

    \noindent\textbf{【初めて学ぶ人】}\par
    \sectionref*{節：数列と漸化式}で記号の読み方を確かめ、
    \sectionref*{節：漸化式を解く前に}から\sectionref*{節：基本4パターン}へ進む。
    その後、\sectionref*{節：追う量を変える}、
    \sectionref*{節：解を組み合わせる見方},\sectionref*{節：応用7パターン}へ進む.
    一つの型を読んだら,その例題と\hyperref[節：演習の読む経路]{「演習の読む経路」}の代表題に取り組む.
    到達目標は,変形後の式を計算できることに加えて,「何に注目し,何を簡単にしたか」を一文で説明し,\bookterm{initial}{初期条件}と元の式への代入で答えを確かめられることである.

    \noindent\textbf{【一度学んだ人】}\par
    \sectionref*{節：解法の決定}と\sectionref*{節：操作を選ぶ練習}から入り,理由を説明できない操作だけ第3章へ戻る.
    \hyperref[節：演習の読む経路]{「演習の読む経路」}の比較題では,別解の手掛かり・計算量・条件の違いを調べる.
    到達目標は,同じ問題の二つの解法を比較し,係数や初項が変わったときにも方法を選び直せることである.

    \noindent\textbf{【数学を深く学びたい人】}\par
    発展パターン・特殊性質型と発展演習を,付録の関連する節と往復して読む.
    到達目標は,異なる型で使った操作の共通点を述べ,\bookterm{general}{一般項}・\bookterm{limit}{極限}・\bookterm{invariant}{不変量}など,問題に合う目標を選べることである.
    \bookterm{matrix}{行列}は付録の冒頭で導入する.微分を使う箇所では接線と導関数の知識が必要であり,\bookterm{trigonometric}{三角関数}・\bookterm{complex-plane}{複素平面}・\bookterm{limit}{極限}は、巻末の参照資料で確認できる。

  \section{本書でのルール}\label{節：本書でのルール}
    左罫線と薄い背景のある囲みは,読み返すための着眼点・条件・結論を示す.
    太字の小見出しで,発想・計算・解釈の切り替わりを示す.
    議論の道筋は、本文の文章と式に記した。
    用語には、巻末の説明へのリンクがある。
    下段の注釈には、条件の確認や短い補足を置いた。
    操作に必要な条件は、本文にも示す。
    PDFの目次・見出し参照・例題の見出しはリンクになっている.
    例題・演習問題の番号から解答へ,方針・解答・解法の番号から問題へ移動できる.

    特に断りが無い限りは,\,$k,l,m,n$は1以上の整数,\,$a,b,c,d,p,q,r,s,t,u,v,x,y,\lambda$は実数として扱う.
    また,\,$\alpha,\beta,\gamma,\delta$は方程式の解として扱う.
    ただし$d$は\bookterm{arithmetic}{等差数列}の公差を表す文字として,\,$0$を含む任意の実数を取りうる.また,添字は$n\ge1$を基本とするが,\,\sectionref{節：虚数解型隣接3項間漸化式}や付録の一部では$a_0$から始まる添字を用いることがある.
    また,\,$z$は複素数として扱う.
    また,\,$\therefore\,,\because\,,\Rightarrow\,,\Leftrightarrow$は左から順に「従って」「なぜならば」「左側が成り立つならば右側が成り立つ」「同値である」を意味する.
    加えて,\,$\cdots$や$\vdots$は途中を省略する時に用いる.

\end{bookcolumns}

% chapter: 計算の基本

\chapter{計算の基本}\label{章：計算の基本}
\begin{multicols*}{2}
  漸化式の解説を行う前に,実際の計算において基本となる内容をこの章では説明する.
  ここで扱う内容は漸化式の一般項を得る過程だけでなく,様々な場面での計算において有用なものである.
  ごく狭い場面で使えるものは避け,多くの場面で使われる手法と,それらに通底する思想を取り扱っている.
  前半では和の計算,式変形,指数・対数を確認し,続いて数学的帰納法と二項定理を扱う.
  後半では三角関数,極形式,複素数の順に,数列の値を角度や平面上の点として表すための知識を準備する.
  複素数の節では,先に説明した三角関数の加法定理と数学的帰納法を用いる.
% section: シグマの計算

\section{シグマの計算}\label{節：シグマの計算}
  \noindent\textbf{【本編で使う場面】}\par
  和の記号と添字の役割を確認する.\sectionref*{節：階差型}で変化を足し戻し,\sectionref*{節：総和型}で総和の式をずらして引くときに使う.

  多くの項を足すとき、すべてを並べる代わりに、一つの記号でまとめられる。
  実数の\bookterm{sequence}{数列}$\{a_k\}$について、$a_1$から$a_n$までの和を
  \[
    \sum_{k=1}^{n} a_k=a_1+a_2+\cdots+ a_n
  \]
  と書く。複素数の\bookterm{sequence}{数列}でも、同じ記号を使う。
  $\sum$は、ギリシャ文字の\bookterm{sigma}{シグマ}である。
  下の$k=1$から上の$n$まで、$k$を一つずつ進めて足す。

  同じ数を繰り返し足せば、項数を掛けた値になる。
  定数倍は和の外へ出せる。二つの和は、項ごとに分けてもよい。
  \begin{align*}
    &\sum_{k=1}^{n} c=nc\\
    &\sum_{k=1}^{n} ca_k=c\sum_{k=1}^{n} a_k\\
    &\sum_{k=1}^{n} (a_k+b_k)=\sum_{k=1}^{n} a_k+\sum_{k=1}^{n} b_k
  \end{align*}
  自然数と、その二乗・三乗の和は、次の式でまとめられる。
  \begin{align*}
    &\sum_{k=1}^{n} k=\frac{1}{2}n(n+1)\\
    &\sum_{k=1}^{n} k^2=\frac{1}{6}n(n+1)(2n+1)\\
    &\sum_{k=1}^{n} k^3=\left\{\frac{1}{2}n(n+1)\right\}^2\\
  \end{align*}

  添字$k$は,和の中で順に動く番号である.同じ範囲なら$\sum_{k=1}^{n}a_k=\sum_{j=1}^{n}a_j$と書き換えても値は変わらない.
  一方,上端の$n$は足す項数を決める.例えば$S_n=\sum_{k=1}^{n}a_k$なら
  \[
    S_{n+1}=S_n+a_{n+1}
  \]
  である.
  \bookterm{recurrence}{漸化式}の添字をずらす場合は,式全体の$n$を置き換える.
  例えば$a_{n+1}=2a_n+n$から$a_{n+2}=2a_{n+1}+n+1$を得る.

% section: 式の変形

\section{式の変形}\label{節：式の変形}
  \noindent\textbf{【本編で使う場面】}\par
  同じ値を保ちながら既知の形を作る操作を確認する.定数の調整は\sectionref*{節：特性方程式型},\bookterm{partial}{部分分数分解}は\sectionref*{節：階差型}で和を計算する際に使う.

  \subsection{\texorpdfstring{$+a-a$をする}{+a-aをする}}\label{節：+a-aをする}
    二次式の中に、二乗の形を作る。
    $a\ne0$として、足りない定数を足し、同じだけ引くと、
    \begin{align*}
      &ax^2+bx+c\\
      &\qquad=a\left(x^2+\frac{b}{a}x\right) +c\\
      &\qquad=a\left(x^2+\frac{b}{a}x+\frac{b^2}{4a^2}-\frac{b^2}{4a^2}\right)+c \tag*{\(\cdots\) (☆)}\\
      &\qquad=a\left(x+\frac{b}{2a}\right)^2-\frac{b^2}{4a}+c
    \end{align*}
    となる。☆の行で、二乗に必要な定数を補った。
    足して引いた量は零なので、式の値は変わらない。
    $x$を含む部分が一つの二乗にまとまり、式の形が見やすくなる。
    この操作を\bookterm{square}{平方完成}という。
    二次関数の最大・最小や、軸の位置を調べるときにも使える。

    値を保ったまま、扱いやすい形を作る。
    同じ発想が、後の定数項を取り除く操作につながる。

  \subsection{部分分数分解}\label{節：部分分数分解}
    次は、積を含む分母を、より簡単な分数の和へ分ける。
    \[
      \frac{1}{(x+a)(x+b)}
    \]
    は、そのままでは和を計算しにくい。
    分母を一次の式に分けるため、次の形を作る。
    ここでは$x+a\ne0$、$x+b\ne0$とする。
    \[
      \frac{1}{(x+a)(x+b)}=\frac{\alpha}{x+a}+\frac{\beta}{x+b}
    \]
    $\alpha,\beta$について求めると,
    \begin{align*}
    \frac{1}{(x+a)(x+b)}
      &= \frac{\alpha}{x+a}+\frac{\beta}{x+b}\\
      &= \frac{(\alpha+\beta)x+\alpha b+\beta a}
      {(x+a)(x+b)}
    \end{align*}
    \[
    \therefore\quad
    \begin{cases}
    \alpha+\beta=0,\\
    \alpha b+\beta a=1
    \end{cases}
    \]
    これを解くと,\,$\beta=-\alpha$であるから,
    \[
      \alpha b-\alpha a=1
      \iff \alpha(b-a)=1
    \]
    したがって,\,$a\ne b$であれば,
    \[
      \alpha=\frac{1}{b-a},
      \qquad
      \beta=-\frac{1}{b-a}
    \]
    が得られる(\,$a=b$のときはそもそも分母が重解$(x+a)^2$になり,この形の分解はできない).
    二つの係数を定めると、一次の分母をもつ分数の和に分かれる。
    この操作を\bookterm{partial}{部分分数分解}という。
    3次の場合も,\,$p,q,r\ne0$で,分母の一次因子$px+a,qx+b,rx+c$の根
    \[
      -\frac{a}{p},\qquad -\frac{b}{q},\qquad -\frac{c}{r}
    \]
    が互いに異なるときは,
\BookMathLayout{\[
      \frac{1}{(px+a)(qx+b)(rx+c)}=\frac{\alpha}{px+a}+\frac{\beta}{qx+b}+\frac{\gamma}{rx+c}
    \]}{
\[
\begin{aligned}
\frac{1}{(px+a)(qx+b)(rx+c)}
&=\frac{\alpha}{px+a}+\frac{\beta}{qx+b}\\
&\quad+\frac{\gamma}{rx+c}
\end{aligned}
\]
}
    と置き,係数$\alpha,\beta,\gamma$を定めることで分解できる.
    言い換えれば,これらの一次因子が互いに定数倍でないことが条件である.
    重複因子があるときは,その重複度に応じて例えば$\dfrac{A}{x+a}+\dfrac{B}{(x+a)^2}$のような項も必要になる.

% section: 等比数列の総和とその導出

\section{等比数列の総和とその導出}\label{節：等比数列の総和とその導出}
  初項が$a$,公比が$r$である等比数列
  \[
    a,\,ar,\,ar^2,\,\ldots,\,ar^{n-1}
  \]
  の初項から第$n$項までの総和を考える.
  その総和を$S_n$とすると,
  \[
    S_n
    =\sum_{k=1}^{n} ar^{k-1}
    =a+ar+ar^2+\cdots+ar^{n-1}
  \]
  と表せる.

  \subsection{総和の公式}
    公比$r$が$1$でないとき,
    \[
      S_n=\frac{a(1-r^n)}{1-r}
    \]
    である.
    分母と分子の符号を同時に変えれば,
    \[
      S_n=\frac{a(r^n-1)}{r-1}
    \]
    とも書ける.
    公比が$1$のときは,すべての項が$a$であるから,
    \[
      S_n=na
    \]
    となる.

  \subsection{公式の導出}
    $r\ne1$として,総和の式の両辺に$r$を掛ける.
    \begin{align*}
      S_n
        &=a+ar+ar^2+\cdots+ar^{n-1},\\
      rS_n
        &=ar+ar^2+ar^3+\cdots+ar^n.
    \end{align*}
    2つの式を項ごとに引くと,右辺では共通する項がすべて消える.
    \begin{align*}
      S_n-rS_n
        &=(a+ar+ar^2+\cdots+ar^{n-1})\\
        &\qquad -(ar+ar^2+ar^3+\cdots+ar^n)\\
        &=a-ar^n.
    \end{align*}
    よって,
    \[
      (1-r)S_n=a(1-r^n)
    \]
    である.
    $r\ne1$より両辺を$1-r$で割ると,
    \[
      S_n=\frac{a(1-r^n)}{1-r}
    \]
    を得る.

    この導出では,総和を公比$r$倍してから元の総和との差を取ることで,途中の項を消している.
    数列の和を求めるときには,このように総和を定数倍して引き,共通する項を消す考え方がよく用いられる.


% section: 指数対数

  \section{指数対数}
    \subsection{指数}
      $a$が$p$回掛けられている状況を考える.
      ここで
      \[
        a\times a\times\cdots\times a
      \]
      のように記すのは煩雑なので,
      \[
        a\times a\times\cdots\times a=a^p
      \]
      と表す.
      このときの$p$を指数といい,
      \[
        f(x)=a^x
      \]
      で表されるような関数を指数関数という.
      ここでは$p$を正の整数として導入したが,\,$a^0=1,\ a^{-p}=\dfrac1{a^p}$のように定義を拡張すれば,指数が$0$や負の整数,有理数,実数であっても同様に$a^x$を考えることができる.
      以降,特に断らない限りこの拡張された意味で指数を用いる.
    \subsection{対数}
      $a^p=q$のとき,\,$p$は
      \[
        \log_a{q}=p
      \]
      のように表される.
      この時$\log$をログと読み,\,$p$を対数という.
      また,\,$a$を底,\,$q$を真数と言う.
      一般に底$a$は$a>0,\,a\ne1$を満たす必要がある.
      また,
      \[
        e=\lim_{x \to \infty} \left(1+\frac{1}{x}\right)^x
      \]
      で定義される数$e$を自然対数の底またはネイピア数という.また,
      \[
        f(x)=\log_a x
      \]
      というような関数を対数関数といい,指数関数の逆関数である.
      逆関数とは
      \[
        y=f(x)
      \]
      を満たすような$x,y$について,
      \[
        x=g(y)
      \]
      となるような関数$g$を指す.
      グラフ上では,逆関数は元の関数と関数$y=x$に対して対称である.
      関数$g$はよく,\,$f^{-1}$と書かれる.

      対数に関して,一般に以下が成り立つ.
      \begin{align*}
        &\log_a{p^q}=q\log_a{p}\\
        &\log_a{MN}=\log_a{M}+\log_a{N}\\
        &\log_a{\frac{M}{N}}=\log_a{M}-\log_a{N}\\
        &\log_a{b}=\frac{\log_c{b}}{\log_c{a}}
      \end{align*}
    

% section: 数学的帰納法

  \section{数学的帰納法}
    自然数$n$についての主張$P(n)$が,すべての$n\ge1$で成り立つことを示したいとする.
    このとき,以下の2つを示せばよい.
    \begin{enumerate}[label=(\roman*),parsep=3pt]
      \item $P(1)$が成り立つ.
      \item ある$n$で$P(n)$が成り立つと仮定すると,\,$P(n+1)$も成り立つ.
    \end{enumerate}
    この2つが示されれば,\,(i)より$P(1)$が成り立ち,\,(ii)より$P(1)\Rightarrow P(2)$,\,$P(2)\Rightarrow P(3)$,\,$\ldots$と次々に成り立つことが分かるので,すべての$n\ge1$について$P(n)$が成り立つ.
    ドミノ倒しに例えられることが多い.\,(i)は最初のドミノを倒すこと,\,(ii)は「あるドミノが倒れれば次のドミノも倒れる」ように並べることに対応する.

    本書では,漸化式から予想した一般項の形を確かめる場面や,三角関数の加法定理から漸化式を導く場面(3.2や3.6.3など)で数学的帰納法を用いる.
% section: 二項定理

  \section{二項定理}
    次に,和の累乗を展開するための公式を確認する.
    自然数$n$と実数$x,y$について,
    \[
      (x+y)^n=\sum_{k=0}^{n}\binom{n}{k}x^{n-k}y^{k}
    \]
    が成り立つ.これを二項定理という.ここで$\dbinom{n}{k}$は二項係数であり,
    \[
      \binom{n}{k}=\frac{n!}{k!(n-k)!}
    \]
    で定義される. ここで,$k$は$0\le k\le n$を満たす整数であり,$0!=1$とする.
    この二項係数は,$n$個から$k$個を選ぶ組合せの数に一致する.
    例えば$n=2$では,
    \[
      (x+y)^2=x^2+2xy+y^2
    \]
    $n=3$では,
    \[
      (x+y)^3=x^3+3x^2y+3xy^2+y^3
    \]
    となる.
    後の節で扱う複素数に対しても,同じ二項定理が成り立つ.
    付録では,この公式を用いて差分を取ると多項式の次数が下がる理由などを説明する.
\section{三角関数}
  ここからは,数列の値を角度や平面上の点として表すための準備を行う.
  まず,極形式や複素数の計算でも用いる三角関数の公式を振り返る.
  まず一般に,
  \begin{align*}
    \sin(\alpha+\beta)
      &=\sin\alpha\cos\beta+\cos\alpha\sin\beta\\
    \cos(\alpha+\beta)
      &=\cos\alpha\cos\beta-\sin\alpha\sin\beta\\
    \tan(\alpha+\beta)
      &=\frac{\tan\alpha+\tan\beta}
              {1-\tan\alpha\tan\beta}
  \end{align*}
  が成り立つ.
  ここで $\alpha=\beta=x$ とすれば,倍角公式
  \begin{align*}
    \sin 2x&=2\sin x\cos x\\
    \cos 2x&=2\cos^2x-1=1-2\sin^2x\\
    \tan 2x&=\frac{2\tan x}{1-\tan^2x}
  \end{align*}
  を得る.また,倍角公式から
  \begin{align*}
    \sin^2x&=\frac{1-\cos 2x}{2}\\
    \cos^2x&=\frac{1+\cos 2x}{2}\\
    \tan^2x&=\frac{1-\cos 2x}{1+\cos 2x}
  \end{align*}
  を導くことができる.また,加法定理から,\,$3x=2x+x$とすることで,
  \begin{align*}
    \sin 3x&=3\sin x-4\sin^3x\\
    \cos 3x&=4\cos^3x-3\cos x
  \end{align*}
  が求められる.
\section{極形式}
  \subsection{極座標系の導入}
    前節で確認した三角関数を用いると,平面上の点を距離と角度によって表すことができる.
    普段用いている直交座標系と比較しながら,この表し方を導入する.
    直交座標系とは,デカルト座標系とも呼ばれ,平面上で直交する2直線の交点を基準として任意の点を表す手法のことである.
    しかしながら円や媒介変数表示された関数などにおいて,一部煩雑な表現をせざるをえない状況があった.
    また,図形の平行移動や対称移動に関しては容易に行えたものの,回転させたり歪ませたりするような操作は非常に複雑な操作が必要だった.
    そこで,直交座標系に加え,新たな概念として極座標系の導入を行う.
    
    極座標系とは,平面上の任意の点を原点からの距離と角度を用いて記述する方法である.
    例えば

    \begin{tikzpicture}[scale=0.8]
      % 座標軸
      \draw[->,>=stealth,semithick] (-1.5,0)--(5.5,0) node[above]{$x$};
      \draw[->,>=stealth,semithick] (0,-1)--(0,5) node[right]{$y$};

      % 原点
      \draw (0,0) node[below right]{O};

      % 点 (1,sqrt(3)) への補助線
      \draw[dashed] (1,0) -- (1,{sqrt(3)});
      \draw[dashed] (0,{sqrt(3)}) -- (1,{sqrt(3)});

      % 座標軸上の値
      \draw (1,0) node[below]{$1$};
      \draw (0,{sqrt(3)}) node[left]{$\sqrt{3}$};

      % 点
      \fill (1,{sqrt(3)}) circle (2pt)
        node[above right]{$(1,\sqrt{3})$};
    \end{tikzpicture}

    のように表される点は,極座標系においては

\begin{tikzpicture}[scale=0.8]

  % 座標軸
  \draw[->,>=stealth,semithick]
    (-1.5,0)--(5.5,0) node[above]{$x$};
  \draw[->,>=stealth,semithick]
    (0,-1)--(0,5) node[right]{$y$};

  % 原点・点A
  \coordinate (O) at (0,0);
  \coordinate (A) at (1,{sqrt(3)});
  \node[below left] at (O) {O};

  % 動径
  \draw[very thick] (O)--(A);

  % 下側の曲線
  \draw[semithick]
    (O)
    .. controls (0.02,0.35) and (0.12,0.62)
    .. (0.3,0.7);

  % 上側の曲線
  % 終点Aでの接線を動径OAと平行にする
  \draw[semithick]
    (0.6,1.4)
    .. controls (0.62,1.58) and
               (0.82,{sqrt(3)-0.10*sqrt(3)})
    .. (A);

  % 動径の長さ
  \node at (0.30,1.05) {$2$};

  % 角度
  \draw[semithick]
    (0.55,0) arc (0:60:0.55);
  \node at (0.68,0.3) {$\frac{\pi}{3}$};

  % 点
  \fill (A) circle (1.5pt);

\end{tikzpicture}
    と表される. この点の原点からの距離は$2$,正の$x$軸から測った角度は$\pi/3$である.
  \subsection{極形式}
    一般に,原点からの距離を$\rho$,正の$x$軸から反時計回りに測った角度を$\theta$とすると,
    点の直交座標$(x,y)$は
    \[
      x=\rho\cos\theta,\qquad y=\rho\sin\theta
    \]
    と表される. ここで$\rho=\sqrt{x^2+y^2}$である.
    原点以外では$\rho>0$であり,$\theta$に$2\pi$の整数倍を加えても同じ点を表す.
    原点では$\rho=0$となり,角度は一意に定まらない.

% section: 複素数の極形式と指数表示

\section{複素数の極形式と指数表示}

前節では,平面上の点を原点からの距離と角度によって表した.
ここでは,その点を複素数として読み替える.
$i^2=-1$を満たす虚数単位$i$を用いると,実数$x,y$に対して複素数は$z=x+iy$と表される.
この複素数を平面上の点$(x,y)$に対応させたものを複素平面という.
前節の$x=\rho\cos\theta,\,y=\rho\sin\theta$を代入すれば,
\[
  z
  =
  \rho(
    \cos\theta+i\sin\theta
  )
\]
と書くことができる. $\rho$は複素数の絶対値であり,原点からの距離を表す.
$z\ne0$のとき,$\theta$を偏角という. 偏角は$2\pi$の整数倍の違いを除いて定まる.

このとき,
\[
  \cos\theta+i\sin\theta
\]
のような複素数を何乗かしたときに,
偏角がどのように変化するかを考える.

まず,三角関数の加法定理を用いる.


\subsection{極形式の積}

  任意の実数$\alpha,\beta$について,
  \begin{align*}
    &(\cos\alpha+i\sin\alpha)
    (\cos\beta+i\sin\beta)\\
    &=
    \cos\alpha\cos\beta
    -\sin\alpha\sin\beta\\
    &\qquad
    +i(
      \sin\alpha\cos\beta
      +\cos\alpha\sin\beta
    ).
  \end{align*}

  加法定理より,
  \[
    \cos\alpha\cos\beta
    -\sin\alpha\sin\beta
    =
    \cos(\alpha+\beta)
  \]
  および
  \[
    \sin\alpha\cos\beta
    +\cos\alpha\sin\beta
    =
    \sin(\alpha+\beta)
  \]
  であるから,
  \[
    (\cos\alpha+i\sin\alpha)
    (\cos\beta+i\sin\beta)
    =
    \cos(\alpha+\beta)
    +i\sin(\alpha+\beta)
  \]
  を得る.

  つまり,
  \[
    \cos\alpha+i\sin\alpha
  \]
  と
  \[
    \cos\beta+i\sin\beta
  \]
  を掛けると,絶対値は1のままで偏角が
  \[
    \alpha+\beta
  \]
  となる.

  この性質を繰り返し用いることで,
  複素数の累乗について重要な公式が得られる.


\subsection{ド・モアブルの定理}

  任意の整数$n$について,
  \[
    (\cos\theta+i\sin\theta)^n
    =
    \cos(n\theta)+i\sin(n\theta)
  \]
  が成り立つ.

  まず,$n$が正の整数の場合を考える.

  $n=1$のとき,
  \[
    (\cos\theta+i\sin\theta)
    =
    \cos\theta+i\sin\theta
  \]
  であるから,成り立つ.

  いま,
  \[
    (\cos\theta+i\sin\theta)^n
    =
    \cos(n\theta)+i\sin(n\theta)
  \]
  が成り立つと仮定する.

  このとき,
  \begin{align*}
    &(\cos\theta+i\sin\theta)^{n+1}\\
    &=
    \{
      \cos(n\theta)+i\sin(n\theta)
    \}
    (\cos\theta+i\sin\theta)\\
    &=
    \cos((n+1)\theta)
    +i\sin((n+1)\theta)
  \end{align*}
  となる.

  よって,数学的帰納法により,
  \[
    (\cos\theta+i\sin\theta)^n
    =
    \cos(n\theta)+i\sin(n\theta)
  \]
  がすべての正の整数$n$について成り立つ.

  $n=0$のときは,
  \[
    (\cos\theta+i\sin\theta)^0
    =
    1
    =
    \cos0+i\sin0
  \]
  である.

  また,
  \[
    \cos(-\theta)=\cos\theta,
    \qquad
    \sin(-\theta)=-\sin\theta
  \]
  より,
  \begin{align*}
    &(\cos\theta+i\sin\theta)
    (\cos(-\theta)+i\sin(-\theta))\\
    &=
    \cos0+i\sin0\\
    &=
    1.
  \end{align*}

  したがって,
  \[
    (\cos\theta+i\sin\theta)^{-1}
    =
    \cos(-\theta)+i\sin(-\theta)
  \]
  である.

  これを用いれば,負の整数についても同じ公式が成り立つ.

  したがって,
  \[
    (\cos\theta+i\sin\theta)^n
    =
    \cos(n\theta)+i\sin(n\theta)
    \qquad(n\in\mathbb Z)
  \]
  が成り立つ.

  これをド・モアブルの定理という.

  例えば,
  \[
    \left(
      \cos\frac{\pi}{6}
      +i\sin\frac{\pi}{6}
    \right)^6
    =
    \cos\pi+i\sin\pi
    =
    -1
  \]
  となる.

  このように,複素数を極形式で表すと,累乗を考えるときには絶対値の累乗と偏角の$n$倍を考えればよいことが分かる.


\subsection{複素数の指数表示}

  次に,実数の指数関数で用いられる
  \[
    e^x
    =
    \lim_{n\to\infty}
    \left(
      1+\frac{x}{n}
    \right)^n
  \]
  という極限表示を複素数に拡張する.

  ここでは,
  \[
    e^{i\theta}
    :=
    \lim_{n\to\infty}
    \left(
      1+\frac{i\theta}{n}
    \right)^n
  \]
  と定める.

  この極限を評価するため,
  \[
    1+\frac{i\theta}{n}
  \]
  を極形式に直す.

  その絶対値は
  \[
    r_n
    =
    \sqrt{
      1+\frac{\theta^2}{n^2}
    }
  \]
  である.

  また,$\phi_n$を
  \[
    \tan\phi_n
    =
    \frac{\theta}{n}
  \]
  を満たす角とする.

  このとき,
  \[
    \cos\phi_n
    =
    \frac{1}{r_n}
  \]
  および
  \[
    \sin\phi_n
    =
    \frac{\theta}{nr_n}
  \]
  であるから,
  \[
    1+\frac{i\theta}{n}
    =
    r_n(
      \cos\phi_n+i\sin\phi_n
    )
  \]
  と表すことができる.

  よって,
  \begin{align*}
    \left(
      1+\frac{i\theta}{n}
    \right)^n
    &=
    r_n^n
    (\cos\phi_n+i\sin\phi_n)^n\\
    &=
    r_n^n
    \{
      \cos(n\phi_n)
      +i\sin(n\phi_n)
    \}.
  \end{align*}

  最後の等式にはド・モアブルの定理を用いた.

  したがって,この極限を求めるためには,
  \[
    r_n^n
  \]
  と
  \[
    n\phi_n
  \]
  の極限を調べればよい.


\subsection{極限の計算}

  まず,
  \[
    r_n^n
    =
    \left(
      1+\frac{\theta^2}{n^2}
    \right)^{n/2}
  \]
  について考える.

  自然対数の基本極限
  \[
    \lim_{x\to0}
    \frac{\log(1+x)}{x}
    =
    1
  \]
  を用いると,
  \begin{align*}
    \log(r_n^n)
    &=
    \frac n2
    \log
    \left(
      1+\frac{\theta^2}{n^2}
    \right)\\
    &=
    \frac{\theta^2}{2n}
    \frac{
      \log
      \left(
        1+\frac{\theta^2}{n^2}
      \right)
    }{
      \frac{\theta^2}{n^2}
    }.
  \end{align*}

  $n\to\infty$のとき,
  \[
    \frac{\theta^2}{2n}\to0
  \]
  であり,
  \[
    \frac{
      \log
      \left(
        1+\frac{\theta^2}{n^2}
      \right)
    }{
      \frac{\theta^2}{n^2}
    }
    \to1
  \]
  であるから,
  \[
    \log(r_n^n)\to0.
  \]

  よって,
  \[
    r_n^n\to1
  \]
  である.

  次に,
  \[
    n\phi_n
  \]
  の極限を考える.

  $\tan\phi_n=\dfrac{\theta}{n}$より,
  \[
    \tan\phi_n\to0
  \]
  であるから,
  \[
    \phi_n\to0
  \]
  である.

  また,
  \[
    \tan x
    =
    \frac{\sin x}{\cos x}
  \]
  および,
  \[
    \lim_{x\to0}\frac{\sin x}{x}=1,
    \qquad
    \lim_{x\to0}\cos x=1
  \]
  より,
  \[
    \lim_{x\to0}\frac{\tan x}{x}=1
  \]
  である.

  したがって,
  \[
    \lim_{x\to0}\frac{x}{\tan x}=1
  \]
  であるから,
  \begin{align*}
    n\phi_n
    &=
    \theta
    \frac{\phi_n}{\tan\phi_n}\\
    &\to
    \theta.
  \end{align*}

  よって,
  \[
    n\phi_n\to\theta
  \]
  が成り立つ.


\subsection{オイラーの公式}

  以上より,
  \[
    r_n^n\to1
  \]
  および
  \[
    n\phi_n\to\theta
  \]
  である.

  $\sin x,\cos x$は連続関数であるから,
  \[
    \cos(n\phi_n)\to\cos\theta
  \]
  および
  \[
    \sin(n\phi_n)\to\sin\theta
  \]
  となる.

  したがって,
  \begin{align*}
    e^{i\theta}
    &=
    \lim_{n\to\infty}
    \left(
      1+\frac{i\theta}{n}
    \right)^n\\
    &=
    \lim_{n\to\infty}
    r_n^n
    \{
      \cos(n\phi_n)
      +i\sin(n\phi_n)
    \}\\
    &=
    \cos\theta+i\sin\theta.
  \end{align*}

  よって,
  \[
    e^{i\theta}
    =
    \cos\theta+i\sin\theta
  \]
  が成り立つ.

  これをオイラーの公式という.

  したがって,
  \[
    \cos\theta+i\sin\theta
    =
    e^{i\theta}
  \]
  と書くことができる.

  これにより,複素数
  \[
    z
    =
    \rho(
      \cos\theta+i\sin\theta
    )
  \]
  は,
  \[
    z
    =
    \rho e^{i\theta}
  \]
  と表せる.


\subsection{ド・モアブルの定理の指数表示}

  オイラーの公式を用いると,
  ド・モアブルの定理は指数関数を使ってさらに簡潔に表すことができる.

  実際,
  \[
    \cos\theta+i\sin\theta
    =
    e^{i\theta}
  \]
  より,
  \begin{align*}
    (\cos\theta+i\sin\theta)^n
    &=
    (e^{i\theta})^n\\
    &=
    e^{in\theta}\\
    &=
    \cos(n\theta)+i\sin(n\theta).
  \end{align*}

  また,
  \[
    z
    =
    \rho e^{i\theta}
  \]
  と表せば,
  \begin{align*}
    z^n
    &=
    \rho^n e^{in\theta}\\
    &=
    \rho^n
    \{
      \cos(n\theta)+i\sin(n\theta)
    \}.
  \end{align*}

  したがって,
  \[
    z^n
    =
    \rho^n
    \{
      \cos(n\theta)+i\sin(n\theta)
    \}
    \qquad(n\in\mathbb Z)
  \]
  となる.

  ここでは,極形式における「絶対値は$n$乗し,偏角は$n$倍する」という性質が,
  指数表示によって簡潔に表されている.


\subsection{実数指数への拡張}

  これまでのド・モアブルの定理では,指数$n$は整数であった.

  ここで,指数を一般の実数にまで広げることを考える.

  まず,複素数
  \[
    z
    =
    \rho(
      \cos\theta+i\sin\theta
    )
    =
    \rho e^{i\theta}
    \qquad(\rho>0)
  \]
  に対して,実数$x$を指数とする場合を考える.

  実数の指数関数では,
  \[
    \rho^x
  \]
  が定義されているので,
  \[
    z^x
  \]
  を
  \[
    z^x
    :=
    \rho^x e^{ix\theta}
  \]
  と定める.

  オイラーの公式より,
  \[
    e^{ix\theta}
    =
    \cos(x\theta)+i\sin(x\theta)
  \]
  であるから,
  \[
    z^x
    =
    \rho^x
    \{
      \cos(x\theta)+i\sin(x\theta)
    \}
  \]
  となる.

  したがって,
  \[
    z^x
    =
    \rho^x
    \{
      \cos(x\theta)+i\sin(x\theta)
    \}
    \qquad(x\in\mathbb R)
  \]
  と表すことができる.

  $x$が整数の場合には,
  これはすでに示したド・モアブルの定理と一致する.

  また,$x=\dfrac{m}{n}$が有理数の場合には,
  \[
    z^{m/n}
    =
    \rho^{m/n}
    \left\{
      \cos\frac{m\theta}{n}
      +i\sin\frac{m\theta}{n}
    \right\}
  \]
  である.

  この数を$n$乗すると,
  \begin{align*}
    \left(
      z^{m/n}
    \right)^n
    &=
    \rho^m
    \left(
      \cos\frac{m\theta}{n}
      +i\sin\frac{m\theta}{n}
    \right)^n\\
    &=
    \rho^m
    \{
      \cos(m\theta)
      +i\sin(m\theta)
    \}\\
    &=
    z^m.
  \end{align*}

  したがって,$z^{m/n}$は$z^m$の$n$乗根となっており,
  この定義は有理数の場合にも通常の指数の意味と整合している.

  以上から,ド・モアブルの定理は,
  \[
    z^x
    =
    \rho^x
    \{
      \cos(x\theta)+i\sin(x\theta)
    \}
    \qquad(x\in\mathbb R)
  \]
  という形に拡張することができる.

  ただし,複素数の偏角は一般には一意ではなく,
  \[
    \theta,\quad
    \theta+2\pi,\quad
    \theta+4\pi,\quad\ldots
  \]
  はすべて同じ複素数を表す.

  このため,一般の複素数の実数乗を考えるときには,
  どの偏角を用いるかをあらかじめ固定しておく必要がある.

  本書では,
  \[
    z
    =
    \rho(
      \cos\theta+i\sin\theta
    )
  \]
  において$\theta$を一つ固定し,
  \[
    z^x
    =
    \rho^x
    \{
      \cos(x\theta)+i\sin(x\theta)
    \}
  \]
  と考えることにする.
\end{multicols*}

% chapter: 漸化式の解説

\chapter{漸化式の解説}\label{章：漸化式の解説}
\begin{bookcolumns}
% section: 漸化式の分類

\section{漸化式の分類}\label{節：漸化式の分類}
  \BookFigClassification

  ここまでに使った方法を、一枚の地図として見直す。
  まず\bookterm{recurrence}{漸化式}は解ける\bookterm{recurrence}{漸化式}と解けない\bookterm{recurrence}{漸化式}に分類される.
  そして解ける\bookterm{recurrence}{漸化式}は「決まった手順で解ける\bookterm{recurrence}{漸化式}」と「特定の性質から解ける\bookterm{recurrence}{漸化式}」に分類される.
  「決まった手順で解ける\bookterm{recurrence}{漸化式}」はさらに分類され,各分類ごとに様々な解法がある.
  ここで言う「解ける」というのは,本書内での定義であり,一般的なものではない.
  ここで解けないとして扱われるような\bookterm{recurrence}{漸化式}は\bookterm{general}{一般項}を表す手段が存在しないわけではない.
  あくまで本書で扱う手法の内では表現できないという意味である.
  「解ける」の定義について述べておく.
  定義は,高校数学程度の範囲内で既知の関数に帰着するか,\,$n$に依存しない有限個の\bookterm{elementary}{初等関数}の演算によって\bookterm{general}{一般項}を表せるような解法が存在するものを意味している.
  このうち「決まった手順で解ける\bookterm{recurrence}{漸化式}」は基本4パターン・応用7パターン・発展パターンの各節で,「特定の性質から解ける\bookterm{recurrence}{漸化式}」は特殊性質の節で扱う.

  基本4パターンでは、差と倍率から項を求めた。
  応用パターンでは、追う量を変えることで、その基本形を取り出した。
  発展パターンと特殊性質型では、追う対象と調べる性質がさらに広がった。

  この分類は,解法を探すための案内である.一つの式に複数の解法が使えることもある.
  本書の方法で\bookterm{closed}{閉じた形}が見つからなくても,\bookterm{sequence}{数列}が定義されないわけではなく,\bookterm{limit}{極限}や整数の性質を調べる価値も失われない.
  例えば\bookproblemref{practice-entrance}{6}{発展問6}では,\bookterm{general}{一般項}を求めずに減り方を調べる.

% section: 漸化式を解く前に

\section{漸化式を解く前に}\label{節：漸化式を解く前に}
  \BookFigApproach

  本章で繰り返し使う問いは,「元の項より,次への変化が簡単になる量はないか」である.
  例えば$a_{n+1}=2a_n+3$では,定数$3$があるため等比型の式をそのまま使えない.
  両辺に$3$を加えると
  \[
    a_{n+1}+3=2(a_n+3)
  \]
  となり,\,$b_n=a_n+3$という同じ形の量が2倍されていると分かる.
  なぜ$3$を選ぶかは\sectionref*{節：特性方程式型}で導く.
  ここでは,変形の目標を具体的に持つことを確認しよう.
  \begin{itemize}
    \item $b_{n+1}-b_n=f(n)$なら,変化を足し戻す（階差型）.
    \item $b_{n+1}=f(n)b_n$なら,倍率を掛け合わせる（階比型）.
    \item $b_{n+1}=b_n$なら,その量は初項から変わらない.
    \item 既知の公式と同じ形なら,その公式の入力に当たる量を追う.
  \end{itemize}
  本書で「きれいな形」と呼ぶのは,こうした更新の仕組みを追いやすい形である.
  和や積の表示が得られた後,それをどこまで計算できるかは別に確認する.

  説明を読むときは,\textbf{観察と発想}と\textbf{式変形}を分けて追ってほしい.
  前者では,どの項を消すか,どの倍率をそろえるかを選ぶ.
  後者では,選んだ量を実際に代入し,初項から一般項を求め,元の量へ戻す.
  逆数・対数・割り算を使う前には,その操作が定義できるかも確かめる.
  候補を探す計算の途中で条件が必要だと分かった場合は,そこで条件を確認してから続ける.

  最初の数項を計算する実験も,変形の候補を見つける助けになる.
  数項の一致は予想の根拠であり,すべての項についての証明にはならない.
  予想した式が初期条件と元の更新式を満たせば,順に次の項が一意に決まる漸化式では,数学的帰納法で一致を示せる.
  \sectionref*{節：数学的帰納法}を参照してほしい.
  帰納法は,変形から得た式を確かめるときにも,実験から得た予想を証明するときにも使える.

  一般項を簡単な式で表すことに加えて,正値・極限・整数性などを調べる目標もある.
  \chapterref{章：実際の解き方}では方法の選択を練習し,発展演習と付録では目標の広がりを見る.

% section: 基本4パターン

\section{基本4パターン}\label{節：基本4パターン}
  \BookFigBasicFour

  差が分かれば、変化を足し戻せる。
  倍率が分かれば、その倍率を掛け合わせられる。
  ここでは、差と倍率が一定の場合から、添字によって変わる場合へ進む。
  後の解法では、元の式を変形し、この四つの仕組みを取り出していく。
  特に等比型は、多くの変形の行き先になる。

% subsection: 等差型 \, $a_{n+1}=a_n+d$

  \subsection{\texorpdfstring{等差型 \, $a_{n+1}=a_n+d$}{等差型}}\label{節：等差型}
  \begin{bookpoint}
    \textbf{差を追う}\par
    一回ごとに足す量が一定なら,初項に増分を足し戻す.項の番号と更新の回数の違いを確かめよう.
  \end{bookpoint}
    まずはもっともシンプルな\bookterm{recurrence}{漸化式}であろう,等差型について説明していく.
    この\bookterm{recurrence}{漸化式}は$a_1$に対して$d$がどんどんと足されていく\bookterm{recurrence}{漸化式}である.
    \bookterm{general}{一般項}は簡単な代入計算によって容易に示せる.
    実際にやってみると
    \begin{align*}
      &a_{n+1}=a_n+d\\
      &a_{n+1}=(a_{n-1}+d)+d\\
      &a_{n+1}=\{(a_{n-2}+d)+d\}+d\\
      &\qquad \vdots\\
      &a_{n+1}=a_1+d+d+\cdots+d\\
      &a_{n+1}=a_1+nd\\
      &\therefore \quad a_n=a_1+(n-1)d
    \end{align*}
    となり\bookterm{general}{一般項}が示される.
    1つ前の項をどんどん代入していくのだ.
    数式の3行目から5行目周辺の式変形が少し心配かもしれないが,\,$d$の個数と$a$の添字部分の和は常に$n+1$となることから考えれば分かるだろう.
    実際,\,$a_{n+1}$に至るまでに何度$d$が足されたかを考えているわけだから,心配であれば実際に適当な$n$で何度か実験してみるのもいい手だろう.

    実際に例題を与えよう.

    \noindent \bookblocklink{example-arithmetic-1}{problem}{answer}{【例題】}\par\nobreak
    \begin{enumerate}[label=(\arabic*),parsep=3pt,format=\bookitemlink{example-arithmetic-1}{problem}{answer}]
      \item $\displaystyle a_1=1,\, a_{n+1}=a_n+1$
      \item $\displaystyle a_1=\frac{1}{2},\, a_{n+1}=a_n+\frac{1}{4}$
      \item $\displaystyle a_1=-\sqrt{3},\, a_{n+1}=a_n-\sqrt{2}$
    \end{enumerate}

    \noindent\bookblocklink{example-arithmetic-1}{hint}{problem}{【方針】}\par\nobreak
    \begin{enumerate}[label=(\arabic*),parsep=3pt,format=\bookitemlink{example-arithmetic-1}{hint}{problem}]
      \item 等差型の\bookterm{general}{一般項}
            $\displaystyle a_n=a_1+(n-1)d$
            を用います.
      \item 公差は$\displaystyle d=\frac14$です.
            等差型の\bookterm{general}{一般項}に代入します.
      \item 公差は$\displaystyle d=-\sqrt2$です.
            負の公差の場合も同じ公式を使います.
    \end{enumerate}

    \noindent\bookblocklink{example-arithmetic-1}{answer}{problem}{【解答】}\par\nobreak
    \begin{enumerate}[label=(\arabic*),parsep=3pt,format=\bookitemlink{example-arithmetic-1}{answer}{problem}]
      \item $\displaystyle a_n=1+(n-1)=n$
      \item $\displaystyle a_n=\frac12+\frac{n-1}{4}
            =\frac{n+1}{4}$
      \item $\displaystyle a_n=-\sqrt3-(n-1)\sqrt2$
    \end{enumerate}

    \noindent\bookblocklink{example-arithmetic-1}{solution}{problem}{【解法】}\par\nobreak
    解法略.

% subsection: 等比型 \, $a_{n+1}=ra_n$

  \subsection{\texorpdfstring{等比型 \, $a_{n+1}=ra_n$}{等比型}}\label{節：等比型}
  \begin{bookpoint}
    \textbf{倍率を追う}\par
    一回ごとの倍率が一定なら,更新の回数だけ掛ける.公比が零の場合も,比で割らず更新の式から読む.
  \end{bookpoint}
    次に説明するのは等比型である.
    一般的な解法も先ほどと同じで,\,$a_{n+1}$から順に$a_1$に向けて書き下して行くことで求めることができる.
    \begin{align*}
      &a_{n+1}=ra_n\\
      &a_{n+1}=r(ra_{n-1})\\
      &a_{n+1}=r\{r(ra_{n-2})\}\\
      &\qquad \vdots\\
      &a_{n+1}=r\cdot r \cdot \cdots \cdot ra_1\\
      &a_{n+1}=a_1r^n\\
      &\therefore \quad a_n=a_1r^{n-1}
    \end{align*}

    \noindent \bookblocklink{example-geometric-1}{problem}{answer}{【例題】}\par\nobreak
    \begin{enumerate}[label=(\arabic*),parsep=3pt,format=\bookitemlink{example-geometric-1}{problem}{answer}]
      \item $\displaystyle a_1=1,\, a_{n+1}=2a_n$
      \item $\displaystyle a_1=1,\, a_{n+1}=-a_n$
      \item $\displaystyle a_1=2,\, a_{n+1}=\frac{1}{2}a_n$
    \end{enumerate}

    \noindent\bookblocklink{example-geometric-1}{hint}{problem}{【方針】}\par\nobreak
    \begin{enumerate}[label=(\arabic*),parsep=3pt,format=\bookitemlink{example-geometric-1}{hint}{problem}]
      \item 等比型の\bookterm{general}{一般項}
            $\displaystyle a_n=a_1r^{n-1}$
            を用います.
      \item 公比は$r=-1$です.
            符号が交互に変化することに注意します.
      \item 初項と公比を\bookterm{general}{一般項}に代入します.
    \end{enumerate}

    \noindent\bookblocklink{example-geometric-1}{answer}{problem}{【解答】}\par\nobreak
    \begin{enumerate}[label=(\arabic*),parsep=3pt,format=\bookitemlink{example-geometric-1}{answer}{problem}]
      \item $\displaystyle a_n=2^{n-1}$
      \item $\displaystyle a_n=(-1)^{n-1}$
      \item $\displaystyle a_n=2\left(\frac12\right)^{n-1}$
    \end{enumerate}

    \noindent\bookblocklink{example-geometric-1}{solution}{problem}{【解法】}\par\nobreak
    解法略.

% subsection: 階差型 \, $a_{n+1}=a_n+f(n)$

  \subsection{\texorpdfstring{階差型 \, $a_{n+1}=a_n+f(n)$}{階差型}}\label{節：階差型}
  \begin{bookpoint}
    \textbf{変化を足す}\par
    差が添字ごとに変わっても,初項からの変化を足せば元の項に戻る.和の範囲は実際の更新の回数に合わせる.
  \end{bookpoint}
    足す量が一定でなくても、初項からの変化を足し戻せる。
    一歩目では$f(1)$、二歩目では$f(2)$を足す。
    直前の項を代入すると、その変化が順に並ぶ。
    \begin{align*}
      &a_{n+1}=a_n+f(n)\\
      &a_{n+1}=\{a_{n-1}+f(n-1)\}+f(n)\\
      &a_{n+1}=\left[\{a_{n-2}+f(n-2)\}+f(n-1)\right]+f(n)\\
      &\qquad \vdots\\
      &a_{n+1}=a_1+f(1)+f(2)+\cdots+f(n)\\
      &a_{n+1}=a_1+\sum_{k=1}^{n} f(k)\\
      &\therefore \quad a_n=a_1+\sum_{k=1}^{n-1} f(k)
    \end{align*}
    変化を足すところまでで、元の\bookterm{sequence}{数列}は和の形に表せた。
    残るのは、その和の計算である。
    $f(n)$の形に合わせて、\bookterm{arithmetic}{等差数列}の和、\bookterm{geometric}{等比数列}の和、\bookterm{partial}{部分分数分解}を使う。

    \noindent \bookblocklink{example-difference-1}{problem}{answer}{【例題】}\par\nobreak
    \begin{enumerate}[label=(\arabic*),parsep=3pt,format=\bookitemlink{example-difference-1}{problem}{answer}]
      \item $\displaystyle a_1=1,\, a_{n+1}=a_n+n$
      \item $\displaystyle a_1=1,\, a_{n+1}=a_n+2^n$
      \item $\displaystyle a_1=0,\, a_{n+1}=a_n+\frac{1}{n(n+1)}$
      \item $\displaystyle a_1=1,\, a_{n+1}=a_n+(n+1)^2$
    \end{enumerate}

    \noindent\bookblocklink{example-difference-1}{hint}{problem}{【方針】}\par\nobreak
    \begin{enumerate}[label=(\arabic*),parsep=3pt,format=\bookitemlink{example-difference-1}{hint}{problem}]
      \item 足す量は$n$です。初項からの変化を、一次式の和にまとめます。
      \item もちろん,\,$f(n)$となる関数には指数関数も含まれます.何にせよやることは\,(1)と同じです.
      \item \bookterm{sigma}{シグマ}は\bookterm{partial}{部分分数分解}で処理しましょう.
      \item 足す量は$(n+1)^2$です。初項を含めて、平方の和としてまとめます。
    \end{enumerate}

    \noindent\bookblocklink{example-difference-1}{answer}{problem}{【解答】}\par\nobreak
    \begin{enumerate}[label=(\arabic*),parsep=3pt,format=\bookitemlink{example-difference-1}{answer}{problem}]
      \item $\displaystyle a_n=1+\frac{n(n-1)}{2}$
      \item $\displaystyle a_n=2^n-1$
      \item $\displaystyle a_n=1-\frac{1}{n}$
      \item $\displaystyle a_n=\frac{1}{6}n(n+1)(2n+1)$
    \end{enumerate}

    \noindent\bookblocklink{example-difference-1}{solution}{problem}{【解法】}\par\nobreak
    \begin{enumerate}[label=(\arabic*),parsep=3pt,format=\bookitemlink{example-difference-1}{solution}{problem}]
      \item $\displaystyle a_1=1,\,a_{n+1}=a_n+n$\\
            \bookterm{difference}{階差}型の\bookterm{general}{一般項}
            \[
              a_n=a_1+\sum_{k=1}^{n-1}f(k)
            \]
            を用いると,
            \begin{align*}
              a_n
                &=1+\sum_{k=1}^{n-1}k\\
                &=1+\frac{n(n-1)}{2}
            \end{align*}

      \item $\displaystyle a_1=1,\,a_{n+1}=a_n+2^n$\\
            同様に,
            \begin{align*}
              a_n
                &=1+\sum_{k=1}^{n-1}2^k\\
                &=1+(2+2^2+\cdots+2^{n-1})\\
                &=1+(2^n-2)\\
                &=2^n-1
            \end{align*}

      \item $\displaystyle a_1=0,\,a_{n+1}=a_n+\frac{1}{n(n+1)}$\\
            \bookterm{partial}{部分分数分解}を用いると,
            \[
              \frac{1}{k(k+1)}
              =\frac{1}{k}-\frac{1}{k+1}
            \]
            であるから,
            \begin{align*}
              a_n
                &=\sum_{k=1}^{n-1}\frac{1}{k(k+1)}\\
                &=\sum_{k=1}^{n-1}
                  \left(\frac{1}{k}-\frac{1}{k+1}\right)\\
                &=\left(1-\frac12\right)
                  +\left(\frac12-\frac13\right)\\
                &\quad +\cdots+\left(\frac{1}{n-1}-\frac1n\right)\\
                &=1-\frac1n
            \end{align*}

      \item $\displaystyle a_1=1,\,a_{n+1}=a_n+(n+1)^2$\\
            \bookterm{difference}{階差}型の\bookterm{general}{一般項}を用いると,
            \begin{align*}
              a_n
                &=1+\sum_{k=1}^{n-1}(k+1)^2\\
                &=1+\sum_{j=2}^{n}j^2\\
                &=\sum_{j=1}^{n}j^2\\
                &=\frac{1}{6}n(n+1)(2n+1)
            \end{align*}
    \end{enumerate}

% subsection: 階比型 \, $a_{n+1}=f(n)a_n$

  \subsection{\texorpdfstring{階比型 \, $a_{n+1}=f(n)a_n$}{階比型}}\label{節：階比型}
  \begin{bookpoint}
    \textbf{倍率を掛ける}\par
    添字ごとの倍率を掛け合わせる方法と,その倍率に合う因子で割る方法を比べよう.零で割らない条件は共通である.
  \end{bookpoint}
    \bookterm{difference}{階差}型では変化を足した.ここでは,各段階の倍率を掛け合わせる.
    難しい点は,積を計算することと,倍率をそろえる因子を見つけることに分かれる.
    まず積の表示を学び,その後に\bookterm{normalize}{正規化}の因子を選ぶ理由を見る.

    具体的な説明の前に1つだけ記号を新たに導入しておく.
    適当な数が並べられた\bookterm{sequence}{数列}$x_n$が存在した時,
    \[
      \prod_{i=1}^n x_i
      =x_1\times x_2\times\cdots\times x_n
    \]
    となるような記号$\prod$（ラージパイ）を定義する.
    これを\bookterm{product}{総積}記号といい,積を表す.
    また,一般に以下が成り立つ.
    \begin{align*}
      &\prod_{i=1}^{n}i=n!\\
      &\prod_{i=1}^{n}k=k^n\\
      &\prod_{i=1}^{n}a_ib_i=\displaystyle\prod_{i=1}^na_i\displaystyle\prod_{i=1}^nb_i\\
      &\prod_{i=1}^{n}a_i=\displaystyle\prod_{i=1+k}^{n+k}a_{i-k}\\
      &\prod_{i=1}^{n}(i+k)=\dfrac{(n+k)!}{k!}
    \end{align*}
    ここでの$a_i$や$b_i$は\bookterm{recurrence}{漸化式}等によって定義されるような,法則性を持った\bookterm{sequence}{数列}である必要はなく,
    「1から順に番号が振られた,数の並び」ぐらいのイメージで問題ない.
    では再び, \,$a_n$の\bookterm{general}{一般項}を求めていく.
    \begin{align*}
      &a_{n+1}=f(n)a_n\\
      &a_{n+1}=f(n)\{f(n-1)a_{n-1}\}\\
      &a_{n+1}=f(n)\left[f(n-1)\{f(n-2)a_{n-2}\}\right]\\
      &\qquad \vdots\\
      &a_{n+1}=f(1)f(2) \cdots f(n)a_1\\
      &a_{n+1}=a_1 \prod_{k=1}^n f(k)\\
      &\therefore \quad a_n=a_1 \prod_{k=1}^{n-1} f(k)
    \end{align*}
    この積による表示は$f(k)=0$となる場合も成り立つ(空積は$1$とする).
    \noindent \bookblocklink{example-ratio-1}{problem}{answer}{【例題】}\par\nobreak
    \begin{enumerate}[label=(\arabic*),parsep=3pt,format=\bookitemlink{example-ratio-1}{problem}{answer}]
      \item $\displaystyle a_1=1,\, a_{n+1}=2na_n$
      \item $\displaystyle a_1=2,\, a_{n+1}=\frac{n+2}{n}a_n$
      \item $\displaystyle a_1=1,\, a_{n+1}=n(n+1)a_n$
    \end{enumerate}

    \noindent\bookblocklink{example-ratio-1}{hint}{problem}{【方針】}\par\nobreak
    \begin{enumerate}[label=(\arabic*),parsep=3pt,format=\bookitemlink{example-ratio-1}{hint}{problem}]
      \item 係数$2n$のうち$n$を階乗で吸収すれば,一定の倍率$2$だけが残ります.
      \item $n(n+1)$を一つのまとまりとして,隣の添字にずらした比を計算してみましょう.倍率$\frac{n+2}{n}$が現れます.
      \item ちょっと頭を捻ってみましょう.先程までは$n!$で割りましたが,今回はうまくいくでしょうか.\,$f(n)=n(n+1)$の形に合わせて,どのように階乗を組み合わせれば「きれい」になるか考えてみてください.
    \end{enumerate}

    \noindent\bookblocklink{example-ratio-1}{answer}{problem}{【解答】}\par\nobreak
    \begin{enumerate}[label=(\arabic*),parsep=3pt,format=\bookitemlink{example-ratio-1}{answer}{problem}]
      \item $\displaystyle a_n=(n-1)!\cdot 2^{n-1}$
      \item $\displaystyle a_n=n(n+1)$
      \item $\displaystyle a_n=n!(n-1)!$
    \end{enumerate}

    \noindent\bookblocklink{example-ratio-1}{solution}{problem}{【解法】}\par\nobreak
    \begin{enumerate}[label=(\arabic*),parsep=3pt,format=\bookitemlink{example-ratio-1}{solution}{problem}]
      \item $\displaystyle a_1=1,\, a_{n+1}=2na_n$\\
            両辺を$n!$で割って,
            \begin{align*}
              &a_{n+1}=2na_n\\
              &\frac{a_{n+1}}{n!}=2\frac{a_n}{(n-1)!}
            \end{align*}
            ここで
            \[
              b_n=\frac{a_n}{(n-1)!}
            \]
            と置くと,
            \begin{align*}
              &b_{n+1}=2b_n\\
              &b_n=b_1\cdot 2^{n-1}
            \end{align*}
            $b_n$を$a_n$に戻して
            \begin{align*}
              &\frac{a_n}{(n-1)!}=\frac{a_1}{(1-1)!}2^{n-1}\\
              &\therefore \quad a_n=a_1 \cdot (n-1)! \cdot 2^{n-1}
            \end{align*}

      \item $\displaystyle a_1=2,\, a_{n+1}=\frac{n+2}{n}a_n$
            \begin{align*}
              &a_{n+1}=\frac{n+2}{n}a_n\\
              &na_{n+1}=(n+2)a_n\\
            \end{align*}
            両辺を
            \[
              n(n+1)(n+2)
            \]
            で割って,
            \begin{align*}
              &\frac{a_{n+1}}{(n+1)(n+2)}=\frac{a_{n}}{n(n+1)}
            \end{align*}
            $a_n$について遡っていくと
            \begin{align*}
              &\frac{a_{n}}{n(n+1)}=\frac{a_{n-1}}{(n-1)(n)}=\\
              &\frac{a_{n-2}}{(n-2)(n-1)}=\cdots=\frac{a_{1}}{1(1+1)}\\
              &\therefore \quad \frac{a_{n}}{n(n+1)}=\frac{a_{1}}{1(1+1)}\\
              &a_n=n(n+1)
            \end{align*}
      \item $\displaystyle a_1=1,\, a_{n+1}=n(n+1)a_n$\\
              両辺を$n!(n+1)!$で割って,
            \begin{align*}
              &a_{n+1}=n(n+1)a_n\\
              &\frac{a_{n+1}}{n!(n+1)!}=\frac{n(n+1)a_n}{n!(n+1)!}\\
              &\frac{a_{n+1}}{n!(n+1)!}=\frac{a_n}{n!(n-1)!}\\
              &\therefore \quad \frac{a_n}{n!(n-1)!}=\frac{a_1}{1!(1-1)!}\\
              &a_n=n!(n-1)!
            \end{align*}
    \end{enumerate}

    因子を見つけにくければ,先に倍率の積を書き,隣り合う分子と分母が相殺するかを調べるとよい.
    \bookterm{normalize}{正規化}はその積の計算を,一歩ずつの更新として表したものである.

    なお,ここでは\bookterm{ratio}{階比}型\bookterm{recurrence}{漸化式}$a_{n+1}=f(n)a_n$について,\,$f(n)$が指数関数でない場合について扱ったが,
    より複雑な関数であったり指数関数を用いたりする場合は\bookterm{logarithm}{対数}を用いることがある.
    それら次節応用7パターンで解説する.


\bookstep{掛ける量を選び、倍率をそろえる}\label{解釈：階比と正規化}
    一方,以下では非零の因子を掛けて倍率をそろえる.これは\sectionref*{節：補助数列・正規化}の\bookterm{normalize}{正規化}に当たる.
    $f(k)\ne0$を仮定し,常に$g(n)\ne0$となるように選ぶ.
    探すのは,\,$\frac{g(n+1)}{g(n)}$が元の倍率$f(n)$の逆数になる因子である.
    与式の両辺に$g(n+1)$を掛けると,
    \[
      g(n+1)a_{n+1}=g(n+1)f(n)a_n
    \]
    ここで$g(n+1)f(n)=g(n)$を満たすような$g(n)$を選べば,
    \begin{align*}
      &g(n+1)a_{n+1}=g(n)a_n\\
      &g(n+1)a_{n+1}=g(n)a_n=g(n-1)a_{n-1}\\
      &=\cdots=g(1)a_1\\
      &\therefore \quad a_n=\frac{g(1)a_1}{g(n)}
    \end{align*}
    この解法では,\,$b_n=g(n)a_n$の更新が$b_{n+1}=b_n$となる.
    項ごとに変わる倍率を取り除いたため,初項の値だけを追えばよくなったのである.
    これを背景に考えると,例えば$f(n)=n$のときは$g(n)=\dfrac1{(n-1)!}$と取れば$g(n+1)f(n)=g(n)$が成り立ち,
    \begin{align*}
      &a_{n+1}=na_n\\
      &\frac{a_{n+1}}{n!}=\frac{a_n}{(n-1)!}\\
      &\frac{a_n}{(n-1)!}=\frac{a_1}{0!}\\
      &\therefore \quad a_n=(n-1)!\cdot a_1
    \end{align*}
    が言える.
    多くはこの場合が出題されるため,演習問題はこの手法を使って解くものを揃えた.


% section: おまけ：定数数列

\section{おまけ：定数数列}\label{節：おまけ：定数数列}
  \bookterm{sequence}{数列}の種類には定数\bookterm{sequence}{数列}というものがある.
  時折,\bookterm{recurrence}{漸化式}を解く中で定数\bookterm{sequence}{数列}が生まれることがある.
  先程の解説から引用しよう.
  \begin{align*}
    &\frac{a_{n}}{n(n+1)}=\frac{a_{n-1}}{(n-1)(n)}=\frac{a_{n-2}}{(n-2)(n-1)}\\
    &\qquad=\cdots=\frac{a_{1}}{1(1+1)}
  \end{align*}

  これはまさに定数\bookterm{sequence}{数列}の良い例である.
  少し式が複雑なので,
  \[
    b_n=\frac{a_{n}}{n(n+1)}
  \]
  とおくと式は
  \[
    b_n=b_{n-1}=b_{n-2}=\cdots=b_1
  \]
  となる.
  ずっと同じ値を取るのだ.
  当たり前過ぎて何を言っているんだ,と思うかもしれないが急に出てきてパニックにならないようにここで念のため触れておく.
  この「一定の値を取り続ける\bookterm{sequence}{数列}」という発想は,\sectionref{節：特性方程式型}で登場する\bookterm{characteristic}{特性方程式}の解$\alpha$(\bookterm{recurrence}{漸化式}$a_{n+1}=pa_n+q$自身を定数\bookterm{sequence}{数列}にしてしまう値)にもつながっている.

% section: 応用7パターン

\section{応用7パターン}\label{節：応用7パターン}
  ここからは、元の式に隠れた差や倍率を取り出す。
  定数を引くと、等比型が現れる。
  逆数を取ると、分数の式が一次の式になる。
  二つの項を組み合わせると、それぞれを一つずつ追えることもある。
  \sectionref{節：漸化式を解く前に}で見た「きれいな形」は、こうした変化を追える形である。
  既知の形へ直す操作を、一つずつ積み重ねていく。

% subsection: 特性方程式型 \, $a_{n+1}=pa_n+q$

\subsection{\texorpdfstring{特性方程式型 \, $a_{n+1}=pa_n+q$}{特性方程式型}}\label{節：特性方程式型}
  \begin{bookpoint}
    \textbf{何を引くか}\par
    更新の前後で変わらない値を基準にすると,そこからの差が等比型になる.倍率が1の場合は先に分ける.
  \end{bookpoint}
  \bookstep{定数項を取り除く}
  \[
    a_{n+1}=pa_n+q
  \]
  では,定数$q$を取り除けば等比型になる.
  \sectionref*{節：補助数列の設計}で述べた「同じ更新をする量を引く」方法を,定数に対して使う.
  ただし$p=1$なら等差型であり,\,$p=0$なら直接$n\ge2$で$a_n=q$と分かる.
  以下では$p\ne0,1$を考える.

  \bookterm{sequence}{数列}全体から同じ数$\alpha$を引き,
  \[
    a_{n+1}-\alpha=p(a_n-\alpha)
  \]
  となることを目指す.右辺を展開すると定数項は$-p\alpha$なので,必要な条件は
  \[
    q-\alpha=-p\alpha
    \quad\Longleftrightarrow\quad
    \alpha=p\alpha+q
  \]
  である.この方程式を本書では\bookterm{characteristic}{特性方程式}と呼ぶ.
  解は$\alpha=\frac{q}{1-p}$であり,これが引くべき値である.
  また,\,$\alpha=p\alpha+q$は,その値から出発すると更新しても変わらないことを表す.
  このような値を\bookterm{fixed}{不動点}という.

  \bookstep{差を求めて、元の項へ戻す}
  元の更新式から\bookterm{fixed}{不動点}の式を引くと,
  \[
    \begin{array}{rrl}
       &a_{n+1}={}&pa_n+q\\
      -\,)&\alpha={}&p\alpha+q\\ \hline
       &a_{n+1}-\alpha={}&p(a_n-\alpha)
    \end{array}
  \]
  となる.$b_n=a_n-\alpha$と置けば
  \[
    b_{n+1}=pb_n,\qquad b_1=a_1-\alpha.
  \]
  等比型の\bookterm{general}{一般項}から
  \[
    b_n=(a_1-\alpha)p^{n-1}
  \]
  となる.$a_n=b_n+\alpha$へ戻すと,
  \[
    a_n=(a_1-\alpha)p^{n-1}+\alpha
  \]
  を得る.大切なのは最終式だけでなく,同じ定数項を持つ二つの式を引いた理由である.
  初項が\bookterm{fixed}{不動点}なら$b_1=0$となり,\,$a_n=\alpha$という定数解もこの式に含まれる.

  \noindent\bookblocklink{check-fixed}{problem}{answer}{【確認問題】}\par\nobreak
  \begin{enumerate}[format=\bookitemlink{check-fixed}{problem}{answer}]
    \item $a_{n+1}=3a_n-4$で,何を引くと等比型になるか.
          $a_1=5$と$a_1=2$の場合の\bookterm{general}{一般項}をそれぞれ求めよ.
  \end{enumerate}
  \noindent\bookblocklink{check-fixed}{answer}{problem}{【確認の解答】}\par\nobreak
  \begin{enumerate}[format=\bookitemlink{check-fixed}{answer}{problem}]
    \item $\alpha=3\alpha-4$より$\alpha=2$を引く.
          $a_1=5$なら$a_n=2+3^n$であり,\,$a_1=2$なら$a_n=2$である.
  \end{enumerate}

\noindent \bookblocklink{example-characteristic-1}{problem}{answer}{【例題】}\par\nobreak
\begin{enumerate}[label=(\arabic*),parsep=3pt,format=\bookitemlink{example-characteristic-1}{problem}{answer}]
  \item $\displaystyle a_1=1,\, a_{n+1}=2a_n+1$
  \item $\displaystyle a_1=1,\, a_{n+1}=\frac{1}{2}a_n+2$
  \item $\displaystyle a_1=1,\, a_{n+1}=a_n+2$
\end{enumerate}

\noindent\bookblocklink{example-characteristic-1}{hint}{problem}{【方針】}\par\nobreak
\begin{enumerate}[label=(\arabic*),parsep=3pt,format=\bookitemlink{example-characteristic-1}{hint}{problem}]
  \item 定数項を消す値$\alpha$を求め、そこからの差を等比型で追います。
  \item 同じ方法です。倍率が分数でも、同じ更新をする定数を引けます。
  \item 倍率が$1$なので、等差型です。定数項を消す値は存在せず、差を足す方法に戻ります。
\end{enumerate}

\noindent\bookblocklink{example-characteristic-1}{answer}{problem}{【解答】}\par\nobreak
\begin{enumerate}[label=(\arabic*),parsep=3pt,format=\bookitemlink{example-characteristic-1}{answer}{problem}]
  \item $\displaystyle a_n=2^n-1$
  \item $\displaystyle a_n=4-3\cdot\left(\frac{1}{2}\right)^{n-1}$
  \item $\displaystyle a_n=1+2(n-1)$
\end{enumerate}

\noindent\bookblocklink{example-characteristic-1}{solution}{problem}{【解法】}\par\nobreak
\begin{enumerate}[label=(\arabic*),parsep=3pt,format=\bookitemlink{example-characteristic-1}{solution}{problem}]
  \item $\displaystyle a_1=1,\,a_{n+1}=2a_n+1$\\
        \bookterm{characteristic}{特性方程式}を考えると,
        \[
          \alpha=2\alpha+1
          \quad\Longleftrightarrow\quad
          \alpha=-1
        \]
        である.

        $b_n=a_n-\alpha=a_n+1$とおくと,
        \begin{align*}
          b_{n+1}
            &=a_{n+1}+1\\
            &=2a_n+2\\
            &=2(a_n+1)\\
            &=2b_n
        \end{align*}
        また,\,$b_1=a_1+1=2$であるから,
        \[
          b_n=2^n
        \]
        よって,
        \[
          a_n=b_n-1=2^n-1
        \]

  \item $\displaystyle a_1=1,\,a_{n+1}=\frac12a_n+2$\\
          \bookterm{characteristic}{特性方程式}を考えると,
          \[
            \alpha=\frac12\alpha+2
            \quad\Longleftrightarrow\quad
            \alpha=4
          \]
          である.

          $b_n=a_n-\alpha=a_n-4$とおくと,
          \begin{align*}
            b_{n+1}
              &=a_{n+1}-4\\
              &=\frac12a_n+2-4\\
              &=\frac12(a_n-4)\\
              &=\frac12b_n
          \end{align*}
          また,\,$b_1=a_1-4=-3$であるから,
          \[
            b_n=-3\left(\frac12\right)^{n-1}
          \]
          よって,
          \[
            a_n=b_n+4
            =4-3\left(\frac12\right)^{n-1}
          \]
   \item $\displaystyle a_1=1,\, a_{n+1}=a_n+2$\\
         $\alpha$は存在しないので,\bookterm{characteristic}{特性方程式}型に帰着させることはできない.
         これが意味することは,\bookterm{fixed}{不動点}を持たないということである.
\end{enumerate}


\bookstep{引いた値を、グラフから見る}\label{解釈：不動点と差分}
  例題で引いた値$\alpha$は、グラフの上にも現れる。
  同じ引き算を、座標の移動として見直す。
  \bookterm{recurrence}{漸化式}を
  \[
    a_{n+1}=f(a_n)
  \]
  と考え,
  \[
    y=f(x)=px+q
  \]
  について,\,$y=x$との交点,すなわち
  \[
    f(\alpha)=\alpha
  \]
  を満たすような点を考える.図示すると

  \begin{tikzpicture}[scale=0.8]
    \draw[->,>=stealth,semithick] (-1.5,0)--(5.5,0)node[above]{$x$};
    \draw[->,>=stealth,semithick] (0,-1)--(0,5)node[right]{$y$};
    \draw (0,0)node[below right]{O};

    \draw[very thick,domain=-1:4]
      plot(\x,\x)node[right]{$y=x$};

    \draw[very thick,domain=-1.5:4]
      plot(\x,\x/2+1) node[right] {$y=px+q$};

    % q
    \node[left=3pt] at (0,1) {$q$};

    % 交点から各軸へ点線
    \draw[dashed] (2,2) -- (2,0);
    \draw[dashed] (2,2) -- (0,2);

    % α
    \node[below=3pt] at (2,0) {$\alpha$};
    \node[left=3pt] at (0,2) {$\alpha$};
  \end{tikzpicture}
  \begin{center}
    \BookFigureCaption{直線の交点として捉える\bookterm{fixed}{不動点}}
  \end{center}

  というような状況である.
  この交点$(\alpha,\alpha)$を原点と思ってグラフをみると,この座標平面は$y=px$と$y=x$のグラフが原点から始まって進んでいくようなものに見える.
  言い換えれば等比型\bookterm{recurrence}{漸化式}に帰着させたということである.
  この交点の$x$座標$\alpha$が,更新$f$の\bookterm{fixed}{不動点}である.

  この\bookterm{fixed}{不動点}を知ることによって,複雑な値の変化をより簡素な値の変化に置き換えることができる.
  また,この話から,\,$|p|>1$のときは一般に項数を重ねるにつれて$p$が生み出す等比的な増大が支配的になる一方,\,$|p|<1$のときは$a_n$は$n\to\infty$で$\alpha=\dfrac{q}{1-p}$に収束し,その\bookterm{limit}{極限}値は$q$(と$p$)によって決まることが分かる.例題(2)(\,$p=\frac12$)が$4$に収束するのはこのためである.

  また,\bookterm{characteristic}{特性方程式}型\bookterm{recurrence}{漸化式}には別の解法が存在する.
  以下に示す.
  \begin{align*}
    &a_{n+1}=pa_n+q \tag*{\(\cdots\) (1)}\\
    &a_{n+2}=pa_{n+1}+q \tag*{\(\cdots\) (2)}\\
  \end{align*}
  $(2)-(1)$をし
  \[
    a_{n+2}-a_{n+1}=p(a_{n+1}-a_n)
  \]
  ここで,\,$b_n=a_{n+1}-a_n$と置くと,
  \begin{align*}
    &b_{n+1}=pb_n
  \end{align*}
  このようにすることで,\bookterm{recurrence}{漸化式}は\bookterm{geometric}{等比数列}に帰着する.
  得られた差を初項から足し戻すと、元の\bookterm{sequence}{数列}が求まる。

  こちらの図形的解釈は先程の交点を考えるものとは違う.傾き$p$の平行な2直線$y=px$と$y=px+q$の間には,常に高さ$q$の差が生じている.\,(2)から(1)を引く操作は,この共通の$q$だけを打ち消す操作に対応している.つまり隣接する2つの式を「ずらして引く」ことで,\,$a_n$に無関係な定数項$q$だけを消してしまうのである.これは付録\sectionref{節：差分から始まる離散的な微積分}で扱う\bookterm{difference-operator}{差分}の考え方と同じ発想である.
  これらの前提を元に図示すると

  \begin{tikzpicture}[scale=0.8]
    \draw[->,>=stealth,semithick] (-1.5,0)--(5.5,0)node[above]{$x$};
    \draw[->,>=stealth,semithick] (0,-1)--(0,5)node[right]{$y$};
    \draw (0,0)node[below right]{O};

    \draw[very thick,domain=-1.5:4]
      plot(\x,\x/2)node[right]{$y=px$};

    \draw[very thick,domain=-1.5:4]
      plot(\x,\x/2+1) node[right] {$y=px+q$};

    % q
    \node[left=3pt] at (0,1) {$q$};

  \end{tikzpicture}
  \begin{center}
    \BookFigureCaption{定数項の違いを表す平行な直線}
  \end{center}

  常に関数の差が$q$と一定になっていることが図からわかる.

  なお,先ほどの等比的な増大について,初項を\bookterm{fixed}{不動点}に選んだ場合は例外になる.
  \bookterm{fixed}{不動点}から出発すれば値は変わらず,\,$a_n=\alpha$という定\bookterm{sequence}{数列}になるからである.
  式で確かめると,\,$p\ne1$のときは
  \[
    a_n-\alpha=(a_1-\alpha)p^{n-1}
  \]
  である.したがって,\,$|p|>1$でも$a_1=\alpha$ならずれは常に$0$であり,\,$a_1\ne\alpha$の場合に\bookterm{fixed}{不動点}からのずれの絶対値が等比的に増大する.

  今後,\bookterm{characteristic}{特性方程式}を用いた解法を多数扱う.
  \bookterm{characteristic}{特性方程式}は\bookterm{recurrence}{漸化式}を解く中で重要な概念なのである.
  その背景には\bookterm{matrix}{行列}や\bookterm{linear}{線形}代数が深く絡んでいる.

% subsection: 階差等比複合型 \, $\displaystyle a_{n+1}=pa_n+f(n)$

\subsection{\texorpdfstring{階差等比複合型 \, $\displaystyle a_{n+1}=pa_n+f(n)$}{階差等比複合型}}\label{節：階差等比複合型}
  \begin{bookpoint}
    \textbf{同じ更新をする量を引く}\par
    付加項を消すことを目標に,引く式の形と係数を決める.候補を見つける計算と,候補が元の式を満たす確認を区別しよう.
  \end{bookpoint}
  \subsubsection{\texorpdfstring{多項式型 \, $\displaystyle a_{n+1}=pa_n+f(n)$}{多項式型}}
    \bookstep{式の形から、追う量を選ぶ}
    \sectionref*{節：特性方程式型}では,同じ更新をする定数を引いた.
    今回は付加項$f(n)$が$n$によって変わるため,引く量も$n$によって変えてみる.
    \sectionref*{節：補助数列の設計}で述べた方針に従い,同じ更新をする既知の量を探そう.
    本文では$f(n)$が多項式の場合を考える.$p=1$なら\bookterm{difference}{階差}型として和を求められるので,まず$p\ne1$とする.

    \bookstep{引く量を選ぶ}
    $b_n=a_n-g(n)$を等比型にするための条件は
    \[
      g(n+1)=pg(n)+f(n)
    \]
    である.これは,\,$g(n)$自身にも同じ更新をさせるという条件である.
    条件が満たされれば,
    \[
      \begin{array}{rrl}
        &a_{n+1}={}&pa_n+f(n)\\
        -\,)&g(n+1)={}&pg(n)+f(n)\\ \hline
        &a_{n+1}-g(n+1)={}&p\{a_n-g(n)\}
      \end{array}
    \]
    となり,付加項が消える.
    \sectionref*{節：非同次と解の組合せ}で見た「一つの既知の解と\bookterm{homogeneous}{同次}の解に分ける」考え方が,
    ここでは$a_n=g(n)+b_n$として現れている.
    $g(n)$を見つける計算と,残る\bookterm{homogeneous}{同次}の解$b_n$を\bookterm{initial}{初期条件}から選ぶ計算を分けているのである.

    $f(n)$が$m$次式なら,\,$p\ne1$のときは$m$次の多項式
    \[
      g(n)=x_0+x_1n+\cdots+x_mn^m
    \]
    を候補にする.$g(n+1)-pg(n)$の最高次の係数は$(1-p)x_m$である.
    これを$f(n)$の最高次の係数に合わせ,次に一段低い係数を合わせる.
    $1-p\ne0$なので,この順に$x_m,x_{m-1},\ldots,x_0$を決められる.
    未知の係数を比較して求めるこの方法を\bookterm{undetermined}{未定係数法}という.

    例えば$a_{n+1}=2a_n+n-1$では,\,$g(n)=un+v$と置くと
    \[
      u=2u+1,\qquad u+v=2v-1
    \]
    となるので,\,$g(n)=-n$である.
    従って$a_n+n$が2倍される量となる.
    この例の計算は\bookproblemref{practice-basic}{4}{標準問4}で確認する.

    \bookstep{変化を求めて、元の項へ戻す}
    $b_{n+1}=pb_n$を解き,\,$a_n=b_n+g(n)$へ戻す.
    $p\ne0,1$なら
    \[
      a_n=\{a_1-g(1)\}p^{n-1}+g(n)
    \]
    となる.$p=0$なら直接$n\ge2$で$a_n=f(n-1)$と求める.
    一方,\,$p=1$で$f(n)$が$m$次なら,付加項を足し戻す多項式は通常$m+1$次になる.
    例えば$a_{n+1}=a_n+n$では,\,$g(n)=\frac{n(n-1)}2$が同じ更新を満たす.

  \subsubsection{\texorpdfstring{指数型  \, $\displaystyle a_{n+1}=pa_n+q^n$}{指数型}}
    前節では$f(n)$が多項式である場合を扱った.
    ここでは$f(n)$が特に指数関数である場合を扱う.
    指数関数である場合にはもう一つ別の解法が存在する.
    \begin{align*}
      a_{n+1}=pa_n+q^n
    \end{align*}
    以下では$q\ne0$とする($q=0$なら$n\ge1$で$q^n=0$なので等比型になる).
    両辺を$q^n$で割って,
    \begin{align*}
      &\frac{a_{n+1}}{q^n}=\frac{p}{q}\cdot\frac{a_n}{q^{n-1}}+1\\
    \end{align*}
    ここで,\,$b_n=\frac{a_n}{q^{n-1}}$と置くと
    \begin{align*}
      &b_{n+1}=\frac{p}{q}b_n+1
    \end{align*}
    あとは\bookterm{characteristic}{特性方程式}型の\bookterm{recurrence}{漸化式}として解くことになる.
    ただし$p=q$のときは$b_{n+1}=b_n+1$(\bookterm{arithmetic}{等差数列})になり,\bookterm{fixed}{不動点}$\alpha$が存在しない(\sectionref{節：特性方程式型}の例題(3)と同じ状況である).
    この場合は$b_n=b_1+(n-1)$として直接求まる.
    例えば$a_1=1,\,a_{n+1}=2a_n+2^n$(\,$p=q=2$)ならば$b_n=\dfrac{a_n}{2^{n-1}}$が\bookterm{arithmetic}{等差数列}となり,\,$a_n=n\cdot2^{n-1}$が得られる.

% subsection: 対数型 \, $a_{n+1}=pa_n^q$

\subsection{対数型 \, $a_{n+1}=pa_n^q$}
  今までは$a_n$に対する操作は定数倍したり関数が掛けられたりするだけだった.
  しかし,\,$a_n$に指数がついたらどうだろうか.
  この場合はかなり難しい発想の変形が必要である.
  それは対数を取るという操作である.
  実際にやってみよう.
  \[
    a_{n+1}=pa_n^q
  \]
  ただし,対数を取るために$a_n>0$とし,底に選ぶ$p$も$p>0,\,p\ne1$を満たすとする.
  この式について,両辺の対数を取ると
  \begin{align*}
    \log_p{a_{n+1}}&=\log_p{pa_n^q}\\
    &=\log_p{p}+\log_p{a_n^q}\\
    &=q\log_p{a_n}+1
  \end{align*}
  $b_n=\log_p{a_n}$とおくと,
  \begin{align*}
    b_{n+1}=qb_n+1
  \end{align*}
  あとはただの特性方程式型の問題である.
  底をいくらにするかは自由だが$p$に関連する数だと扱いやすい.
  たとえば$p=k^m$の時は$k$でとっておくことが多い.
  作問者側は解答がシンプルな形になるようにするため,より小さい底を取ることで答えが求めやすくなる.

  また,対数の強みは積を和に,商を差にできるところにある.
  漸化式が全て積や商のみで表される場合にその強みが現れる.
  これはここで解説するよりも演習問題を通して経験するほうがわかりやすいと判断したため,この内容は(2)で取り扱っている.

  \noindent\textbf{【例題】}
  \begin{enumerate}[label=(\arabic*),parsep=3pt]
    \item $\displaystyle a_1=2,\quad a_{n+1}=2a_n^2$
    \item $\displaystyle a_1=\frac{1}{2},\quad a_{n+1}=\frac{1}{a_n}$
    \item $\displaystyle a_1=1,\quad a_{n+1}=2^na_n$
    \item $\displaystyle a_1=5,\quad a_{n+1}=a_n^2-2a_n+2$
    \item $\displaystyle a_1=2,\quad a_{n+1}=\frac{a_n^2}{2^n}$
    \item $\displaystyle a_1=1,\quad a_{n+1}=(n+1)a_n$
  \end{enumerate}

  \noindent\textbf{【方針】}
  \begin{enumerate}[label=(\arabic*),parsep=3pt]
    \item 両辺の底を$2$とする対数を取り,\,$b_n=\log_2a_n$とおいて一次の漸化式に変形します.その結果,漸化式は特性方程式型になります.
    \item 一見すると分数型(後述)のように見えますが,対数型として解くことができます.
    \item $f(n)=2^n$が指数関数なので,対数を取ると和の形に帰着できます.
    \item うまく変形すると対数型になります.
    \item 両辺の対数を取ると,二次の漸化式が一次の漸化式に変わり,さらに等差数列の和を使って解くことができます.
    \item 階比型に見えますが,対数型としても解けます.
  \end{enumerate}

  \noindent\textbf{【解答】}
  \begin{enumerate}[label=(\arabic*),parsep=3pt]
    \item $\displaystyle a_n=2^{2^{n}-1}$
    \item $a_n=2^{(-1)^{n}}$
    \item $a_n=2^{\frac{n(n-1)}{2}}$
    \item $a_n=2^{2^n}+1$
    \item $\displaystyle a_n=2^{n-1}-n+1$
    \item $a_n=n!$
  \end{enumerate}

  \noindent\textbf{【解法】}
  \begin{enumerate}[label=(\arabic*),parsep=3pt]
    \item $\displaystyle a_1=2,\quad a_{n+1}=2a_n^2$\\
          両辺の底を$2$とする対数を取ると,
          \begin{align*}
            \log_2a_{n+1}
              &=\log_2{2a_n^2}\\
              &=\log_2{2}+\log_2{a_n^2}\\
              &=2\log_2{a_n}+1
          \end{align*}
          ここで
          \[
            b_n=\log_2a_n
          \]
          とおくと,
          \begin{align*}
            &b_{n+1}=2b_n+1\\
            &b_{n+1}+1=2(b_n+1)
          \end{align*}
          ここで
          \[
            c_n=b_n+1
          \]
          とおくと,
          \[
            c_{n+1}=2c_n
          \]
          $c_1$は,
          \[
            c_1=b_1+1=\log_2{a_1}+1=\log_2{2}+1=1+1=2
          \]
          従って,
          \begin{align*}
            &c_n=2\cdot2^{n-1}=2^n\\
            &b_n+1=2^n\\
            &b_n=2^n-1\\
            &a_n=2^{2^{n}-1}
          \end{align*}
    \item $\displaystyle a_1=\frac12,\quad a_{n+1}=\frac{1}{a_n}$\\
          与式の両辺の底を2とする対数を取ると,
          \begin{align*}
            &\log_2{a_{n+1}}=\log_2{\frac{1}{a_n}}\\
            &\log_2{a_{n+1}}=\log_2{1}-\log_2{a_n}\\
            &\log_2{a_{n+1}}=-\log_2{a_n}\\
            &\log_2{a_n}=\log_2{a_1}(-1)^{n-1}\\
            &\log_2{a_n}=\log_2{\frac{1}{2}}(-1)^{n-1}\\
            &\log_2{a_n}=(-1)^{n}\\
            &\therefore \quad a_n=2^{(-1)^{n}}
          \end{align*}
    \item $\displaystyle a_1=1,\quad a_{n+1}=2^na_n$\\
          両辺の自然対数を取ると
          \[
            \log a_{n+1}=n\log2+\log a_n
          \]
          したがって
          \begin{align*}
            \log a_n&=\log a_1+\sum_{k=1}^{n-1} k\log2\\
            &=\frac{(n-1)n}{2}\log2\\
            &\therefore \quad a_n=2^{\frac{n(n-1)}{2}}
          \end{align*}
    \item $\displaystyle a_1=5,\quad a_{n+1}=a_n^2-2a_n+2$\\
          与式を変形して
          \[
            a_{n+1}-1=(a_n-1)^2
          \]
          両辺の底を2とする対数を取ると,
          \begin{align*}
            &\log_2{(a_{n+1}-1)}=\log_2{(a_n-1)^2}\\
            &\log_2{(a_{n+1}-1)}=2\log_2{(a_n-1)}\\
            &\log_2{(a_n-1)}=2^{n-1}\cdot\log_2{(a_1-1)}\\
            &\log_2{(a_n-1)}=2^{n-1}\cdot\log_2{4}\\
            &\log_2{(a_n-1)}=2^{n}\\
            &a_n-1=2^{2^{n}}\\
            &\therefore \quad a_n=2^{2^{n}}+1
          \end{align*}
    \item $\displaystyle a_1=2,\quad a_{n+1}=\frac{a_n^2}{2^n}$\\
          両辺を底を2とする自然体数を取って
          \[
            equation
          \]
    \item $\displaystyle a_1=1,\quad a_{n+1}=(n+1)a_n$\\
          両辺の自然対数を取って,
          \[
            \log a_{n+1}=\log{n+1}+\log a_n
          \]
          したがって
          \begin{align*}
            \log a_n&=\log a_1 + \log{2}+\cdots+ \log{n}\\
            &=\log 1 + \log{2}+\cdots +\log{n}\\
            &=\log n!\\
            &\therefore \quad a_n=n!
          \end{align*}
  \end{enumerate}
  \noindent\textbf{【補足】}

  (3)では底を自然対数としたが,底は任意の値で構わない.


% parent file: 分数型
\subsection{分数型}\label{節：分数型}
\begin{bookpoint}
  \textbf{分母を簡単にする量を選ぶ}\par
  逆数,または\bookterm{fixed}{不動点}からの二つの差の比を追う.
  分母が零になる初期値と定数解の分岐を確かめてから,置き換えを使う.
\end{bookpoint}
\input{第3章 漸化式の解説/05_応用7パターン/04_分数型/01_逆数型/逆数型.tex}
\input{第3章 漸化式の解説/05_応用7パターン/04_分数型/02_1次分数型/1次分数型.tex}
\input{第3章 漸化式の解説/05_応用7パターン/04_分数型/03_極形式による図形的解釈/極形式による図形的解釈.tex}

% subsection: 隣接3項間漸化式型 \, $\displaystyle a_{n+2}+pa_{n+1}+qa_n=0$

\subsection{\texorpdfstring{隣接3項間漸化式型 \, $\displaystyle a_{n+2}+pa_{n+1}+qa_n=0$}{隣接3項間漸化式型}}\label{節：隣接3項間漸化式型}
  \BookFigThreeTerm

  \subsubsection{\texorpdfstring{二次特性方程式型 \, $\displaystyle a_{n+2}+pa_{n+1}+qa_n=0$}{二次特性方程式型}}\label{節：二次特性方程式型}
    節は変われど,特性方程式を使って考える.
    結局,大方の場合,等比数列に変形できればいいのだから,これを等比数列に変形できないかといくつかの試行錯誤を重ねるのが王道の解き方である.
    結果から言ってしまえば,特性方程式型のところで少し触れた式を生み出すことで等比数列を作ることができる.
    \[
      a_{n+2}+pa_{n+1}+qa_n=0
    \]
    について,
    \[
      a_{n+2}-\alpha a_{n+1}=\beta(a_{n+1}-\alpha a_{n})
    \]
    となるような$\alpha,\beta$の組を探せば良い.
    変形すると
    \begin{align*}
      &a_{n+2}-\alpha a_{n+1}=\beta(a_{n+1}-\alpha a_{n})\\
      &\therefore \quad a_{n+2}-(\alpha+\beta)a_{n+1}+\alpha\beta a_n=0
    \end{align*}
    である.これが元の漸化式
    \[
      a_{n+2}+pa_{n+1}+qa_n=0
    \]
    と一致するには,
    \[
      \alpha+\beta=-p,\quad \alpha\beta=q
    \]
    であればよい.これは解と係数の関係そのものであるから,\,$\alpha,\beta$は
    \[
      x^2+px+q=0
    \]
    の2解として得られる.

    \,$x^2+px+q=0$の2解を$\alpha,\beta$とすると,元の漸化式は
    \[
      \begin{cases}
        a_{n+1}-\alpha a_n=\beta(a_n-\alpha a_{n-1})\\
        a_{n+1}-\beta a_n=\alpha(a_n-\beta a_{n-1})
      \end{cases}
    \]
    という2つの等比数列を生む(添字は1つずれているが本質的には先程と同じ議論である).

    $\alpha\ne\beta$の時,
    \[
      b_n=a_{n+1}-\alpha a_n
    \]
    は公比$\beta$の等比数列であるから,
    \[
      b_n=b_1\beta^{n-1}
    \]
    である.ここで,
    \[
      b_1=a_2-\alpha a_1
    \]
    なので,
    \[
      a_{n+1}-\alpha a_n
      =(a_2-\alpha a_1)\beta^{n-1}
    \]
    を得る.

    同様に,
    \[
      c_n=a_{n+1}-\beta a_n
    \]
    は公比$\alpha$の等比数列であるから,
    \[
      c_n=c_1\alpha^{n-1}
    \]
    である.また,
    \[
      c_1=a_2-\beta a_1
    \]
    なので,
    \[
      a_{n+1}-\beta a_n
      =(a_2-\beta a_1)\alpha^{n-1}
    \]
    となる.

    ここで,この2式から$a_{n+1}$を消去するために,上の式から下の式を引くと,
    \begin{align*}
      &(a_{n+1}-\alpha a_n)-(a_{n+1}-\beta a_n)\\
      &\qquad=(a_2-\alpha a_1)\beta^{n-1}
      -(a_2-\beta a_1)\alpha^{n-1}
    \end{align*}
    である.左辺を整理すると,
    \[
      (\beta-\alpha)a_n
      =(a_2-\alpha a_1)\beta^{n-1}
      -(a_2-\beta a_1)\alpha^{n-1}.
    \]
    $\alpha\ne\beta$より$\beta-\alpha\ne0$であるから,
    \[
      a_n
      =\frac{a_2-\alpha a_1}{\beta-\alpha}\beta^{n-1}
      -\frac{a_2-\beta a_1}{\beta-\alpha}\alpha^{n-1}
    \]
    となる.

    ここで,\,$a_1,a_2,\alpha,\beta$は$n$によらない定数であるから,係数をまとめれば,
    \[
      a_n=A\alpha^{n-1}+B\beta^{n-1}
    \]
    と書くことができる.
    これは\sectionref*{節：解を組み合わせる見方}で学んだ,二つの解の線形結合である.
    非零の異なる根なら,$\alpha^{n-1}$と$\beta^{n-1}$の初めの二項は$(1,\alpha)$と$(1,\beta)$であり,
    一次独立なので基底になる.係数$A,B$を初期条件で選ぶことは,2次元の解全体から一つを選ぶことに当たる.
    \,$\alpha=\beta$(重解)の場合は上の議論がそのままは使えないため,別に扱う.
    同じ冪を二つ並べても一次独立にはならないので,別の形の解が必要になる.
    このときは$b_n=a_{n+1}-\alpha a_n$のみを考え,\,$b_{n+1}=\alpha b_n$という等比数列に帰着させてから,さらに階差型または等比複合型として解く(具体的な手順は例題(2)を参照).

    \noindent \bookblocklink{example-three-term-1}{problem}{answer}{【例題】}
    \begin{enumerate}[label=(\arabic*),parsep=3pt,format=\bookitemlink{example-three-term-1}{problem}{answer}]
      \item $\displaystyle a_1=1,\,a_2=1,\, a_{n+2}=5a_{n+1}-6a_n$
      \item $\displaystyle a_1=1,\,a_2=4,\, a_{n+2}=4a_{n+1}-4a_n$
      \item $\displaystyle a_1=2,\,a_2=1,\, a_{n+2}=a_{n+1}+6a_n$
    \end{enumerate}

    \noindent\bookblocklink{example-three-term-1}{hint}{problem}{【方針】}
    \begin{enumerate}[label=(\arabic*),parsep=3pt,format=\bookitemlink{example-three-term-1}{hint}{problem}]
      \item 特性方程式$x^2-5x+6=0$の解は異なる2つの実数解です.素直に2本の等比数列を作りましょう.
      \item 特性方程式$x^2-4x+4=0$を解くと重解になります.等比数列は1本しか作れないことに注意しましょう.
      \item 特性方程式の解の1つが負になるパターンです.符号に注意して計算しましょう.
    \end{enumerate}

    \noindent\bookblocklink{example-three-term-1}{answer}{problem}{【解答】}
    \begin{enumerate}[label=(\arabic*),parsep=3pt,format=\bookitemlink{example-three-term-1}{answer}{problem}]
      \item $\displaystyle a_n=2^n-3^{n-1}$
      \item $\displaystyle a_n=n\cdot2^{n-1}$
      \item $\displaystyle a_n=(-2)^{n-1}+3^{n-1}$
    \end{enumerate}

    \noindent\bookblocklink{example-three-term-1}{solution}{problem}{【解法】}
    \begin{enumerate}[label=(\arabic*),parsep=3pt,format=\bookitemlink{example-three-term-1}{solution}{problem}]
      \item $\displaystyle a_1=1,\,a_2=1,\, a_{n+2}=5a_{n+1}-6a_n$\\
            特性方程式
            \[
              x^2-5x+6=0
            \]
            を解くと,
            \[
              (x-2)(x-3)=0
            \]
            より$\alpha=2,\beta=3$である.
            \begin{align*}
              &a_{n+2}-2a_{n+1}=3(a_{n+1}-2a_n)\\
              &a_{n+2}-3a_{n+1}=2(a_{n+1}-3a_n)
            \end{align*}
            なので
            \[
              b_n=a_{n+1}-2a_n,\,c_n=a_{n+1}-3a_n
            \]
            と置くと,
            \begin{align*}
              &b_n=b_1\cdot3^{n-1},\quad b_1=a_2-2a_1=-1\\
              &c_n=c_1\cdot2^{n-1},\quad c_1=a_2-3a_1=-2\\
              &\therefore \quad b_n=-3^{n-1},\quad c_n=-2^n
            \end{align*}
            したがって
            \[
              b_n-c_n=(a_{n+1}-2a_n)-(a_{n+1}-3a_n)=a_n
            \]
            であるから,
            \[
              a_n=-3^{n-1}-(-2^n)=2^n-3^{n-1}
            \]
      \item $\displaystyle a_1=1,\,a_2=4,\, a_{n+2}=4a_{n+1}-4a_n$\\
            特性方程式
            \[
              x^2-4x+4=0
            \]
            を解くと,
            \[
              (x-2)^2=0
            \]
            より$\alpha=2$である.
            \[
              b_n=a_{n+1}-2a_n
            \]
            と置くと,
            \begin{align*}
              b_{n+1}&=a_{n+2}-2a_{n+1}\\
              &=(4a_{n+1}-4a_n)-2a_{n+1}\\
              &=2(a_{n+1}-2a_n)\\
              &=2b_n
            \end{align*}
            また,
            \begin{align*}
              b_1=a_2-2a_1=2
            \end{align*}
            から,
            \[
              b_n=2\cdot2^{n-1}=2^n
            \]
            したがって
            \[
              a_{n+1}-2a_n=2^n
            \]
            である.両辺を$2^{n+1}$で割ると,
            \[
              \frac{a_{n+1}}{2^{n+1}}-\frac{a_n}{2^n}=\frac12
            \]
            ここで
            \[
              c_n=\dfrac{a_n}{2^n}
            \]
            と置くと,
            \begin{align*}
              &c_{n+1}=c_n+\frac12,\quad c_1=\frac{a_1}{2}=\frac12\\
              &\therefore \quad c_n=\frac12+(n-1)\cdot\frac12=\frac n2\\
              &\therefore \quad a_n=2^n\cdot c_n=n\cdot2^{n-1}
            \end{align*}
      \item $\displaystyle a_1=2,\,a_2=1,\, a_{n+2}=a_{n+1}+6a_n$\\
            特性方程式
            \[
              x^2-x-6=0
            \]
            を解くと,
            \[
              (x-3)(x+2)=0
            \]
            より$\alpha=3,\beta=-2$である.
            \begin{align*}
              &a_{n+2}-3a_{n+1}=-2(a_{n+1}-3a_n)\\
              &a_{n+2}+2a_{n+1}=3(a_{n+1}+2a_n)
            \end{align*}
            なので,
            \[
              b_n=a_{n+1}-3a_n,\ c_n=a_{n+1}+2a_n
            \]
            と置くと,
\BookMathLayout{\begin{align*}
              &b_n=b_1\cdot(-2)^{n-1},\quad b_1=a_2-3a_1=-5\\
              &c_n=c_1\cdot3^{n-1},\quad c_1=a_2+2a_1=5\\
              &\therefore \quad b_n=-5(-2)^{n-1},\quad c_n=5\cdot3^{n-1}
            \end{align*}}{
\[
\begin{aligned}
b_n&=b_1\cdot(-2)^{n-1},\\ b_1&=a_2-3a_1=-5\\
c_n&=c_1\cdot3^{n-1},\\ c_1&=a_2+2a_1=5\\
\therefore\quad b_n&=-5(-2)^{n-1},\\ c_n&=5\cdot3^{n-1}
\end{aligned}
\]
}
            \[
              c_n-b_n=(a_{n+1}+2a_n)-(a_{n+1}-3a_n)=5a_n
            \]
            であるから,
            \begin{align*}
              5a_n&=5\cdot3^{n-1}-\{-5(-2)^{n-1}\}\\
              &=5\cdot3^{n-1}+5(-2)^{n-1}\\
              \therefore\quad a_n&=(-2)^{n-1}+3^{n-1}
            \end{align*}
    \end{enumerate}

  \subsubsection{\texorpdfstring{等差二次特性方程式型 \, $\displaystyle a_{n+2}+pa_{n+1}+qa_n=r$}{等差二次特性方程式型}}\label{節：等差二次特性方程式型}
    今度は右辺が0ではなく$r$となった.
    つまり$r$だけズレたのである.
    ズレたらもとに戻す,つまり不動点を探すのである.
    \[
      b_n=a_n-\alpha
    \]
    となる$\alpha$を求めると,
    \begin{align*}
      &(a_{n+2}-\alpha)+p(a_{n+1}-\alpha)+q(a_n-\alpha)=0\\
      &a_{n+2}+pa_{n+1}+qa_n=\alpha+p\alpha+q\alpha\\
      &\therefore \quad r=\alpha +p\alpha +q\alpha\\
      &\Leftrightarrow \quad \alpha=\frac{r}{p+q+1}\quad(p+q+1\ne0)
    \end{align*}
    ただし$p+q+1=0$かつ$r\ne0$のときは定数解$\alpha$が存在しない.\,$p+q+1=0$かつ$r=0$なら任意の定数が条件を満たし,元の漸化式はすでに前節の形である.定数解が存在しない場合は$d_n=a_{n+1}-a_n$とおくと,元の漸化式から$d_{n+1}=qd_n+r$となり,特性方程式型または階差型に帰着できる.
    あとは前節で述べた隣接3項間漸化式の問題である.

    \noindent \bookblocklink{example-three-term-2}{problem}{answer}{【例題】}
    \begin{enumerate}[label=(\arabic*),parsep=3pt,format=\bookitemlink{example-three-term-2}{problem}{answer}]
      \item $\displaystyle a_1=2,\,a_2=2,\, a_{n+2}=5a_{n+1}-6a_n+2$
      \item $\displaystyle a_1=1,\,a_2=2,\, a_{n+2}=4a_{n+1}-4a_n+3$
    \end{enumerate}

    \noindent\bookblocklink{example-three-term-2}{hint}{problem}{【方針】}
    \begin{enumerate}[label=(\arabic*),parsep=3pt,format=\bookitemlink{example-three-term-2}{hint}{problem}]
      \item まず不動点$\alpha$を求め,\,$b_n=a_n-\alpha$と置いて前節の二次特性方程式型に帰着させましょう.
      \item こちらは重解のパターンです.同じ手順で$\alpha$を求めてから前節の重解の場合の処理を行いましょう.
    \end{enumerate}

    \noindent\bookblocklink{example-three-term-2}{answer}{problem}{【解答】}
    \begin{enumerate}[label=(\arabic*),parsep=3pt,format=\bookitemlink{example-three-term-2}{answer}{problem}]
      \item $\displaystyle a_n=1+2^n-3^{n-1}$
      \item $\displaystyle a_n=3+2^{n-2}(3n-7)$
    \end{enumerate}

    \noindent\bookblocklink{example-three-term-2}{solution}{problem}{【解法】}
    \begin{enumerate}[label=(\arabic*),parsep=3pt,format=\bookitemlink{example-three-term-2}{solution}{problem}]
      \item $\displaystyle a_1=2,\,a_2=2,\, a_{n+2}=5a_{n+1}-6a_n+2$
            \begin{align*}
              &a_{n+2}=5a_{n+1}-6a_n+2\\
              &a_{n+2}-5a_{n+1}+6a_n=2
            \end{align*}
            ここで
            \[
              (a_{n+2}-\alpha)-5(a_{n+1}-\alpha)+6(a_n-\alpha)=0 
            \]
            を満たす$\alpha$を求めると,
            \[
              \alpha=1
            \]
            $b_n=a_n-1$と置くと,
            \begin{align*}
              &b_{n+2}=5b_{n+1}-6b_n\\
              &b_1=1,\, b_2=1
            \end{align*}
            （答案中略。先ほどの二次特性方程式型の例題(1)と同じ漸化式に帰着する.全く同じ計算を$b_n$に対して行う）
            \[
              b_n=2^n-3^{n-1}
            \]
            を得るから,
            \[
              a_n=1+b_n=1+2^n-3^{n-1}
            \]
      \item $\displaystyle a_1=1,\,a_2=2,\, a_{n+2}=4a_{n+1}-4a_n+3$
            \begin{align*}
              &a_{n+2}=4a_{n+1}-4a_n+3\\
              &a_{n+2}-4a_{n+1}+4a_n=3
            \end{align*}
            ここで
            \[
              (a_{n+2}-\alpha)-4(a_{n+1}-\alpha)+4(a_n-\alpha)=0 
            \]
            を満たす$\alpha$を求めると,
            \[
              \alpha=3
            \]
            なので,\,$b_n=a_n-3$と置くと,
            \begin{align*}
              &b_{n+2}=4b_{n+1}-4b_n\\
              &b_1=a_1-3=-2\\
              &b_2=a_2-3=-1                
            \end{align*}
            特性方程式$x^2-4x+4=0$は重解$\alpha'=2$を持つから,\,$c_n=b_{n+1}-2b_n$と置くと,
            \begin{align*}
              &c_{n+1}=2c_n\\
              &c_1=b_2-2b_1=-1-2(-2)=3\\
              &\therefore \quad c_n=3\cdot2^{n-1}
            \end{align*}
            したがって$b_{n+1}-2b_n=3\cdot2^{n-1}$である.両辺を$2^{n+1}$で割ると,
            \begin{align*}
              &\frac{b_{n+1}}{2^{n+1}}-\frac{b_n}{2^n}=\frac34
            \end{align*}
            $f_n=\dfrac{b_n}{2^n}$と置くと,
            \begin{align*}
              &f_{n+1}=f_n+\frac34,\quad f_1=\frac{b_1}{2}=-1\\
              &\therefore \quad f_n=-1+(n-1)\cdot\frac34=\frac{3n-7}{4}\\
              &\therefore \quad b_n=2^n\cdot f_n=2^{n-2}(3n-7)
            \end{align*}
            よって,
            \[
              a_n=3+b_n=3+2^{n-2}(3n-7)
            \]
    \end{enumerate}

% subsection: 総和型

\subsection{総和型}
  \subsubsection{総和と一般項の漸化式型 \, $S_n=pa_n+f(n)$}
    ある漸化式によって定められる数列$a_n$によって,
    \[
      S_n=\sum_{k=1}^{n} a_k
    \]
    として定められる数列$S_n$について考えよう.
    $S_n$と$a_n$の関係を表した
    \[
      S_n=pa_n+f(n)
    \]
    のような漸化式を考える.
    このような漸化式は今まで取り扱った形にはない,未知のものである.
    であれば既知の形にするのが数学である.
    これまでのことをきちんと理解していれば,\,$a_{n+1}$と$a_n$の関係が分かっている時には解くことができる.
    $S_n$をうまく変形して$a_{n}$関連だけの式を産むには,\,$p \neq 1$の時,
    \begin{align*}
      &S_n=pa_n+f(n) \tag*{\(\cdots\) (1)}\\
      &S_{n+1}=pa_{n+1}+f(n+1) \tag*{\(\cdots\) (2)}
    \end{align*}
    $(2)-(1)$をして
    \begin{align*}
      &a_{n+1}=pa_{n+1}-pa_n+f(n+1)-f(n)\\
      &a_{n+1}-pa_{n+1}=-pa_n+f(n+1)-f(n)\\
      &a_{n+1}(1-p)=-pa_n+f(n+1)-f(n)\\
      &a_{n+1}=\frac{-pa_n+f(n+1)-f(n)}{(1-p)}\\
      &a_{n+1}=\frac{p}{p-1}a_n+\frac{1}{1-p}\{f(n+1)-f(n)\}
    \end{align*}
    $\displaystyle g(n)=\frac{1}{1-p}\{f(n+1)-f(n)\}$とおくと,
    \begin{align*}
      &a_{n+1}=\frac{p}{p-1}a_n+g(n)
    \end{align*}
    従って漸化式は階差等比複合型(\ref{節：階差等比複合型})に帰着する(係数$\frac{p}{p-1}$は$p=p-1$となることがないため,決して$1$にはならない).
    \,$p=1$のときはこの割り算ができないが,\,$S_n=a_n+f(n)$と$S_n=S_{n-1}+a_n$($n\ge2$)を見比べることで,直接$a_n=f(n+1)-f(n)$が得られる.

    \noindent \textbf{【例題】}
    \begin{enumerate}[label=(\arabic*),parsep=3pt]
      \item $\displaystyle S_n=2a_n-n$
      \item $\displaystyle S_n=\frac12a_n+n$
    \end{enumerate}

    \noindent\textbf{【方針】}
    \begin{enumerate}[label=(\arabic*),parsep=3pt]
      \item まず$n=1$を代入して$a_1$を求めてから,本文の変形をなぞって$a_{n+1}$と$a_n$だけの漸化式を作りましょう.
      \item 同様です.できた漸化式は特性方程式型に帰着します.
    \end{enumerate}

    \noindent\textbf{【解答】}
    \begin{enumerate}[label=(\arabic*),parsep=3pt]
      \item $\displaystyle a_n=2^n-1$
      \item $\displaystyle a_n=1+(-1)^{n-1}$
    \end{enumerate}

    \noindent\textbf{【解法】}
    \begin{enumerate}[label=(\arabic*),parsep=3pt]
    \item $\displaystyle S_n=2a_n-n$

        $n=1$を代入すると,
        \[
          S_1=a_1=2a_1-1
        \]
        より
        \[
          a_1=1
        \]
        また,
        \begin{align*}
          a_{n+1}
            &=S_{n+1}-S_n\\
            &=\{2a_{n+1}-(n+1)\}-(2a_n-n)\\
            &=2a_{n+1}-2a_n-1
        \end{align*}
        したがって,
        \[
          a_{n+1}=2a_n+1
        \]
        ここで $b_n=a_n+1$ とおくと,
        \begin{align*}
          b_{n+1}
            &=a_{n+1}+1\\
            &=2a_n+2\\
            &=2(a_n+1)\\
            &=2b_n
        \end{align*}
        また,
        \[
          b_1=a_1+1=2
        \]
        であるから,
        \[
          b_n=2^n
        \]
        よって,
        \[
          a_n=2^n-1
        \]
    \item $\displaystyle S_n=\frac12a_n+n$
        $n=1$を代入すると,
        \[
          S_1=a_1=\frac12a_1+1
        \]
        より
        \[
          a_1=2
        \]
        \begin{align*}
          a_{n+1}
            &=S_{n+1}-S_n\\
            &=\left\{\frac12a_{n+1}+(n+1)\right\}-\left(\frac12a_n+n\right)\\
            &=\frac12a_{n+1}-\frac12a_n+1
        \end{align*}
        したがって,
        \[
          a_{n+1}=-a_n+2
        \]

        ここで $b_n=a_n-1$ とおくと,
        \begin{align*}
          b_{n+1}
            &=a_{n+1}-1\\
            &=-a_n+1\\
            &=-(a_n-1)\\
            &=-b_n
        \end{align*}
        また,
        \[
          b_1=a_1-1=1
        \]
        であるから,
        \[
          b_n=(-1)^{n-1}
        \]
        よって,
        \[
          a_n=1+(-1)^{n-1}
        \]

    \end{enumerate}

  \subsubsection{純総和型 \, $S_{n+1}=pS_n+f(n) $}
    両辺に総和がきてもやることは一緒で,1つずらした総和と元の式を足し引きしてやることで,\,$a_n$や$a_{n+1}$の項を作れば良い.
    \begin{align*}
      &S_{n+1}=pS_n+f(n) \tag*{\(\cdots\) (1)}\\
      &S_{n}=pS_{n-1}+f(n-1) \tag*{\(\cdots\) (2)}\\
    \end{align*}
      $(1)-(2)$をして,
    \begin{align*}
      &a_{n+1}=pa_{n}+f(n)-f(n-1)
    \end{align*}
    ただし(2)は(1)の$n$を$n-1$に置き換えたものなので,この式は$n\ge2$の範囲でのみ成り立つ.\,$a_1$と$a_2$の関係は,別途$S_1,S_2$から直接求める必要がある.
    あとはただの階差数列型の計算である.

    \noindent \textbf{【例題】}
    \begin{enumerate}[label=(\arabic*),parsep=3pt]
      \item $\displaystyle S_1=1,\ S_{n+1}=2S_n+n$
    \end{enumerate}

    \noindent\textbf{【方針】}
    \begin{enumerate}[label=(\arabic*),parsep=3pt]
      \item 本文の式は$n\ge2$でしか使えないことに注意し,\,$a_2$だけは$S_2=2S_1+f(1)$から別に求めましょう.
    \end{enumerate}

    \noindent\textbf{【解答】}
    \begin{enumerate}[label=(\arabic*),parsep=3pt]
      \item $\displaystyle a_1=1,\quad a_n=3\cdot2^{n-2}-1\ (n\ge2)$
    \end{enumerate}

    \noindent\textbf{【解法】}
    \begin{enumerate}[label=(\arabic*),parsep=3pt]
    \item $\displaystyle S_1=1,\ S_{n+1}=2S_n+n$

        $a_1=S_1=1$である.
        本文で示した公式
        \[
          a_{n+1}=pa_{n}+f(n)-f(n-1)
        \]
        に$p=2,\ f(n)=n$を代入すると,
        \[
          a_{n+1}=2a_n+\{n-(n-1)\}=2a_n+1\quad(n\ge2)
        \]
        を得る(この式は$n\ge2$でのみ成り立つことに注意する).

        実際,\,$n=1$のときは
        \[
          a_2=S_2-S_1=(2S_1+1)-S_1=2
        \]
        であるが,\,$2a_1+1=2\cdot1+1=3\ne a_2$であるから,この漸化式は$a_1$から$a_2$へは使えない.

        そこで,\,$n\ge2$の範囲で $b_n=a_n+1$ とおくと,
        \begin{align*}
          b_{n+1}
            &=a_{n+1}+1\\
            &=2a_n+2\\
            &=2(a_n+1)\\
            &=2b_n\qquad(n\ge2)
        \end{align*}
        また,
        \[
          b_2=a_2+1=3
        \]
        であるから,
        \[
          b_n=3\cdot2^{n-2}\quad(n\ge2)
        \]
        よって,
        \[
          a_1=1,\quad a_n=3\cdot2^{n-2}-1\quad(n\ge2)
        \]
    \end{enumerate}


% subsection: 連立型

\subsection{連立型}\label{節：連立型}
  \subsubsection{加減法型}
    一般形が
    \[
      \begin{cases}
        a_{n+1}=pa_{n}+qb_{n} \\
        b_{n+1}=qa_{n}+pb_{n}
      \end{cases}
    \]
    の形で表される漸化式の解き方について考えよう.
    二つの式では係数$p,q$の位置が入れ替わっている.
    足すと両方の係数が$p+q$になり,引くと$p-q$とその符号反転が現れる.
    そこで,和と差がそれぞれ単独で更新されないか調べる.
    実際にすすめていくと,
    \begin{align*}
      &a_{n+1}+b_{n+1}=(p+q)(a_{n}+b_{n})\\
      &a_n+b_n=(a_1+b_1)(p+q)^{n-1}\tag*{\(\cdots\) (1)}
    \end{align*}
    また,
    \begin{align*}
      &a_{n+1}-b_{n+1}=(p-q)(a_{n}-b_{n})\\
      &a_n-b_n=(a_1-b_1)(p-q)^{n-1}\tag*{\(\cdots\) (2)}
    \end{align*}
    $(1)+(2)$をすると
\BookMathLayout{\begin{align*}
      &2a_n=(a_1+b_1)(p+q)^{n-1}+(a_1-b_1)(p-q)^{n-1}\\
      &a_n=\frac{(a_1+b_1)(p+q)^{n-1}+(a_1-b_1)(p-q)^{n-1}}{2}\\
    \end{align*}}{
\[
\begin{aligned}
2a_n&=(a_1+b_1)(p+q)^{n-1}\\
&\quad+(a_1-b_1)(p-q)^{n-1}\\
a_n&=\frac{1}{2}\bigl\{(a_1+b_1)(p+q)^{n-1}\\
&\qquad+(a_1-b_1)(p-q)^{n-1}\bigr\}
\end{aligned}
\]
}
    $(1)-(2)$をすると
\BookMathLayout{\begin{align*}
      &2b_n=(a_1+b_1)(p+q)^{n-1}-(a_1-b_1)(p-q)^{n-1}\\
      &b_n=\frac{(a_1+b_1)(p+q)^{n-1}-(a_1-b_1)(p-q)^{n-1}}{2}
    \end{align*}}{
\[
\begin{aligned}
2b_n&=(a_1+b_1)(p+q)^{n-1}\\
&\quad-(a_1-b_1)(p-q)^{n-1}\\
b_n&=\frac{1}{2}\bigl\{(a_1+b_1)(p+q)^{n-1}\\
&\qquad-(a_1-b_1)(p-q)^{n-1}\bigr\}
\end{aligned}
\]
}
    \sectionref*{節：一次独立・基底・次元}で,初期条件を独立な二つの解に分けたことを思い出そう.
    ここでは初期の組$(a_1,b_1)$を
    \[
      \begin{aligned}
        (a_1,b_1)
        &=\frac{a_1+b_1}{2}(1,1)\\
        &\quad+\frac{a_1-b_1}{2}(1,-1)
      \end{aligned}
    \]
    と分けた.$(1,1)$から始まる部分は$p+q$倍,$(1,-1)$から始まる部分は$p-q$倍される.
    上の一般項は,二つの解を線形結合して元の初期条件に合わせた式である.
    和と差という補助量を取り出すことが,解を独立な二つの部分に分けることにつながっていた.
    別解及び演習は一般型の方で行う.
  \subsubsection{一般型}
    連立一般型の漸化式は,
    \[
      \begin{cases}
        a_{n+1}=pa_{n}+qb_{n} \\
        b_{n+1}=ra_{n}+sb_{n}
      \end{cases}
    \]
    で表される漸化式である.
    いつものように等比数列型への帰着を目指そう.
    $a_n$と$b_n$が互いに影響し合っているため,独立させて解くことは難しそうだ.
    ならばそれぞれについて求めるのではなく,いっそのこと合わせて求めてしまおうというのが解法の思想である.
    その流れは,
\BookMathLayout{\[
      \begin{array}{rrl}
        & a_{n+1}={} &pa_{n}+qb_{n}\\
        +\,) & b_{n+1}={} & ra_{n}+sb_{n}\\ \hline
        & a_{n+1}+b_{n+1}={} & (p+r)a_{n}+(q+s)b_{n}
      \end{array}
    \]}{
\[
\setlength{\arraycolsep}{2pt}
\begin{array}{rrl}
&a_{n+1}={}&pa_n+qb_n\\
+\,)&b_{n+1}={}&ra_n+sb_n\\\hline
&a_{n+1}+b_{n+1}={}&(p+r)a_n+(q+s)b_n
\end{array}\]
}
    について
    \[
      \alpha a_{n+1}+\beta b_{n+1}
      =\gamma(\alpha a_n+\beta b_n)
    \]
    を満たすようにしてやればよい.
    実際に左辺を計算すると,
\BookMathLayout{\begin{align*}
      \alpha a_{n+1}+\beta b_{n+1}&=\alpha(pa_n+qb_n)+\beta(ra_n+sb_n)\\
      &=(p\alpha+r\beta)a_n+(q\alpha+s\beta)b_n
    \end{align*}}{
\[
\begin{aligned}
\alpha a_{n+1}+\beta b_{n+1}
&=\alpha(pa_n+qb_n)\\ &\quad+\beta(ra_n+sb_n)\\
&=(p\alpha+r\beta)a_n\\ &\quad+(q\alpha+s\beta)b_n
\end{aligned}
\]
}
    一方,
    \[
      \gamma(\alpha a_n+\beta b_n)=\gamma\alpha a_n+\gamma\beta b_n
    \]
    である.
    したがって,
    \[
    \begin{cases}
      p\alpha+r\beta=\gamma\alpha\\
      q\alpha+s\beta=\gamma\beta
    \end{cases}
    \]
    となるように$\alpha,\beta,\gamma$を決めればよい.

    ここで,いきなり3つの文字をすべて求めようとすると少し複雑に見える.
    そこで,まず$\gamma$について考える.

    上の2式を変形すると,
    \[
    \begin{cases}
      (p-\gamma)\alpha+r\beta=0\\
      q\alpha+(s-\gamma)\beta=0
    \end{cases}
    \]
    である.
    $\alpha,\beta$は同時に$0$ではないので,この連立方程式が$\alpha,\beta$について解を持つためには,
    \[
      (p-\gamma)(s-\gamma)-qr=0
    \]
    でなければならない.
    整理すると,
    \[
      \gamma^2-(p+s)\gamma+(ps-qr)=0
    \]
    となる.
    この二次方程式を,この連立漸化式の特性方程式と呼ぶ.
    特性方程式の解を$\gamma_1,\gamma_2$とする.
    まず$\gamma=\gamma_1$を上の連立方程式に代入し,
    \[
      \begin{cases}
        (p-\gamma_1)\alpha+r\beta=0\\
        q\alpha+(s-\gamma_1)\beta=0
      \end{cases}
    \]
    を満たす$\alpha,\beta$を1組求める.
    すると,
    \[
      c_n=\alpha a_n+\beta b_n
    \]
    と置くことで,
    \begin{align*}
      c_{n+1}&=\alpha a_{n+1}+\beta b_{n+1}\\
      &=\gamma_1(\alpha a_n+\beta b_n)\\
      &=\gamma_1c_n
    \end{align*}
    となる.
    したがって,
    \[
      c_n=c_1\gamma_1^{\,n-1}
    \]
    という等比数列が得られる.
    同様に,\,$\gamma=\gamma_2$に対しても$\alpha,\beta$を求めれば,
    \[
      d_n=\alpha' a_n+\beta' b_n
    \]
    について
    \[
      d_n=d_1\gamma_2^{\,n-1}
    \]
    という等比数列が得られる.
    なお,ここで得られた2つの等比数列から$a_n,b_n$を求めるには,
    \[
      \begin{cases}
      c_n=\alpha a_n+\beta b_n\\
      d_n=\alpha'a_n+\beta'b_n
      \end{cases}
    \]
    を$a_n,b_n$について解けばよい.
    ただし,この連立方程式から$a_n,b_n$を一意に求めるには,係数の組$(\alpha,\beta)$と$(\alpha',\beta')$が比例しないこと,すなわち$\alpha\beta'-\alpha'\beta\ne0$が必要である.
    $\gamma_1\ne\gamma_2$ならこの条件を満たすが,これは十分条件であって必要条件ではない.
    重解の場合でも,例えば$a_{n+1}=2a_n,\,b_{n+1}=2b_n$では$(1,0),(0,1)$を選べる.
    一方,重解に対して比例しない2組を取れない場合は,別の工夫が必要になる(\sectionref{節：二次特性方程式型}の重解の場合と同様の考え方を用いる).
    また,$\,\gamma_1,\gamma_2$が虚数になる場合もこの手法は形式的には成り立つが,実数の形に書き直すには\sectionref{節：虚数解型隣接3項間漸化式}で扱う手法が必要になる.
    このように,\,$a_n,b_n$をそれぞれ直接求めるのではなく,\,$a_n,b_n$の和の形を作ることで,連立していた漸化式を等比数列へと分離することができる.

    \noindent \bookblocklink{example-system-1}{problem}{answer}{【例題】}
    \begin{enumerate}[label=(\arabic*),parsep=3pt,format=\bookitemlink{example-system-1}{problem}{answer}]
      \item $\displaystyle a_1=1,\,b_1=1.\ \begin{cases}a_{n+1}=3a_n+b_n\\ b_{n+1}=a_n+3b_n\end{cases}$
      \item $\displaystyle a_1=1,\,b_1=0,\ \begin{cases}a_{n+1}=3a_n+b_n\\ b_{n+1}=2a_n+2b_n\end{cases}$
    \end{enumerate}

    \noindent\bookblocklink{example-system-1}{hint}{problem}{【方針】}
    \begin{enumerate}[label=(\arabic*),parsep=3pt,format=\bookitemlink{example-system-1}{hint}{problem}]
      \item 前節で述べた問題です.一般型,加減法型,どちらのやり方でも問題ありません.
      \item これはこの節で述べたものです.
    \end{enumerate}

    \noindent\bookblocklink{example-system-1}{answer}{problem}{【解答】}
    \begin{enumerate}[label=(\arabic*),parsep=3pt,format=\bookitemlink{example-system-1}{answer}{problem}]
      \item $\displaystyle a_n=4^{n-1},\qquad b_n=a_n=4^{n-1}$
      \item $\displaystyle a_n=\frac{2\cdot4^{n-1}+1}{3},\qquad b_n=\frac{2\cdot4^{n-1}-2}{3}$
    \end{enumerate}

    \noindent\bookblocklink{example-system-1}{solution}{problem}{【解法】}
    \begin{enumerate}[label=(\arabic*),parsep=3pt,format=\bookitemlink{example-system-1}{solution}{problem}]
    \item $\displaystyle a_1=1,\,b_1=1.\ \begin{cases}a_{n+1}=3a_n+b_n\\ b_{n+1}=a_n+3b_n\end{cases}$\\
          \begin{enumerate}
            \item 【解法1】\\
                  $\alpha a_{n+1}+\beta b_{n+1}$を計算すると,
                  \begin{align*}
                    \alpha a_{n+1}+\beta b_{n+1}
                      &=\alpha(3a_n+b_n)\\
                      &\qquad +\beta(a_n+3b_n)\\
                      &=(3\alpha+\beta)a_n\\
                      &\qquad +(\alpha+3\beta)b_n
                  \end{align*}
                  これが
                  \[
                    \gamma(\alpha a_n+\beta b_n)
                    =\gamma\alpha a_n+\gamma\beta b_n
                  \]
                  と一致するためには,
                  \[
                  \begin{cases}
                    (3-\gamma)\alpha+\beta=0\\
                    \alpha+(3-\gamma)\beta=0
                  \end{cases}
                  \]
                  でなければならない.

                  $\alpha,\beta$が同時に$0$ではないため,
                  \[
                    (3-\gamma)^2-1=0
                  \]
                  より,
                  \begin{align*}
                    &\gamma^2-6\gamma+8=0\\
                    &(\gamma-2)(\gamma-4)=0
                  \end{align*}
                  したがって,
                  \[
                    \gamma_1=2,\qquad\gamma_2=4
                  \]
                  である.

                  まず$\gamma=2$とすると,
                  \[
                  \begin{cases}
                    \alpha+\beta=0\\
                    \alpha+\beta=0
                  \end{cases}
                  \]
                  となるので,
                  \[
                    \alpha=1,\qquad\beta=-1
                  \]
                  と取ることができる.

                  そこで,
                  \[
                    c_n=a_n-b_n
                  \]
                  とおくと,
                  \begin{align*}
                    c_{n+1}
                      &=a_{n+1}-b_{n+1}\\
                      &=(3a_n+b_n)-(a_n+3b_n)\\
                      &=2(a_n-b_n)\\
                      &=2c_n
                  \end{align*}
                  また,
                  \[
                    c_1=a_1-b_1=0
                  \]
                  であるから,
                  \[
                    c_n=0
                  \]
                  したがって,
                  \[
                    a_n-b_n=0
                  \]

                  次に$\gamma=4$とすると,
                  \[
                  \begin{cases}
                    -\alpha+\beta=0\\
                    \alpha-\beta=0
                  \end{cases}
                  \]
                  となるので,
                  \[
                    \alpha=1,\qquad\beta=1
                  \]
                  と取ることができる.

                  そこで,
                  \[
                    d_n=a_n+b_n
                  \]
                  とおくと,
                  \begin{align*}
                    d_{n+1}
                      &=a_{n+1}+b_{n+1}\\
                      &=(3a_n+b_n)+(a_n+3b_n)\\
                      &=4(a_n+b_n)\\
                      &=4d_n
                  \end{align*}
                  また,
                  \[
                    d_1=a_1+b_1=2
                  \]
                  であるから,
                  \[
                    d_n=2\cdot4^{n-1}
                  \]

                  以上より,
                  \[
                  \begin{cases}
                    a_n-b_n=0\\
                    a_n+b_n=2\cdot4^{n-1}
                  \end{cases}
                  \]
                  となる.

                  2式を辺々加えると,
                  \[
                    2a_n=2\cdot4^{n-1}
                  \]
                  であるから,
                  \[
                    a_n=4^{n-1}
                  \]
                  また,
                  \[
                    b_n=a_n=4^{n-1}
                  \]
            \item 【解法2】\\
                  2つの式を辺々加えると,
                  \begin{align*}
                    a_{n+1}+b_{n+1}
                      &=(3a_n+b_n)\\
                      &\quad +(a_n+3b_n)\\
                      &=4(a_n+b_n)
                  \end{align*}
                  よって,
                  \[
                    a_n+b_n=(a_1+b_1)4^{n-1}=2\cdot4^{n-1}. \tag*{\(\cdots\) (1)}
                  \]

                  また,2つの式を辺々引くと,
                  \begin{align*}
                    a_{n+1}-b_{n+1}
                      &=(3a_n+b_n)\\
                      &\quad -(a_n+3b_n)\\
                      &=2(a_n-b_n)
                  \end{align*}
                  よって,
                  \[
                    a_n-b_n=(a_1-b_1)2^{n-1}=0. \tag*{\(\cdots\) (2)}
                  \]

                  \((1)+(2)\)をすると,
                  \[
                    2a_n=2\cdot4^{n-1}
                  \]
                  なので,
                  \[
                    a_n=4^{n-1}
                  \]

                  また,\((1)-(2)\)をすると,
                  \[
                    2b_n=2\cdot4^{n-1}
                  \]
                  なので,
                  \[
                    b_n=4^{n-1}
                  \]
          \end{enumerate}
    \item $\displaystyle a_1=1,\,b_1=0,\ \begin{cases}a_{n+1}=3a_n+b_n\\ b_{n+1}=2a_n+2b_n\end{cases}$\\
          $\alpha a_{n+1}+\beta b_{n+1}$を計算すると,
          \begin{align*}
            \alpha a_{n+1}+\beta b_{n+1}
              &=\alpha(3a_n+b_n)\\
              &\quad +\beta(2a_n+2b_n)\\
              &=(3\alpha+2\beta)a_n\\
              &\quad +(\alpha+2\beta)b_n
          \end{align*}
          これが
          \[
            \gamma(\alpha a_n+\beta b_n)
            =\gamma\alpha a_n+\gamma\beta b_n
          \]
          と一致するためには,
          \[
          \begin{cases}
            (3-\gamma)\alpha+2\beta=0\\
            \alpha+(2-\gamma)\beta=0
          \end{cases}
          \]
          でなければならない.

          $\alpha,\beta$が同時に$0$ではないため,
          \[
            (3-\gamma)(2-\gamma)-2=0
          \]
          より,
          \begin{align*}
            &\gamma^2-5\gamma+4=0\\
            &(\gamma-1)(\gamma-4)=0
          \end{align*}
          したがって,
          \[
            \gamma_1=1,\qquad\gamma_2=4
          \]
          である.

          まず$\gamma=1$とすると,
          \[
          \begin{cases}
            2\alpha+2\beta=0\\
            \alpha+\beta=0
          \end{cases}
          \]
          となるので,
          \[
            \alpha=1,\qquad\beta=-1
          \]
          と取ることができる.

          そこで,
          \[
            c_n=a_n-b_n
          \]
          とおくと,
          \begin{align*}
            c_{n+1}
              &=a_{n+1}-b_{n+1}\\
              &=(3a_n+b_n)-(2a_n+2b_n)\\
              &=a_n-b_n\\
              &=c_n
          \end{align*}
          また,
          \[
            c_1=a_1-b_1=1
          \]
          であるから,
          \[
            c_n=1
          \]
          したがって,
          \[
            a_n-b_n=1
          \]

          次に$\gamma=4$とすると,
          \[
          \begin{cases}
            -\alpha+2\beta=0\\
            \alpha-2\beta=0
          \end{cases}
          \]
          となるので,
          \[
            \alpha=2,\qquad\beta=1
          \]
          と取ることができる.

          そこで,
          \[
            d_n=2a_n+b_n
          \]
          とおくと,
          \begin{align*}
            d_{n+1}
              &=2a_{n+1}+b_{n+1}\\
              &=2(3a_n+b_n)+(2a_n+2b_n)\\
              &=8a_n+4b_n\\
              &=4(2a_n+b_n)\\
              &=4d_n
          \end{align*}
          また,
          \[
            d_1=2a_1+b_1=2
          \]
          であるから,
          \[
            d_n=2\cdot4^{n-1}
          \]

          以上より,
          \[
          \begin{cases}
            a_n-b_n=1\\
            2a_n+b_n=2\cdot4^{n-1}
          \end{cases}
          \]
          となる.

          2式を辺々加えると,
          \[
            3a_n=1+2\cdot4^{n-1}
          \]
          であるから,
          \[
            a_n=\frac{2\cdot4^{n-1}+1}{3}
          \]
          また,
          \[
            b_n=a_n-1
          \]
          より,
          \[
            b_n=\frac{2\cdot4^{n-1}-2}{3}
          \]
    \end{enumerate}



% section: 発展パターン

\section{発展パターン}\label{節：発展パターン}
ここでは更に発展的なパターンについて紹介する.
扱われる漸化式は,一部の難関大において時々出題されるようなものである.
基本的に誘導がついた状態で出題されるため,ひと目見て解けるようになる必要はそこまで高くないが,知っていて損はない.
ここで追加する変数係数型は階比型と階差型の組合せ,
二重数列は階差と帰納法を二方向へ広げる見方である.
新しい型名を覚えるだけでなく,既習の操作がどこで働くかを確かめよう.
\subsection{変数係数の一次漸化式}\label{節：変数係数の一次漸化式}
  \sectionref*{節：階比型}では添字ごとの倍率を掛け合わせ,
  \sectionref*{節：階差等比複合型}では同じ更新をする量を引いた.
  ここでは,倍率と付加項がどちらも添字に依存する場合を,この二つの操作で解こう.

  \subsubsection{倍率をそろえて,変化を足す}
  まず
  \BookMathLayout{\[
    a_1=2,\qquad
    a_{n+1}=\frac{n+2}{n}a_n+(n+1)(n+2)
    \quad(n\ge1)
  \]}{\[
    \begin{aligned}
      a_1&=2,\\
      a_{n+1}&=\frac{n+2}{n}a_n+(n+1)(n+2)\\
      &\quad(n\ge1)
    \end{aligned}
  \]}
  を考える.係数の分子・分母に注目すると,
  \[
    \frac{(n+1)(n+2)}{n(n+1)}=\frac{n+2}{n}
  \]
  である.\textbf{隣り合う項の比が元の倍率になる量}として$h_n=n(n+1)$を選べる.
  これは\sectionref*{節：補助数列の設計}で目標から割る量を選んだ方法である.
  $b_n=\frac{a_n}{n(n+1)}$と置いて実際に計算すると,
  \[
    \frac{a_{n+1}}{(n+1)(n+2)}
      =\frac{a_n}{n(n+1)}+1
  \]
  となる.従って$b_1=1$, $b_{n+1}=b_n+1$より$b_n=n$であり,
  \[
    a_n=n^2(n+1)
  \]
  を得る.分母は$n\ge1$で非零である.初項と元の式へ代入しても確かめられる.

  \subsubsection{一般の手順と,割り算の条件}
  $a_{n+1}=p_na_n+q_n$ $(n\ge1)$に対して,
  $h_1=1$, $h_{n+1}=p_nh_n$とする.
  すべての$p_n$が非零なら,$h_n$も非零であり,
  \[
    h_n=\prod_{k=1}^{n-1}p_k
  \]
  である.空積は$1$とする.両辺を$h_{n+1}$で割ると,
  \[
    b_n=\frac{a_n}{h_n},\qquad
    b_{n+1}-b_n=\frac{q_n}{h_{n+1}}
  \]
  だから,階差を足して
  \[
    a_n=h_n\left(a_1+\sum_{k=1}^{n-1}\frac{q_k}{h_{k+1}}\right)
  \]
  を得る.先ほどは定数倍して$h_n=n(n+1)$を選んだが,
  その場合の$b_1$を正しく合わせれば同じ答えになる.

  \begin{bookpoint}
    \textbf{正規化の役割}\par
    倍率を取り除くと,残った付加項を足す問題になる.
    積・和の表示を得ることと,最後まで簡単な式で計算できることは別である.
  \end{bookpoint}
  もし$p_j=0$なら,$a_{j+1}=q_j$はそれ以前の項によらず決まる.
  零を含む$h_n$で割らず,零となる更新の前後で区間を分け,
  非零の倍率が続く範囲で同じ方法を使う.

  \subsubsection{同じ更新をする量を引く別解}
  先ほどの例では,$v_n=n^2(n+1)$が元の更新を満たすことを直接確かめられる.
  この候補は,正規化後の$b_n=n$から見つけたものである.
  $d_n=a_n-v_n$と置くと,
  \[
    d_{n+1}=\frac{n+2}{n}d_n
  \]
  となり,階比型へ戻る.初項を一般の$a_1=s$に変えると,
  \[
    a_n=n(n+1)\left(n+\frac{s-2}{2}\right)
  \]
  である.一つの解$v_n$と,同次の解の定数倍に分かれている.
  これは\sectionref*{節：非同次と解の組合せ}の見方そのものである.

  \begin{bookpoint}
    \textbf{二つの方法の選び方}\par
    隣り合う比が分かるなら割る因子を探す.
    同じ更新をする既知の量が見えるならそれを引く.
    どちらも,付加項と倍率を分けて扱うための操作である.
  \end{bookpoint}
  操作だけの比較は\sectionref*{節：操作を選ぶ練習},
  一般項まで求める練習は\bookproblemref{practice-connections}{1}{接続演習問1}へ進もう.

\subsection{ガウス型}\label{節：ガウス型}
  ここまでの漸化式では,差を取る,定数を引く,別の数列に置き換えるなどの操作から,既知の形へ帰着させてきた.
  ガウス記号を含む場合も,この方針は変わらない.
  ただし,変形の手掛かりとして,整数性,不等式,余り,偶奇,数え上げなどが必要になる.
  ガウス型では,整数問題で用いる考え方が前面に現れるのである.

  ガウス記号があるというだけで,一つの解法が決まるわけではない.
  本項では,\textbf{ガウス記号が式の中で何をしているか}に注目して,10種類の使い方を整理する.
  各解説の直後に代表例題を一つ置いた.
  同じ問題が複数の使い方に関係することもあるので,分類は排他的な区分ではなく,着眼点を選ぶための見取り図として使おう.

  \subsubsection{下準備と分類の見取り図}\label{節：ガウス型の下準備}
    ガウス記号$\lfloor x\rfloor$または$[x]$は,実数$x$を超えない最大の整数を表し,床関数と呼ばれる.
    整数$k$について,
    \[
      \lfloor x\rfloor=k\iff k\le x<k+1
    \]
    である.\,$x\ge0$では小数点以下を切り捨てた値に一致するが,
    負数では0に近づく切り捨てと異なる.例えば$\lfloor-1.5\rfloor=-2$である.
    天井関数$\lceil x\rceil$は$x$以上の最小の整数であり,
    \[
      \lceil x\rceil=k\iff k-1<x\le k
    \]
    で定義される.

    \BookFigFloorGraph

    式を読むときは,床関数が\textbf{項の値,添字,付加項のどれに作用するか}を最初に確かめよう.
    そのうえで,以下の使い方を考える.
    \begin{enumerate}[label=(\arabic*)]
      \item 整数性・不等式で整数部分を確定する.
      \item 同じ値を取る区間を群としてまとめる.
      \item 余りを取り出し,有限個の値と周期を調べる.
      \item 床関数の差で境界・条件・繰り上がりを検出する.
      \item 丸め誤差と係数の関係を使う.
      \item 整数への丸めで不動点と停止を調べる.
      \item 添字を二つに分割する.
      \item 商と余りで整数の桁を読む.
      \item 小数部分を残し,回転として読む.
      \item 条件を満たす整数や格子点の個数を数える.
    \end{enumerate}
    最初の6種類は主に値やその変化,次の2種類は添字,
    最後の2種類は小数部分や個数という意味に注目するものである.
    一般項を求めるだけでなく,周期や止まる値,何を数えているかを調べることにも意味がある.

    なお,整数$k$について$\lfloor x+k\rfloor=\lfloor x\rfloor+k$だが,
    実数同士では$\lfloor x+y\rfloor=\lfloor x\rfloor+\lfloor y\rfloor$とは限らない.
    数項を計算する実験は着眼点を探すために使い,予想は不等式や帰納法などで確かめよう.

  \input{第3章 漸化式の解説/06_発展パターン/01_ガウス型/01_整数部分の確定.tex}
  \input{第3章 漸化式の解説/06_発展パターン/01_ガウス型/02_群数列.tex}
  \input{第3章 漸化式の解説/06_発展パターン/01_ガウス型/03_余りと周期.tex}
  \input{第3章 漸化式の解説/06_発展パターン/01_ガウス型/04_境界と条件.tex}
  \input{第3章 漸化式の解説/06_発展パターン/01_ガウス型/05_丸め誤差.tex}
  \input{第3章 漸化式の解説/06_発展パターン/01_ガウス型/06_不動点と停止.tex}
  \input{第3章 漸化式の解説/06_発展パターン/01_ガウス型/07_添字の分割.tex}
  \input{第3章 漸化式の解説/06_発展パターン/01_ガウス型/08_桁の抽出.tex}
  \input{第3章 漸化式の解説/06_発展パターン/01_ガウス型/09_小数部分と回転.tex}
  \input{第3章 漸化式の解説/06_発展パターン/01_ガウス型/10_個数と格子点.tex}

  \medskip
  \noindent \textbf{【着眼点の振り返り】}\\
    ガウス記号の値を不等式で確定する方法,群や余りで値を整理する方法,
    添字を分割する方法,小数部分や個数として読む方法を見てきた.
    そのどれを使うかは,記号の位置と,式の中で果たす役割から考えよう.
    余りと桁,群数列と数え上げのように,一つの問題から別の見方へつながることもある.

\subsection{隣接4項間漸化式型}\label{節：隣接4項間漸化式型}
隣接2項間や隣接3項間の漸化式を今まで扱った.
これをより広い範囲に適用してみよう.
例えば
\[
  a_{n+3}+pa_{n+2}+qa_{n+1}+ra_n=0
\]
のような漸化式を考える.
これは三階の同次線形漸化式なので,\sectionref*{節：一次独立・基底・次元}で述べたように,
初めの三項を自由に選ぶ解全体は3次元である.
ここでは一つの補助数列を先に求めることで,残りの計算を二階の式へ分ける.
\sectionref*{節：補助数列の設計}で述べたように,先に単純な更新を決め,それをする量を探そう.
3項をまとめた$d_n=a_{n+2}+\alpha a_{n+1}+\beta a_n$を作り,
$d_{n+1}=\gamma d_n$となることを目指す.
その条件は
\BookMathLayout{\[
  a_{n+3}+\alpha a_{n+2}+\beta a_{n+1}=\gamma \left(a_{n+2}+\alpha a_{n+1}+\beta a_{n}\right)
\]}{
\[
\begin{aligned}
a_{n+3}+\alpha a_{n+2}+\beta a_{n+1}\\
{}=\gamma\left(a_{n+2}+\alpha a_{n+1}+\beta a_n\right)
\end{aligned}
\]
}
である.両辺を左側に集め,元の漸化式と係数を比べると,
\[
  \alpha-\gamma=p,\qquad
  \beta-\gamma\alpha=q,\qquad
  -\gamma\beta=r
\]
となる.最初の二式から
\[
  \alpha=p+\gamma,\qquad
  \beta=q+p\gamma+\gamma^2
\]
なので,最後の式へ代入すると
\[
  \gamma^3+p\gamma^2+q\gamma+r=0
\]
を得る.この三次方程式の実数解を一つ選べば,$\alpha,\beta$も決まる.
三次方程式は少なくとも一つ実数解を持つが,計算しやすい値を見つけられるかは別の問題である.
以下の例題では因数分解や階差の形を手掛かりに選ぶ.
これによって,\,$d_n=a_{n+2}+\alpha a_{n+1}+\beta a_n$とおくと$d_{n+1}=\gamma d_n$が成り立つので,
\[
  a_{n+2}+\alpha a_{n+1}+\beta a_n=d_1\gamma^{\,n-1}
\]
という3項間漸化式が得られる.
ただし,\,\sectionref{節：等差二次特性方程式型}で扱ったのは右辺が定数$r$の場合であり,ここでは右辺が指数関数$d_1\gamma^{n-1}$になる点が異なる.
この形は\sectionref{節：階差等比複合型}の考え方を3項間に応用することで解くことができる.
ここでさらに同じ更新をする既知の量を引くと,右辺の指数関数を消せる場合がある.
例えば第1問では,$p=-3,q=3,r=-1$なので上の方程式は$(\gamma-1)^3=0$となり,
$\gamma=1,\alpha=-2,\beta=1$を得る.
補助数列が$d_n=a_{n+2}-2a_{n+1}+a_n$となる理由は,この係数の一致にある.
元の式が三階差の形であると気付けば,同じ量を階差からも作れる.

  次の数列の一般項を求めてみよう.

  \noindent \bookblocklink{example-four-term-1}{problem}{answer}{【例題】}
  \begin{enumerate}[label=(\arabic*),parsep=3pt,format=\bookitemlink{example-four-term-1}{problem}{answer}]
    \item $\displaystyle a_1=1,\,a_2=2,\,a_3=4$\\
          $a_{n+3}=3a_{n+2}-3a_{n+1}+a_n\quad(n\ge1)$
    \item $\displaystyle a_1=2,\,a_2=1,\,a_3=5$\\
          $a_{n+3}=2a_{n+2}+a_{n+1}-2a_n\quad(n\ge1)$
    \item $\displaystyle a_1=1,\,a_2=4,\,a_3=11$\\
          $a_{n+3}=5a_{n+2}-7a_{n+1}+3a_n\quad(n\ge1)$
  \end{enumerate}

  \noindent\bookblocklink{example-four-term-1}{hint}{problem}{【方針】}
  \begin{enumerate}[label=(\arabic*),parsep=3pt,format=\bookitemlink{example-four-term-1}{hint}{problem}]
    \item $d_n=a_{n+2}-2a_{n+1}+a_n$とおき,階差をさらにとると一定になることを利用します.
    \item $d_n=a_{n+2}-a_{n+1}-2a_n$とおくと$d_{n+1}=d_n$となります.初期条件から$d_n=0$を導き,3項間漸化式として解きます.
    \item $d_n=a_{n+2}-2a_{n+1}+a_n$とおくと$d_{n+1}=3d_n$となります.得られた右辺を満たす数列を引き,階差が一定の形に帰着させます.
  \end{enumerate}

  \noindent\bookblocklink{example-four-term-1}{answer}{problem}{【解答】}
  \begin{enumerate}[label=(\arabic*),parsep=3pt,format=\bookitemlink{example-four-term-1}{answer}{problem}]
    \item $\displaystyle a_n=\frac{n^2-n+2}{2}$
    \item $\displaystyle a_n=2^{n-1}+(-1)^{n-1}$
    \item $\displaystyle a_n=n-1+3^{n-1}$
  \end{enumerate}

  \noindent\bookblocklink{example-four-term-1}{solution}{problem}{【解法】}
  \begin{enumerate}[label=(\arabic*),parsep=3pt,format=\bookitemlink{example-four-term-1}{solution}{problem}]
    \item 漸化式を移項すると,
          \[
            a_{n+3}-3a_{n+2}+3a_{n+1}-a_n=0
          \]
          である.ここで
          \[
            d_n=a_{n+2}-2a_{n+1}+a_n
          \]
          とおくと,
          \[
            d_{n+1}-d_n=a_{n+3}-3a_{n+2}+3a_{n+1}-a_n=0
          \]
          より$d_n$は一定である.初期条件から$d_1=4-2\cdot2+1=1$なので,
          \[
            a_{n+2}-2a_{n+1}+a_n=1
          \]
          となる.さらに$e_n=a_{n+1}-a_n$とおくと$e_{n+1}-e_n=1$であり,$e_1=1$だから$e_n=n$である.したがって,
          \[
            a_n=a_1+\sum_{k=1}^{n-1}e_k
                =1+\sum_{k=1}^{n-1}k
                =\frac{n^2-n+2}{2}
          \]
          を得る.
    \item
          \[
            d_n=a_{n+2}-a_{n+1}-2a_n
          \]
          とおくと,
          \[
            d_{n+1}-d_n=a_{n+3}-2a_{n+2}-a_{n+1}+2a_n=0
          \]
          である.また,$d_1=5-1-2\cdot2=0$だから$d_n=0$であり,
          \[
            a_{n+2}-a_{n+1}-2a_n=0
          \]
          となる.特性方程式$x^2-x-2=0$の解は$2,-1$なので,
          \[
            a_n=A\,2^{n-1}+B(-1)^{n-1}
          \]
          とおける.初期条件より$A+B=2,\,2A-B=1$であるから$A=B=1$を得る.
    \item
          \[
            d_n=a_{n+2}-2a_{n+1}+a_n
          \]
          とおくと,与えられた漸化式から$d_{n+1}=3d_n$である.初期条件より$d_1=11-2\cdot4+1=4$なので,
          \[
            a_{n+2}-2a_{n+1}+a_n=4\cdot3^{n-1}
          \]
          となる.一方,$3^{n-1}$の2階差も$4\cdot3^{n-1}$である.そこで$b_n=a_n-3^{n-1}$とおくと,
          \[
            b_{n+2}-2b_{n+1}+b_n=0
          \]
          となる.$b_1=0,\,b_2=1$より$b_n=n-1$であるから,
          \[
            a_n=b_n+3^{n-1}=n-1+3^{n-1}
          \]
          を得る.
  \end{enumerate}

\subsection{虚数解型隣接3項間漸化式}\label{節：虚数解型隣接3項間漸化式}
  \BookFigComplexRoots

  今までの隣接3\,項間漸化式(\sectionref{節：隣接3項間漸化式型})では,特性方程式の解$\alpha,\beta$が異なれば$a_n=A\alpha^{n-1}+B\beta^{n-1}$という式自体は虚数解の場合でもそのまま成り立つ.
  ここでは,特性方程式が虚数解をもつ場合に,この式を三角関数を用いた実数の形に書き直す方法を考える.

  本項では,実数解で使った「特性根の冪を組み合わせる」流れを複素数へ広げる.
  \sectionref*{節：解を組み合わせる見方}で学んだ線形結合と初期条件の一意性を使い,
  実数の形への書き直しを確かめる.その後に,補助数列と加法定理から同じ形を見る.
  実数の係数$p,q$について,まず,
  \[
    a_{n+2}+pa_{n+1}+qa_n=0
  \]
  を考える.特性方程式は
  \[
    x^2+px+q=0
  \]
  である.

  これまで,特性方程式の解が異なる実数
  \[
    x=\alpha,\beta
  \]
  をもつ場合には,
  \[
    a_n=A\alpha^{n-1}+B\beta^{n-1}
  \]
  という形で考えた.本項の非実の根は零でないので,指数を一つずらした分を係数に吸収できる.
  以下では指数を$n$にした表し方も使う.

  では,特性方程式が実数解をもたず,
  \[
    x=\alpha,\overline{\alpha}
  \]
  という互いに共役な解をもつ場合には,一般項はどのような形になるだろうか.

  ここで「二つの冪で解全体を表せる」という理由を,前の共通説明につなげておこう.
  $p^2-4q<0$だから$q>0$であり,$\alpha,\overline\alpha$は異なる非零の根である.
  $u_n=\alpha^{n-1}$を元の式へ代入すると,
  \[
    u_{n+2}+pu_{n+1}+qu_n
    =\alpha^{n-1}(\alpha^2+p\alpha+q)=0
  \]
  となる.$v_n=\overline\alpha^{\,n-1}$も同じ式を満たす.
  それらの初めの二項は$(1,\alpha)$と$(1,\overline\alpha)$である.
  $Cu_n+Dv_n=0$なら$n=1,2$から
  \[
    C+D=0,\qquad C\alpha+D\overline\alpha=0
  \]
  を得る.$\alpha\ne\overline\alpha$なので$C=D=0$であり,二つの解は一次独立である.
  また,どんな初期条件もこの二つの係数で一意に合わせられる.
  従って,\sectionref*{節：一次独立・基底・次元}で述べた2次元の解全体の基底になっている.
  ここでは係数$C,D$も複素数として考えている.
  初めの二項を合わせれば,同じ更新からすべての項の一致が分かる.
  非実の根でも,実数解のときと同じ理由で冪の線形結合を使えるのである.

  例えば,
  \[
    x^2-2x+2=0
  \]
  を考えると,
  \[
    x=1\pm i
  \]
  である.したがって,漸化式
  \[
    a_{n+2}=2a_{n+1}-2a_n
  \]
  の特性方程式は虚数解をもつ.

  したがって,この例の一般項を
  \[
    a_n=A(1+i)^n+B(1-i)^n
  \]
  と書くことができる.しかし,\,$a_n$が実数列である場合,このままでは一般項の実数性が分かりにくい.そこで,複素数の極形式を利用する.

  一般に,特性方程式の解を
  \[
    \alpha=\rho(\cos\theta+i\sin\theta)=\rho e^{i\theta}
  \]
  とすると,共役解は
  \[
    \overline{\alpha}=\rho(\cos\theta-i\sin\theta)=\rho e^{-i\theta}
  \]
  である.したがって,
  \[
    \alpha^n=\rho^n e^{in\theta}
    =\rho^n(\cos n\theta+i\sin n\theta)
  \]
  および
  \[
    \overline{\alpha}^{\,n}
    =\rho^n e^{-in\theta}
    =\rho^n(\cos n\theta-i\sin n\theta)
  \]
  となる.

  ここで,$a_n$が実数列なら,共役を取っても値が変わらないので,
  \[
    A\alpha^n+B\overline\alpha^{\,n}
    =\overline B\alpha^n+\overline A\,\overline\alpha^{\,n}
  \]
  である.先ほど確かめた一次独立性から,係数の表し方は一意だから,
  $B=\overline A$となる.実数の解を選ぶ条件が,係数を共役にすることとして現れた.
  実際,\,$A=\dfrac{C-iD}{2},\ B=\dfrac{C+iD}{2}$(\,$C,D$は実数)とおくと,
\BookMathLayout{\begin{align*}
    A\alpha^n+B\overline{\alpha}^{\,n}
    &=\frac{C-iD}{2}\rho^n(\cos n\theta+i\sin n\theta)\\
    &\qquad+\frac{C+iD}{2}\rho^n(\cos n\theta-i\sin n\theta)\\
    &=\rho^n(C\cos n\theta+D\sin n\theta)
  \end{align*}}{
\[
\begin{aligned}
A\alpha^n+B\overline\alpha^{\,n}\\
{}=\frac{C-iD}{2}\rho^n(\cos n\theta+i\sin n\theta)\\
\quad+\frac{C+iD}{2}\rho^n(\cos n\theta-i\sin n\theta)\\
{}=\rho^n(C\cos n\theta+D\sin n\theta)
\end{aligned}
\]
}
  のように,虚部が打ち消し合って
  \[
    a_n=\rho^n(C\cos n\theta+D\sin n\theta)
  \]
  という実数の形に書き直せることが分かる.

  したがって,虚数解をもつ隣接3\,項間漸化式では,
  \[
    a_n=\rho^n(C\cos n\theta+D\sin n\theta)
  \]
  という形が基本的な一般項となる.ここで,\,$\rho$は特性根の絶対値,\,$\theta$はその偏角である.

  では,一般項
  \[
    a_n=\rho^n(C\cos n\theta+D\sin n\theta)
  \]
  に含まれる定数$C,D$をどのように決定すればよいだろうか.

  これまでの漸化式と同様に,初期条件を代入することで$C,D$を決定することができる.例えば,\,$a_1,a_2$が与えられているとする.一般項に$n=1,2$を代入すると,
  \[
    a_1=\rho(C\cos\theta+D\sin\theta)
  \]
  および
  \[
    a_2=\rho^2(C\cos2\theta+D\sin2\theta)
  \]
  を得る.

  これは$C,D$についての連立方程式である.したがって,この2つの方程式を解けば$C,D$が決定される.

  特に,$a_0,a_1$から始め,漸化式を$n\ge0$で使う場合には,より簡単になる.一般項に$n=0$を代入すると,
  \[
    a_0=C
  \]
  となる.したがって,
  \[
    C=a_0
  \]
  が直ちに得られる.

  さらに,\,$n=1$を代入すると,
  \[
    a_1=\rho(C\cos\theta+D\sin\theta)
  \]
  であるから,
  \[
    D=
    \cfrac{\frac{a_1}{\rho}-C\cos\theta}{\sin\theta}
  \]
  となる.ここで,\,$\sin\theta\neq0$であることに注意する.特性根が実数ではなく虚部をもつ場合には,\,$\theta$は$0$や$\pi$ではないため,
  \[
    \sin\theta\neq0
  \]
  が成り立つ.

  したがって,\,$a_0,a_1$が与えられている場合には,
  \[
    C=a_0,\qquad
    D=\cfrac{\frac{a_1}{\rho}-a_0\cos\theta}{\sin\theta}
  \]
  として$C,D$を求めることができる.

  例えば,先ほどの
  \[
    a_{n+2}=2a_{n+1}-2a_n
  \]
  を考える.特性根は
  \[
    1+i
    =\sqrt{2}
    \left(
      \cos\frac{\pi}{4}
      +i\sin\frac{\pi}{4}
    \right)
  \]
  であるから,
  \[
    \rho=\sqrt{2},\qquad
    \theta=\frac{\pi}{4}
  \]
  である.したがって一般項は
  \[
    a_n=(\sqrt{2})^n
    \left(
      C\cos\frac{n\pi}{4}
      +D\sin\frac{n\pi}{4}
    \right)
  \]
  となる.

  例えば初期条件
  \[
    a_0=1,\qquad a_1=2
  \]
  が与えられているとする.まず,
  \[
    C=a_0=1
  \]
  である.また,
  \[
    2=\sqrt{2}
    \left(
      C\cos\frac{\pi}{4}
      +D\sin\frac{\pi}{4}
    \right)
  \]
  であるから,
  \[
    2=\sqrt{2}
    \left(
      \frac{1}{\sqrt{2}}
      +\frac{D}{\sqrt{2}}
    \right)
    =1+D
  \]
  より,
  \[
    D=1
  \]
  を得る.

  よって,この場合の一般項は
  \[
    a_n=(\sqrt{2})^n
    \left(
      \cos\frac{n\pi}{4}
      +\sin\frac{n\pi}{4}
    \right)
  \]
  となる.

  このように,虚数解をもつ場合でも,基本的な流れは実数解の場合と変わらない.まず特性根から$\rho,\theta$を求め,
  \[
    a_n=\rho^n(C\cos n\theta+D\sin n\theta)
  \]
  とおき,初期条件を代入して$C,D$を決定すればよい.

  \noindent\textbf{【発想のつながり：倍率をそろえると加法定理が現れる】}\par
  特性根を極形式で書いたのは,一歩ごとの倍率と回転を読み取るためだった.
  解と係数の関係から$p=-2\rho\cos\theta,q=\rho^2$なので,元の式は
  \[
    a_{n+2}-2\rho\cos\theta\,a_{n+1}+\rho^2a_n=0
  \]
  である.今度はこの係数の形から補助数列を設計してみよう.
  初項$a_1$に合わせて指数を$n-1$にし,以下では別の係数$c,d$を使う.
  係数に$\rho$と$\rho^2$があるので,添字一つにつき$\rho$倍される量を基準にする.
  \sectionref*{節：補助数列の設計}の正規化を使い,
  \[
    b_n=\frac{a_n}{\rho^{n-1}}
  \]
  と置く.元の式を$\rho^{n+1}$で割ると
  \[
    b_{n+2}-2\cos\theta\,b_{n+1}+b_n=0
  \]
  となる.大きさの倍率を除いたことで,加法定理と同じ関係が残った.
  実際,
  \[
    \cos(x+\theta)+\cos(x-\theta)=2\cos\theta\cos x
  \]
  と,余弦を正弦に替えた同じ公式に$x=n\theta$を入れると,
  \[
    u_n=\cos((n-1)\theta),\qquad
    v_n=\sin((n-1)\theta)
  \]
  はどちらも$w_{n+2}-2\cos\theta\,w_{n+1}+w_n=0$を満たすと分かる.
  \sectionref*{節：線形結合}の重ね合わせの原理から,$cu_n+dv_n$も同じ関係を満たす.
  初めの二項は$(1,\cos\theta)$と$(0,\sin\theta)$であり,$\sin\theta\ne0$なので一次独立である.
  先ほどの複素数の冪に代わり,正弦・余弦が実数の解全体の基底になる.
  そこで
  \[
    b_n=c\cos((n-1)\theta)+d\sin((n-1)\theta)
  \]
  を候補にしよう.

  \noindent\textbf{【初期条件を合わせ,すべての項で確かめる】}\par
  候補の初めの二項は$c$と$c\cos\theta+d\sin\theta$である.
  元の初期条件から$b_1=a_1,b_2=\frac{a_2}{\rho}$なので,
  \[
    c=a_1,\qquad
    d=\cfrac{\frac{a_2}{\rho}-a_1\cos\theta}{\sin\theta}
  \]
  と選べば一致する.$\sin\theta\ne0$だから,この二つの係数は必ず決められる.
  候補も元の$b_n$も,直前の二項から同じ式で次の項を作る.
  従って初めの二項が一致すれば第三項も一致し,同じことを繰り返すとすべての項で一致する.
  初期条件と加法定理だけで,この候補が一般項であることを確かめられた.
  元へ戻すと,
  \BookMathLayout{\[
    a_n=\rho^{n-1}\{c\cos((n-1)\theta)+d\sin((n-1)\theta)\}
  \]}{\[
    \begin{aligned}
      a_n=\rho^{n-1}\bigl\{&c\cos((n-1)\theta)\\
                           &+d\sin((n-1)\theta)\bigr\}
    \end{aligned}
  \]}
  を得る.倍率をそろえ,加法定理に合う量を作り,初めの二項を合わせるのが解法の流れである.

  \noindent\textbf{【振動と周期性を区別する】}\par
  $C=D=0$なら零数列であり,以下では少なくとも一方が零でないとする.
  三角関数の部分には振動する形が現れるが,整数の添字で見た数列が周期的になるとは限らない.
  $\rho^{n}$は全体の大きさの倍率なので,
  \[
    \begin{cases}
      0<\rho<1 & \text{振幅が減衰する},\\
      \rho=1 & \text{振幅が一定である},\\
      \rho>1 & \text{振幅が増大する}
    \end{cases}
  \]
  と区別できる.
  $\rho=1$で$\frac{\theta}{2\pi}$が有理数なら,正整数$T$について$T\theta$を$2\pi$の整数倍にでき,
  加法定理から$a_{n+T}=a_n$となる.
  非零の解については,逆も成り立つ.
  複素数の冪の式で周期$T$を仮定すると,
  \[
    A(\alpha^T-1)\alpha^{n}
    +\overline A(\overline\alpha^{\,T}-1)\overline\alpha^{\,n}=0
  \]
  を得る.$n=1,2$の二つの式を使うと,$\alpha\ne\overline\alpha$だから両方の係数が零になる.
  非零の実数解では$A,\overline A$も零でないので,$\alpha^T=\overline\alpha^{\,T}=1$である.
  従って周期性には$\rho=1$と$\frac{\theta}{2\pi}\in\mathbb Q$の両方を確認する.
  三角関数があるというだけで周期性を結論しないようにしよう.

  次の数列の一般項を求めてみよう.

  \noindent \bookblocklink{example-complex-1}{problem}{answer}{【例題】}
  \begin{enumerate}[label=(\arabic*),parsep=3pt,format=\bookitemlink{example-complex-1}{problem}{answer}]
    \item $\displaystyle a_1=1,\,a_2=0,\quad a_{n+2}=-a_n\quad(n\ge1)$
\item \BookMathLayout{$\displaystyle a_1=0,\,a_2=1,\quad a_{n+2}=2a_{n+1}-2a_n\quad(n\ge1)$}{$\displaystyle a_1=0,\,a_2=1$\\ $\displaystyle a_{n+2}=2a_{n+1}-2a_n$\\ $\displaystyle(n\ge1)$}
    \item $\displaystyle a_1=1,\,a_2=2,\quad a_{n+2}=-4a_n\quad(n\ge1)$
\item \BookMathLayout{$\displaystyle a_1=1,\,a_2=0,\quad a_{n+2}=2a_{n+1}-4a_n\quad(n\ge1)$}{$\displaystyle a_1=1,\,a_2=0$\\ $\displaystyle a_{n+2}=2a_{n+1}-4a_n$\\ $\displaystyle(n\ge1)$}
  \end{enumerate}

  \noindent\bookblocklink{example-complex-1}{hint}{problem}{【方針】}
  \begin{enumerate}[label=(\arabic*),parsep=3pt,format=\bookitemlink{example-complex-1}{hint}{problem}]
    \item 初項から数項を計算して周期性を見つけ,一般項を予想します.その後,数学的帰納法で証明します.
    \item 特性根を極形式で表し,$\rho=\sqrt2,\,\theta=\dfrac{\pi}{4}$を用いて一般項を置きます.
    \item 特性根は$\pm2i$です.$\rho=2,\,\theta=\dfrac{\pi}{2}$として初期条件を代入します.
    \item 特性根を極形式で表し,$\rho=2,\,\theta=\dfrac{\pi}{3}$として係数を決定します.
  \end{enumerate}

  \noindent\bookblocklink{example-complex-1}{answer}{problem}{【解答】}
  \begin{enumerate}[label=(\arabic*),parsep=3pt,format=\bookitemlink{example-complex-1}{answer}{problem}]
    \item $\displaystyle a_n=\cos\frac{(n-1)\pi}{2}$
    \item $\displaystyle a_n=(\sqrt2)^{n-1}\sin\frac{(n-1)\pi}{4}$
    \item $\displaystyle a_n=2^{n-1}\left(\cos\frac{(n-1)\pi}{2}+\sin\frac{(n-1)\pi}{2}\right)$
\item \BookMathLayout{$\displaystyle a_n=2^{n-1}\left(\cos\frac{(n-1)\pi}{3}-\frac1{\sqrt3}\sin\frac{(n-1)\pi}{3}\right)$}{$\displaystyle a_n=2^{n-1}\Bigl(\cos\frac{(n-1)\pi}{3}$\\
$\displaystyle\hspace{3zw}-\frac1{\sqrt3}\sin\frac{(n-1)\pi}{3}\Bigr)$}
  \end{enumerate}

  \noindent\bookblocklink{example-complex-1}{solution}{problem}{【解法】}
  \begin{enumerate}[label=(\arabic*),parsep=3pt,format=\bookitemlink{example-complex-1}{solution}{problem}]
    \item 初項から数項を計算すると,
          \begin{align*}
            &a_1=1,\quad a_2=0,\quad a_3=-1,\quad a_4=0,\\
            &a_5=1,\quad a_6=0,\quad a_7=-1,\quad a_8=0
          \end{align*}
          となる.この並びから,
          \[
            a_n=\cos\frac{(n-1)\pi}{2}
          \]
          と予想できる.これを数学的帰納法で証明する.
          $n=1,2$では
          \[
            a_1=1=\cos0,\qquad a_2=0=\cos\frac{\pi}{2}
          \]
          となり,主張が成り立つ.
          ある$k\ge1$で,$n=k,k+1$に対して
          \[
            a_k=\cos\frac{(k-1)\pi}{2},\qquad
            a_{k+1}=\cos\frac{k\pi}{2}
          \]
          が成り立つと仮定する.漸化式と$-\cos x=\cos(x+\pi)$より,
          \[
            \begin{aligned}
              a_{k+2}
              &=-a_k\\
              &=-\cos\frac{(k-1)\pi}{2}\\
              &=\cos\left(\frac{(k-1)\pi}{2}+\pi\right)\\
              &=\cos\frac{(k+1)\pi}{2}
            \end{aligned}
          \]
          である.したがって,$n=k+1,k+2$でも主張が成り立つ.初めの2項で成り立つことと合わせ,数学的帰納法により
          \[
            a_n=\cos\frac{(n-1)\pi}{2}
          \]
          がすべての$n\ge1$で成り立つ.
    \item 特性方程式は
          \[
            x^2-2x+2=0
          \]
          であり,特性根は$1\pm i$である.したがって
          \[
            \rho=\sqrt2,\qquad\theta=\frac{\pi}{4}
          \]
          として,
\BookMathLayout{\[
            a_n=(\sqrt2)^{n-1}
            \left(
              C\cos\frac{(n-1)\pi}{4}
              +D\sin\frac{(n-1)\pi}{4}
            \right)
          \]}{
\[
\begin{aligned}
a_n&=(\sqrt2)^{n-1}\Bigl(C\cos\frac{(n-1)\pi}{4}\\
&\hspace{4zw}+D\sin\frac{(n-1)\pi}{4}\Bigr)
\end{aligned}
\]
}
          とおける.$a_1=0$より$C=0$であり,$a_2=1$より$D=1$である.よって,
          \[
            a_n=(\sqrt2)^{n-1}\sin\frac{(n-1)\pi}{4}
          \]
          を得る.
    \item 特性方程式は$x^2+4=0$であり,特性根は$2i,-2i$である.したがって
          \[
            a_n=2^{n-1}
            \left(
              C\cos\frac{(n-1)\pi}{2}
              +D\sin\frac{(n-1)\pi}{2}
            \right)
          \]
          とおける.$a_1=1$より$C=1$であり,$a_2=2$より$2D=2$,すなわち$D=1$である.よって,
          \[
            a_n=2^{n-1}
            \left(
              \cos\frac{(n-1)\pi}{2}
              +\sin\frac{(n-1)\pi}{2}
            \right)
          \]
          を得る.
    \item 特性方程式は
          \[
            x^2-2x+4=0
          \]
          であり,特性根は$1\pm\sqrt3\,i$である.したがって
          \[
            \rho=2,\qquad\theta=\frac{\pi}{3}
          \]
          として,
          \[
            a_n=2^{n-1}
            \left(
              C\cos\frac{(n-1)\pi}{3}
              +D\sin\frac{(n-1)\pi}{3}
            \right)
          \]
          とおける.$a_1=1$より$C=1$であり,$a_2=0$より
          \[
            0=2\left(\frac12+D\frac{\sqrt3}{2}\right)
          \]
          となるので$D=-\dfrac1{\sqrt3}$である.よって,
          \[
            a_n=2^{n-1}
            \left(
              \cos\frac{(n-1)\pi}{3}
              -\frac1{\sqrt3}\sin\frac{(n-1)\pi}{3}
            \right)
          \]
          を得る.
  \end{enumerate}

\subsection{非隣接2項間漸化式型}
  よく見られる非隣接二項間型の漸化式は
  \[
    a_{n+2}=f(a_{n})
  \]
  のような見た目のものである.
  このような場合には大きく分けて2つの解放がある.
  一つは$n=2k$のときと$n=2k-1$のときで場合分けして一般項を求める手法で,
  もう一方は0倍の$a_{n+1}$が存在すると考えて,隣接三項間漸化式に持ち込む方法である.

\subsection{二重数列と二方向の階差}\label{節：二重数列と二方向の階差}
  一つの番号に沿って項を並べる代わりに,行と列の二つの番号で値を並べてみよう.
  ここでは$n,m\ge0$を整数とし,$a_{n,m}$を数表の位置$(n,m)$の値とする.
  このように二つの添字をもつものを\textbf{二重数列}という.
  新しい万能の解法ではなく,階差と帰納法を二方向で使う見方を学ぶ.

  \subsubsection{距離を二乗し,二方向へ差を取る}
  原点から点$(n,m)$までの距離は
  \[ a_{n,m}=\sqrt{n^2+m^2} \]
  である.
  根号を取り除くため$b_{n,m}=a_{n,m}^2$とすると,
  一方向の差は$b_{n+1,m}-b_{n,m}=2n+1$となり,$m$によらない.
  従って,さらにもう一方向へ差を取れば零になる.

  この計算を広げ,境界の値
  \[
    b_{n,0}=n^2,\qquad b_{0,m}=m^2
  \]
  と,一定の実数$c$について
  \BookMathLayout{\[
    b_{n+1,m+1}-b_{n+1,m}-b_{n,m+1}+b_{n,m}=c
  \]}{\[
    \begin{aligned}
      b_{n+1,m+1}-b_{n+1,m}\\
      {}-b_{n,m+1}+b_{n,m}=c
    \end{aligned}
  \]}
  が与えられているとする.内側の値はこの式と境界から順に決まる.
  右辺を$b_{n+1,m+1}$について解けば,ほかの三つの位置は添字の和が小さい.
  $n+m$についての帰納法により一意性を確かめられる.

  横方向の差$d_{n,m}=b_{n+1,m}-b_{n,m}$を追うと,
  \[
    d_{n,m+1}-d_{n,m}=c,\qquad d_{n,0}=2n+1
  \]
  なので$d_{n,m}=2n+1+cm$である.これを横に足し戻せば,
  \[
    b_{n,m}=n^2+m^2+cnm
  \]
  となる.\textbf{二方向に一つずつ階差を取り,逆順に足し戻した}のである.
  $c$が一定という条件が,この短い計算を可能にしている.

  \subsubsection{図形と数え上げに戻る}
  二つの単位ベクトル$\boldsymbol u,\boldsymbol v$のなす角を$\theta$とすると,
  $c=2\cos\theta$の場合には
  \[
    b_{n,m}=|n\boldsymbol u+m\boldsymbol v|^2
  \]
  と読める.この図形的な解釈には$-2\le c\le2$が必要である.
  単に数表を作る式としてなら$c$は任意の実数でよいが,
  その値を距離の二乗と読むためには非負性も確認する.

  一方,格子点$(0,0)$から$(n,m)$へ,右と上へ一歩ずつ進む道の数を$A_{n,m}$とする.
  最後の一歩を分ければ,
  \[
    A_{n,m}=A_{n-1,m}+A_{n,m-1}\quad(n,m\ge1)
  \]
  であり,境界では$A_{n,0}=A_{0,m}=1$である.
  \BookFigLatticeRecurrence
  下の数表は,上と左の値を足して作ったものである.
  \[
    \begin{array}{c|rrrr}
      m\backslash n&0&1&2&3\\\hline
      0&1&1&1&1\\
      1&1&2&3&4\\
      2&1&3&6&10\\
      3&1&4&10&20
    \end{array}
  \]
  全$n+m$歩から上へ進む$m$歩の位置を選ぶので,
  \[
    A_{n,m}=\binom{n+m}{m}
  \]
  である.最後の一歩で分ける説明は,二項係数の加法公式の説明でもある.
  この公式と境界を使った帰納法による確認は,
  \bookproblemref{practice-connections}{3}{接続演習問3}で行おう.

  \begin{bookpoint}
    \textbf{添字を増やしても,確認することは同じ}\par
    どこが初期条件に当たる境界か,どの値から次が決まるかを先に確かめる.
    二方向の関係を別々に与える場合は,同じ位置を二通りに計算して矛盾しないことも必要である.
  \end{bookpoint}
  格子経路と二項係数の資料は参考文献のNIST DLMF §26.3へ.


% section: 特殊性質型

\section{特殊性質型}\label{節：特殊性質型}

1つ面白い話をしよう.
まず,次の3つの漸化式を見てほしい.
\[
  a_{n+1}=a_n^2,
  \quad
  a_{n+1}=a_n^2-1,
  \quad
  a_{n+1}=a_n^2-2
\]
1つ目はきちんとここまで本書を読んできた諸君らならば,容易に解くことができるだろう.
実際,繰り返し代入すれば
\[
  a_n=a_1^{2^{n-1}}
\]
と分かる.\,$a_1>0$なら,両辺の対数を取って等比数列に帰着させることもできる.
一方,\,$a_1=0$や$a_1<0$では実数の対数を最初から取ることはできないが,元の漸化式そのものは定義されている.

2つ目はどうだろうか.\,$a_n$に指数がついているから,対数を使って処理したいが,\,$-1$の項があるためそのままではうまくいかない.
ただし,不動点が存在しないわけではない.
不動点$\alpha$は
\[
  \alpha=\alpha^2-1
  \quad\Longleftrightarrow\quad
  \alpha=\frac{1\pm\sqrt5}{2}
\]
である.
したがって,初項をどちらかの不動点に選べば,\,$a_n=\alpha$という定数解が得られる.
これは,後で扱う\sectionref{節：初項依存型}にもつながる考え方である.

しかし,不動点が見つかることと,一般の初項に対して漸化式がこれまでの型へ帰着できることは別である.
実際,\,$b_n=a_n-\alpha$と置くと,不動点の関係$\alpha^2-\alpha-1=0$より
\[
  b_{n+1}=b_n^2+2\alpha b_n
\]
となり,二次の項が残る.
つまり,不動点を引くだけでは等比型などの線形な漸化式にはならない.
少々意地悪であったかもしれないが,これは(\sectionref{節：漸化式の分類}で述べた意味において)本書でここまで扱った手法だけから一般の初項に対する閉じた一般項を得ることが難しい漸化式である.

3つ目も一見先の例のように一般項が得難い見た目をしている.
しかしながらこれはある特殊な条件の噛み合いによって偶然にも初等的な解法が存在する.

ここまで,等差型,等比型,階差型,階比型など,漸化式を一定の型に
帰着して解く方法を学んできた.しかし,右辺が少し複雑になると,
同じ道具が急に働かなくなる.

ここで扱うのは,一般の漸化式がいつでも解けるという話ではない.
一見すると解けそうにない漸化式でも,既知の数学的対象と一致する構造を持つことがある.
その一致を見抜くことで,初等的な型に還元できない漸化式を解いていく.
本章はそういう話である.
\subsection{基本対称式型}\label{節：基本対称式型}
導入に続いて,先程の\bookterm{recurrence}{漸化式}を考えていこう.
\[
  a_1=3, \quad a_{n+1}=a_n^2-2
\]
この\bookterm{recurrence}{漸化式}を解くに当たって邪魔をしているのは$-2$の項である.
これをうまく消せるような操作を考えたい.
また,\,2乗があることからうまく因数分解できないか考えてみよう.

この$-2$を打ち消す$2$が生まれるような対称式を考える.

2乗の項は正だから,
\[
  (a+b)^2=a^2+2ab+b^2
\]
で言うところの$2ab$が$2$になってくれるように調整すればうまくいきそうだ.
そのようなふうに考えていくと,
\[
  \left(\alpha+\frac{1}{\alpha}\right)^2=\alpha^2+2+\frac{1}{\alpha^2}
\]
という式が思い当たる.
実際,ここに$-2$をすると
\[
  \alpha^2+\frac{1}{\alpha^2}
\]
が残るから,成り立ちそうなことが分かる.
この実験を元に解答に移る.

\noindent \bookblocklink{example-symmetric-1}{problem}{answer}{【問題】}\par\nobreak
\begin{smallquote}
  次の\bookterm{recurrence}{漸化式}で定義される\bookterm{sequence}{数列}の\bookterm{general}{一般項}を求めよ.
  \[
    a_1=3, \quad a_{n+1}=a_n^2-2
  \]
\end{smallquote}
\par\noindent\bookblocklink{example-symmetric-1}{answer}{problem}{【解答】}\par\nobreak
  \begin{smallquote}
  \bookterm{general}{一般項}が
    \[
      b_n=\alpha^{2^{n-1}}+\alpha^{-2^{n-1}}
    \]
    で表される\bookterm{sequence}{数列}$b_n$を考える.
    \begin{align*}
      b_{n+1}&=\alpha^{2^{n}}+\alpha^{-2^{n}}\\
      &=\left(\alpha^{2^{n-1}}\right)^2+\left(\alpha^{-2^{n-1}}\right)^2\\
      &=\left(\alpha^{2^{n-1}}+\alpha^{-2^{n-1}}\right)^2-2\\
      &=b_n^2-2
    \end{align*}
    これは与式の\bookterm{recurrence}{漸化式}と一致する.また,
    \[
      b_1=a_1=\alpha+\alpha^{-1}
    \]
    となる$\alpha$を考えると,
    \begin{align*}
      &\alpha+\alpha^{-1}=3\\
      &\alpha^2-3\alpha+1=0\\
      &\therefore \quad \alpha=\frac{3\pm \sqrt{5}}{2}
    \end{align*}
    得られた$\alpha$の2解は逆数の関係にある.
    以上より,
    \[
      a_n=\left(\frac{3+ \sqrt{5}}{2}\right)^{2^{n-1}}+\left(\frac{3- \sqrt{5}}{2}\right)^{2^{n-1}}
    \]
  \end{smallquote}

ここで,定数を $-2$ から $-1$ に変えてみよう.
\[
  a_{n+1}=a_n^2-1
\]
同じ置換を試すと,
\begin{align*}
  a_n^2-1
    &=\left(u+u^{-1}\right)^2-1\\
    &=u^2+u^{-2}+1,
\end{align*}
となる.しかし,次の項として欲しい形は $u^2+u^{-2}$ である.
余分な $1$ が残るため,同じ対称式の族の中に戻ってこない.

\bookterm{recurrence}{漸化式}の $-2$ と $-1$ は,見た目には小さな違いである.
しかし,対称式を保つ恒等式にとっては,その違いが決定的になる.
「ほとんど同じ式だから,ほとんど同じ方法で解ける」とは限らない.
解法が成立するかどうかは,係数の大きさではなく,恒等式が完全に一致
するかどうかで決まる.

同じ考え方は,三次式にも現れる.
\[
  \left(\alpha+\frac{1}{\alpha}\right)^3
  -3\left(\alpha+\frac{1}{\alpha}\right)
  =\alpha^3+\frac{1}{\alpha^3}
\]
このように,対称な組み合わせを選ぶと,べき乗の複雑な展開が
再び同じ種類の式へ戻る.ただし,恒等式が少しでもずれると,その閉じた
構造は失われる.これが基本対称式型の面白さである.

三乗の恒等式に合う\bookterm{recurrence}{漸化式}も作れる。

\noindent\bookblocklink{example-symmetric-2}{problem}{answer}{【問題】}\par\nobreak
\begin{smallquote}
  $\displaystyle a_1=\frac{8}{3},\quad a_{n+1}=a_n^3+3a_n$
\end{smallquote}
\par\noindent\bookblocklink{example-symmetric-2}{answer}{problem}{【解答】}\par\nobreak
\begin{smallquote}
  \bookterm{general}{一般項}が
    \[
      b_n=\alpha^{3^{n-1}}-\alpha^{-3^{n-1}}
    \]
    で表される\bookterm{sequence}{数列}$b_n$を考える.
    \begin{align*}
      b_{n+1}&=\alpha^{3^{n}}-\alpha^{-3^{n}}\\
      &=\left(\alpha^{3^{n-1}}\right)^3-\left(\alpha^{-3^{n-1}}\right)^3\\
      &=\left(\alpha^{3^{n-1}}-\alpha^{-3^{n-1}}\right)^3\\
      &\qquad +3\left(\alpha^{3^{n-1}}-\alpha^{-3^{n-1}}\right)\\
      &=b_n^3+3b_n
    \end{align*}
    これは与式の\bookterm{recurrence}{漸化式}と一致する.また,
    \[
      b_1=a_1=\alpha-\alpha^{-1}=\frac{8}{3}
    \]
    となる$\alpha$を考えると,
    \begin{align*}
      &\alpha-\alpha^{-1}=\frac{8}{3}\\
      &3\alpha^2-8\alpha-3=0\\
      &\therefore \quad \alpha=3,\,-\frac13
    \end{align*}
    どちらを選んでも結果は変わらないので,\,$\alpha=3$とすると,
    \[
      a_n=3^{3^{n-1}}-3^{-3^{n-1}}
    \]
\end{smallquote}

\subsection{三角関数型}\label{節：三角関数型}
  \subsubsection{解法}\label{節：解法}
    またもや似た見た目の漸化式を考える.
    \[
      a_1=\frac{\sqrt{3}}{2},\quad a_{n+1}=2a_n^2-1
    \]
    
    この漸化式も前節と同様に,今までの方法ではうまく解くことができない.
    
    今回の鍵は三角関数である.
    三角関数と置くことで,数列の漸化式が角度の漸化式へ変わることがある.
    \[
      a_1=\frac{\sqrt{3}}{2},\quad a_{n+1}=2a_n^2-1
    \]
    右辺は $\cos 2x=2\cos^2x-1$ と一致している.さらに,
    初項は
    \[
      a_1=\cos\frac{\pi}{6}
    \]
    と書ける.そこで $a_n=\cos\theta_n$ と置くと,
    \[
      \theta_1=\frac{\pi}{6},
      \quad
      \theta_{n+1}=2\theta_n
    \]
    と選べばよい.ただし,$\cos$の値から角度は一意に決まらないため,
    $\theta_{n+1}=2\theta_n$は元の漸化式から必ず従う式ではなく,
    その漸化式を満たす数列を作るための角度の選び方である.従って候補は
    \[
      a_n=\cos\left(2^{n-1}\frac{\pi}{6}\right)
    \]
    である.
    
    この候補を帰納法で確認する.\,$n=1$ では初項に一致する.ある $n$ で
    \[
      a_n=\cos\left(2^{n-1}\frac{\pi}{6}\right)
    \]
    が成り立つとすると,
    \begin{align*}
      a_{n+1}
        &=2a_n^2-1\\
        &=2\cos^2\left(2^{n-1}\frac{\pi}{6}\right)-1\\
        &=\cos\left(2^n\frac{\pi}{6}\right)
    \end{align*}
    これは $n+1$ の式である.よって,
    \[
      a_n=\cos\left(2^{n-1}\frac{\pi}{6}\right)
    \]
    が示された.
    
    この解法の本質は,二次式を無理に展開したり因数分解したりすることではない.
    数列の値を角度の関数として表すことで,二次式の計算を角度の2倍という単純な操作に翻訳しているのである.
  
  \subsubsection{倍角を回転として見る}\label{節：倍角を回転として見る}
    この「角度の2倍」は,単位円上の点の動きとしても捉えられる.
    原点を$O$とし,正の$x$軸からの角が$\theta$である点を
    \[
      P(\theta)=(\cos\theta,\sin\theta)
    \]
    とする.点$P(\theta)$を原点のまわりにさらに$\theta$だけ回転すると,
    正の$x$軸からの角は$\theta+\theta=2\theta$となり,$P(2\theta)$へ移る.
    つまり倍角は,\textbf{現在の角度と同じ角度だけ,さらに回す操作}と考えられる.
    
    \begin{center}
      \begin{tikzpicture}[scale=2]
        \draw[->] (-1.15,0) -- (1.3,0) node[right] {$x$};
        \draw[->] (0,-1.15) -- (0,1.3) node[above] {$y$};
        \draw (0,0) circle (1);
        \node[below left] at (0,0) {$O$};
        \draw (0,0) -- (35:1);
        \draw (0,0) -- (70:1);
        \fill (35:1) circle (0.025) node[right] {$P(\theta)$};
        \fill (70:1) circle (0.025) node[above right] {$P(2\theta)$};
        \draw[->] (0.38,0) arc (0:35:0.38);
        \node at (17.5:0.55) {$\theta$};
        \draw[->] (35:0.72) arc (35:70:0.72);
        \node at (52.5:0.88) {$\theta$};
      \end{tikzpicture}
    \end{center}
    
    先ほどの数列では,$a_n$は点$P(\theta_n)$の横座標である.
    回転後の横座標が
    \[
      \cos(2\theta_n)=2\cos^2\theta_n-1=2a_n^2-1
    \]
    となるので,円周上の点の動きを横座標だけで記録したものが,
    漸化式$a_{n+1}=2a_n^2-1$なのである.
    例えば,初めの角度は
    \[
      \frac{\pi}{6}\ \longrightarrow\ \frac{\pi}{3}
      \ \longrightarrow\ \frac{2\pi}{3}
      \ \longrightarrow\ \frac{4\pi}{3}
    \]
    と変わり,横座標は
    \[
      \frac{\sqrt{3}}{2}\ \longrightarrow\ \frac12
      \ \longrightarrow\ -\frac12
      \ \longrightarrow\ -\frac12
    \]
    と変わる.最後の二点は円周上では異なる点だが,横座標は同じである.
    このように,数列の値だけでは見えない点の動きが,角度を使うと見える.
    
    ただし,毎回一定の角$\phi$だけ回す操作では
    $\theta_{n+1}=\theta_n+\phi$となるのに対し,
    ここでは$\theta_{n+1}=2\theta_n$となり,回す角度もそのつど変わる.
    円周上の角度は$2\pi$の整数倍だけ違っても同じ点を表すので,
    点の動きを考えるときは一周ごとに角度を読み替えてよい.
    また,$P(\theta)$と$P(\theta+\pi)$は倍角によって同じ点へ移る.
    したがって,円周全体への写像としての倍角は,一つの一定角の回転とは異なる.
    
    \,\sectionref{節：極形式と複素平面}の見方を使えば,この操作はさらに短く書ける.
    \[
      z_n=\cos\theta_n+i\sin\theta_n
    \]
    と置くと,ド・モアブルの定理より
    \[
      z_{n+1}=z_n^2,\qquad z_n=z_1^{\,2^{n-1}}
    \]
    となる.$|z_n|=1$なので,$z_n$を掛けることは角$\theta_n$の回転である.
    $z_n^2=z_n\cdot z_n$は,現在の点にその点自身を掛けることで,
    角度をもう一度加えている.$a_n$はその実部に当たる.
    
    \subsubsection{正接では直線の傾きを記録する}\label{節：正接では直線の傾きを記録する}
    今度は
    \[
      a_1=\tan\frac{\pi}{12}=2-\sqrt{3},
      \quad
      a_{n+1}=\frac{2a_n}{1-a_n^2}
    \]
    を考える.右辺は$\tan$の倍角公式そのものである.したがって,
    \[
      a_n=\tan\left(2^{n-1}\frac{\pi}{12}\right)
    \]
    と予想できる.
    
    実際,\,$n=1$ で初項に一致し,帰納法の段階では
    \begin{align*}
      a_{n+1}
        &=\frac{2a_n}{1-a_n^2}\\
        &=\cfrac{2\tan\left(2^{n-1}\cfrac{\pi}{12}\right)}
                {1-\tan^2\left(2^{n-1}\cfrac{\pi}{12}\right)}\\
        &=\tan\left(2^n\frac{\pi}{12}\right)
    \end{align*}
    となる.ただし,分母が $0$ になると$\tan$は定義できない.
    三角関数型では,置換を思いつくだけでなく,途中で定義できなくなる項がないかも確認しなければならない.
    
    この場合,$a_n=\tan\theta_n$は,原点と$P(\theta_n)$を通る直線の傾きである.
    直線をさらに$\theta_n$だけ回すと,その傾きが$\tan(2\theta_n)$となる.
    従って,今度は同じ倍角の操作を横座標ではなく傾きとして記録している.
    直線の向きは$\pi$の整数倍だけ違っても同じであり,
    縦向きになると傾きは定義できない.
    $a_n=\pm1$のときは$\cos(2\theta_n)=0$となり,
    この縦向きの状態が漸化式の分母$1-a_n^2=0$に対応する.
    なお,この例の角度は,$n\geq2$では
    \[
      \theta_n=\frac{2^{n-2}\pi}{6}
    \]
    である.$\theta_n=\pi/2+k\pi$となるには$2^{n-2}=3+6k$が必要だが,
    2の累乗は3の倍数ではないので,そのような整数$k$は存在しない.
    初項の角度$\pi/12$でも傾きは定義できるため,この数列はすべての項で定義できる.
    
  \subsubsection{整数倍の公式へ}\label{節：整数倍の公式へ}
    $\cos$の例と$\tan$の例は,同じ方針を共有している.すなわち,
    \[
      a_n=\varphi(\theta_n),
      \quad
      \varphi(m\theta)=F_m(\varphi(\theta))
    \]
    となる関数 $\varphi$ と整数倍公式を探し,角度の漸化式を解いている.
    いわば媒介変数表示や置換のようなことをしている.
    
    また,ここで現れた多項式
    \[
      2x^2-1,
      \quad
      4x^3-3x
    \]
    は偶然の産物ではない.チェビシェフ多項式 $T_m(x)$ を
    \[
      T_m(\cos\theta)=\cos(m\theta)
    \]
    で定義すれば,
    \[
      T_0(x)=1,  T_1(x)=x,
    \]
    \[
      T_{m+1}(x)=2xT_m(x)-T_{m-1}(x)
    \]
    となる.例えば $T_2(x)=2x^2-1$ である.
    したがって,$\cos$の例は
    \[
      a_{n+1}=T_2(a_n),
      \quad
      a_n=T_{2^{n-1}}(a_1)
    \]
    と書き直せる.三角関数の公式から始めて多項式の族へ進む道が,
    付録のチェビシェフ多項式へつながっている.

\subsection{ニュートン法型}\label{節：ニュートン法型}

\noindent\textbf{【前提と学ぶこと】}\par
微分係数$f'(a)$が接線の傾きを表すことを使う.
微分が未習なら,接線と$x$軸の交点から次の近似値を作る図形的な発想と,得られた漸化式の変形を先に読んでよい.
閉じた一般項を求める計算と,近似がどれほど速く改善するかを調べることは別の目標である.
後者は\sectionref*{節：差分から始まる離散的な微積分}で見直す.

少々離れたところから話を始める.
ある関数$f(x)$について,\,$f(a)$が既知の時,その近傍の点$x=a+h$で
\[
  f(a+h) \approx f(a)+f'(a)h
\]
が成り立つ.
つまり,\,$f(a)$の値を元に,\,$f(a+h)$を近似したのである.
方程式
\[
  f(x)=x^2-2=0
\]
の正の解 $\sqrt{2}$ を求めたいとする.いまの近似値を $x=a_n$ とし,
点 $(a_n,f(a_n))$ における接線を引く.その接線の傾きは
$f'(a_n)=2a_n$ だから,接線の方程式は
\[
  y=f(a_n)+2a_n(x-a_n)
\]
である.\,$y=0$ と置いて $x$ 軸との交点を求めると,
\begin{align*}
  a_{n+1}
    &=a_n-\frac{f(a_n)}{f'(a_n)}\\
    &=a_n-\frac{a_n^2-2}{2a_n}\\
    &=\frac{a_n^2+2}{2a_n}
\end{align*}
これが
\[
  a_{n+1}=\frac{a_n^2+2}{2a_n}
\]
という,ニュートン法を$f(x)=x^2-2$に適用したときの漸化式である.

例えば $a_1=2$ とすると,
\[
  a_2=\frac{3}{2},
  \quad
  a_3=\frac{17}{12},
  \quad
  a_4=\frac{577}{408}
\]
となり,これらは $\sqrt{2}$ に急速に近づく.この速さは,誤差を
$e_n=a_n-\sqrt{2}$ と置くと見える.
\begin{align*}
  e_{n+1}
    &=\frac{a_n^2+2}{2a_n}-\sqrt{2}\\
    &=\frac{(a_n-\sqrt{2})^2}{2a_n}\\
    &=\frac{e_n^2}{2a_n}
\end{align*}
誤差がほぼ二乗されるため,正しい近似に近づくと桁数が一気に増える.

実は,この漸化式は閉じた式で表すことができる.
\[
  a_{n+1}-\sqrt{2}=\frac{(a_n-\sqrt{2})^2}{2a_n},
  \qquad
  a_{n+1}+\sqrt{2}=\frac{(a_n+\sqrt{2})^2}{2a_n}
\]
であることに注目すると,両辺の比を取ることで$2a_n$が消え,
\[
  \frac{a_{n+1}-\sqrt{2}}{a_{n+1}+\sqrt{2}}
  =\left(\frac{a_n-\sqrt{2}}{a_n+\sqrt{2}}\right)^{2}
\]
という関係が得られる.つまり
\[
  b_n=\frac{a_n-\sqrt{2}}{a_n+\sqrt{2}}
\]
と置くと,\,$b_{n+1}=b_n^{2}$という,対数型(\sectionref{節：対数型})と同じ構造の漸化式になる.先ほど見た誤差の2乗のふるまいは,この$b_{n+1}=b_n^2$そのものだったのである.

\,$a_1=2$のとき,\,$b_1=\dfrac{2-\sqrt2}{2+\sqrt2}=3-2\sqrt2=\dfrac1{3+2\sqrt2}$であるから,\,$u=3+2\sqrt2$とおくと$b_n=u^{-2^{n-1}}$である.これを$b_n=\dfrac{a_n-\sqrt2}{a_n+\sqrt2}$について$a_n$で解けば,
\[
  a_n=\sqrt{2}\,\frac{u^{2^{n-1}}+1}{u^{2^{n-1}}-1},\qquad u=3+2\sqrt2
\]
という閉じた式が得られる.

このように,今回の$f(x)=x^2-2$に対するニュートン法型漸化式では,特別な変形によって一般項を閉じた式で表すことができる.
しかし,一般のニュートン法型漸化式が同じように閉じた式で表せるとは限らない.
閉じた式が得られなくても,収束先や誤差の減り方を調べることには十分な価値がある.
この解法の主眼は閉じた式を得ることだけにあるのではない.
「接線を使って,次の近似値をどう設計するか」を考え,その設計が
どのような速さで誤差を減らすかを調べることこそが本質である.付録では,この考えを
漸化式の収束,固定点,誤差評価,微分方程式の数値解法へ広げる.

\subsection{初項依存型}\label{節：初項依存型}
  \BookFigCobweb

  図では$f(x)=\displaystyle\frac{x^2-1}{2}$として,
  横軸上の$a_1=2$から曲線$y=f(x)$へ縦に進み,
  直線$y=x$へ横に進む操作を繰り返している.
  縦の移動は$a_{n+1}=f(a_n)$という計算そのものであり,
  横の移動で得られた値を次の入力として読む.
  最初の三つの入力は$a_1=2,\,a_2=\displaystyle\frac32,\,
  a_3=\displaystyle\frac58$に対応する.
  初項を変えれば出発点が変わり,同じ更新規則でも別の軌道になる.



最後に,\bookterm{recurrence}{漸化式}の形だけでなく,初項の選び方そのものが問題を左右する
例を見よう.
\[
  a_1=c,
  \quad
  a_{n+1}=\frac{a_n^2-1}{2}
\]
初項を記号 $c$ のままにすると,
\begin{align*}
  a_2&=\frac{c^2-1}{2},\\
  a_3&=\frac{c^4-2c^2-3}{8},\\
  a_4&=\frac{c^8-4c^6-2c^4+12c^2-55}{128}
\end{align*}
一般に $a_n$ は $c$ の次数 $2^{n-1}$ の多項式になる.初項を一般の
文字のまま扱おうとすると,項を一つ進めるたびに次数が倍になり,
\bookterm{trigonometric}{三角関数}型や基本対称式型のような固定された表現が見えてこない.

特に $c=2$ と固定すると,
\[
  a_1=2,
  \quad
  a_2=\frac{3}{2},
  \quad
  a_3=\frac{5}{8},
  \quad
  a_4=-\frac{39}{128}
\]
となる.しかし,この列は直前の例のように角度を2倍する形にも,
指数の差を保つ形にも,接線近似の形にもならない.

このように$c=2$では,ここまでに試した変換から直ちに簡単な表現を見出せない.ただし,数項の計算だけで,他の初項についても簡単な表現が存在しないとは結論できない.
一方で,初項をうまく選べば特別な構造が現れることもある.
例えば定\bookterm{sequence}{数列}を探すだけなら,
\[
  c=1+\sqrt{2},
  \quad
  c=1-\sqrt{2}
\]
では $a_{n+1}=a_n=c$ となる.大切なのは,ある初項を選んだときだけ
特別な恒等式や置換が働くことがあり,初項を一般化するとその偶然の
構造が失われる,という事実である.

初項によって解法の見通しが変わる例は,先に挙げた基本対称式型にもある.
もちろん,一般の初項での解法は先に示したとおりであるが,より狭義の解き方ではこのようなものがある.
\[
  a_1=-1,\quad a_{n+1}=a_n^2-2
\]
のような場合であれば,第5項までが以下のようになる.
\[
  -1,-1,-1,-1,-1
\]
\bookterm{recurrence}{漸化式}が表す\bookterm{sequence}{数列}が定数\bookterm{sequence}{数列}を成したのだ.
このような解はシンプルに\bookterm{fixed}{不動点}を求めることで作り出せる.
つまりは
\[
  \alpha=\alpha^2-2
\]
を$\alpha$について解けばよい.
この二次方程式の解は$\alpha=-1,2$である.
\,$\alpha=2$を初項に選ぶと,\bookterm{recurrence}{漸化式}から常に$a_n=2$となる.

実は,この定\bookterm{sequence}{数列}$a_n=-1$は,\,\sectionref{節：基本対称式型}で扱った$a_{n+1}=a_n^2-2$型の特別な場合と見ることができる.
\,$\alpha+\alpha^{-1}=a_1$を満たす$\alpha$が単位円上,すなわち$\alpha=e^{i\theta}$となるように($-2\le a_1\le2$のように)初項を選べば,
\[
  a_n=\alpha^{2^{n-1}}+\alpha^{-2^{n-1}}=2\cos(2^{n-1}\theta)
\]
という\bookterm{trigonometric}{三角関数}の形になる.今回の$a_1=-1$はまさにこの場合で,\,$2\cos\theta=-1$すなわち$\theta=\dfrac{2\pi}3$ととれる.このとき$2^{n-1}\theta$は$2$倍するごとに$\dfrac{2\pi}3$と$\dfrac{4\pi}3$を行き来し(どちらも余弦は$-\dfrac12$),\,$a_n=2\cos(2^{n-1}\theta)=-1$が常に成り立つ.つまり$a_1=-1$は,角度を$2$倍しても余弦の値が変わらないという特別な初項を選んだ場合に他ならない.

この例では,少なくとも本節で扱った初等的な道具からは\bookterm{general}{一般項}を得る
決まった方法が出てこない.それでも\bookterm{recurrence}{漸化式}は\bookterm{sequence}{数列}を定めており,
\[
  a_n=f^{\,\circ(n-1)}(c),
  \quad
  f(x)=\frac{x^2-1}{2}
\]
と書くことはできる.これは解答の終わりではなく,「\bookterm{general}{一般項}を求める」
という目標自体を見直す入口である.付録では,閉じた式が得られない
場合にも,軌道,\bookterm{limit}{極限},近似,\bookterm{matrix}{行列},\bookterm{generating}{母関数}など別の言葉で\bookterm{sequence}{数列}を理解する
方法を考える.


\end{bookcolumns}

% chapter: 実際の解き方

\chapter{実際の解き方}\label{章：実際の解き方}
\begin{bookcolumns}

  差を取る、倍率をそろえる、複数の項をまとめる。
  これまでの操作を、式の特徴と結び付けて見直す。
  一つの問題に複数の方法が使えるときは、その条件と残る計算を比べる。
% section: 解法の決定

\section{解法の決定}\label{節：解法の決定}
  以下にここまでで扱った解法を体系的にまとめあげ,実際に問題と相対した時にどこに注目すればよいか,なにを考えればよいかをまとめた.
  \chapterref{章：漸化式の解説}で目指してきたように,式の特徴を手掛かりとして,既知の「きれい」な形に帰着させることを考えよう.

  チャートは「式全体」「数列の項」「足されている数」の3つの観点から,着目点と試せる変形を整理したものである.
  すべてを順番に確認する必要はなく,式の中で目立つ特徴から必要な項目を使えばよい.
  一つの式に複数の特徴が当てはまることもある.どの特徴を使い,どの順番で変形するかは,この後の問題演習で練習していこう.
  変形した後にも,新しい式について同じ着目点を確認するとよい.

  \noindent\textbf{【観察から変形まで】}\par
  次の4段階を,必要に応じて繰り返す.チャートの出口は候補であり,一つの正解に固定するものではない.
  \begin{enumerate}
    \item \textbf{観察}：総和,分母,倍率,付加項のうち,何が目立つかを具体的な式で指す.
    \item \textbf{候補}：その特徴が変形後に消えるか,一定になるかを考える.複数の候補があってよい.
    \item \textbf{条件}：割る量の非零,対数の真数の正値,初項による定数解の分岐を確かめる.
    \item \textbf{変形と評価}：新しい量の更新を一行計算し,何が簡単になったかを見る.新しい式を観察し直し,必要なら別の候補へ戻る.
  \end{enumerate}
  候補が正しくても計算が長くなる場合はある.条件の見やすさと計算量も比較しよう.
  一般項が得られたら,初期条件と元の更新式を確認する.

  \begingroup
  \clubpenalty=10000
  \widowpenalty=10000
  \medskip
  \noindent \textbf{【チャート1：式全体について】}
  \begin{enumerate}
    \item \textbf{総和$S_n$があるか}

      $S_n=\sum_{k=1}^{n}a_k$なら,\,$S_{n+1}-S_n=a_{n+1}$である.
      添字を一つずらした式との差を取り,総和を消去できないか考える（総和型）.

    \item \textbf{数列の項を含む式が分母にあるか}

      逆数を取って簡単にならないか考える.
      定数係数の1次式同士の商なら,定数を引いて逆数型へ帰着させる方法もある（分数型）.
      分母が$0$にならないことを確認する.

    \item \textbf{何項間か.添字は隣接しているか}

      例えば$a_{n+2}$と$a_n$の関係なら,奇数番目と偶数番目を別々の数列として考えられる（非隣接2項間漸化式型）.
      また,\,$a_{n+2}=qa_n$は
      \[
        a_{n+2}=0\cdot a_{n+1}+qa_n
      \]
      と書ける.定数係数で,各項が1次の和として現れるなら,欠けた項を係数$0$で補い,隣接3項間・4項間の解法を検討する.

    \item \textbf{展開・因数分解・通分で簡単になるか}

      見た目が複雑でも,整理すると基本的な形が現れることがある.
      分数を約分する際には,消す因子が$0$になる場合を別に確認する.

    \item \textbf{同じ式のまとまりが繰り返し現れるか}

      例えば$a_n+n$と$a_{n+1}+(n+1)$が現れるなら,\,$b_n=a_n+n$と置いて関係を調べる.

    \item \textbf{既知の公式に似ているか}

      $a_n^2-2$のような形なら,積が$1$となる二つの数の和を使う置換を考える（基本対称式型）.
      $2a_n^2-1$や$\displaystyle \frac{2a_n}{1-a_n^2}$は,三角関数の倍角公式を思い出す手掛かりになる（三角関数型）.
      $\displaystyle \frac12\left(a_n+\frac{c}{a_n}\right)$のような平均の形なら,ニュートン法型を検討する.
      見た目に加えて,初項や置換の条件も確認する.

    \begin{samepage}
    \item \textbf{複数の数列の関係が与えられているか}\\*[3pt]
      定数係数の線形連立漸化式なら,数列の和・差などを使って一つずつの関係に分けられないか考える（連立型）.
    \par
    \end{samepage}

    \item \textbf{整数部分を表す記号があるか}

      まず,整数部分の記号が項の値・添字・付加項のどこに作用しているかを見る（ガウス型）.
      項の値なら整数性や余り,添字なら偶奇や群への分割,付加項なら同じ値を取る範囲が手掛かりになる.
      床関数を外せるかを調べるときは,整数性と上下の不等式を確かめる.
  \end{enumerate}

  \medskip
  \noindent \textbf{【チャート2：数列の項について】}

  係数については,主に$a_{n+1}=p_na_n+f(n)$の形に整理できる場合を考える.
  \begin{enumerate}
    \item \textbf{係数は定数か,\,$n$に依存するか}

      \begin{itemize}[beginpenalty=10000]
        \item 定数の場合は,等比型や特性方程式型などを考える.係数が$1$なら,階差に注目する.
        \item $n$に依存する場合は,その係数の積が求めやすいか調べる（階比型）.
              足されている数もある場合は,適切な数列で割って係数をそろえられないか考える.
              例えば係数が$\displaystyle \frac{n+1}{n}$なら,\,$b_n=\displaystyle \frac{a_n}{n}$という置換が候補になる.
      \end{itemize}
      係数だけで判断せず,足されている数についてもチャート3を確認する.

    \item \textbf{累乗の指数はどうなっているか}

      \begin{itemize}[beginpenalty=10000]
        \item すべて1次か.2次以上なら,因数分解や$a_n^2$などのまとまりを使う置換を考える.
        \item 指数は定数か,\,$n$に依存するか.対数を取ると,指数を係数として取り出せる場合がある（対数型）.
      \end{itemize}
      実数の対数を取る各量は,正であることを確認する.

    \item \textbf{数列の項同士の積や商があるか}

      $a_na_{n+1}$などの積があるなら,因数分解や置換でまとまりを作れないか考える.
      積・商・累乗を対数で和・差・定数倍に変えられるかも調べる.
      ただし,和の対数を各項の対数の和に分けることはできない.
  \end{enumerate}

  \medskip
  \noindent \textbf{【チャート3：足されている数について】}

  ここでは,数列の項を含まずに足されている部分を「付加項」と呼ぶ.
  以下では$a_{n+1}=p_na_n+f(n)$の形を考え,\,$f(n)$を調べよう.
  \begin{enumerate}
    \item \textbf{付加項はないか,定数か}

      付加項がなければ,係数の積から求める方法を考える（等比型・階比型）.
      付加項も係数も定数なら,定数を引いて消せないか考える（特性方程式型）.
      ただし$a_{n+1}=a_n+d$なら,そのまま等差型として解ける.

    \item \textbf{$n$に依存するなら,どのような形か}

      \begin{itemize}[beginpenalty=10000]
        \item \textbf{指数の形}：$cq^n$などなら,\,$q\ne0$を確認し,\,$q^n$などで割って指数の部分を取り除くことを考える.
        \item \textbf{$n$の多項式}：$n^2+2n$などで係数が定数なら,多項式を引いて付加項を消す方法を考える（階差等比複合型）.
        \item \textbf{多項式と指数の積}：$n2^n$などなら,まず指数の部分で割り,残った多項式をどう扱うか考える.
      \end{itemize}
      例えば$n+2^n$のように複数の形が混ざっていれば,それぞれの部分を見る.

    \item \textbf{総和を求めやすい形にできるか}

      階差型へ帰着できたら,付加項の総和を考える.
      $n(n+1)$のような連続整数の積や,隣り合う項の差として表せないか調べる.
      差の形にできれば,総和を取ったときに途中の項が消える.
  \end{enumerate}

  \medskip
  \noindent \textbf{【補足：初項も手掛かりになる】}

  例えば初項が$\displaystyle \frac{1}{\sqrt{6}-\sqrt{2}}$だとしよう.
  わざわざこのような数字が設定されたということはきっと裏がある.
  実際これは$\cos 15^\circ$または$\sin 75^\circ$の値である.
  そこから考えて,この漸化式は三角関数型のものではないか,と予想するのである.
  ただし,初項だけで解法が決まるわけではない.漸化式の形と合わせて判断しよう.
  \par
  \endgroup

% 解法の選択を、計算の前に確認する短い練習。
\subsection{操作を選ぶ練習}\label{節：操作を選ぶ練習}
  ここでは\bookterm{general}{一般項}まで一気に計算する前に,新しい量を選び,その更新だけを確かめる.
  引く・割る・\bookterm{logarithm}{対数}を取るという操作が,それぞれ何を簡単にするかを比較しよう.

  \noindent\bookblocklink{operation-choice}{problem}{answer}{【問題】}\par\nobreak
  \begin{enumerate}[format=\bookitemlink{operation-choice}{problem}{answer}]
    \item $a_1=1$,\,$a_{n+1}=3a_n+2n+1$とする.
          定数を引くだけで付加項を消せるか.
          $g(n)=un+v$を引く場合の$u,v$を求め,\,$b_n=a_n-g(n)$の更新式と初項を示せ.
    \item $a_1=2$,\,$a_{n+1}=\frac{n+2}{n}a_n+(n+1)(n+2)$とする.
          $b_n=\frac{a_n}{n}$と$c_n=\frac{a_n}{n(n+1)}$を候補とし,それぞれの更新式を求めよ.
          どちらで$a_n$の係数を$1$にそろえられるか.
    \item $a_1=-2$,\,$a_{n+1}=2^{2n+1}a_n$とする.
          $b_n=\log_2a_n$は使えるか.
          \bookterm{logarithm}{対数}を使わず,\,$c_n=\frac{a_n}{2^{n^2}}$の更新式と初項を求めよ.
  \end{enumerate}

  第2問は,\sectionref*{節：階比型}で倍率をそろえた操作を,付加項がある場合に使う練習である.
  一般に
  \[
    a_{n+1}=p_na_n+q_n
  \]
  に対して,既知の非零の\bookterm{sequence}{数列}$h_n$が$h_{n+1}=p_nh_n$を満たすなら,
  \[
    b_n=\frac{a_n}{h_n},\qquad
    b_{n+1}-b_n=\frac{q_n}{h_{n+1}}
  \]
  となる.倍率をそろえる操作の後に,\bookterm{difference}{階差}を足し戻す操作が続く.
  $p_n\ne0$なら,\,$h_1=1$,\,$h_n=\prod_{k=1}^{n-1}p_k$を候補にできる（空積は$1$）.
  係数が零になる場合はこの割り算を使わず,その段階の更新を別に確認する.
  また,残った和がいつでも簡単に計算できるとは限らない.

  \noindent\bookblocklink{operation-choice}{answer}{problem}{【解答と選択の理由】}\par\nobreak
  \begin{enumerate}[format=\bookitemlink{operation-choice}{answer}{problem}]
    \item $a_n-\alpha$では付加項の$2n$が残る.
          $g(n+1)=3g(n)+2n+1$から
          \[
            u=3u+2,\qquad u+v=3v+1
          \]
          となり,\,$u=v=-1$である.
          従って$b_n=a_n+n+1$により$b_{n+1}=3b_n$,\,$b_1=3$を得る.
          引く量を定数から一次式へ変えたことで,付加項を消せた.

    \item 候補ごとに一歩だけ更新すると,
          \[
            b_{n+1}=\frac{n+2}{n+1}b_n+n+2
          \]
          と
          \[
            c_{n+1}=c_n+1,\qquad c_1=1
          \]
          を得る.
          $h_n=n(n+1)$なら
          \[
            \frac{h_{n+1}}{h_n}=\frac{n+2}{n}
          \]
          が元の倍率と一致する.
          この一致を手掛かりに割る因子を選ぶ.
          $n\ge1$なので分母は非零であり,\,$c_n=n$から$a_n=n^2(n+1)$も得られる.

    \item 正の倍率で更新するので,各項は負であり,\,$\log_2a_n$は実数として定義できない.
          一方,\,$2n+1=(n+1)^2-n^2$から
          \[
            c_{n+1}=c_n,\qquad c_1=-1
          \]
          となる.従って$a_n=-2^{n^2}$を得る.
          \bookterm{normalize}{正規化}で割る$2^{n^2}$は常に非零であり,元の項が正である必要はない.
          初項が正の比較例は\bookproblemref{practice-basic}{42}{標準問42}へ.
  \end{enumerate}


 \section{解法の変更}\label{節：解法の変更}

  先に挙げた着目点から試した変形でも簡単にならない時は,一度最初からやり直してみるとよい.
  この際に今までと同じ考え方ではなく,なるべく別のアプローチができないか検討すると解法を見つけやすい.

  また,それ以外の場合でも次のどれかが起こったら,今の方法を一度止め,解法の再検討を行うと良い.
  \begin{enumerate}
    \item 対数を取っても,数列の項を1次の和として表せない.
    \item 定数を引いて付加項を消そうとしたが,条件を満たす定数が求まらない.
    \item 置換後に新しい種類の項が増え続ける.
  \end{enumerate}
  このようなときには$a_2$から$a_5$ぐらいまでの値を求めてみると法則性に気づき,帰納法で証明できる場合がある.
  

  それでも見つけられない時は,特殊関数型の置換や,まだ試していない式変形も候補に入れよう.
  一旦ご飯でも食べて改めて取り組もう.
  ご飯の時間でないのなら,今までの数学の知識を総動員して似たものを探そう.
  
  \begin{samepage}
  この手順を身につけると,個別のパターンを暗記するだけでなく,
  見慣れない漸化式にも最初の一手を選べるようになる.
  \par
  \end{samepage}


\section{計算と検算の確認}\label{節：読む前の確認}
ここまでに使った、添字・和・式変形・検算を確かめる五題である。
解答には、計算の理由を読み返す場所も示した。

    \noindent\bookblocklink{reading-check}{problem}{answer}{【確認問題】}\par\nobreak
    \begin{enumerate}[format=\bookitemlink{reading-check}{problem}{answer}]
      \item $a_{n+1}=2a_n+n$の添字$n$を$n+1$に置き換えた式を書け.
      \item $\displaystyle\frac1{x(x+1)}=\frac1x-\frac1{x+1}$を確かめ,使える$x$の範囲を述べよ.
      \item $\displaystyle\sum_{k=1}^{3}k$と$\displaystyle\sum_{k=1}^{3}2^{k-1}$を計算せよ.
      \item $b_n=a_n+3$,\,$b_1=4$,\,$b_{n+1}=2b_n$から$a_n$を求めよ.
      \item $a_1=1$,\,$a_{n+1}=a_n+2$について,最初の3項が$1,3,5$であることだけで$a_n=2n-1$を証明できるか.すべての項で示すために何を確かめるか.
    \end{enumerate}

    \noindent\bookblocklink{reading-check}{answer}{problem}{【解答と戻り先】}\par\nobreak
    \begin{enumerate}[format=\bookitemlink{reading-check}{answer}{problem}]
      \item $a_{n+2}=2a_{n+1}+n+1$.
            項と付加項の添字を両方ずらす.\sectionref*{節：シグマの計算}では添字の役割を確認する.
      \item 通分すると右辺は$\displaystyle\frac{(x+1)-x}{x(x+1)}$となる.
            $x\ne0,-1$が必要である.\sectionref*{節：式の変形}へ.
      \item $6$と$7$.
            \sectionref*{節：シグマの計算},\sectionref*{節：等差数列とその総和の導出},\sectionref*{節：等比数列の総和とその導出}へ.
      \item $b_n=4\cdot2^{n-1}$より$a_n=4\cdot2^{n-1}-3$.
            元へ戻す際は$a_n=b_n-3$である.\sectionref*{節：等比型}へ.
      \item 3項の一致は予想の手掛かりである.
            初項の一致と,任意の$n\ge1$について$a_n=2n-1$なら$a_{n+1}=2(n+1)-1$となることを確かめる.
            \sectionref*{節：数学的帰納法}へ.
    \end{enumerate}


% 問題演習を第4章の節として読み込む。
% chapter: 問題演習

  \section{問題演習}\label{節：問題演習}
    ここでは,これまでに扱った漸化式の解法を実際の問題を通して確認する.
    標準難易度55題では,基本・応用11パターンを横断し,式の構造から解法を選ぶ.
    応用難易度10題では,複数の解法をどの順番で組み合わせるかを考える.
    発展難易度10題では,漸化式を立てて,さまざまな数学の対象を調べる.
    方針・解答・解法は問題と同じ問番号に対応し,問番号から対応箇所へ移動できる.

% 第4章 標準難易度：基本4パターン・応用7パターンを各5題。
% 主分類・出題意図・別解の対応は docs/chapter4-standard-design-2026-10-04.md。

    \subsection{標準難易度}\label{節：標準難易度}\label{節：基本と応用から出題されるもの}
      ここでは,第3章の基本4パターンと応用7パターンから,それぞれ5題ずつ,計55題を出題する.
      同じ解法で解ける問題も,式の並び方や係数が変わると違って見える.
      まず移項や約分によって式を整理し,どの数量の差や比に注目すればよいかを考えよう.
      問題は型ごとにまとめずに並べてあるので,式の構造を見て解法を選んでほしい.
      方針と解法は問題と同じ問番号に対応する.
      別解のある問題では,最初の解法を理解したあとに,別の見方でも同じ答えになることを確認しよう.
      特に階比型と対数型には共通して使える見方がある.
      以下,\,$n$は正の整数とし,必要なときは
      \[
        S_n=\sum_{k=1}^{n}a_k
      \]
      とおく.

      \subsubsection{問題}\label{節：標準難易度：問題}\label{節：基本と応用から出題されるもの：問題}
        \begin{enumerate}[label=\textbf{問\arabic*.},leftmargin=3.8em,format=\bookitemlink{practice-basic}{problem}{answer}]
          % Q1
          \item $\displaystyle a_1=3$ とし,
                \[
                  (2a_n+3)a_{n+1}=3a_n
                \]
                を満たす数列$\{a_n\}$の一般項を求めよ.
          % S1
          \item $\displaystyle S_n=2a_n-n^2$
                                を満たす数列$\{a_n\}$の一般項を求めよ.
          % R1
          \item $\displaystyle a_1=3,\qquad a_{n+1}=2(n+1)a_n$
                で定められる数列の一般項を求めよ.
          % M1
          \item $\displaystyle a_1=2,\qquad a_{n+1}=2a_n+n-1$
                で定められる数列$\{a_n\}$の一般項を求めよ.
          % G1
          \item $\displaystyle a_1=6,\qquad a_{n+1}=-\frac12a_n$
                で定められる数列の一般項を求めよ.
          % T1
          \item $\displaystyle a_1=2,\,a_2=5$ とし,
                \[
                  a_{n+2}+6a_n=5a_{n+1}
                \]
                を満たす数列$\{a_n\}$の一般項を求めよ.
          % U1
          \item $\displaystyle a_1=3,\qquad b_1=1,$
                            \[
                              \begin{cases}
                                a_{n+1}=2a_n-b_n,\\
                                b_{n+1}=-a_n+2b_n
                              \end{cases}
                            \]
                                で定められる2つの数列$\{a_n\},\{b_n\}$の一般項を求めよ.
          % L1
          \item $\displaystyle a_1=4,\qquad a_{n+1}=\frac{a_n^3}{4}$
                で定められる数列$\{a_n\}$の一般項を求めよ.
          % C1
          \item $\displaystyle a_1=5,\qquad a_{n+1}=3a_n-4$
                で定められる数列$\{a_n\}$の一般項を求めよ.
          % D1
          \item $\displaystyle a_1=\frac56,\qquad 3a_{n+1}=3a_n-2$
                で定められる数列の一般項を求めよ.
          % F1
          \item $\displaystyle a_1=2,\qquad a_{n+1}=a_n+3n-1$
                で定められる数列の一般項を求めよ.
          % T2
          \item $\displaystyle a_1=3,\,a_2=1$ とし,
                \[
                  a_{n+2}+a_{n+1}=6a_n
                \]
                を満たす数列$\{a_n\}$の一般項を求めよ.
          % F2
          \item $\displaystyle a_1=1,\qquad a_{n+1}=a_n+3n^2+3n+1$
                で定められる数列の一般項を求めよ.
          % Q2
          \item $\displaystyle a_1=2,\,a_{n+1}=-\frac{2a_n}{3a_n+2}$
                で定められる数列$\{a_n\}$の一般項を求めよ.
          % C2
          \item $\displaystyle a_1=1,\qquad 2a_{n+1}+a_n=9$
                で定められる数列$\{a_n\}$の一般項を求めよ.
          % U2
          \item $\displaystyle a_1=3,\qquad b_1=1,$
                            \[
                              \begin{cases}
                                a_{n+1}=4a_n-2b_n,\\
                                b_{n+1}=a_n+b_n
                              \end{cases}
                            \]
                                で定められる2つの数列$\{a_n\},\{b_n\}$の一般項を求めよ.
          % D2
          \item $\displaystyle a_1=-4,\qquad 2a_{n+1}+a_n=3a_n+6$
                で定められる数列の一般項を求めよ.
          % G2
          \item $\displaystyle a_1=3,\qquad 2a_{n+1}-5a_n=a_{n+1}+a_n$
                で定められる数列の一般項を求めよ.
          % R2
          \item $\displaystyle a_1=6,\qquad
                 (n+1)(n+2)a_{n+1}=n(n+3)a_n$
                で定められる数列の一般項を求めよ.
          % M2
          \item $\displaystyle a_1=3,\qquad a_{n+1}+a_n=2n^2+2n+1$
                で定められる数列$\{a_n\}$の一般項を求めよ.
          % S2
          \item $\displaystyle S_n=a_n+(n-1)2^{n-1}$
                                を満たす数列$\{a_n\}$の一般項を求めよ.
          % L2
          \item $\displaystyle a_1=5,\qquad a_{n+1}=a_n^2-4a_n+6$
                で定められる数列$\{a_n\}$の一般項を求めよ.
          % T3
          \item $\displaystyle a_1=1,\,a_2=0$ とし,
                \[
                  a_{n+2}-4a_{n+1}+4a_n=0
                \]
                を満たす数列$\{a_n\}$の一般項を求めよ.
          % M3
          \item $\displaystyle a_1=1,\qquad a_{n+1}=3a_n-2^n$
                で定められる数列$\{a_n\}$の一般項を求めよ.
          % U3
          \item $\displaystyle a_1=1,\qquad b_1=2,$
                            \[
                              \begin{cases}
                                a_{n+1}=2a_n+b_n,\\
                                b_{n+1}=2b_n
                              \end{cases}
                            \]
                                で定められる2つの数列$\{a_n\},\{b_n\}$の一般項を求めよ.
          % Q3
          \item $\displaystyle a_1=4,\,a_{n+1}=\frac{3a_n+2}{a_n+2}$
                で定められる数列$\{a_n\}$の一般項を求めよ.
          % G3
          \item $\displaystyle a_1=4,\qquad (n+1)a_{n+1}=\frac{n+1}{3}a_n$
                で定められる数列の一般項を求めよ.
          % D3
          \item $\displaystyle a_1=1-\sqrt2,\qquad a_{n+1}+n=a_n+n+\sqrt2$
                で定められる数列の一般項を求めよ.
          % L3
          \item $\displaystyle a_1=2,\qquad a_{n+1}=a_n^{\frac{n+2}{n}}$
                で定められる正の数列$\{a_n\}$の一般項を求めよ.
          % S3
          \item $\displaystyle S_1=1,\qquad S_{n+1}=3S_n+2$
                                を満たす数列$\{a_n\}$の一般項を求めよ.
          % F3
          \item $\displaystyle a_1=3,\qquad a_{n+1}=a_n+2\cdot3^n$
                で定められる数列の一般項を求めよ.
          % R3
          \item $\displaystyle a_1=4,\qquad (n+1)^2a_{n+1}=n^2a_n$
                で定められる数列の一般項を求めよ.
          % C3
          \item $\displaystyle a_1=0,\qquad a_{n+1}-a_n=\frac{6-a_n}{3}$
                で定められる数列$\{a_n\}$の一般項を求めよ.
          % S4
          \item $\displaystyle a_1=1,\qquad a_{n+1}=2S_n+2^n$
                                で定められる数列$\{a_n\}$の一般項を求めよ.
          % Q4
          \item $\displaystyle a_1=1,\,a_{n+1}=2+\frac3{a_n}$
                で定められる数列$\{a_n\}$の一般項を求めよ.
          % G4
          \item 数列$\{a_n\}$は,ある実数$r$について
                \[
                  a_{n+1}=ra_n
                \]
                を満たし,\,$a_2=-6$,\, $a_5=48$である.\,$r$および一般項$a_n$を求めよ.
          % F4
          \item $\displaystyle a_1=0,\qquad a_{n+1}=a_n+\frac{2}{(n+1)(n+3)}$
                で定められる数列の一般項を求めよ.
          % T4
          \item $\displaystyle a_1=3,\,a_2=5$ とし,
                \[
                  a_{n+2}=2a_{n+1}+3a_n-8
                \]
                を満たす数列$\{a_n\}$の一般項を求めよ.
          % U4
          \item $\displaystyle a_1=0,\qquad b_1=1,$
                            \[
                              \begin{cases}
                                a_{n+1}=2a_n+b_n+1,\\
                                b_{n+1}=a_n+2b_n-1
                              \end{cases}
                            \]
                                で定められる2つの数列$\{a_n\},\{b_n\}$の一般項を求めよ.
          % L4
          \item $\displaystyle a_1=1,\qquad a_{n+1}=2\sqrt{a_n}$
                で定められる数列$\{a_n\}$の一般項を求めよ.
          % C4
          \item 実数$t$に対し,
                \[
                  a_1=t,\qquad a_{n+1}=4a_n-9
                \]
                で数列$\{a_n\}$を定める.一般項を求め,この数列が定数数列となる$t$の値を求めよ.
          % R4
          \item $\displaystyle a_1=2,\qquad a_{n+1}=2^{2n+1}a_n$
                で定められる数列の一般項を求めよ.
          % D4
          \item 数列$\{a_n\}$は,ある実数$d$について
                \[
                  a_{n+1}=a_n+d
                \]
                を満たし,\,$a_3=8$,\, $a_8=-7$である.\,$d$,\, $a_1$および一般項$a_n$を求めよ.
          % M4
          \item $\displaystyle a_1=2,\qquad a_{n+1}=2a_n+3\cdot2^n$
                で定められる数列$\{a_n\}$の一般項を求めよ.
          % U5
          \item $\displaystyle a_1=1,\qquad b_1=0,$
                            \[
                              \begin{cases}
                                a_{n+1}=a_n+2b_n,\\
                                b_{n+1}=a_n
                              \end{cases}
                            \]
                                で定められる2つの数列$\{a_n\},\{b_n\}$の一般項を求めよ.
          % R5
          \item $\displaystyle a_1=1,\qquad a_{n+1}=\frac{n(n+1)}2a_n$
                で定められる数列の一般項を求めよ.
          % Q5
          \item $\displaystyle a_1=2,\,a_{n+1}=2-\frac1{a_n}$
                で定められる数列$\{a_n\}$の一般項を求めよ.
          % G5
          \item $\displaystyle a_1=0,\qquad a_{n+1}=7a_n$
                で定められる数列の一般項を求めよ.
          % L5
          \item $\displaystyle a_1=4,\qquad a_{n+1}=\frac{a_n^2}{2^{n-1}}$
                で定められる数列$\{a_n\}$の一般項を求めよ.
          % F5
          \item $\displaystyle a_1=1,\qquad a_{n+1}=a_n+(-1)^n$
                で定められる数列の一般項を求めよ.
          % C5
          \item 実数$p,q$を用いて,
                \[
                  a_{n+1}=pa_n+q
                \]
                で定められる数列$\{a_n\}$について,
                \[
                  a_1=8,\qquad a_2=5,\qquad a_3=\frac72
                \]
                であることが分かっている.\,$p,q$の値と一般項を求めよ.
          % T5
          \item $\displaystyle a_1=1,\,a_2=2$ とし,
                \[
                  a_{n+2}-3a_{n+1}+2a_n=1
                \]
                を満たす数列$\{a_n\}$の一般項を求めよ.
          % D5
          \item $\displaystyle a_1=\frac74,\qquad 4a_{n+1}-4a_n=0$
                で定められる数列の一般項を求めよ.
          % M5
          \item $\displaystyle a_1=0,\qquad a_{n+1}=2a_n+3(-1)^n$
                で定められる数列$\{a_n\}$の一般項を求めよ.
          % S5
          \item $\displaystyle a_1=2,\qquad S_n=\frac{n+2}{3}a_n$
                                を満たす数列$\{a_n\}$の一般項を求めよ.
        \end{enumerate}

      \subsubsection{方針}\label{節：標準難易度：方針}\label{節：基本と応用から出題されるもの：方針}
        \begin{enumerate}[label=\textbf{問\arabic*.},leftmargin=3.8em,format=\bookitemlink{practice-basic}{hint}{problem}]
          % Q1
          \item $a_{n+1}$について解き,両辺の逆数を取る.\,$\displaystyle \frac{1}{a_n}$の増え方に注目する.
          % S1
          \item $n=1$から初項を求める. 次に$S_{n+1}-S_n=a_{n+1}$を使い, 総和を消す. 得られた漸化式では$a_n$に$n$の一次式を加えて等比数列を作る.
          % R1
          \item $(n+1)!=(n+1)n!$を使い,\,$a_n$を$n!$で割る.
          % M1
          \item $a_n$から適切な一次式を引き,残りが等比数列になるようにする.
          % G1
          \item 初項から第$n$項までに,$-\dfrac12$を掛ける回数を考える.
          % T1
          \item 特性方程式を解き,\,$a_{n+1}-2a_n$と$a_{n+1}-3a_n$を求める.二式の差で$a_n$が残る.
          % U1
          \item 2つの式を足すと和は変わらず, 引くと差は3倍になる. 和と差を求めてから$a_n,b_n$に戻す.
          % L1
          \item 全項が正であることを確かめ,底を$2$とする対数を取る.
          % C1
          \item $a_n$の値が変わらないとしたときの値を求め,その値との差を考える.
          % D1
          \item 両辺を$3$で割り,\,$a_{n+1}-a_n$を求める.
          % F1
          \item $a_{k+1}-a_k=3k-1$を$k=1$から$n-1$まで加える.
          % T2
          \item 特性方程式の解は$2$と$-3$である.負の解に対応する補助数列と$(-3)^{n-1}$の符号に注意する.
          % F2
          \item $3n^2+3n+1$が$(n+1)^3-n^3$であることを使う.
          % Q2
          \item 逆数を取って一次の漸化式に帰着させる.最初の数項を計算すると,より短い解法も見つかる.
          % C2
          \item まず$a_{n+1}$について解く.不動点との差は負の公比をもつ等比数列になる.
          % U2
          \item $a_n+t b_n$が等比数列になるように$t$を選ぶ. $t=-1,-2$で2本の等比数列を得られる. 片方を消して隣接3項間の漸化式にする別解もある.
          % D2
          \item 両辺から$a_n$を引き,その後で$2$で割る.
          % G2
          \item $a_{n+1}$を左辺に,\,$a_n$を右辺にまとめる.
          % R2
          \item 倍率を$\displaystyle \frac{n}{n+1}\frac{n+3}{n+2}$と分解し,順に掛ける.
          % M2
          \item 右辺を$(n+1)^2+n^2$と見れば,引くべき二次式が分かる.
          % S2
          \item $S_{n+1}-S_n=a_{n+1}$に代入する. 両辺の$a_{n+1}$が消え,\,$a_n$が直接求まる. $n=1$の与式だけでは初項が決まらないことにも注意する.
          % L2
          \item 平方完成すると$a_{n+1}-2=(a_n-2)^2$となる.正の量$a_n-2$の対数を取る.
          % T3
          \item 特性方程式は重解$2$を持つ.\,$a_{n+1}-2a_n$を等比数列にした後,\,$\displaystyle \frac{a_n}{2^{n-1}}$を考える.
          % M3
          \item $2^n$に関係する等比数列を引くか,$2^{n-1}$で割って一次の漸化式にする.
          % U3
          \item $b_n$は$a_n$と関係なく先に解ける. その結果を第1式に代入し,\,$\displaystyle \frac{a_n}{2^{n-1}}$の階差を調べる.
          % Q3
          \item 不動点は$2$と$-1$である.\,$\displaystyle \frac{a_n-2}{a_n+1}$を考える.\,$\displaystyle \frac{1}{a_n-2}$を考える別解もある.
          % G3
          \item $n$は正整数なので,\,$n+1$で割ることができる.
          % D3
          \item 両辺に現れる$n$を消去してから,隣り合う項の差を見る.
          % L3
          \item 対数を取った数列の階比が$\displaystyle \frac{n+2}{n}$となる.積の途中の因子が消えることに注目する.
          % S3
          \item $S_n+1$を等比数列として解いてから,\,$a_n=S_n-S_{n-1}$を使う. この差の公式を使う範囲と$a_1$を区別する.
          % F3
          \item 階差を加えた後,\,$3+3^2+\cdots+3^{n-1}$を求める.
          % R3
          \item 左辺は第$n+1$項に$(n+1)^2$を掛けた形,右辺は第$n$項に$n^2$を掛けた形である.
          % C3
          \item 右辺が$0$になる値$6$との差を取ると,等比型になる.
          % S4
          \item $S_{n+1}=S_n+a_{n+1}$から$S_n$の漸化式を作る. $S_n+2^n$が等比数列になる. 1つ前の式との差から$a_n$の漸化式を作る別解もある.
          % Q4
          \item 右辺を一つの分数にまとめる.不動点$3,-1$からの比を取ると,公比が負の等比数列になる.
          % G4
          \item 第$2$項から第$5$項までには$r$を$3$回掛ける.
          % F4
          \item $\displaystyle \frac{2}{(k+1)(k+3)}=\frac1{k+1}-\frac1{k+3}$と分解する.
          % T4
          \item 定数解$2$を引く.その後の特性方程式の一つの解$3$を使うと,初項が$0$の補助数列が現れる.
          % U4
          \item 定数項にも注意して2つの式を足し引きする. 和は3倍され, 差には2が加えられる.
          % L4
          \item 平方根を$\displaystyle a_n^{\frac{1}{2}}$と書き,底を$2$とする対数を取る.
          % C4
          \item 不動点との差を取る.その差の初項が$0$となる場合も,分けて確認する.
          % R4
          \item $2$の累乗を掛けると指数が加わる.対数を取って和に直してもよい.
          % D4
          \item 第$3$項から第$8$項までには,公差を何回加えるかを数える.
          % M4
          \item $2^{n-1}$で割った数列を考える.変形後は等差型になる.
          % U5
          \item 第1式の添字を1つ進め, 第2式を代入して$b_n$を消す. $a_2$はもとの連立から求める. 最後に$b_n$を求めるときは初項も確認する.
          % R5
          \item $(n-1)!n!$の次の項との比が$n(n+1)$であることを使う.
          % Q5
          \item 不動点の方程式は重解$1$を持つ.二つの差の比の代わりに,\,$\displaystyle \frac{1}{a_n-1}$を考える.
          % G5
          \item $0$に$7$を掛けても$0$のままである.項で割らずに考える.
          % L5
          \item 対数を取ると$b_{n+1}=2b_n-n+1$となる.この式を満たす一次式を探す.
          % F5
          \item $(-1)^k$の和は等比数列の和として求められる.最初の数項も調べる.
          % C5
          \item 最初の2回の更新式を引けば$q$が消える.係数を求めた後,不動点との差を取る.
          % T5
          \item 定数を引くだけでは右辺を$0$にできない.\,$d_n=a_{n+1}-a_n$を考え,一次の漸化式に帰着させる.
          % D5
          \item $a_{n+1}=a_n$なら,項の値は変わらない.
          % M5
          \item $(-1)^n$も指数の形である.$(-1)^{n-1}$で割った数列を考える.
          % S5
          \item $S_{n+1}-S_n=a_{n+1}$から$\displaystyle \frac{a_{n+1}}{a_n}$を求める. 比を順に掛けると, 分子と分母の大部分が打ち消し合う.
        \end{enumerate}

      \subsubsection{解答}\label{節：標準難易度：解答}\label{節：基本と応用から出題されるもの：解答}
        \begin{enumerate}[label=\textbf{問\arabic*.},leftmargin=3.8em,format=\bookitemlink{practice-basic}{answer}{problem}]
          % Q1
          \item \[
                  a_n=\frac{3}{2n-1}
                \]
          % S1
          \item \[
                  a_n=6\cdot2^{n-1}-2n-3
                \]
          % R1
          \item \[
                  a_n=3\cdot2^{n-1}n!
                \]
          % M1
          \item \[
                  a_n=3\cdot2^{n-1}-n
                \]
          % G1
          \item \[
                  a_n=6\left(-\frac12\right)^{n-1}
                \]
          % T1
          \item \[
                  a_n=2^{n-1}+3^{n-1}
                \]
          % U1
          \item \[
                  a_n=2+3^{n-1},\qquad b_n=2-3^{n-1}
                \]
          % L1
          \item \[
                  a_n=2^{\,1+3^{n-1}}
                \]
          % C1
          \item \[
                  a_n=2+3^n
                \]
          % D1
          \item \[
                  a_n=\frac56-\frac23(n-1)=\frac{9-4n}{6}
                \]
          % F1
          \item \[
                  a_n=\frac{3n^2-5n+6}{2}
                \]
          % T2
          \item \[
                  a_n=2^n+(-3)^{n-1}
                \]
          % F2
          \item \[
                  a_n=n^3
                \]
          % Q2
          \item \[
                  a_n=\frac4{5(-1)^{n-1}-3}
                \]
          % C2
          \item \[
                  a_n=3-2\left(-\frac12\right)^{n-1}
                \]
          % U2
          \item \[
                  a_n=4\cdot3^{n-1}-2^{n-1},\qquad b_n=2\cdot3^{n-1}-2^{n-1}
                \]
          % D2
          \item \[
                  a_n=3n-7
                \]
          % G2
          \item \[
                  a_n=3\cdot6^{n-1}
                \]
          % R2
          \item \[
                  a_n=\frac{2(n+2)}n
                \]
          % M2
          \item \[
                  a_n=n^2+2(-1)^{n-1}
                \]
          % S2
          \item \[
                  a_n=(n+1)2^{n-1}
                \]
          % L2
          \item \[
                  a_n=2+3^{\,2^{n-1}}
                \]
          % T3
          \item \[
                  a_n=(2-n)2^{n-1}
                \]
          % M3
          \item \[
                  a_n=2^n-3^{n-1}
                \]
          % U3
          \item \[
                  a_n=n2^{n-1},\qquad b_n=2^n
                \]
          % Q3
          \item \[
                  a_n=\frac{10\cdot4^{n-1}+2}{5\cdot4^{n-1}-2}
                \]
          % G3
          \item \[
                  a_n=4\left(\frac13\right)^{n-1}
                \]
          % D3
          \item \[
                  a_n=1+(n-2)\sqrt2
                \]
          % L3
          \item \[
                  a_n=2^{\frac{n(n+1)}2}
                \]
          % S3
          \item \[
                  a_1=1,\qquad a_n=4\cdot3^{n-2}\quad(n\ge2)
                \]
          % F3
          \item \[
                  a_n=3^n
                \]
          % R3
          \item \[
                  a_n=\frac4{n^2}
                \]
          % C3
          \item \[
                  a_n=6\left\{1-\left(\frac23\right)^{n-1}\right\}
                \]
          % S4
          \item \[
                  a_n=2\cdot3^{n-1}-2^{n-1}
                \]
          % Q4
          \item \[
                  a_n=\frac{3^n+(-1)^n}{3^{n-1}-(-1)^n}
                \]
          % G4
          \item \[
                  r=-2,\qquad a_n=3(-2)^{n-1}
                \]
          % F4
          \item \[
                  a_n=\frac56-\frac1{n+1}-\frac1{n+2}
                \]
          % T4
          \item \[
                  a_n=2+3^{n-1}
                \]
          % U4
          \item \[
                  a_n=\frac{3^{n-1}+2n-3}{2},\qquad b_n=\frac{3^{n-1}-2n+3}{2}
                \]
          % L4
          \item \[
                  a_n=2^{\,2-2^{2-n}}
                \]
          % C4
          \item \[
                  a_n=3+(t-3)4^{n-1},\qquad
                \text{定数数列となるのは }t=3
                \]
          % R4
          \item \[
                  a_n=2^{n^2}
                \]
          % D4
          \item \[
                  d=-3,\qquad a_1=14,\qquad a_n=17-3n
                \]
          % M4
          \item \[
                  a_n=(3n-1)2^{n-1}
                \]
          % U5
          \item \[
                  a_n=\frac{2^n+(-1)^{n-1}}{3},\qquad b_n=\frac{2^{n-1}-(-1)^{n-1}}{3}
                \]
          % R5
          \item \[
                  a_n=\frac{(n-1)!\,n!}{2^{n-1}}
                \]
          % Q5
          \item \[
                  a_n=\frac{n+1}{n}
                \]
          % G5
          \item \[
                  a_n=0
                \]
          % L5
          \item \[
                  a_n=2^{\,n+2^{n-1}}
                \]
          % F5
          \item \[
                  a_n=\frac{1+(-1)^{n-1}}2
                \]
          % C5
          \item \[
                  p=\frac12,\qquad q=1,\qquad
                a_n=2+6\left(\frac12\right)^{n-1}
                \]
          % T5
          \item \[
                  a_n=2^n-n
                \]
          % D5
          \item \[
                  a_n=\frac74
                \]
          % M5
          \item \[
                  a_n=(-1)^{n-1}-2^{n-1}
                \]
          % S5
          \item \[
                  a_n=n(n+1)
                \]
        \end{enumerate}

      \subsubsection{解法}\label{節：標準難易度：解法}\label{節：基本と応用から出題されるもの：解法}
        \begin{enumerate}[label=\textbf{問\arabic*.},leftmargin=3.8em,format=\bookitemlink{practice-basic}{solution}{problem}]
          % Q1
          \item $a_1>0$であり,\,$a_n>0$なら
                \[
                  a_{n+1}=\frac{3a_n}{2a_n+3}>0
                \]
                である.したがって,すべての項は正で,分母$2a_n+3$も$0$ではない.
                両辺の逆数を取ると,
                \[
                  \frac1{a_{n+1}}=\frac{2a_n+3}{3a_n}
                  =\frac1{a_n}+\frac23.
                \]
                $\displaystyle b_n=\frac{1}{a_n}$と置けば,初項$\displaystyle \frac{1}{3}$,公差$\displaystyle \frac{2}{3}$の等差数列になる.よって,
                \[
                  b_n=\frac13+\frac23(n-1)=\frac{2n-1}{3},
                  \qquad
                  a_n=\frac3{2n-1}.
                \]
                求めた式もすべての正整数$n$で正であり,逆数を戻す際の分母も$0$にならない.

                \noindent \textbf{【別解】}\\
                最初の数項は
                \[
                  a_1=\frac31,\quad a_2=\frac33,\quad a_3=\frac35
                \]
                なので,\,$\displaystyle a_n=\frac{3}{2n-1}$と予想できる.
                この式は$n=1$で成り立つ.\,$\displaystyle a_k=\frac{3}{2k-1}$と仮定すると,
                \[
                  a_{k+1}
                  =\frac{3\cdot\frac3{2k-1}}{2\cdot\frac3{2k-1}+3}
                  =\frac3{2k+1}
                  =\frac3{2(k+1)-1}.
                \]
                したがって,数学的帰納法によって一般項が確かめられる.
          % S1
          \item $n=1$を代入すると$a_1=2a_1-1$であるから,\,$a_1=1$である.
                                また,
                            \begin{align*}
                              a_{n+1}
                                &=S_{n+1}-S_n\\
                                &=2a_{n+1}-(n+1)^2\\
                                &\quad{}-2a_n+n^2
                            \end{align*}
                                より,
                            \[
                              a_{n+1}=2a_n+2n+1.
                            \]
                                右辺に$n$があるので, 定数を足すだけでは等比数列にならない.
                                $b_n=a_n+2n+3$とおくと,
                            \begin{align*}
                              b_{n+1}
                                &=2a_n+2n+1+2(n+1)+3\\
                                &=2(a_n+2n+3)=2b_n.
                            \end{align*}
                                $b_1=1+2+3=6$より$b_n=6\cdot2^{n-1}$である.
                                したがって,
                            \[
                              a_n=6\cdot2^{n-1}-2n-3.
                            \]
          % R1
          \item 係数$n+1$は階乗に含まれるので,漸化式の両辺を$(n+1)!$で割ると,
                \[
                  \frac{a_{n+1}}{(n+1)!}=2\frac{a_n}{n!}
                \]
                となる.そこで
                \[
                  b_n=\frac{a_n}{n!}
                \]
                と置くと,\,$b_{n+1}=2b_n$,\, $b_1=3$である.
                よって,\,$b_n=3\cdot2^{n-1}$であり,
                \[
                  a_n=3\cdot2^{n-1}n!
                \]
                を得る.分母を$(n-1)!$にするのではなく,係数$n+1$と添字が合う$n!$を選ぶことがポイントである.

                \noindent \textbf{【別解】}\\
                倍率を初項から順に掛けると,
                \[
                  a_n
                    =3\prod_{k=1}^{n-1}2(k+1)
                    =3\cdot2^{n-1}(2\cdot3\cdots n)
                    =3\cdot2^{n-1}n!
                \]
                である.この積には$2$が$n-1$個含まれる.
                $n=1$では空積を$1$とするので,同じ式が成り立つ.
          % M1
          \item 定数を引くだけでは$n-1$を消せないので,同じ漸化式を満たす一次式$g(n)=un+v$を探す.
                必要な条件は,
                \[
                  g(n+1)=2g(n)+n-1
                \]
                である.これに$g(n)=un+v$を代入して係数を比べると,
                \[
                  u=2u+1,\qquad u+v=2v-1
                \]
                なので,\,$u=-1,v=0$を得る.つまり$g(n)=-n$である.
                \,$b_n=a_n-g(n)=a_n+n$とおくと,
                \[
                  b_{n+1}=a_{n+1}+n+1=2a_n+2n=2b_n
                \]
                であり,\,$b_1=2+1=3$だから,
                \[
                  b_n=3\cdot2^{n-1},\qquad
                  a_n=3\cdot2^{n-1}-n
                \]
                となる.最終式に$n=1$を代入すると$2$である.

                \noindent \textbf{【別解：階差を先に求める】}\\
                隣り合う更新式を引くと,
                \[
                  d_n=a_{n+1}-a_n,\qquad d_{n+1}=2d_n+1
                \]
                となる.\,$a_2=4$より$d_1=2$なので,
                \[
                  d_n+1=3\cdot2^{n-1}
                \]
                である.従って,\,$n\ge2$では,
                \[
                  a_n=2+\sum_{k=1}^{n-1}(3\cdot2^{k-1}-1)
                     =3\cdot2^{n-1}-n
                \]
                となる.\,$n=1$でも最終式は成り立つ.
          % G1
          \item 次の項は前の項に$-\dfrac12$を掛けて得られる.
                したがって,初項$6$,\,公比$-\dfrac12$の等比数列であるから,
                \[
                  a_n=6\left(-\frac12\right)^{n-1}
                \]
                である.最初の数項は$6,\,-3,\,\dfrac32,\,-\dfrac34,\ldots$となる.
                公比が負であるため符号は交互に変わり,その絶対値は毎回$\dfrac12$倍になる.
          % T1
          \item 漸化式を$a_{n+2}-5a_{n+1}+6a_n=0$と整理する.
                特性方程式は
                \[
                  x^2-5x+6=(x-2)(x-3)=0
                \]
                なので,二つの解は$2$と$3$である.
                そこで,
                \[
                  b_n=a_{n+1}-2a_n,\qquad c_n=a_{n+1}-3a_n
                \]
                と置くと,
                \begin{align*}
                  b_{n+1}
                  &=a_{n+2}-2a_{n+1}
                  \\
                  &=3(a_{n+1}-2a_n)=3b_n,\\
                  c_{n+1}
                  &=a_{n+2}-3a_{n+1}
                  \\
                  &=2(a_{n+1}-3a_n)=2c_n.
                \end{align*}
                初項は$b_1=5-2\cdot2=1,\,c_1=5-3\cdot2=-1$だから,
                \[
                  b_n=3^{n-1},\qquad c_n=-2^{n-1}.
                \]
                二式の差を取ると$b_n-c_n=a_n$である.よって,
                \[
                  a_n=2^{n-1}+3^{n-1}.
                \]
                二つの等比数列から,共通して含まれている$a_{n+1}$を消去することが要点である.
          % U1
          \item $s_n=a_n+b_n,\ d_n=a_n-b_n$とおく.
                                2つの式を足し引きすると,
                            \[
                              s_{n+1}=s_n,\qquad d_{n+1}=3d_n.
                            \]
                                $s_1=3+1=4,\ d_1=3-1=2$より,
                            \[
                              s_n=4,\qquad d_n=2\cdot3^{n-1}.
                            \]
                                和と差からもとの2数列を復元すると,
                            \begin{align*}
                              a_n&=\frac{s_n+d_n}{2}=2+3^{n-1},\\
                              b_n&=\frac{s_n-d_n}{2}=2-3^{n-1}.
                            \end{align*}

                \noindent \textbf{【別解】}\\
                和が変わらないことから$a_n+b_n=4$である.
                                したがって$b_n=4-a_n$を第1式に代入して,
                            \[
                              a_{n+1}=2a_n-(4-a_n)=3a_n-4.
                            \]
                                $a_{n+1}-2=3(a_n-2)$と$a_1-2=1$より$a_n-2=3^{n-1}$である.
                                $b_n=4-a_n$を使えば, 2数列を求められる.
          % L1
          \item $a_1=4>0$であり,\,$a_n>0$なら$\displaystyle a_{n+1}=\frac{a_n^3}{4}>0$である.
                従って数学的帰納法により全ての項が正となり,対数を取ることができる.
                \,$b_n=\log_2a_n$とおくと,
                \[
                  b_{n+1}=3b_n-2,\qquad b_1=2
                \]
                となる.不動点$1$との差を取れば,
                \[
                  b_{n+1}-1=3(b_n-1)
                \]
                であり,\,$b_1-1=1$だから,
                \[
                  b_n=1+3^{n-1}
                \]
                を得る.\,$b_n=\log_2a_n$をもとの数列に戻して,
                \[
                  a_n=2^{\,1+3^{n-1}}
                \]
                となる.

                \noindent \textbf{【別解：指数を直接追う】}\\
                \,$\displaystyle c_n=\frac{a_n}{2}$とおくと,
                \[
                  c_{n+1}=\frac{a_n^3}{8}=c_n^3,\qquad c_1=2
                \]
                となる.1回進むたびに指数が$3$倍されるので,
                \[
                  c_n=2^{3^{n-1}}
                \]
                である.実際,\,$n=1$では$2$となり,この式の両辺を3乗すれば$n+1$の式が得られる.
                従って,\,$a_n=2c_n=2^{1+3^{n-1}}$となる.

                \noindent \textbf{【別解：階比を補助数列にする】}\\
                全項が正なので,\,$\displaystyle r_n=\frac{a_{n+1}}{a_n}=\frac{a_n^2}{4}$とおける.
                この階比について,
                \[
                  r_{n+1}=\frac{a_{n+1}^2}{4}
                         =\frac{a_n^6}{64}=r_n^3,\qquad r_1=4
                \]
                が成り立つ.指数を追うと$r_n=4^{3^{n-1}}$となるから,\,$n\ge2$では,
                \[
                  a_n=4\prod_{k=1}^{n-1}r_k
                     =4^{\,1+\sum_{k=1}^{n-1}3^{k-1}}
                     =4^{\,\frac{1+3^{n-1}}{2}}
                     =2^{\,1+3^{n-1}}
                \]
                である.ここでは階比の積を,指数の和に読み替えている.
                最終式は$n=1$でも成り立つ.
          % C1
          \item 定数項$-4$を消すため,同じ漸化式を満たす定数$\alpha$を探す.
                \[
                  \alpha=3\alpha-4\quad\Longleftrightarrow\quad\alpha=2
                \]
                この式をもとの漸化式から引くと,
                \[
                  a_{n+1}-2=3(a_n-2)
                \]
                となる.そこで,\,$b_n=a_n-2$とおくと,
                \[
                  b_{n+1}=3b_n,\qquad b_1=5-2=3
                \]
                である.初項が$3$,公比が$3$の等比数列なので,
                \[
                  b_n=3\cdot3^{n-1}=3^n
                \]
                従って,
                \[
                  a_n=2+3^n
                \]
                となる.\,$n=1$を代入すると$5$となり,初期条件も満たす.

                \noindent \textbf{【別解：階差に注目する】}\\
                漸化式を1つずらして引くと,
                \[
                  a_{n+2}-a_{n+1}=3(a_{n+1}-a_n)
                \]
                である.\,$d_n=a_{n+1}-a_n$とおくと,\,$a_2=11$より$d_1=6$であり,
                \[
                  d_n=6\cdot3^{n-1}
                \]
                となる.従って,\,$n\ge2$では,
                \[
                  a_n=5+6\sum_{k=1}^{n-1}3^{k-1}
                     =5+3(3^{n-1}-1)=2+3^n
                \]
                を得る.得られた式は$n=1$でも成り立つ.
          % D1
          \item 両辺を$3$で割ると,
                \[
                  a_{n+1}=a_n-\frac23
                \]
                となる.隣り合う項の差が常に$-\dfrac23$であるから,これは公差$-\dfrac23$の等差数列である.
                初項から第$n$項までにこの公差を$n-1$回加えるので,
                \[
                  a_n=\frac56-\frac23(n-1)=\frac{9-4n}{6}
                \]
                である.この式に$n=1$を代入すると,\,$a_1=\dfrac56$も確かめられる.
          % F1
          \item 隣り合う項の差は$3n-1$なので,初項にこれまでの増加分を加えればよい.
                $n\ge2$のとき,
                \begin{align*}
                  a_n
                    &=2+\sum_{k=1}^{n-1}(3k-1)\\
                    &=2+3\frac{n(n-1)}2-(n-1)\\
                    &=\frac{3n^2-5n+6}{2}
                \end{align*}
                となる.最後の式は$n=1$でも$2$となるから,すべての正整数$n$について成り立つ.
                和の上限を$n$にしてしまうと,第$n+1$項を求めることになるので注意する.
          % T2
          \item 特性方程式は
                \[
                  x^2+x-6=(x-2)(x+3)=0
                \]
                なので,二つの解は$2$と$-3$である.
                \[
                  b_n=a_{n+1}-2a_n,\qquad c_n=a_{n+1}+3a_n
                \]
                と置くと,
                \begin{align*}
                  b_{n+1}
                  &=a_{n+2}-2a_{n+1}
                  \\
                  &=-3(a_{n+1}-2a_n)=-3b_n,\\
                  c_{n+1}
                  &=a_{n+2}+3a_{n+1}
                  \\
                  &=2(a_{n+1}+3a_n)=2c_n.
                \end{align*}
                $b_1=1-2\cdot3=-5,\,c_1=1+3\cdot3=10$だから,
                \[
                  b_n=-5(-3)^{n-1},\qquad c_n=10\cdot2^{n-1}.
                \]
                $c_n-b_n=5a_n$なので,
                \[
                  5a_n=10\cdot2^{n-1}+5(-3)^{n-1}.
                \]
                よって,
                \[
                  a_n=2^n+(-3)^{n-1}.
                \]
                負の公比も等比数列の公比としてそのまま扱える.
          % F2
          \item 二項定理より,
                \[
                  3n^2+3n+1=(n+1)^3-n^3
                \]
                である.したがって,\,$n\ge2$のとき,
                \begin{align*}
                  a_n
                    &=1+\sum_{k=1}^{n-1}\{(k+1)^3-k^3\}\\
                    &=1+(2^3-1^3)+(3^3-2^3)\\
                    &\quad{}+\cdots+\{n^3-(n-1)^3\}\\
                    &=1+n^3-1\\
                    &=n^3
                \end{align*}
                となる.途中の立方数が次々に相殺するため,平方和を別に計算する必要はない.
                $n=1$でも$a_1=1=1^3$なので,結論はすべての正整数$n$で成り立つ.

                \noindent \textbf{【別解1】}\\
                最初の数項を計算すると$1,\,8,\,27,\ldots$となり,\,$a_n=n^3$が予想できる.
                これを数学的帰納法で確かめる.
                $n=1$では$a_1=1=1^3$である.
                $a_n=n^3$を仮定すると,
                \[
                  a_{n+1}=n^3+3n^2+3n+1=(n+1)^3
                \]
                となるから,予想はすべての正整数$n$について正しい.

                \noindent \textbf{【別解2】}\\
                漸化式を
                \[
                  a_{n+1}-(n+1)^3=a_n-n^3
                \]
                と書けば,\,$a_n-n^3$は一定である.
                その値は$a_1-1^3=0$だから,\,$a_n=n^3$となる.
          % Q2
          \item 初項から計算すると,
                \[
                  a_1=2,\qquad a_2=-\frac12,\qquad a_3=2.
                \]
                $2$と$\displaystyle -\frac{1}{2}$は互いに移り合い,これらの値で分母$3a_n+2$はそれぞれ$8$と$\displaystyle \frac{1}{2}$になる.
                したがって,数列はすべての正整数$n$について定まり,どの項も$0$ではない.
                一般項を一つの式で表すために,\,$\displaystyle b_n=\frac{1}{a_n}$と置くと,
                \[
                  b_{n+1}=-b_n-\frac32,\qquad b_1=\frac12.
                \]
                定数項を消すために$\displaystyle \frac{3}{4}$を加えれば,
                \[
                  b_{n+1}+\frac34=-\left(b_n+\frac34\right).
                \]
                よって,
                \[
                  b_n+\frac34=\left(\frac12+\frac34\right)(-1)^{n-1}
                  =\frac54(-1)^{n-1}.
                \]
                したがって,
                \[
                  b_n=\frac{5(-1)^{n-1}-3}{4},
                  \qquad
                  a_n=\frac4{5(-1)^{n-1}-3}.
                \]
                この分母は$n$が奇数なら$2$,偶数なら$-8$なので,\,$0$にはならない.

                \noindent \textbf{【別解】}\\
                漸化式の右辺を$\displaystyle f(x)=\frac{-2x}{3x+2}$と書くと,
                \[
                  f(2)=-\frac12,\qquad f\left(-\frac12\right)=2.
                \]
                したがって,初項からこの二つの値を交互に取ることが分かる.
                \[
                  a_n=
                  \begin{cases}
                    2 &(n\text{が奇数}),\\
                    -\dfrac12 &(n\text{が偶数}).
                  \end{cases}
                \]
                これも一般項の表し方である.逆数を用いた式は,この二つの場合を$(-1)^{n-1}$によってまとめたものである.
          % C2
          \item まず次の項だけを左辺に残すと,
                \[
                  a_{n+1}=-\frac12a_n+\frac92
                \]
                である.不動点は,
                \[
                  \alpha=-\frac12\alpha+\frac92
                  \quad\Longleftrightarrow\quad \alpha=3
                \]
                となる.\,$b_n=a_n-3$とおけば,
                \[
                  b_{n+1}=-\frac12b_n,\qquad b_1=1-3=-2
                \]
                なので,
                \[
                  b_n=-2\left(-\frac12\right)^{n-1},\qquad
                  a_n=3-2\left(-\frac12\right)^{n-1}
                \]
                を得る.負の公比により,\,$a_n-3$の符号が1項ごとに入れ替わる.
                例えば,\,$a_1=1<3$であり,\,$a_2=4>3$である.
          % U2
          \item $a_n+t b_n$が等比数列になるように$t$を探す.
                            \[
                              a_{n+1}+t b_{n+1}
                                =(4+t)a_n+(t-2)b_n.
                            \]
                                これが$(4+t)(a_n+t b_n)$と一致するためには,
                            \[
                              t-2=t(4+t),\qquad
                              (t+1)(t+2)=0
                            \]
                                であればよい.
                                $t=-1,-2$をそれぞれ用いると,
                            \[
                              a_{n+1}-b_{n+1}=3(a_n-b_n),
                            \]
                            \[
                              a_{n+1}-2b_{n+1}=2(a_n-2b_n).
                            \]
                                初項から,
                            \[
                              a_n-b_n=2\cdot3^{n-1},\qquad
                              a_n-2b_n=2^{n-1}.
                            \]
                                この2式の差を取ると$b_n=2\cdot3^{n-1}-2^{n-1}$である.
                                さらに$a_n=b_n+2\cdot3^{n-1}$より,
                            \[
                              a_n=4\cdot3^{n-1}-2^{n-1},\qquad
                              b_n=2\cdot3^{n-1}-2^{n-1}.
                            \]

                \noindent \textbf{【別解】}\\
                第1式から$2b_n=4a_n-a_{n+1}$である.
                                第2式を使うと,
                            \begin{align*}
                              a_{n+2}
                                &=4a_{n+1}-2b_{n+1}\\
                                &=4a_{n+1}-2(a_n+b_n)\\
                                &=5a_{n+1}-6a_n.
                            \end{align*}
                                $a_1=3,\ a_2=4\cdot3-2\cdot1=10$であり, 特性方程式は
                            \[
                              r^2-5r+6=(r-2)(r-3)=0.
                            \]
                                $a_n=C\cdot2^{n-1}+D\cdot3^{n-1}$とおくと,
                            \[
                              C+D=3,\qquad 2C+3D=10
                            \]
                                より$C=-1,\ D=4$である.
                                $\displaystyle b_n=\frac{4a_n-a_{n+1}}{2}$から$b_n$も求められる.
          % D2
          \item $a_n$を右辺にまとめると,
                \begin{align*}
                  2a_{n+1}&=2a_n+6,\\
                  a_{n+1}&=a_n+3
                \end{align*}
                となる.したがって,初項$-4$,\,公差$3$の等差数列であるから,
                \[
                  a_n=-4+3(n-1)=3n-7
                \]
                である.係数が両辺に分かれていても,まず隣り合う項の差を取り出せばよい.
          % G2
          \item 両辺から$a_{n+1}$を引き,さらに$5a_n$を加えると,
                \[
                  a_{n+1}=6a_n
                \]
                となる.定数項はなく,毎回同じ倍率で増える数列である.
                よって,初項$3$,\,公比$6$の等比数列として,
                \[
                  a_n=3\cdot6^{n-1}
                \]
                を得る.
          % R2
          \item $n\ge1$より$(n+1)(n+2)\ne0$であるから,
                \[
                  a_{n+1}=\frac{n}{n+1}\frac{n+3}{n+2}a_n
                \]
                と書ける.この2つの分数の積を分けて掛ければ,\,$n\ge2$のとき,
                \begin{align*}
                  a_n
                    &=6\prod_{k=1}^{n-1}\frac{k}{k+1}
                        \prod_{k=1}^{n-1}\frac{k+3}{k+2}\\
                    &=6\left(\frac12\frac23\cdots\frac{n-1}{n}\right)
                        \left(\frac43\frac54\cdots\frac{n+2}{n+1}\right)\\
                    &=6\cdot\frac1n\cdot\frac{n+2}{3}\\
                    &=\frac{2(n+2)}n
                \end{align*}
                となる.分子と次の分母が相殺するので,大きな積を計算する必要はない.
                $n=1$でも右辺は$6$であり,初項と一致する.

                \noindent \textbf{【別解】}\\
                両辺が同じ形になるように,
                \[
                  b_n=\frac{n}{n+2}a_n
                \]
                と置く.もとの漸化式より,
                \[
                  b_{n+1}
                    =\frac{n+1}{n+3}
                      \frac{n(n+3)}{(n+1)(n+2)}a_n
                    =\frac{n}{n+2}a_n
                    =b_n
                \]
                である.初項は$b_1=\dfrac13\cdot6=2$だから,\,$b_n=2$であり,
                \[
                  a_n=\frac{2(n+2)}n
                \]
                となる.
          % M2
          \item 二次式を引いて等比型にすることを考える.
                右辺は,
                \[
                  2n^2+2n+1=(n+1)^2+n^2
                \]
                と書ける.このため与式は,
                \[
                  a_{n+1}-(n+1)^2=-(a_n-n^2)
                \]
                となる.\,$b_n=a_n-n^2$とおけば,
                \[
                  b_{n+1}=-b_n,\qquad b_1=3-1=2
                \]
                なので,
                \[
                  b_n=2(-1)^{n-1},\qquad a_n=n^2+2(-1)^{n-1}
                \]
                を得る.二次式の部分$n^2$を引いた残りは,\,$2,-2,2,-2,\ldots$と交互に変わる.
          % S2
          \item $n=1$のとき与式は$S_1=a_1$となるだけなので, この式だけでは$a_1$は決まらない.
                                ただし, すべての正整数$n$で与式が成り立つことを使えば, 初項も決められる.
                            \begin{align*}
                              a_{n+1}
                                &=S_{n+1}-S_n\\
                                &=a_{n+1}+n2^n\\
                                &\quad{}-a_n-(n-1)2^{n-1}.
                            \end{align*}
                                両辺の$a_{n+1}$を消すと,
                            \begin{align*}
                              a_n
                                &=n2^n-(n-1)2^{n-1}\\
                                &=(n+1)2^{n-1}.
                            \end{align*}
                                この計算は$n\ge1$で成り立つので,\,$a_1=2$も含めて求められた.
          % L2
          \item 項の加減を含む式のまま対数を取っても簡単にならないので,まず平方完成する.
                \[
                  a_{n+1}-2=(a_n-2)^2
                \]
                ここで$a_1-2=3>0$であり,正の数の2乗は正なので,
                数学的帰納法により全ての$n$で$a_n-2>0$となる.
                そこで,\,$b_n=\log_3(a_n-2)$とおくと,
                \[
                  b_{n+1}=2b_n,\qquad b_1=\log_3 3=1
                \]
                である.従って,
                \[
                  b_n=2^{n-1},\qquad a_n-2=3^{\,2^{n-1}}
                \]
                なので,
                \[
                  a_n=2+3^{\,2^{n-1}}
                \]
                を得る.\,$n=1$でも$a_1=5$となる.

                \noindent \textbf{【別解：繰り返し2乗する】}\\
                \,$c_n=a_n-2$とおけば,\,$c_{n+1}=c_n^2,c_1=3$である.
                従って指数が$1,2,4,8,\ldots$と変わり,
                \[
                  c_n=3^{2^{n-1}}
                \]
                が予想できる.この式は$n=1$で成り立ち,
                \[
                  (3^{2^{n-1}})^2=3^{2^n}
                \]
                より次の項でも成り立つので,数学的帰納法で確認できる.
          % T3
          \item 特性方程式は
                \[
                  x^2-4x+4=(x-2)^2=0
                \]
                であり,重解$2$を持つ.二つの異なる補助数列は作れないので,
                \[
                  b_n=a_{n+1}-2a_n
                \]
                だけをまず調べる.
                \[
                  b_{n+1}=a_{n+2}-2a_{n+1}
                  =2(a_{n+1}-2a_n)=2b_n.
                \]
                $b_1=0-2\cdot1=-2$だから,
                \[
                  b_n=-2^n,\qquad a_{n+1}-2a_n=-2^n.
                \]
                両辺を$2^n$で割ると,
                \[
                  \frac{a_{n+1}}{2^n}-\frac{a_n}{2^{n-1}}=-1.
                \]
                $\displaystyle c_n=\frac{a_n}{2^{n-1}}$と置けば,\,$c_1=1,\,c_{n+1}=c_n-1$となる.よって,
                \[
                  c_n=1-(n-1)=2-n,
                  \qquad a_n=(2-n)2^{n-1}.
                \]
                割ったものは$2^n$であり,\,$a_2=0$でもこの操作に問題はない.

                \noindent \textbf{【別解】}\\
                初めから$\displaystyle c_n=\frac{a_n}{2^{n-1}}$と置き,もとの漸化式を$2^{n+1}$で割ると,
                \[
                  c_{n+2}-2c_{n+1}+c_n=0.
                \]
                したがって,
                \[
                  c_{n+2}-c_{n+1}=c_{n+1}-c_n.
                \]
                つまり,\,$\{c_n\}$は等差数列である.\,$c_1=1,\,c_2=0$より公差は$-1$なので,
                \[
                  c_n=2-n,\qquad a_n=(2-n)2^{n-1}.
                \]
                重解に対応する指数の部分を除くと,等差数列が残ることを直接確かめられる.
          % M3
          \item 加わる項が$-2^n$なので,同じ漸化式を満たす$g(n)=c\,2^n$を探す.
                \[
                  c\,2^{n+1}=3c\,2^n-2^n
                \]
                両辺を$2^n$で割ると$2c=3c-1$であり,\,$c=1$である.
                従って,\,$b_n=a_n-2^n$とおけば,
                \[
                  b_{n+1}=a_{n+1}-2^{n+1}=3(a_n-2^n)=3b_n
                \]
                となる.\,$b_1=1-2=-1$より,
                \[
                  b_n=-3^{n-1},\qquad a_n=2^n-3^{n-1}
                \]
                を得る.一般項は$n=1$で$1$となる.

                \noindent \textbf{【別解：指数列で割る】}\\
                \,$\displaystyle c_n=\frac{a_n}{2^{n-1}}$とおくと,
                \[
                  c_{n+1}=\frac32c_n-1,\qquad c_1=1
                \]
                である.不動点は$2$なので,
                \[
                  c_n-2=-\left(\frac32\right)^{n-1}
                \]
                となる.従って,
                \[
                  a_n=2^{n-1}\left\{2-\left(\frac32\right)^{n-1}\right\}
                     =2^n-3^{n-1}
                \]
                を得る.
          % U3
          \item 第2式はそれだけで解くことができる.
                                $b_1=2,\ b_{n+1}=2b_n$より,
                            \[
                              b_n=2^n.
                            \]
                                これを第1式に代入すると,
                            \[
                              a_{n+1}=2a_n+2^n.
                            \]
                                両辺を$2^n$で割り,
                            \[
                              \frac{a_{n+1}}{2^n}
                                =\frac{a_n}{2^{n-1}}+1
                            \]
                                と書き直す.
                                $\displaystyle c_n=\frac{a_n}{2^{n-1}}$とおけば$c_{n+1}=c_n+1,\ c_1=1$であるから,\,$c_n=n$である.
                                したがって,
                            \[
                              a_n=n2^{n-1},\qquad b_n=2^n.
                            \]

                \noindent \textbf{【別解】}\\
                $b_n=a_{n+1}-2a_n$を使い,\,$b_{n+1}=2b_n$から$b_n$を消去すると,
                            \[
                              a_{n+2}-2a_{n+1}
                                =2(a_{n+1}-2a_n).
                            \]
                                よって$a_{n+2}=4a_{n+1}-4a_n$である.
                                特性方程式は$(r-2)^2=0$である.
                                この形で$\displaystyle c_n=\frac{a_n}{2^{n-1}}$とおくと,
                            \[
                              c_{n+2}-2c_{n+1}+c_n=0
                            \]
                                となる. つまり$c_{n+1}-c_n$は一定なので,\,$c_n=C+D(n-1)$と書ける.
                                したがって,
                            \[
                              a_n=\{C+D(n-1)\}2^{n-1}.
                            \]
                                $a_1=1,\ a_2=2a_1+b_1=4$より$C=D=1$であり,\,$a_n=n2^{n-1}$を得る.
          % Q3
          \item 不動点の方程式は
                \[
                  x=\frac{3x+2}{x+2}
                  \quad\Longleftrightarrow\quad
                  x^2-x-2=(x-2)(x+1)=0
                \]
                なので,不動点は$2$と$-1$である.
                また,\,$a_1>2$であり,
                \[
                  a_{n+1}-2=\frac{a_n-2}{a_n+2}
                \]
                から,\,$a_n>2$なら$a_{n+1}>2$である.
                よって,すべての項で$a_n>2$となり,\,$a_n+2$と$a_n+1$は$0$ではない.
                そこで,
                \[
                  b_n=\frac{a_n-2}{a_n+1}
                \]
                と置くと,
                \begin{align*}
                  b_{n+1}
                  &=\cfrac{\cfrac{3a_n+2}{a_n+2}-2}
                           {\cfrac{3a_n+2}{a_n+2}+1}\\
                  &=\frac{a_n-2}{4a_n+4}
                  =\frac14b_n.
                \end{align*}
                $\displaystyle b_1=\frac{2}{5}$より,
                \[
                  b_n=\frac25\left(\frac14\right)^{n-1}.
                \]
                $b_n(a_n+1)=a_n-2$を$a_n$について解いて,
                \[
                  a_n=\frac{2+b_n}{1-b_n}
                  =\frac{10\cdot4^{n-1}+2}{5\cdot4^{n-1}-2}.
                \]
                $\displaystyle 0<b_n\leq\frac{2}{5}$なので,最後の変形で割った$1-b_n$も$0$ではない.

                \noindent \textbf{【別解】}\\
                不動点$2$からの差について,
                \[
                  a_{n+1}-2=\frac{a_n-2}{a_n+2}
                \]
                が成り立つ.\,$a_n>2$なので,この差の逆数を取ることができる.
                $\displaystyle c_n=\frac{1}{a_n-2}$と置くと,
                \[
                  c_{n+1}=\frac{a_n+2}{a_n-2}
                  =1+4c_n,\qquad c_1=\frac12.
                \]
                よって,
                \[
                  c_{n+1}+\frac13=4\left(c_n+\frac13\right),
                  \qquad
                  c_n=\frac56\cdot4^{n-1}-\frac13.
                \]
                したがって,
                \[
                  a_n=2+\frac1{c_n}
                  =2+\frac6{5\cdot4^{n-1}-2}
                  =\frac{10\cdot4^{n-1}+2}{5\cdot4^{n-1}-2}.
                \]
                二つの不動点からの比を使う方法と,一つの不動点からの差の逆数を使う方法は,同じ漸化式を別の形で単純にしている.
          % G3
          \item $n\ge1$より$n+1\ne0$である.両辺を$n+1$で割れば,
                \[
                  a_{n+1}=\frac13a_n
                \]
                となる.係数に$n$が含まれていても,整理した後の倍率は一定である.
                したがって,
                \[
                  a_n=4\left(\frac13\right)^{n-1}
                \]
                である.
          % D3
          \item 両辺から$n$を引けば,
                \[
                  a_{n+1}=a_n+\sqrt2
                \]
                となる.もとの式には$n$が現れるが,増加する量は$n$によらず一定である.
                よって,公差$\sqrt2$の等差数列として,
                \begin{align*}
                  a_n
                    &=1-\sqrt2+(n-1)\sqrt2\\
                    &=1+(n-2)\sqrt2
                \end{align*}
                を得る.公差が無理数でも,等差数列の一般項の求め方は同じである.
          % L3
          \item 全項が正であるので,\,$b_n=\log_2a_n$とおくことができる.
                すると,
                \[
                  b_{n+1}=\frac{n+2}{n}b_n,\qquad b_1=1
                \]
                となる.これは対数を取った後の数列が階比型となる例である.
                \,$n\ge2$では,
                \[
                  b_n=\prod_{k=1}^{n-1}\frac{k+2}{k}
                     =\frac{3\cdot4\cdots(n+1)}{1\cdot2\cdots(n-1)}
                     =\frac{n(n+1)}2
                \]
                である.最終式は$n=1$でも$b_1=1$を与える.
                従って,
                \[
                  a_n=2^{b_n}=2^{\frac{n(n+1)}2}
                \]
                となる.分数の指数が含まれていても,対数を取れば一次の係数になる.

                \noindent \textbf{【別解：対数を取らず指数をそろえる】}\\
                階比型の解法で現れた因子$n(n+1)$を利用し,
                \[
                  c_n=a_n^{\frac{2}{n(n+1)}}
                \]
                とおく.全項が正なので指数法則を使うことができ,
                \[
                  c_{n+1}
                  =\left(a_n^{\frac{n+2}{n}}\right)^{\frac{2}{(n+1)(n+2)}}
                  =a_n^{\frac{2}{n(n+1)}}=c_n
                \]
                となる.\,$c_1=2$だから$c_n=2$であり,
                \[
                  a_n=2^{\frac{n(n+1)}2}
                \]
                を得る.
          % S3
          \item 総和$S_n$も1つの数列なので, まずこれについて解く.
                            \[
                              S_{n+1}+1=3(S_n+1),\qquad S_1+1=2
                            \]
                                より,
                            \[
                              S_n=2\cdot3^{n-1}-1.
                            \]
                                $n\ge2$に対しては,
                            \begin{align*}
                              a_n
                                &=S_n-S_{n-1}\\
                                &=(2\cdot3^{n-1}-1)
                                  -(2\cdot3^{n-2}-1)\\
                                &=4\cdot3^{n-2}.
                            \end{align*}
                                一方, 初項は$a_1=S_1=1$である.
                                上で求めた式に$n=1$を代入すると$\displaystyle \frac{4}{3}$となり, 初項とは一致しない.
                                したがって,
                            \[
                              a_1=1,\qquad a_n=4\cdot3^{n-2}\quad(n\ge2).
                            \]

                \noindent \textbf{【別解】}\\
                $n\ge2$において,
                            \begin{align*}
                              a_{n+1}
                                &=S_{n+1}-S_n\\
                                &=(3S_n+2)-(3S_{n-1}+2)=3a_n.
                            \end{align*}
                                ただし,\,$S_n=3S_{n-1}+2$を使ったので, この式の導出は$n\ge2$で行っている.
                                $S_2=3S_1+2=5$より$a_2=5-1=4$であるから,
                            \[
                              a_n=4\cdot3^{n-2}\quad(n\ge2).
                            \]
                                $a_1=1$は別に付け加える.
          % F3
          \item $n\ge2$のとき,階差を順に加えると,
                \begin{align*}
                  a_n
                    &=3+2\sum_{k=1}^{n-1}3^k\\
                    &=3+2\frac{3(3^{n-1}-1)}{3-1}\\
                    &=3^n
                \end{align*}
                となる.最後の式は$n=1$でも初項$3$に一致する.
                階差が指数関数であっても,まず差を加え,その後で和を計算すればよい.

                \noindent \textbf{【別解1】}\\
                $2\cdot3^n=3^{n+1}-3^n$に気づけば,漸化式は
                \[
                  a_{n+1}-3^{n+1}=a_n-3^n
                \]
                と書ける.したがって,\,$a_n-3^n$は初項での値$a_1-3=0$に等しい.
                よって,\,$a_n=3^n$である.
                増加分が隣り合う指数関数の差になっているため,相殺を使っても解ける.

                \noindent \textbf{【別解2】}\\
                $a_n=3^n$を予想して帰納法を使うこともできる.
                初項は予想と一致し,\,$a_n=3^n$とすると,
                \[
                  a_{n+1}=3^n+2\cdot3^n=3^{n+1}
                \]
                となる.したがって,予想はすべての正整数$n$について成り立つ.
          % R3
          \item もとの式の両辺は,添字に合わせてその平方を掛けた同じ形になっている.
                そこで
                \[
                  b_n=n^2a_n
                \]
                と置くと,\,$b_{n+1}=b_n$である.
                初項は$b_1=1^2\cdot4=4$だから,すべての項で$b_n=4$となる.
                よって,
                \[
                  n^2a_n=4,\qquad a_n=\frac4{n^2}
                \]
                である.\,$n$は正整数なので,\,$n^2$で割ることができる.

                \noindent \textbf{【別解】}\\
                最初の数項は$4,\,1,\,\dfrac49,\,\dfrac14,\ldots$なので,\,$a_n=\dfrac4{n^2}$が予想できる.
                初項はこの式と一致する.
                $a_n=\dfrac4{n^2}$を仮定すると,
                \[
                  a_{n+1}
                    =\frac{n^2}{(n+1)^2}\frac4{n^2}
                    =\frac4{(n+1)^2}
                \]
                となる.したがって,数学的帰納法により一般項が確認できる.
          % C3
          \item 左辺が階差であっても,右辺には$a_n$が含まれる.
                そのまま既知の式の和を取ることはできないので,まず式を整理する.
                \[
                  a_{n+1}=\frac23a_n+2
                \]
                この式を満たす定数は$6$であり,
                \[
                  a_{n+1}-6=\frac23(a_n-6)
                \]
                となる.\,$b_n=a_n-6$とおくと,\,$b_1=-6$であるから,
                \[
                  b_n=-6\left(\frac23\right)^{n-1}
                \]
                よって,
                \[
                  a_n=6-6\left(\frac23\right)^{n-1}
                \]
                である.\,$n=1$では$0$となる.

                \noindent \textbf{【別解：増加量を求める】}\\
                与式を1つずらした式からもとの式を引くと,
                \[
                  (a_{n+2}-a_{n+1})-(a_{n+1}-a_n)
                  =-\frac13(a_{n+1}-a_n)
                \]
                となる.従って,\,$d_n=a_{n+1}-a_n$は,
                \[
                  d_{n+1}=\frac23d_n,\qquad d_1=2
                \]
                を満たす.これより,\,$n\ge2$では,
                \[
                  a_n=\sum_{k=1}^{n-1}2\left(\frac23\right)^{k-1}
                     =6\left\{1-\left(\frac23\right)^{n-1}\right\}
                \]
                を得る.この式は$n=1$でも正しい.
          % S4
          \item $S_1=a_1=1$であり,
                            \[
                              S_{n+1}=S_n+a_{n+1}=3S_n+2^n.
                            \]
                                そこで$c_n=S_n+2^n$とおくと,
                            \begin{align*}
                              c_{n+1}
                                &=3S_n+2^n+2^{n+1}\\
                                &=3(S_n+2^n)=3c_n.
                            \end{align*}
                                $c_1=1+2=3$より$c_n=3^n$となるので,
                            \[
                              S_n=3^n-2^n.
                            \]
                                よって$n\ge2$では,
                            \begin{align*}
                              a_n
                                &=S_n-S_{n-1}\\
                                &=2\cdot3^{n-1}-2^{n-1}.
                            \end{align*}
                                $n=1$でも右辺は$1=a_1$であるから,
                            \[
                              a_n=2\cdot3^{n-1}-2^{n-1}\quad(n\ge1).
                            \]

                \noindent \textbf{【別解】}\\
                $n\ge2$において, 与式から1つ前の式を引くと,
                            \[
                              a_{n+1}-a_n
                                =2(S_n-S_{n-1})+2^n-2^{n-1}
                                =2a_n+2^{n-1}.
                            \]
                                したがって$a_{n+1}=3a_n+2^{n-1}$である.
                                $a_2=2S_1+2=4=3a_1+1$より, この関係は$n=1$でも成り立つ.
                                $b_n=a_n+2^{n-1}$とおけば,
                            \[
                              b_{n+1}=3b_n,\qquad b_1=2.
                            \]
                                よって$b_n=2\cdot3^{n-1}$となり, 同じ答えを得る.
          % Q4
          \item $a_1>0$であり,\,$a_n>0$なら$\displaystyle a_{n+1}=2+\frac{3}{a_n}>0$なので,すべての項は正である.
                したがって,もとの分母$a_n$は$0$ではない.
                右辺を一つの分数にまとめると,
                \[
                  a_{n+1}=\frac{2a_n+3}{a_n}.
                \]
                不動点の方程式は
                \[
                  x=\frac{2x+3}{x}
                  \quad\Longleftrightarrow\quad
                  (x-3)(x+1)=0
                \]
                なので,\,$3$と$-1$が不動点である.
                $a_n+1>0$に注意して,
                \[
                  b_n=\frac{a_n-3}{a_n+1}
                \]
                と置けば,
                \begin{align*}
                  b_{n+1}
                  &=\frac{2a_n+3-3a_n}{2a_n+3+a_n}\\
                  &=-\frac{a_n-3}{3(a_n+1)}
                  =-\frac13b_n.
                \end{align*}
                $b_1=-1$だから,
                \[
                  b_n=-\left(-\frac13\right)^{n-1}
                  =\frac{(-1)^n}{3^{n-1}}.
                \]
                これを$a_n$に戻すと,
                \[
                  a_n=\frac{3+b_n}{1-b_n}
                  =\frac{3^n+(-1)^n}{3^{n-1}-(-1)^n}.
                \]
                最後の分母は,\,$n$が奇数なら$3^{n-1}+1$,偶数なら$3^{n-1}-1$である.
                偶数の場合は$n\geq2$なので,いずれも正であり$0$ではない.

                \noindent \textbf{【別解】}\\
                数列$\{u_n\}$を
                \[
                  u_1=u_2=1,\qquad u_{n+2}=2u_{n+1}+3u_n
                \]
                で定める.初めの二項が正で,右辺の係数も正なので,すべての$u_n$は正である.
                このとき,
                \[
                  \frac{u_{n+2}}{u_{n+1}}
                  =2+\cfrac{3}{\frac{u_{n+1}}{u_n}},
                  \qquad \frac{u_2}{u_1}=1.
                \]
                したがって,\,$\displaystyle a_n=\frac{u_{n+1}}{u_n}$となる.
                隣接3項間漸化式の特性方程式は
                \[
                  x^2-2x-3=(x-3)(x+1)=0
                \]
                である.\,$u_{n+2}+u_{n+1}=3(u_{n+1}+u_n)$に対応する等比数列と,
                $u_{n+2}-3u_{n+1}=-(u_{n+1}-3u_n)$に対応する等比数列を初項から求めると,
                \[
                  u_{n+1}+u_n=2\cdot3^{n-1},\qquad
                  u_{n+1}-3u_n=-2(-1)^{n-1}
                \]
                を得る.二式の差から,
                \[
                  u_n=\frac{3^{n-1}+(-1)^{n-1}}2.
                \]
                よって,
                \[
                  a_n=\frac{u_{n+1}}{u_n}
                  =\frac{3^n+(-1)^n}{3^{n-1}-(-1)^n}.
                \]
                項そのものではなく,隣り合う項の比として考えると,分数型と隣接3項間型がつながる.
          % G4
          \item 第$2$項から第$5$項までには$r$を$3$回掛けるから,
                \[
                  a_5=r^3a_2
                \]
                である.よって,\,$48=-6r^3$より$r^3=-8$となる.
                $r$は実数なので,\,$r=-2$である.
                さらに,\,$a_2=ra_1$より,
                \[
                  a_1=\frac{-6}{-2}=3
                \]
                となる.したがって,
                \[
                  a_n=3(-2)^{n-1}
                \]
                である.公比が負なので,正負が交互に現れることにも注意する.
          % F4
          \item 分母の2つの因数が$2$だけ離れていることに注目すると,
                \[
                  \frac{2}{(k+1)(k+3)}
                    =\frac1{k+1}-\frac1{k+3}
                \]
                と部分分数分解できる.よって,\,$n\ge2$のとき,
                \[
                  a_n=\sum_{k=1}^{n-1}
                       \left(\frac1{k+1}-\frac1{k+3}\right)
                \]
                である.和の範囲を分けて書くと,
                \begin{align*}
                  a_n
                    &=\sum_{j=2}^{n}\frac1j-\sum_{j=4}^{n+2}\frac1j\\
                    &=\frac12+\frac13-\frac1{n+1}-\frac1{n+2}\\
                    &=\frac56-\frac1{n+1}-\frac1{n+2}
                \end{align*}
                となる.2つの和は添字が$2$だけずれているので,両端にそれぞれ2項が残る.
                $n=2$では直接計算しても$a_2=\dfrac14$となり,得られた式と一致する.
                また,\,$n=1$を代入すると$0$となるから,この一般項は初項も含めて成り立つ.
          % T4
          \item 定数項を消すための定数解は,
                \[
                  \alpha=2\alpha+3\alpha-8
                  \quad\Longleftrightarrow\quad \alpha=2
                \]
                である.\,$b_n=a_n-2$と置けば,
                \[
                  b_{n+2}=2b_{n+1}+3b_n,
                  \qquad b_1=1,\quad b_2=3.
                \]
                特性方程式は
                \[
                  x^2-2x-3=(x-3)(x+1)=0.
                \]
                そこで,\,$c_n=b_{n+1}-3b_n$と置くと,
                \[
                  c_{n+1}=-c_n,\qquad c_1=3-3\cdot1=0.
                \]
                初項が$0$なので,\,$c_n=0$がすべての$n$で成り立つ.
                つまり,
                \[
                  b_{n+1}=3b_n.
                \]
                よって,\,$b_1=1$から$b_n=3^{n-1}$となり,
                \[
                  a_n=2+3^{n-1}.
                \]
                この問題では一方の補助数列が常に$0$になるため,二本を両方求める前に元の数列を等比数列に帰着できる.

                \noindent \textbf{【別解】}\\
                一般項を直接予想する前に,\,$a_{n+1}=3a_n-4$という短い漸化式が成り立つことを示す方法もある.
                まず$a_2=5=3a_1-4$である.
                ある$k$について$a_{k+1}=3a_k-4$が成り立つとき,
                \[
                  a_{k+2}-(3a_{k+1}-4)
                  =-a_{k+1}+3a_k-4
                  =0.
                \]
                したがって,数学的帰納法から$a_{n+1}=3a_n-4$が成り立つ.
                この一次の漸化式を
                \[
                  a_{n+1}-2=3(a_n-2)
                \]
                と変形すれば,\,$a_n=2+3^{n-1}$を得る.
          % U4
          \item $s_n=a_n+b_n,\ d_n=a_n-b_n$とおく.
                                2つの式を足すと定数項$1$と$-1$が消え,
                            \[
                              s_{n+1}=3s_n.
                            \]
                                一方, 引いたときには定数項の差は$2$となるので,
                            \[
                              d_{n+1}=d_n+2.
                            \]
                                $s_1=1,\ d_1=-1$より,
                            \[
                              s_n=3^{n-1},\qquad d_n=-1+2(n-1)=2n-3.
                            \]
                                したがって,
                            \begin{align*}
                              a_n&=\frac{s_n+d_n}{2}
                                   =\frac{3^{n-1}+2n-3}{2},\\
                              b_n&=\frac{s_n-d_n}{2}
                                   =\frac{3^{n-1}-2n+3}{2}.
                            \end{align*}
          % L4
          \item $a_1=1>0$であり,\,$a_n>0$なら$2\sqrt{a_n}>0$である.
                従って全ての項が正となる.
                \,$b_n=\log_2a_n$とおくと,
                \[
                  b_{n+1}=1+\frac12b_n,\qquad b_1=0
                \]
                である.不動点$2$との差を取れば,
                \[
                  b_{n+1}-2=\frac12(b_n-2)
                \]
                なので,
                \[
                  b_n=2-2\left(\frac12\right)^{n-1}=2-2^{2-n}
                \]
                となる.従って,
                \[
                  a_n=2^{\,2-2^{2-n}}
                \]
                を得る.\,$n=1$では指数が$0$となり,\,$a_1=1$である.

                \noindent \textbf{【別解：正の不動点で割る】}\\
                正の不動点は$4=2\sqrt4$を満たす$4$である.
                \,$\displaystyle c_n=\frac{a_n}{4}$とおくと,
                \[
                  c_{n+1}=\sqrt{c_n},\qquad c_1=\frac14
                \]
                となる.従って,
                \[
                  c_n=\left(\frac14\right)^{(\frac12)^{n-1}},\qquad
                  a_n=4\left(\frac14\right)^{(\frac12)^{n-1}}
                     =2^{\,2-2^{2-n}}
                \]
                を得る.1回進むたびに指数が半分になることを利用した.
          % C4
          \item 不動点を求めると,
                \[
                  \alpha=4\alpha-9\quad\Longleftrightarrow\quad\alpha=3
                \]
                である.\,$b_n=a_n-3$とおけば,
                \[
                  b_{n+1}=4b_n,\qquad b_1=t-3
                \]
                なので,
                \[
                  a_n=3+(t-3)4^{n-1}
                \]
                を得る.この計算では$t-3$で割っていないので,\,$t=3$の場合も含まれる.
                実際,\,$t=3$なら全ての$n$について$a_n=3$である.
                逆に定数数列なら$a_2=a_1$が必要であり,
                \[
                  4t-9=t\quad\Longleftrightarrow\quad t=3
                \]
                となる.従って求める値は$t=3$だけである.
          % R4
          \item 初項は正で,各段階の倍率$2^{2n+1}$も正なので,すべての項が正である.
                倍率を順に掛けると,\,$n\ge2$のとき,
                \begin{align*}
                  a_n
                    &=2\prod_{k=1}^{n-1}2^{2k+1}\\
                    &=2^{\,1+\sum_{k=1}^{n-1}(2k+1)}\\
                    &=2^{\,1+n(n-1)+(n-1)}\\
                    &=2^{n^2}
                \end{align*}
                となる.最後の式は$n=1$でも初項$2$に一致する.
                積の計算そのものではなく,指数の和を計算している点が大切である.

                \noindent \textbf{【別解1】}\\
                すべての項が正であるから,
                \[
                  b_n=\log_2 a_n
                \]
                と置くことができる.漸化式の両辺の対数を取ると,
                \[
                  b_{n+1}=b_n+2n+1,\qquad b_1=1
                \]
                となる.これは階差型であり,
                \[
                  b_n=1+\sum_{k=1}^{n-1}(2k+1)=n^2
                \]
                を得る.したがって,\,$a_n=2^{b_n}=2^{n^2}$である.
                指数の和を直接計算する方法と,対数を取る方法は,同じ計算を異なる形で表している.

                \noindent \textbf{【別解2】}\\
                $2n+1=(n+1)^2-n^2$を使うと,
                \[
                  \frac{a_{n+1}}{2^{(n+1)^2}}
                    =\frac{a_n}{2^{n^2}}
                \]
                となる.この比は初項で$\displaystyle \frac{2}{2}=1$だから,\,$a_n=2^{n^2}$を得る.
          % D4
          \item 第$3$項から第$8$項までには,公差$d$を$5$回加える.したがって,
                \[
                  a_8-a_3=5d
                \]
                であるから,\,$-7-8=5d$より$d=-3$を得る.
                また,\,$a_3=a_1+2d$であるから,
                \[
                  a_1=8-2(-3)=14
                \]
                である.よって,
                \[
                  a_n=14-3(n-1)=17-3n
                \]
                となる.求めた式に$n=3,\,8$を代入すれば,指定された2つの値をともに満たす.
          % M4
          \item $a_n$に掛かる係数と,加わる指数列の底がともに$2$である.
                そこで,\,$\displaystyle b_n=\frac{a_n}{2^{n-1}}$とおく.
                与式を$2^n$で割ると,
                \[
                  b_{n+1}=b_n+3,\qquad b_1=2
                \]
                となる.これは初項$2$,公差$3$の等差数列だから,
                \[
                  b_n=2+3(n-1)=3n-1
                \]
                である.よって,
                \[
                  a_n=(3n-1)2^{n-1}
                \]
                となる.変形後の式では不動点を求める必要がない.
                係数と底が等しいため,指数列で割った残りに$n$の一次式が現れる.
          % U5
          \item 第1式の添字を1つ進め, 第2式を代入すると,
                            \[
                              a_{n+2}=a_{n+1}+2b_{n+1}
                                =a_{n+1}+2a_n.
                            \]
                                $a_1=1,\ a_2=a_1+2b_1=1$である.
                                特性方程式は,
                            \[
                              r^2-r-2=(r-2)(r+1)=0.
                            \]
                                よって$a_n=C\cdot2^{n-1}+D(-1)^{n-1}$とおくと,
                            \[
                              C+D=1,\qquad 2C-D=1
                            \]
                                から$C=\frac{2}{3},\ D=\frac{1}{3}$となる.
                                したがって,
                            \[
                              a_n=\frac{2^n+(-1)^{n-1}}{3}.
                            \]
                                $n\ge2$では$b_n=a_{n-1}$であるから,
                            \[
                              b_n=\frac{2^{n-1}+(-1)^{n-2}}{3}
                                =\frac{2^{n-1}-(-1)^{n-1}}{3}.
                            \]
                                最後の式は$n=1$でも$0=b_1$となるので,
                            \[
                              b_n=\frac{2^{n-1}-(-1)^{n-1}}{3}\quad(n\ge1).
                            \]

                \noindent \textbf{【別解】}\\
                $c_n=a_n+b_n,\ d_n=a_n-2b_n$とおくと,
                            \[
                              c_{n+1}=2c_n,\qquad d_{n+1}=-d_n.
                            \]
                                $c_1=d_1=1$より$c_n=2^{n-1},\ d_n=(-1)^{n-1}$である.
                                $c_n-d_n=3b_n,\ 2c_n+d_n=3a_n$を用いると,\,$a_n,b_n$を直接求められる.
          % R5
          \item 係数$n$と$n+1$をそれぞれ階乗に吸収するため,
                \[
                  b_n=\frac{a_n}{(n-1)!\,n!}
                \]
                と置く.実際,
                \begin{align*}
                  b_{n+1}
                    &=\frac{a_{n+1}}{n!(n+1)!}\\
                    &=\frac{n(n+1)a_n}{2n!(n+1)!}\\
                    &=\frac12\frac{a_n}{(n-1)!\,n!}\\
                    &=\frac12b_n
                \end{align*}
                となる.ここでは$n!=n(n-1)!$,\, $(n+1)!=(n+1)n!$を用いた.
                $0!=1$より$b_1=1$なので,\,$b_n=2^{-(n-1)}$である.
                元に戻せば,
                \[
                  a_n=\frac{(n-1)!\,n!}{2^{n-1}}
                \]
                を得る.

                \noindent \textbf{【別解】}\\
                倍率を分けて掛ければ,
                \begin{align*}
                  a_n
                    &=\prod_{k=1}^{n-1}\frac{k(k+1)}2\\
                    &=\frac1{2^{n-1}}
                        \left(\prod_{k=1}^{n-1}k\right)
                        \left(\prod_{k=1}^{n-1}(k+1)\right)\\
                    &=\frac{(n-1)!\,n!}{2^{n-1}}
                \end{align*}
                となる.$n=1$では2つの積を空積$1$として扱うため,初項も同じ式に含まれる.
          % Q5
          \item 不動点の方程式は
                \[
                  x=2-\frac1x
                  \quad\Longleftrightarrow\quad
                  (x-1)^2=0
                \]
                であり,重解$1$を持つ.
                異なる二つの不動点を使った比は作れないので,\,$a_n-1$の変化を調べる.
                \[
                  a_{n+1}-1=1-\frac1{a_n}=\frac{a_n-1}{a_n}.
                \]
                $a_1>1$であり,\,$a_n>1$なら右辺は正なので$a_{n+1}>1$である.
                したがって,すべての項で$a_n>1$となり,もとの分母$a_n$も,差$a_n-1$も$0$ではない.
                両辺の逆数を取って,
                \[
                  \frac1{a_{n+1}-1}
                  =\frac{a_n}{a_n-1}
                  =1+\frac1{a_n-1}.
                \]
                $\displaystyle b_n=\frac{1}{a_n-1}$と置けば,\,$b_1=1,\,b_{n+1}=b_n+1$となる.
                よって,
                \[
                  b_n=n,\qquad a_n=1+\frac1n=\frac{n+1}{n}.
                \]
                不動点が重なる場合には,差の逆数から等差数列が現れた.

                \noindent \textbf{【別解】}\\
                数項を計算すると,
                \[
                  a_1=\frac21,\quad a_2=\frac32,\quad a_3=\frac43
                \]
                なので,\,$\displaystyle a_n=\frac{n+1}{n}$と予想できる.
                この式は$n=1$で成り立つ.\,$\displaystyle a_k=\frac{k+1}{k}$を仮定すると,
                \[
                  a_{k+1}=2-\frac{k}{k+1}
                  =\frac{k+2}{k+1}.
                \]
                したがって,数学的帰納法によって一般項を確かめられる.
          % G5
          \item $a_1=0$より,\,$a_2=7a_1=0$である.
                同様に,ある項が$0$なら次の項も$0$になるので,すべての項が$0$である.
                よって,
                \[
                  a_n=0
                \]
                を得る.等比型の一般項を直接用いても,\,$a_n=0\cdot7^{n-1}=0$となる.
                この例では$a_n=0$なので,\,$\displaystyle \frac{a_{n+1}}{a_n}$を作るために$a_n$で割ってはいけない.
          % L5
          \item 初項は正であり,正の項からは次の項も正となるので,全項について対数を取ることができる.
                \,$b_n=\log_2a_n$とおけば,
                \[
                  b_{n+1}=2b_n-n+1,\qquad b_1=2
                \]
                となる.項ごとに変わる$-n+1$を消すため,同じ式を満たす一次式$g(n)=un+v$を探す.
                \[
                  g(n+1)=2g(n)-n+1
                \]
                係数を比べると,
                \[
                  u=2u-1,\qquad u+v=2v+1
                \]
                なので,\,$u=1,v=0$であり,\,$g(n)=n$となる.
                従って,\,$c_n=b_n-n$とおけば,
                \[
                  c_{n+1}=2c_n,\qquad c_1=2-1=1
                \]
                である.よって,
                \[
                  b_n=n+2^{n-1},\qquad a_n=2^{\,n+2^{n-1}}
                \]
                を得る.指数に$n=1$を代入すると$2$なので,初項は$4$となる.

                \noindent \textbf{【別解：指数列で割ってから指数を追う】}\\
                対数を取った解法で$n$を引くことは,もとの項を$2^n$で割ることに対応する.
                \,$\displaystyle d_n=\frac{a_n}{2^n}$とおけば,
                \[
                  d_{n+1}=\frac{a_{n+1}}{2^{n+1}}
                         =\frac{a_n^2}{2^{2n}}=d_n^2,\qquad d_1=2
                \]
                となる.従って,
                \[
                  d_n=2^{2^{n-1}},\qquad
                  a_n=2^nd_n=2^{\,n+2^{n-1}}
                \]
                を得る.
          % F5
          \item $n\ge2$のとき,階差を加えると,
                \[
                  a_n=1+\sum_{k=1}^{n-1}(-1)^k
                \]
                となる.右辺の和は初項$-1$,\,公比$-1$の等比数列の和なので,
                \begin{align*}
                  a_n
                    &=1-\frac{1-(-1)^{n-1}}2\\
                    &=\frac{1+(-1)^{n-1}}2
                \end{align*}
                である.この式は$n=1$でも$1$となる.
                したがって,奇数番目の項は$1$,\,偶数番目の項は$0$である.

                \noindent \textbf{【別解】}\\
                2回分の変化をまとめると,
                \[
                  a_{n+2}
                    =a_n+(-1)^n+(-1)^{n+1}
                    =a_n
                \]
                となる.また,\,$a_1=1$,\, $a_2=0$であるから,数列は$1,\,0,\,1,\,0,\ldots$と繰り返す.
                この偶奇による違いを1つの式で表せば,
                \[
                  a_n=\frac{1+(-1)^{n-1}}2
                \]
                となる.
          % C5
          \item 最初の2回の更新より,
                \[
                  5=8p+q,\qquad \frac72=5p+q
                \]
                である.2式を引くと,
                \[
                  -\frac32=-3p
                \]
                なので,\,$p=\frac12$であり,\,$q=5-8\cdot\frac12=1$となる.
                従って,
                \[
                  a_{n+1}=\frac12a_n+1
                \]
                である.不動点$2$との差を取り,\,$b_n=a_n-2$とおくと,
                \[
                  b_{n+1}=\frac12b_n,\qquad b_1=6
                \]
                より,
                \[
                  a_n=2+6\left(\frac12\right)^{n-1}
                \]
                を得る.これに$n=1,2,3$を代入すると,順に$8,5,\frac72$となる.

                \noindent \textbf{【別解：階差から一般項を求める】}\\
                係数を決めた後,漸化式を1つずらして引くと,
                \[
                  d_n=a_{n+1}-a_n,\qquad
                  d_{n+1}=\frac12d_n,\qquad d_1=-3
                \]
                となる.従って,\,$n\ge2$では,
                \[
                  a_n=8-3\sum_{k=1}^{n-1}\left(\frac12\right)^{k-1}
                     =2+6\left(\frac12\right)^{n-1}
                \]
                を得る.最終式は$n=1$でも成り立つ.
          % T5
          \item すべての項を同じ定数$\alpha$に置き換えると,左辺は
                \[
                  \alpha-3\alpha+2\alpha=0
                \]
                となり,右辺の$1$とは一致しない.
                したがって,この問題には定数解がなく,定数を引くだけでは右辺を$0$にできない.
                そこで,階差
                \[
                  d_n=a_{n+1}-a_n
                \]
                を考える.漸化式を整理すると,
                \[
                  a_{n+2}-a_{n+1}=2(a_{n+1}-a_n)+1
                \]
                なので,
                \[
                  d_{n+1}=2d_n+1,\qquad d_1=2-1=1.
                \]
                これを$d_{n+1}+1=2(d_n+1)$と変形すれば,
                \[
                  d_n+1=2\cdot2^{n-1}=2^n,
                  \qquad d_n=2^n-1.
                \]
                よって,\,$n\geq2$では,
                \begin{align*}
                  a_n
                  &=a_1+\sum_{k=1}^{n-1}(a_{k+1}-a_k)\\
                  &=1+\sum_{k=1}^{n-1}(2^k-1)\\
                  &=1+(2^n-2)-(n-1)=2^n-n.
                \end{align*}
                この式は$n=1$でも$a_1=1$を与える.したがって,
                \[
                  a_n=2^n-n.
                \]
                階差もすぐに等比数列になるわけではないが,さらに$1$を加えることで等比数列になった.

                \noindent \textbf{【別解】}\\
                定数解の代わりに,\,$-n$が漸化式を満たすことに気付けば,
                \[
                  -(n+2)-3\{-(n+1)\}+2(-n)=1
                \]
                なので,\,$b_n=a_n+n$と置くことができる.
                \[
                  b_{n+2}-3b_{n+1}+2b_n=0,
                  \qquad b_1=2,\quad b_2=4.
                \]
                ここで,\,$c_n=b_{n+1}-2b_n$と置くと,
                \[
                  c_{n+1}=c_n,\qquad c_1=4-2\cdot2=0.
                \]
                したがって,\,$b_{n+1}=2b_n$となり,\,$b_n=2^n$を得る.
                よって,
                \[
                  a_n=b_n-n=2^n-n.
                \]
                右辺の定数を消すために,項ごとに変わる一次式を引くこともできる.
          % D5
          \item 両辺を$4$で割ると,\,$a_{n+1}-a_n=0$である.つまり,次の項は前の項と等しい.
                初項が$\dfrac74$なので,
                \[
                  a_1=a_2=a_3=\cdots=\frac74
                \]
                となり,\,$a_n=\dfrac74$である.
                これは公差$0$の等差数列であり,公式を使っても
                \[
                  a_n=\frac74+(n-1)\cdot0=\frac74
                \]
                と求まる.
          % M5
          \item $(-1)^n$は$0$にならないので,\,$\displaystyle b_n=\frac{a_n}{(-1)^{n-1}}$とおける.
                与式を$(-1)^n$で割ると,
                \[
                  b_{n+1}=-2b_n+3,\qquad b_1=0
                \]
                となる.この式の不動点は$1$であり,
                \[
                  b_{n+1}-1=-2(b_n-1)
                \]
                なので,
                \[
                  b_n=1-(-2)^{n-1}
                \]
                を得る.もとの数列に戻すと,
                \[
                  a_n=(-1)^{n-1}\{1-(-2)^{n-1}\}
                     =(-1)^{n-1}-2^{n-1}
                \]
                となる.最後の等号には$(-1)^{n-1}(-2)^{n-1}=2^{n-1}$を用いた.

                \noindent \textbf{【別解：交代する項を直接引く】}\\
                \,$g_n=(-1)^{n-1}$とおくと,
                \[
                  g_{n+1}=2g_n+3(-1)^n
                \]
                が成り立つ.与式からこれを引けば,
                \[
                  a_{n+1}-g_{n+1}=2(a_n-g_n)
                \]
                となる.\,$a_1-g_1=-1$だから,
                \[
                  a_n-g_n=-2^{n-1}
                \]
                であり,同じ一般項を得る.
          % S5
          \item 与式の$n$を$n+1$に置き換えた式との差を取ると,
                            \[
                              a_{n+1}
                                =\frac{n+3}{3}a_{n+1}
                                 -\frac{n+2}{3}a_n.
                            \]
                                整理して,
                            \[
                              na_{n+1}=(n+2)a_n,\qquad
                              a_{n+1}=\frac{n+2}{n}a_n.
                            \]
                                $a_1=2$よりすべての項は正であり,\,$n\ge2$では,
                            \begin{align*}
                              a_n
                                &=2\prod_{k=1}^{n-1}\frac{k+2}{k}\\
                                &=2\frac{3\cdot4\cdots(n+1)}{1\cdot2\cdots(n-1)}\\
                                &=n(n+1).
                            \end{align*}
                                この式は$n=1$でも$a_1=2$と一致する.
                                したがって,
                            \[
                              a_n=n(n+1).
                            \]

                \noindent \textbf{【別解】}\\
                $na_{n+1}=(n+2)a_n$を
                            \[
                              \frac{a_{n+1}}{(n+1)(n+2)}
                                =\frac{a_n}{n(n+1)}
                            \]
                                と書き直す. よって$\displaystyle \frac{a_n}{n(n+1)}$は一定であり, 初項からその値は$\displaystyle \frac{2}{1\cdot2}=1$である.
                                したがって$a_n=n(n+1)$となる.
        \end{enumerate}


    \subsection{応用難易度}\label{節：応用難易度}\label{節：高難易度のもの}
      ここでは,標準問題で学んだ解法を複数組み合わせる10題を扱う.
      見た目の複雑さをそのまま計算するのではなく,共通の分母,和と差,積や総和の構造を探そう.
      どの数量に変換するかと,変換をどの順番で使うかを考えてほしい.
      以下,\,$n$は正の整数とし,必要なときは$S_n=\sum_{k=1}^{n}a_k$とおく.

      \subsubsection{問題}\label{節：応用難易度：問題}\label{節：高難易度のもの：問題}
        \begin{enumerate}[label=\textbf{問\arabic*.},leftmargin=3.8em,format=\bookitemlink{practice-advanced}{problem}{answer}]
          % A1: 階乗で重みをそろえた総和
          \item $\displaystyle a_1=1$ とし,
                \[
                  a_{n+1}
                  =4(n+1)a_n
                    -(n+1)!\sum_{k=1}^{n}\frac{a_k}{k!}
                \]
                で数列$\{a_n\}$を定める.一般項を求めよ.
          % A2: 添字を含む分母と不動点
          \item $\displaystyle a_1=2$ とし,
                \[
                  a_{n+1}
                   =\frac{(n+1)a_n+1-n}{na_n+2-n}
                \]
                で数列$\{a_n\}$を定める.一般項を求めよ.
          % A3: 三角形状の重みをもつ積
          \item $\displaystyle a_1=1$ とし,
                \[
                  a_{n+1}
                   =2^n\prod_{k=1}^{n}a_k^{\,n-k+1}
                \]
                で正の数列$\{a_n\}$を定める.一般項を求めよ.
          % A4: 非線形3項間式から階比を作る
          \item $\displaystyle a_1=1,\qquad a_2=2$ とし,
                \[
                  (na_n+a_{n+1})a_{n+2}
                    =(n+1)a_{n+1}^{\,2}
                \]
                で正の数列$\{a_n\}$を定める.一般項を,階乗を用いて表せ.
          % A5: 和と差で異なる構造が現れる連立
          \item $\displaystyle a_1=\frac32,\qquad b_1=\frac12$ とし,
                \begin{align*}
                 a_{n+1}
                  &=(n+1)(a_n+b_n)\\
                  &\quad+\frac{3(a_n-b_n)}
                      {2\{3+(n+1)(a_n-b_n)\}},\\
                 b_{n+1}
                  &=(n+1)(a_n+b_n)\\
                  &\quad-\frac{3(a_n-b_n)}
                      {2\{3+(n+1)(a_n-b_n)\}}
                \end{align*}
                で2つの数列を定める.それぞれの一般項を求めよ.
          % A6: 積の3項間関係と二重の階差
          \item $\displaystyle a_1=2,\qquad a_2=16$ とし,
                \[
                  a_{n+2}a_n=2^{2n+1}a_{n+1}^{\,2}
                \]
                を満たす正の数列$\{a_n\}$の一般項を求めよ.
          % A7: 添字とともに動く基準点
          \item $\displaystyle a_1=3$ とし,
                \[
                  a_{n+1}
                   =\frac{(n+1)^2a_n+n(n+1)}
                          {a_n+n(n+1)}
                \]
                で数列$\{a_n\}$を定める.一般項を求めよ.
          % A8: 係数の和が零になる3項間関係
          \item $\displaystyle a_1=0,\qquad a_2=1$ とし,
                \begin{align*}
                 &(n+2)(n+4)a_{n+2}\\
                 &\quad-(2n^2+10n+11)a_{n+1}\\
                 &\quad+(n+1)(n+3)a_n=0
                \end{align*}
                を満たす数列$\{a_n\}$の一般項を求めよ.
          % A9: 掛け算で結び付いた二つの数列
          \item $\displaystyle a_1=b_1=2$ とし,
                \[
                 \begin{cases}
                  a_{n+1}=2^{2n+1}a_n^{\,2}b_n,\\
                  b_{n+1}=\dfrac{a_nb_n^{\,2}}{2^{2n+1}}
                 \end{cases}
                \]
                で2つの正の数列を定める.それぞれの一般項を求めよ.
          % A10: 二つの平方が作る分数式
          \item $\displaystyle a_1=\frac97$ とし,
                \[
                  a_{n+1}
                   =\frac{(a_n+1)^2+2^n(a_n-1)^2}
                          {(a_n+1)^2-2^n(a_n-1)^2}
                \]
                で数列$\{a_n\}$を定める.一般項を求めよ.
        \end{enumerate}

      \subsubsection{方針}\label{節：応用難易度：方針}\label{節：高難易度のもの：方針}
        \begin{enumerate}[label=\textbf{問\arabic*.},leftmargin=3.8em,format=\bookitemlink{practice-advanced}{hint}{problem}]
          % A1: 階乗で重みをそろえた総和
          \item 総和の中の$a_k/k!$と,前項の係数$n+1$を同時に見よう.まず$b_n=a_n/n!$とおき,その部分和を消去する.得られる3項間漸化式では,さらに$2^{n-1}$で割る.
          % A2: 添字を含む分母と不動点
          \item 係数に$n$があるので,初めから一定の公比を探してもうまくいかない.まず$a_{n+1}-1$を計算し,その逆数を取る.最後に$n$の一次式を加えて等比型にする.
          % A3: 三角形状の重みをもつ積
          \item 対数を取ると,係数が$n-k+1$の重み付き総和になる.隣接する式の差を取ると重みが$1$になり,もう一度差を取ると総和が消える.
          % A4: 非線形3項間式から階比を作る
          \item 項そのものではなく,\, $r_n=a_{n+1}/a_n$を導入する.次に$1/r_n$の漸化式を作り,定数$1$を引くと積の相殺が見える.
          % A5: 和と差で異なる構造が現れる連立
          \item 2式に共通する部分と,符号だけが違う部分を分ける.和は係数が$n$に依存する階比型で,差は逆数を取ると階差型になる.
          % A6: 積の3項間関係と二重の階差
          \item 対数を取った後,\, $x_{n+1}-x_n$を新しい数列として扱う.また,対数を取る前に$a_{n+1}/a_n$を作るルートもある.
          % A7: 添字とともに動く基準点
          \item $a_n=n$を代入すると次項は$n+1$になる.一定の不動点を探す前に,\, $c_n=a_n/n$とおいて基準をそろえよう.その後$c_n=1,-1$からの距離の比を作る.
          % A8: 係数の和が零になる3項間関係
          \item 3つの係数を符号も含めて足すと$0$になる.式を隣接する項の差でまとめ,その階差数列を階比型として解く.最後の和には部分分数分解を使う.
          % A9: 掛け算で結び付いた二つの数列
          \item 対数を取ると,対称な連立漸化式になる.和と差を取れば,指数の増え方を分離できる.対数を取る前に$a_nb_n$と$a_n/b_n$を作る方法も考えよう.
          % A10: 二つの平方が作る分数式
          \item 分子と分母の和・差を使って$a_{n+1}\pm1$を計算する.それらの比を取ると平方が現れるが,係数$2^n$も残る.その次に対数を使う.
        \end{enumerate}

      \subsubsection{解答}\label{節：応用難易度：解答}\label{節：高難易度のもの：解答}
        \begin{enumerate}[label=\textbf{問\arabic*.},leftmargin=3.8em,format=\bookitemlink{practice-advanced}{answer}{problem}]
          % A1: 階乗で重みをそろえた総和
          \item 一般項は
                \[
                  a_n=n!(n+1)2^{n-2}
                \]
                である.この式は$n=1$でも成り立つ.
          % A2: 添字を含む分母と不動点
          \item 一般項は
                \[
                  a_n
                  =1+\frac{1}{3\cdot2^{n-1}-n-1}
                \]
                である.
          % A3: 三角形状の重みをもつ積
          \item $\displaystyle
                 \alpha=\frac{3+\sqrt5}{2},\quad
                 \beta=\frac{3-\sqrt5}{2}$
                とおくと,一般項は
                \[
                  a_n=2^{\frac{\alpha^{n-1}-\beta^{n-1}}{\sqrt5}}
                \]
                である.
          % A4: 非線形3項間式から階比を作る
          \item 一般項は
                \[
                  a_n
                   =\frac{4^{n-1}\{(n-1)!\}^2}{(2n-2)!}
                \]
                である.ここでは$0!=1$とする.
          % A5: 和と差で異なる構造が現れる連立
          \item 一般項は
                \begin{align*}
                 a_n&=2^{n-1}n!+\frac3{n^2+n+4},\\
                 b_n&=2^{n-1}n!-\frac3{n^2+n+4}
                \end{align*}
                である.
          % A6: 積の3項間関係と二重の階差
          \item $\displaystyle
                 E_n=\frac{n(n-1)(2n-1)}6+2n-1$
                とおくと,
                \[
                  a_n=2^{E_n}
                \]
                である.
          % A7: 添字とともに動く基準点
          \item 一般項は
                \[
                  a_n
                   =n\,\frac{n^2+n+1}{n^2+n-1}
                \]
                である.
          % A8: 係数の和が零になる3項間関係
          \item 一般項は
                \[
                  a_n
                   =\frac{10}{3}-\frac4{n+1}-\frac4{n+2}
                \]
                である.
          % A9: 掛け算で結び付いた二つの数列
          \item 一般項は
                \begin{align*}
                 a_n&=2^{\,3^{n-1}+n^2-1},\\
                 b_n&=2^{\,3^{n-1}-n^2+1}
                \end{align*}
                である.
          % A10: 二つの平方が作る分数式
          \item $\displaystyle E_n=n+1+2^{n-1}$ とおくと,一般項は
                \[
                  a_n=\frac{2^{E_n}+1}{2^{E_n}-1}
                \]
                である.
        \end{enumerate}

      \subsubsection{解法}\label{節：応用難易度：解法}\label{節：高難易度のもの：解法}
        \begin{enumerate}[label=\textbf{問\arabic*.},leftmargin=3.8em,format=\bookitemlink{practice-advanced}{solution}{problem}]
          % A1: 階乗で重みをそろえた総和
          \item まず総和を消そうとすると,添字によって階乗も係数も変わる.そこで,総和の各項に合わせて
                \[
                  b_n=\frac{a_n}{n!},\qquad
                  T_n=\sum_{k=1}^{n}b_k
                \]
                とおく.両辺を$(n+1)!$で割れば,
                \[
                  b_{n+1}=4b_n-T_n,\qquad b_1=1.
                \]
                これより$b_2=3$である.次の式では
                \[
                  b_{n+2}=4b_{n+1}-T_{n+1}
                          =3b_{n+1}-T_n
                \]
                となるので,\, $T_n=4b_n-b_{n+1}$を代入して
                \[
                  b_{n+2}-4b_{n+1}+4b_n=0
                \]
                を得る.特性方程式は$(r-2)^2=0$であり,同じ等比数列を2本作っても情報は増えない.ここでは
                \[
                  c_n=\frac{b_n}{2^{n-1}}
                \]
                とおくと,
                \[
                  c_{n+2}-2c_{n+1}+c_n=0.
                \]
                したがって$c_{n+1}-c_n$は一定である.初めの2項は$c_1=1$,\, $c_2=3/2$だから,
                \[
                  c_n=1+\frac{n-1}{2}=\frac{n+1}{2}.
                \]
                よって$b_n=(n+1)2^{n-2}$となり,\, $a_n=n!b_n$から一般項が得られる.階乗による正規化を最初に行うと,その後は本編の総和型と重解型の手順で解ける.
          % A2: 添字を含む分母と不動点
          \item 分子と分母の係数を個別に追うより,差を計算しよう.
                \[
                  a_{n+1}-1
                  =\frac{a_n-1}{2+n(a_n-1)}.
                \]
                $a_1>1$であり,\, $a_n>1$ならば右辺も正である.したがって帰納的に$a_n>1$が成り立ち,元の分母も常に正となる.
                
                ここで
                \[
                  b_n=\frac{1}{a_n-1}
                \]
                とおくと,
                \[
                  b_{n+1}=2b_n+n,\qquad b_1=1.
                \]
                逆数を取ると分数型は消えるが,まだ等比型ではない.右辺の$n$を消すため,
                \[
                  c_n=b_n+n+1
                \]
                とおけば,
                \begin{align*}
                  c_{n+1}
                    &=2b_n+n+(n+1)+1\\
                    &=2(b_n+n+1)=2c_n.
                \end{align*}
                $c_1=3$だから$c_n=3\cdot2^{n-1}$であり,
                \[
                  b_n=3\cdot2^{n-1}-n-1.
                \]
                逆数を戻して一般項を得る.なお,\, $b_1=1$で,\, $b_{n+1}=2b_n+n>0$なので,得られた式の分母も$0$にならない.
                
                この問題では,不動点への移動,逆数化,一次式の移動を別々の目的で使っている.一度の変換で解ける形にならなくても,次に消すべき項を見極めればよい.

                \noindent 【別解】\\
                $b_n$の式を$2^n$で割ると,
                \[
                 \frac{b_{n+1}}{2^n}
                  -\frac{b_n}{2^{n-1}}=\frac{n}{2^n}.
                \]
                したがって
                \[
                 \frac{b_n}{2^{n-1}}
                 =1+\sum_{k=1}^{n-1}\frac{k}{2^k}.
                \]
                ここで
                \[
                 \frac{k}{2^k}
                 =\frac{k+1}{2^{k-1}}-\frac{k+2}{2^k}
                \]
                だから総和は相殺し,
                \[
                 \sum_{k=1}^{n-1}\frac{k}{2^k}
                  =2-\frac{n+1}{2^{n-1}}.
                \]
                これから同じ$b_n$を得る.一次式を加える解法は,この相殺を補助数列の形で先に組み込んだものといえる.
          % A3: 三角形状の重みをもつ積
          \item 各項は正なので,
                \[
                  x_n=\log_2a_n
                \]
                とおける.対数を取ると積が和になり,
                \[
                  x_{n+1}
                   =n+\sum_{k=1}^{n}(n-k+1)x_k,
                  \qquad x_1=0
                \]
                となる.この式から$x_2=1$である.
                
                重みのある総和を一度に計算する必要はない.同じ形の式を1つずらして引くと,\, $n\ge2$について
                \[
                  x_{n+1}-x_n
                  =1+\sum_{k=1}^{n}x_k.
                \]
                $n=1$でも$x_2-x_1=1+x_1$が成り立つので,この関係はすべての$n\ge1$で使える.さらに隣接する式の差を取ると,
                \[
                  x_{n+2}-2x_{n+1}+x_n=x_{n+1}.
                \]
                したがって,
                \[
                  x_{n+2}-3x_{n+1}+x_n=0.
                \]
                特性方程式の2根を
                \[
                  \alpha=\frac{3+\sqrt5}{2},\qquad
                  \beta=\frac{3-\sqrt5}{2}
                \]
                とすると,\, $\alpha+\beta=3$,\, $\alpha\beta=1$である.よって
                \[
                  x_{n+2}-\beta x_{n+1}
                   =\alpha(x_{n+1}-\beta x_n).
                \]
                初期値$x_1=0$,\, $x_2=1$に合う形は
                \[
                  x_n
                   =\frac{\alpha^{n-1}-\beta^{n-1}}{\alpha-\beta}
                   =\frac{\alpha^{n-1}-\beta^{n-1}}{\sqrt5}.
                \]
                実際,この式は上の3項間関係と2つの初期値を満たす.その関係から次項は一意に決まるので,これが$x_n$である.最後に$a_n=2^{x_n}$と戻せばよい.
                
                指数は$0,1,3,8,21,\ldots$となる.巨大な積を計算する難しさは,重み付き総和の差を2回取ることで消えている.
          % A4: 非線形3項間式から階比を作る
          \item 元の式は項の積を含むが,次数はすべて$2$である.各項は正なので,
                \[
                  r_n=\frac{a_{n+1}}{a_n}
                \]
                とおき,元の式を$a_na_{n+1}$で割ると,
                \[
                  (n+r_n)r_{n+1}=(n+1)r_n.
                \]
                したがって,
                \[
                  r_{n+1}=\frac{(n+1)r_n}{n+r_n},
                  \qquad r_1=2.
                \]
                ここで逆数$c_n=1/r_n$を取れば,
                \[
                  c_{n+1}=\frac{nc_n+1}{n+1}.
                \]
                各$n$で$c_n=1$が変わらないことに注目して,
                \[
                  d_n=c_n-1
                \]
                とおくと,
                \[
                  d_{n+1}=\frac{n}{n+1}d_n,\qquad d_1=-\frac12.
                \]
                よって積の相殺から
                \[
                  d_n=-\frac1{2n},\qquad
                  c_n=\frac{2n-1}{2n},\qquad
                  r_n=\frac{2n}{2n-1}.
                \]
                初項$a_1=1$に戻って,
                \[
                  a_n=\prod_{k=1}^{n-1}\frac{2k}{2k-1}.
                \]
                $m=n-1$とおくと分子は$2^m m!$である.分母の奇数の積については,
                \[
                  (2m)!
                   =(1\cdot3\cdots(2m-1))\,2^m m!
                \]
                なので,
                \[
                  a_n=\frac{4^m(m!)^2}{(2m)!}.
                \]
                $m=0$では空積を$1$とし,\, $a_1=1$も含められる.分数型を解く前に,階比を導入することが決め手である.

                \noindent 【別解】\\
                逆数の式で,平行移動せず直接両辺に$n+1$を掛けてもよい.
                \[
                  (n+1)c_{n+1}-nc_n=1.
                \]
                したがって$e_n=nc_n$は公差$1$の等差数列であり,\, $e_1=1/2$から
                \[
                  e_n=n-\frac12,\qquad c_n=1-\frac1{2n}
                \]
                となる.不動点を引く方法と,添字を掛けて階差を一定にする方法の両方が使える.
          % A5: 和と差で異なる構造が現れる連立
          \item 2式で共通する部分は,引けば消える.符号だけが違う部分は,足せば消える.そこで
                \[
                  s_n=a_n+b_n,\qquad d_n=a_n-b_n
                \]
                とおくと,
                \[
                  s_{n+1}=2(n+1)s_n,\qquad s_1=2,
                \]
                \[
                  d_{n+1}
                   =\frac{3d_n}{3+(n+1)d_n},\qquad d_1=1.
                \]
                $d_n>0$なら$d_{n+1}>0$であるから,帰納的に差は常に正となり,元の分母は$0$にならない.
                
                和の式は,階乗で割ると
                \[
                  \frac{s_{n+1}}{(n+1)!}=2\frac{s_n}{n!}
                \]
                となる.よって
                \[
                  s_n=2^n n!.
                \]
                一方,差の式の逆数を取って$c_n=1/d_n$とおけば,
                \[
                  c_{n+1}=c_n+\frac{n+1}{3},\qquad c_1=1.
                \]
                したがって,
                \begin{align*}
                 c_n
                  &=1+\frac13\sum_{k=1}^{n-1}(k+1)\\
                  &=1+\frac{(n-1)(n+2)}6
                   =\frac{n^2+n+4}{6}.
                \end{align*}
                これより$d_n=6/(n^2+n+4)$である.最後に
                \[
                  a_n=\frac{s_n+d_n}{2},\qquad
                  b_n=\frac{s_n-d_n}{2}
                \]
                と戻して一般項を得る.
                
                連立式全体を一つの方法で解こうとする必要はない.和と差を作った後,それぞれに適した解法を選ぶ.
          % A6: 積の3項間関係と二重の階差
          \item 正値の条件を使って$x_n=\log_2a_n$とおくと,
                \[
                  x_{n+2}-2x_{n+1}+x_n=2n+1,
                  \qquad x_1=1,\quad x_2=4.
                \]
                この式は$x_n$自体の階差ではなく,階差の階差を表している.そこで
                \[
                  d_n=x_{n+1}-x_n
                \]
                とおけば,
                \[
                  d_{n+1}-d_n=2n+1,\qquad d_1=3.
                \]
                $2n+1=(n+1)^2-n^2$なので,
                \[
                  d_n=3+(n^2-1)=n^2+2.
                \]
                さらにもう一度階差を加えて,
                \begin{align*}
                  x_n
                   &=1+\sum_{k=1}^{n-1}(k^2+2)\\
                   &=1+\frac{n(n-1)(2n-1)}6+2(n-1)\\
                   &=\frac{n(n-1)(2n-1)}6+2n-1.
                \end{align*}
                最後に$a_n=2^{x_n}$と戻せばよい.
                
                本編の隣接3項間漸化式と似た形でも,定数項が$n$に依存し,しかも左辺が二重の階差そのものになっている.この構造を見れば,特性方程式の計算に進む前に短い解法が見つかる.

                \noindent 【別解】\\
                $r_n=a_{n+1}/a_n$とおき,元の式を$a_na_{n+1}$で割ると,
                \[
                  r_{n+1}=2^{2n+1}r_n,\qquad r_1=8.
                \]
                これは階比型なので,
                \begin{align*}
                 r_n
                  &=2^{\,3+\sum_{k=1}^{n-1}(2k+1)}\\
                  &=2^{\,n^2+2}.
                \end{align*}
                もう一度比を掛け合わせて,
                \[
                 a_n
                  =2\prod_{k=1}^{n-1}2^{k^2+2}
                  =2^{\,1+\sum_{k=1}^{n-1}(k^2+2)}
                \]
                を得る.対数で差を取る方法と,元の数列で比を取る方法は,同じ構造を異なる形で使っている.
          % A7: 添字とともに動く基準点
          \item $a_n=n$なら$a_{n+1}=n+1$となるので,基準になる値も$n$に応じて変わっている.そこで
                \[
                  c_n=\frac{a_n}{n}
                \]
                とおくと,
                \[
                  c_{n+1}
                   =\frac{(n+1)c_n+1}{c_n+n+1},
                  \qquad c_1=3.
                \]
                $c_n>1$なら分母は正であり,
                \[
                  c_{n+1}-1=\frac{n(c_n-1)}{c_n+n+1}>0.
                \]
                したがって全項$c_n>1$が成り立ち,元の分母も正となる.
                
                この式では$c_n=1$と$c_n=-1$がともに動かない.両方の距離を計算すると,
                \[
                  c_{n+1}+1
                   =\frac{(n+2)(c_n+1)}{c_n+n+1}.
                \]
                同じ分母が消えるように,
                \[
                  b_n=\frac{c_n-1}{c_n+1}
                \]
                とおけば,
                \[
                  b_{n+1}=\frac{n}{n+2}b_n,\qquad b_1=\frac12.
                \]
                よって,
                \[
                  b_n
                   =\frac12\prod_{k=1}^{n-1}\frac{k}{k+2}
                   =\frac{1}{n(n+1)}.
                \]
                $c_n=(1+b_n)/(1-b_n)$だから,
                \[
                  c_n=\frac{n(n+1)+1}{n(n+1)-1}.
                \]
                最後に$a_n=nc_n$と戻す.ここでは正規化によって分数型を見つけ,さらにその分数型を階比型へ変えている.

                \noindent 【別解】\\
                $c_n$の式から$u_n=1/(c_n-1)$とおいてもよい.
                \[
                 u_{n+1}=\frac{n+2}{n}u_n+\frac1n.
                \]
                定数$1/2$を加えると,
                \[
                 u_{n+1}+\frac12
                   =\frac{n+2}{n}\left(u_n+\frac12\right).
                \]
                $u_1=1/2$から,\, $u_n+1/2=n(n+1)/2$を得る.二基準点の比を取る方法と,一基準点との差の逆数を取ってから定数を加える方法が対応している.
          % A8: 係数の和が零になる3項間関係
          \item 係数は$n$に依存するため,定数係数の特性方程式型ではない.一方,
                \[
                  2n^2+10n+11
                   =(n+2)(n+4)+(n+1)(n+3)
                \]
                なので,元の式は
                \begin{align*}
                 &(n+2)(n+4)(a_{n+2}-a_{n+1})\\
                 &\qquad=(n+1)(n+3)(a_{n+1}-a_n)
                \end{align*}
                とまとめられる.ここで$d_n=a_{n+1}-a_n$とおけば,
                \[
                  d_{n+1}
                   =\frac{(n+1)(n+3)}{(n+2)(n+4)}d_n,
                  \qquad d_1=1.
                \]
                2つの比に分けて掛け合わせると,
                \begin{align*}
                 d_n
                  &=\prod_{k=1}^{n-1}
                      \frac{k+1}{k+2}\frac{k+3}{k+4}\\
                  &=\frac2{n+1}\frac4{n+3}
                   =\frac8{(n+1)(n+3)}.
                \end{align*}
                次に元の数列を求めるには,この階差を足す.
                \[
                 d_k=4\left(\frac1{k+1}-\frac1{k+3}\right)
                \]
                だから,\, $n\ge2$について
                \begin{align*}
                 a_n
                  &=4\sum_{k=1}^{n-1}
                       \left(\frac1{k+1}-\frac1{k+3}\right)\\
                  &=4\left(\frac12+\frac13
                          -\frac1{n+1}-\frac1{n+2}\right).
                \end{align*}
                これを整理すれば一般項が得られる.得られた式は$n=1$でも$a_1=0$を満たす.階差を作って終わるのでなく,差の階比を解いた後,総和で戻るところまでが一連の手順である.

                \noindent 【別解】\\
                階差の式に両辺の係数を掛ける形のまま,
                \[
                 (n+2)(n+4)d_{n+1}
                   =(n+1)(n+3)d_n
                \]
                を眺めると,\, $e_n=(n+1)(n+3)d_n$が定数数列である.初項$e_1=8$から直ちに$d_n=8/\{(n+1)(n+3)\}$を得る.相殺する積を全部書く代わりに,相殺後に残る因子を補助数列へ組み込める.
          % A9: 掛け算で結び付いた二つの数列
          \item 各項は正なので,
                \[
                  x_n=\log_2a_n,\qquad y_n=\log_2b_n
                \]
                とおくと,
                \[
                 \begin{cases}
                  x_{n+1}=2x_n+y_n+2n+1,\\
                  y_{n+1}=x_n+2y_n-2n-1
                 \end{cases}
                \]
                となる.ここからは本編の連立型と同じように,
                \[
                  s_n=x_n+y_n,\qquad d_n=x_n-y_n
                \]
                を使う.足し算では$n$の項が消えて,
                \[
                  s_{n+1}=3s_n,\qquad s_1=2,
                \]
                引き算では
                \[
                  d_{n+1}=d_n+4n+2,\qquad d_1=0
                \]
                となる.よって$s_n=2\cdot3^{n-1}$である.また
                \[
                 4n+2=2\{(n+1)^2-n^2\}
                \]
                なので,
                \[
                 d_n=2(n^2-1).
                \]
                したがって,
                \[
                 x_n=\frac{s_n+d_n}{2}=3^{n-1}+n^2-1,
                \]
                \[
                 y_n=\frac{s_n-d_n}{2}=3^{n-1}-n^2+1.
                \]
                指数に戻して答えを得る.対数型と連立型は別々の分類に見えるが,同じ問題で順に使うことができる.

                \noindent 【別解】\\
                対数を取らずに,積と商を作っても分離できる.
                \[
                  P_n=a_nb_n,\qquad R_n=\frac{a_n}{b_n}
                \]
                とおくと,
                \[
                  P_{n+1}=P_n^3,\quad P_1=4,
                  \qquad R_{n+1}=2^{4n+2}R_n,\quad R_1=1.
                \]
                よって,
                \[
                  P_n=2^{\,2\cdot3^{n-1}},\qquad
                  R_n=2^{\,2(n^2-1)}.
                \]
                正値を使って$a_n=\sqrt{P_nR_n}$,\, $b_n=\sqrt{P_n/R_n}$と戻せば同じ答えになる.対数の和差は,元の数列の積と商に対応する.
          % A10: 二つの平方が作る分数式
          \item 分子だけを展開すると構造が見えにくくなる.分母を$D_n$と書けば,
                \[
                 a_{n+1}-1=\frac{2^{n+1}(a_n-1)^2}{D_n},
                 \qquad
                 a_{n+1}+1=\frac{2(a_n+1)^2}{D_n}.
                \]
                共通する分母を消すため,
                \[
                  b_n=\frac{a_n-1}{a_n+1}
                \]
                とおくと,
                \[
                  b_{n+1}=2^n b_n^2,\qquad b_1=\frac18.
                \]
                ここで,次項が定義できることも帰納的に確かめておこう.
                \[
                  a_n>1,\qquad 0<b_n\le2^{-n-2}
                \]
                が$n=1$で成り立つ.この評価がある$n$で成り立つと仮定すると,まず
                \[
                  D_n=(a_n+1)^2(1-2^nb_n^2)>0
                \]
                なので,\, $a_{n+1}$が定義できる.また$a_{n+1}-1$の式は正であり,\, $a_{n+1}>1$である.その後,比の式を使えば
                \[
                  0<b_{n+1}\le2^{-n-4}<2^{-(n+1)-2}
                \]
                を得る.したがって,元の数列はすべての$n$で定義でき,逆数や対数の操作も有効である.
                
                次に$x_n=-\log_2b_n$とおけば,
                \[
                  x_{n+1}=2x_n-n,\qquad x_1=3.
                \]
                $n$の一次式を消すため$y_n=x_n-n-1$とおくと,
                \[
                  y_{n+1}=2y_n,\qquad y_1=1.
                \]
                よって$x_n=n+1+2^{n-1}=E_n$であり,
                \[
                  b_n=2^{-E_n}.
                \]
                $a_n=(1+b_n)/(1-b_n)$と戻して答えを得る.分数の変形だけでも,対数だけでも解き切れない.平方の構造を作ってから対数を使う順序が大切である.
        \end{enumerate}


    \subsection{発展難易度}\label{節：発展難易度}\label{節：大学過去問ベースのもの}
      ここでは,漸化式を場合の数,確率,整数,図形,積分,極限の問題に使う10題を扱う.
      まず何を数列として追うかを選び,元の条件から漸化式を立て,最後に元の問いへ戻そう.
      一般項を求めることだけが目的ではない. 求めたい量の性質に応じて,和,差,剰余,不変量にも注目する.
      問1から問9は大学入試問題の構造・発想を参考に,問いや誘導を本書向けに再構成した演習である.
      問10は著者の掲載問題を改稿したもので,完全順列の古典的題材を扱う.
      参考資料は各問の末尾に示す. 問5は数学IIIの置換積分・部分積分・断面積による体積計算を前提とする.

      \subsubsection{問題}\label{節：発展難易度：問題}\label{節：大学過去問ベースのもの：問題}
        \begin{enumerate}[label=\textbf{問\arabic*.},leftmargin=3.8em,format=\bookitemlink{practice-entrance}{problem}{answer}]
          % F1: 正方形上の確率移動と分布の収束
          \item 正方形の頂点を時計回りに$A,B,C,D$とする.点$P$は初め$A$にあり,1秒ごとに,その頂点にとどまる,時計回りに隣の頂点へ移る,反時計回りに隣の頂点へ移る,の3つの操作から1つを選ぶ.どの操作も選ばれる確率は$\dfrac13$であり,各回の選択は独立である.
                $n$秒後に$P$が$A,C$にいる確率をそれぞれ$p_n,q_n$,\,$B,D$のいずれかにいる確率を$r_n$とする.ただし,\,$n=0,1,2,\ldots$である.
                
                \smallskip
                \noindent (1)~$p_{n+1},q_{n+1}$を$p_n,q_n,r_n$で表せ.
                
                \smallskip
                \noindent (2)~$u_n=p_n+q_n$,\,$v_n=p_n-q_n$とおく.$u_n,v_n$の一般項を求め,各頂点にいる確率を求めよ.
                
                \smallskip
                \noindent (3)~各頂点にいる確率の$n\to\infty$における極限を求めよ.
                \par\smallskip{\small\textit{\href{https://www.u-tokyo.ac.jp/content/400263415.pdf}{東京大学2024年理科第3問の状態分類を参考に再構成}}}
          % F2: 奇数が続く条件と2の因子
          \item 正の整数$a_0$から,数列$\{a_n\}$を
                \[
                  a_{n+1}=\begin{cases}
                    \dfrac{a_n}{2}&(a_n\text{が偶数}),\\[3pt]
                    \dfrac{3a_n+1}{2}&(a_n\text{が奇数})
                  \end{cases}
                \]
                によって定める.
                
                \smallskip
                \noindent (1)~$a_0,a_1,\ldots,a_m$がすべて奇数であるとき,$a_j+1$を$a_0$と$j$で表せ.ただし,\,$0\le j\le m$とする.
                
                \smallskip
                \noindent (2)~正整数$m$に対し,$a_0,a_1,\ldots,a_m$がすべて奇数となるような$a_0$をすべて求め,その最小値を答えよ.
                
                \smallskip
                \noindent (3)~$a_0+1=2^s t$と表す.ただし,\,$s$は$0$以上の整数,$t$は正の奇数である.初項$a_0$から連続して現れる奇数項の個数を求めよ.
                \par\smallskip{\small\textit{\href{https://www.kyoto-u.ac.jp/ja/admissions/undergrad/past-eq}{京都大学2024年理系第4問の偶奇分岐を参考に再構成}}}
          % F3: 畳み込み和が一定になる数列
          \item 数列$\{c_n\}$を
                \[
                  c_0=1,\qquad
                  c_{n+1}=\frac{2n+1}{2n+2}c_n
                  \quad(n=0,1,2,\ldots)
                \]
                によって定め,
                \[
                  T_n=\sum_{k=0}^{n}c_kc_{n-k}
                \]
                とする.
                
                \smallskip
                \noindent (1)~$c_n$を階乗を用いて表せ.ただし,\,$0!=1$とする.
                
                \smallskip
                \noindent (2)~正整数$n$に対し,
                \[
                  nT_n=\sum_{j=0}^{n-1}(2j+1)c_jc_{n-1-j}
                \]
                を示せ.さらに,この等式を用いて$T_n$を求めよ.
                
                \smallskip
                \noindent (3)~正整数$n$に対し,
                \[
                  \sum_{k=0}^{n}\binom{2k}{k}
                     \binom{2n-2k}{n-k}=4^n
                \]
                が成り立つことを示せ.
                \par\smallskip{\small\textit{\href{https://www.osaka-u.ac.jp/ja/admissions/faculty/general/files/9o1vy3/@@download/file}{大阪大学2025年文系第2問の畳み込み構造を参考に再構成}}}
          % F4: 内分を繰り返す平行四辺形
          \item 平行四辺形$A_0B_0C_0D_0$の頂点を反時計回りに並べる.
                一定の実数$t$が$0<t<1$を満たすとする.
                四角形$A_nB_nC_nD_n$の各辺上に,
                \[
                 \begin{aligned}
                 A_nA_{n+1}:A_{n+1}B_n&=t:(1-t),\\
                 B_nB_{n+1}:B_{n+1}C_n&=t:(1-t),\\
                 C_nC_{n+1}:C_{n+1}D_n&=t:(1-t),\\
                 D_nD_{n+1}:D_{n+1}A_n&=t:(1-t)
                 \end{aligned}
                \]
                となる点を取り,この操作を繰り返す.\,$n$は$0$以上の整数とする.
                \par\noindent
                (1) 各四角形は平行四辺形であり,対角線の交点が変わらないことを示せ.
                \par\noindent
                (2) $A_nB_nC_nD_n$の面積を$S_n$とする.\,$S_n$を$S_0,t,n$で表せ.
                \par\noindent
                (3) $n\to\infty$のとき,すべての頂点が最初の対角線の交点$O$に近づくことを示せ.
                \par\noindent
                (4) 初めの図形が辺長$\ell$の正方形で,\,$t=1/2$の場合を考える.
                操作を$n$回行った正方形の辺長を求めよ.
                また,有向線分$A_nB_n$の向きは,最初の向きからどれだけ回転するか.
                \par\smallskip{\small\textit{\href{https://www.densu.jp/tsukuba/22tsukubaspass.pdf}{筑波大学2022前期 第3問の内分操作を参考に独自に再構成}}}
          % F5: 積分を数列にして円錐を切る
          \item この問題では,数学IIIの部分積分・置換積分と,断面積による体積計算を用いる.
                正整数$n$に対して,
                \[
                 I_n=\int_0^{\pi/3}\frac{d\theta}{\cos^n\theta}
                \]
                とする.
                \par\noindent
                (1) $I_1,I_2$を求めよ.
                \par\noindent
                (2) $n\geq3$について$\,I_n$と$I_{n-2}$の関係式を求め,\,$I_3$を計算せよ.
                \par\noindent
                (3) 座標空間において,原点を頂点とし,平面$z=2$上の半径$2$の円を底面とする円錐を考える.
                その内部と表面は,
                \[
                 0\leq z\leq2,\qquad x^2+y^2\leq z^2
                \]
                で表される.この円錐のうち$x\geq1$を満たす部分の体積を求めよ.
                \par\smallskip{\small\textit{\href{https://www.nagoya-u.ac.jp/admissions/exam/upload/06_h31_sugakurikakei-seikai.pdf}{名古屋大学2019理系 第1問の構造を参考に独自に再構成}}}
          % F6: 一般項を求めずに減衰の速さを決める
          \item $\displaystyle a_1=1,\qquad
                a_{n+1}=\frac{a_n}{(1+2a_n)^2}$
                で数列$\{a_n\}$を定め,$b_n=1/a_n$とする.
                \par\noindent
                (1) すべての項が正であることを確かめ,$n\geq2$について$b_n>4n-3$を示せ.
                \par\noindent
                (2) 次の極限を求めよ.
                \[
                 \lim_{n\to\infty}\frac1n\sum_{k=1}^n a_k
                \]
                (3) $\displaystyle\lim_{n\to\infty}na_n$を求めよ.
                \par\smallskip{\small\textit{\href{https://toudainyuushi.com/wp-content/uploads/2020/06/2006S.pdf}{東京大学2006理科 第5問を参考に再構成した既存題を維持・増補}}}
          % F7: 1往復で整列できる札の配置
          \item $n$枚の札には$1,2,\ldots,n$の数字が1つずつ書かれている.
                これらを横一列に並べ, 隣り合う2枚を左から順に調べる.
                左の数字が右の数字より大きいときだけ交換し, 右端まで進む.
                続いて, 隣り合う2枚を右から順に調べ, 同じ規則で交換して左端まで戻る.
                この1往復で$1,2,\ldots,n$の順になる初期配置を\textbf{成功配置}と呼び, その個数を$c_n$とする.
                ただし,\,$n=1$では操作を行わず,\,$c_1=1$とする. 以下の問に答えよ.
                
                \smallskip
                \noindent (1)~$c_2,c_3$を求めよ. また,\,$n\ge2$として成功配置の最初の2枚の数字を$A_1,A_2$とするとき,
                少なくとも一方が$2$以下であることを示せ.
                
                \smallskip
                \noindent (2)~$n\ge3$とする. 最初の2枚に$1$を含む場合と,\,$1$を含まず$2$を含む場合に分けて数え,
                $c_n,c_{n-1},c_{n-2}$の関係式を求めよ.
                後者では, 最初の比較を済ませてから$2$の札を取り除くことを考えてよい.
                
                \smallskip
                \noindent (3)~$c_n$を求めよ.
                \par\smallskip{\small\textit{\href{https://www.u-tokyo.ac.jp/content/400239118.pdf}{東京大学2025年度理科第5問を参考に再構成}}}
          % F8: 巨大な項の整除と, 添字の互除法
          \item $\displaystyle a_1=1,\quad a_{n+1}=a_n^2+1$
                で定まる数列$\{a_n\}$について, 以下の問に答えよ.
                整数$x,y$の最大公約数を$\gcd(x,y)$と表す.
                
                \smallskip
                \noindent (1)~$a_n$が$5$の倍数となる正整数$n$をすべて求めよ.
                
                \smallskip
                \noindent (2)~任意の正整数$m,n$に対して,\,$a_m$が$a_n$を割り切ることと,
                $m$が$n$を割り切ることは同値であることを示せ.
                必要なら, 補助的に$a_0=0$と定めてよい.
                
                \smallskip
                \noindent (3)~任意の正整数$m,n$について
                \[
                 \gcd(a_m,a_n)=a_{\gcd(m,n)}
                \]
                が成り立つことを示し,\,$\gcd(a_{2026},a_{3040})$を求めよ.
                \par\smallskip{\small\textit{\href{https://www.densu.jp/tokyo/22tokyospass.pdf}{東京大学2022年度理科第2問を参考に再構成}}}
          % F9: 整数方程式の解を生み出す漸化式
          \item 正の整数$x,y,z$についての方程式
                \[
                 x^2+y^2+z^2=xyz
                \]
                を考える. 以下の問に答えよ.
                
                \smallskip
                \noindent (1)~$x=3$を満たす解$(3,y,z)$について,\,$y\le z$ならば
                $(3,z,3z-y)$も正の整数解であり,\,$3z-y>z$となることを示せ.
                
                \smallskip
                \noindent (2)~数列$\{u_n\}$を
                \[
                 u_0=u_1=3,\qquad
                 u_{n+2}=3u_{n+1}-u_n\quad(n\ge0)
                \]
                によって定める. 一般項を求めよ.
                
                \smallskip
                \noindent (3)~$(3,u_n,u_{n+1})$がすべての$n\ge0$について上の方程式の正の整数解となることを示せ.
                さらに, 異なる正の整数解が無限個あることを示し,
                $\displaystyle\lim_{n\to\infty}u_{n+1}/u_n$を求めよ.
                \par\smallskip{\small\textit{\href{https://toudainyuushi.com/wp-content/uploads/2020/06/2006S.pdf}{東京大学2006年度理科第4問を参考に独自に作成}}}
          % F10: 元のペアをすべて解消する組み替え
          \item $n$を正の整数とする. 互いに異なる$2n$個の要素
                \[
                 A_1,\ldots,A_n,\quad B_1,\ldots,B_n
                \]
                があり, 初めは$A_i$と$B_i$がペアを組んでいる.
                各ペアを$A$側と$B$側から1個ずつ選んで作り, すべての要素をちょうど1回ずつ使って$n$組のペアを作り直す.
                完成した組合せを数え, ペアを並べる順序や組み替える操作の手順は区別しない.
                元のペア$(A_i,B_i)$が1組も残らない組合せの数を$S(n)$とする.
                以下,\,$S(0)=1$と定める.
                \par\smallskip
                (1) $S(1),S(2),S(3)$を求めよ.
                \par\smallskip
                (2) $A_1$の相手を$B_j$とする.
                $A_j$の相手が$B_1$である場合と, そうでない場合に分け,
                $S(n),S(n-1),S(n-2)$の関係式を求めよ.
                \par\smallskip
                (3) $S(n)-nS(n-1)$を求め,\,$S(n)$を有限和で表せ.
                \par\smallskip
                (4) $r$を$0\le r\le n$を満たす整数とする.
                元のペアがちょうど$r$組残る組合せの数を求めよ.
                また, 全$n!$通りの組合せについて, 元のペアが残る組数の平均値を求めよ.
                \par\smallskip{\small\textit{\href{https://jukenmath.net/problems/798}{著者作問「ペアの入れ替えパターン数」}}}
        \end{enumerate}

      \subsubsection{方針}\label{節：発展難易度：方針}\label{節：大学過去問ベースのもの：方針}
        \begin{enumerate}[label=\textbf{問\arabic*.},leftmargin=3.8em,format=\bookitemlink{practice-entrance}{hint}{problem}]
          % F1: 正方形上の確率移動と分布の収束
          \item $A$への移動は,$A$にとどまる場合と,$B,D$から移る場合に分かれる.$p_n+q_n+r_n=1$を使って和の式から$r_n$を消去する.差の式では$r_n$が相殺する.$B,D$にいる確率が等しいことも,操作の対称性から確かめる.
          % F2: 奇数が続く条件と2の因子
          \item 奇数の分岐だけを使える間は,$a_{n+1}+1=\dfrac32(a_n+1)$となる.(2)では,$a_m+1$が偶数になる必要条件を求めた後,その条件から実際に各項が奇数になることを確認する.(3)では,$a_j+1$から2の因子が1つずつ減ることに注目する.
          % F3: 畳み込み和が一定になる数列
          \item (1)は階比型として係数を掛ける.(2)では,$n=k+(n-k)$と分け,\,$k$を$n-k$に入れ替えて同じ和が現れることを使う.漸化式で$kc_k$を$c_{k-1}$に直した後も,添字を逆順にした2つの式を加えると係数が一定になる.(3)は(1)の式を$T_n$に代入する.
          % F4: 内分を繰り返す平行四辺形
          \item 隣接する辺を$\boldsymbol p_n=\overrightarrow{A_nB_n}$,
                \,$\boldsymbol q_n=\overrightarrow{A_nD_n}$とする.
                面積は,内側の四角形の周囲に残る四つの三角形を引けばよい.
                収束を示すときは,\,$|\boldsymbol p_n|^2+|\boldsymbol q_n|^2$の漸化式を作る.
                正方形の場合は,辺を表す複素数への掛け算で回転を読み取る.
          % F5: 積分を数列にして円錐を切る
          \item $I_n$を$\cos^{-(n-2)}\theta$と$\cos^{-2}\theta$の積として部分積分し,
                $\sin^2\theta=1-\cos^2\theta$で整理する.
                体積は平面$z=$一定で切る.断面は半径$z$の円の弓形になるので,
                弓形の半分の中心角を$\theta$とし,$z\cos\theta=1$で置換する.
          % F6: 一般項を求めずに減衰の速さを決める
          \item 逆数を取ると,$b_{n+1}-b_n=4+4a_n$になる.
                階差を足し合わせた式を使い,まず$a_n\to0$を示す.
                平均の極限は,有限個の初期項と,十分小さくなった残りの項に分けて評価する.
                最後に$b_n/n$を調べ,$na_n$へ戻す.
          % F7: 1往復で整列できる札の配置
          \item 最初の比較の後, 左端には$\min(A_1,A_2)$が残る. この札がその後で触れられるのは, 復路の最後の比較だけである.
                個数を数える際には, 最初の2枚のどちらに$1$または$2$があるかで4つに分ける.
                $1$を取り除く場合と$2$を取り除く場合は, 小さい配置との対応が異なることに注意する.
          % F8: 巨大な項の整除と, 添字の互除法
          \item (1)では項そのものを計算する代わりに,\,$5$で割った余りに漸化式を適用する.
                (2)では法を$a_m$へ変える. $a_0=0$とおけば,\,$a_m$から先の余りは$a_0$からの余りを繰り返す.
                (3)では$n$を$m$で割った余りへ置き換える操作が, 項の最大公約数を変えないことを示す.
          % F9: 整数方程式の解を生み出す漸化式
          \item (1)では$9+z^2+(3z-y)^2-3z(3z-y)$を整理する.
                この整理で保たれる式を,\,(3)では隣接する2項に適用する.
                (2)は実数解型の3項間漸化式である. (3)の整数性は一般項を展開するよりも, 初期値と漸化式から示す方が短い.
          % F10: 元のペアをすべて解消する組み替え
          \item (2) 後者では,\,$B_1$の相手を$A_k$とする.
                $A_1,B_1$を除き,\,$A_k$と$B_j$をペアにして,\,$n-1$組の問題との対応を考える.
                (3) 3項間の関係式を,\,$S(n)-nS(n-1)$について書き直す.
                その後,\,$n!$で割って階差を作る.
                (4) 残すペアを選んでから, 残りを完全に組み替える.
                平均値は, 各元のペアが何通りの組合せに含まれるかを数えると短い.
        \end{enumerate}

      \subsubsection{解答}\label{節：発展難易度：解答}\label{節：大学過去問ベースのもの：解答}
        \begin{enumerate}[label=\textbf{問\arabic*.},leftmargin=3.8em,format=\bookitemlink{practice-entrance}{answer}{problem}]
          % F1: 正方形上の確率移動と分布の収束
          \item (1)~
                \[
                  p_{n+1}=\frac{p_n+r_n}{3},\qquad
                  q_{n+1}=\frac{q_n+r_n}{3}.
                \]
                (2)~
                \[
                \begin{aligned}
                  u_n&=\frac12+\frac12\left(-\frac13\right)^n,\\
                  v_n&=\left(\frac13\right)^n.
                \end{aligned}
                \]
                したがって
                \[
                \begin{aligned}
                  p_n&=\frac14+\frac12\left(\frac13\right)^n
                       +\frac14\left(-\frac13\right)^n,\\
                  q_n&=\frac14-\frac12\left(\frac13\right)^n
                       +\frac14\left(-\frac13\right)^n.
                \end{aligned}
                \]
                $B,D$にいる確率はそれぞれ
                \[
                  \frac14-\frac14\left(-\frac13\right)^n
                \]
                である.(3)~極限はいずれも$\dfrac14$である.
          % F2: 奇数が続く条件と2の因子
          \item (1)~
                \[
                  a_j+1=\frac{3^j}{2^j}(a_0+1)
                  \quad(0\le j\le m).
                \]
                (2)~求める初期値は
                \[
                  a_0=2^{m+1}\ell-1
                  \quad(\ell=1,2,3,\ldots)
                \]
                であり,最小値は$2^{m+1}-1$である.
                
                (3)~奇数項の個数は$s$である.$s\ge1$のとき,$a_0,\ldots,a_{s-1}$が奇数であり,$a_s=3^s t-1$は偶数である.$s=0$のときは初項から偶数である.
          % F3: 畳み込み和が一定になる数列
          \item (1)~
                \[
                  c_n=\frac{(2n)!}{4^n(n!)^2}
                     =\frac1{4^n}\binom{2n}{n}.
                \]
                (2)~設問の等式から$T_n=T_{n-1}$が得られ,$T_0=1$より
                \[
                  T_n=1\quad(n=0,1,2,\ldots).
                \]
                (3)~$c_kc_{n-k}$を(1)の式で表し,$T_n=1$の両辺を$4^n$倍すれば示される.
          % F4: 内分を繰り返す平行四辺形
          \item (1) 対角線の交点は常に$O$である.
                \par\noindent
                (2) $\rho=(1-t)^2+t^2$と置くと,
                \[
                 S_n=S_0\rho^n.
                \]
                (3) 四つの頂点はすべて$O$に近づく.
                \par\noindent
                (4) 辺長は$\ell\,2^{-n/2}$であり,有向線分の向きは反時計回りに$n\pi/4$だけ回転する.
                向きの角は$2\pi$の整数倍の違いを同一視する.
          % F5: 積分を数列にして円錐を切る
          \item (1)
                \[
                 I_1=\log(2+\sqrt3),\qquad I_2=\sqrt3.
                \]
                (2)
                \[
                 (n-1)I_n=(n-2)I_{n-2}+2^{n-2}\sqrt3,
                \]
                \[
                 I_3=\sqrt3+\frac12\log(2+\sqrt3).
                \]
                (3) 求める体積は
                \[
                 \frac{8\pi}{9}-\frac{4\sqrt3}{3}
                 +\frac13\log(2+\sqrt3).
                \]
                ここで$\log$は自然対数である.
          % F6: 一般項を求めずに減衰の速さを決める
          \item (1) すべての項は正であり,
                \[
                 b_n=4n-3+4\sum_{k=1}^{n-1}a_k>4n-3
                 \qquad(n\geq2).
                \]
                (2) 極限は$0$である.
                \par\noindent
                (3) 極限は$1/4$である.
          % F7: 1往復で整列できる札の配置
          \item (1)~$c_2=2,\ c_3=6$である.
                
                (2)~$n\ge3$に対して
                \[
                  c_n=4c_{n-1}-2c_{n-2}.
                \]
                (3)~したがって,
                \[
                  \boxed{\displaystyle
                  c_n=\frac{(2+\sqrt2)^{n-1}+(2-\sqrt2)^{n-1}}2}
                  \quad(n\ge1).
                \]
          % F8: 巨大な項の整除と, 添字の互除法
          \item (1)~$n=3,6,9,\ldots$, すなわち$3$の正の倍数である.
                
                (2)~$a_m\mid a_n\ \Longleftrightarrow\ m\mid n$である.
                
                (3)~$\gcd(a_m,a_n)=a_{\gcd(m,n)}$であり,
                \[
                 \boxed{\gcd(a_{2026},a_{3040})=a_2=2}.
                \]
          % F9: 整数方程式の解を生み出す漸化式
          \item $\displaystyle\alpha=\frac{3+\sqrt5}2,\quad
                \beta=\frac{3-\sqrt5}2$とおくと,
                \[
                 \boxed{\displaystyle
                 u_n=\frac32\left\{
                 \alpha^n+\beta^n-
                 \frac{\alpha^n-\beta^n}{\sqrt5}
                 \right\}}.
                \]
                $(3,u_n,u_{n+1})$は正の整数解であり,\,$n\ge1$で$u_{n+1}>u_n$だから異なる解が無限個得られる.
                また,
                \[
                 \boxed{\displaystyle
                 \lim_{n\to\infty}\frac{u_{n+1}}{u_n}
                 =\frac{3+\sqrt5}2}.
                \]
          % F10: 元のペアをすべて解消する組み替え
          \item (1)
                \[
                 S(1)=0,\quad S(2)=1,\quad S(3)=2.
                \]
                (2)
                \[
                 S(n)=(n-1)\{S(n-1)+S(n-2)\}
                 \quad(n\ge2).
                \]
                (3)
                \[
                 S(n)-nS(n-1)=(-1)^n,
                \]
                \[
                 \boxed{S(n)=n!\sum_{k=0}^{n}\frac{(-1)^k}{k!}}.
                \]
                (4) 組合せの数は
                \[
                 \binom{n}{r}S(n-r).
                \]
                元のペアが残る組数の平均値は$1$である.
        \end{enumerate}

      \subsubsection{解法}\label{節：発展難易度：解法}\label{節：大学過去問ベースのもの：解法}
        \begin{enumerate}[label=\textbf{問\arabic*.},leftmargin=3.8em,format=\bookitemlink{practice-entrance}{solution}{problem}]
          % F1: 正方形上の確率移動と分布の収束
          \item \textbf{(1)~直前の位置で場合を分ける.}
                $n+1$秒後に$A$にいるためには,$n$秒後に$A$にいてとどまるか,$B$にいて反時計回りに移るか,$D$にいて時計回りに移ればよい.これらは同時には起こらないから,確率を加えられる.
                $B,D$のどちらにいても,次に$A$へ移る確率は$\dfrac13$である.よって,その2頂点にいる確率の合計$r_n$を用いて
                \[
                  p_{n+1}=\frac13p_n+\frac13r_n
                \]
                と書ける.同様に
                \[
                  q_{n+1}=\frac13q_n+\frac13r_n
                \]
                である.状態をまとめてよい理由は,$B,D$から$A,C$への移動確率がそれぞれ同じだからである.
                
                \smallskip
                \textbf{(2)~和と差で連立を解く.}
                初めは$A$にいるので,
                \[
                  p_0=1,\quad q_0=0,\quad r_0=0.
                \]
                また,必ず4頂点のどれかにいるから
                \[
                  p_n+q_n+r_n=1.
                \]
                2本の漸化式を加え,$r_n=1-u_n$を代入すると
                \[
                \begin{aligned}
                  u_{n+1}&=\frac{u_n+2r_n}{3}\\
                         &=\frac23-\frac13u_n.
                \end{aligned}
                \]
                これは本編の特性方程式型である.定数の解$\dfrac12$を引くと
                \[
                  u_{n+1}-\frac12=-\frac13\left(u_n-\frac12\right).
                \]
                $u_0=1$より
                \[
                  u_n=\frac12+\frac12\left(-\frac13\right)^n.
                \]
                一方,2本の漸化式を引くと,$r_n$の項が消え,
                \[
                  v_{n+1}=\frac13v_n,\qquad v_0=1
                \]
                となる.したがって
                \[
                  v_n=\left(\frac13\right)^n.
                \]
                本編の連立漸化式で用いたように,和と差を選ぶことで,2本の数列を別々に解ける形にしている.
                $\displaystyle p_n=\frac{u_n+v_n}{2}$,\,$\displaystyle q_n=\frac{u_n-v_n}{2}$より,(2)の式を得る.
                
                $B,D$にいる確率は等しい.実際,時計回りと反時計回りをすべて入れ替えると,$A,C$を固定して$B,D$を交換する経路になる.もとの経路と確率は同じであり,この対応は逆向きにも戻せる.よって,それぞれの確率は
                \[
                  \frac{r_n}{2}=\frac{1-u_n}{2}
                  =\frac14-\frac14\left(-\frac13\right)^n
                \]
                となる.求めた4つの確率に$n=0$を代入すると,初期条件も満たす.
                
                \smallskip
                \textbf{(3)~出発点の影響が消える.}
                公比の絶対値がいずれも$\dfrac13<1$なので,
                \[
                  \left(\frac13\right)^n\to0,\qquad
                  \left(-\frac13\right)^n\to0.
                \]
                したがって,各頂点にいる確率はすべて$\dfrac14$に近づく.
                対称な図形であっても,初めは$A$に確率が集中している.その偏りが時間とともに消えることを,漸化式から示せたのである.また,$u_n-\dfrac12$は符号を交互に変える.対頂点の組$\{A,C\}$にいる確率は,$\dfrac12$の上下を行き来しながら近づく.
          % F2: 奇数が続く条件と2の因子
          \item \textbf{(1)~使われる分岐を固定する.}
                偶奇で式が変わるので,初めから1本の式で全項を求めようとすると難しい.ここでは,奇数が続く区間に限って考える.
                $a_j$が奇数である間は
                \[
                  a_{j+1}=\frac32a_j+\frac12.
                \]
                本編の特性方程式型と同じく,定数の解$-1$を引けば
                \[
                  a_{j+1}+1=\frac32(a_j+1)
                \]
                となる.したがって,等比数列の式から
                \[
                  a_j+1=\frac{3^j}{2^j}(a_0+1)
                  \quad(0\le j\le m)
                \]
                を得る.この式を使える範囲は,奇数の分岐が使われる範囲である.
                
                \smallskip
                \textbf{(2)~必要条件を求め,逆向きに確認する.}
                まず,$a_0,\ldots,a_m$がすべて奇数であると仮定する.(1)より
                \[
                  3^m(a_0+1)=2^m(a_m+1).
                \]
                $a_m+1$は偶数なので,右辺は$2^{m+1}$で割り切れる.$3^m$と$2^{m+1}$は互いに素だから,$a_0+1$も$2^{m+1}$で割り切れる.よって
                \[
                  a_0=2^{m+1}\ell-1
                  \quad(\ell\text{は正整数})
                \]
                が必要である.
                
                次に,この形の$a_0$から始めれば条件を満たすことを示す.
                \[
                  b_j=3^j2^{m+1-j}\ell-1
                  \quad(0\le j\le m)
                \]
                とおく.$m+1-j\ge1$なので,$b_j$はすべて正の奇数である.また,$b_0=a_0$であり,$0\le j<m$で
                \[
                \begin{aligned}
                  \frac{3b_j+1}{2}
                    &=3^{j+1}2^{m-j}\ell-1\\
                    &=b_{j+1}
                \end{aligned}
                \]
                が成立する.したがって,各段階で実際に奇数の分岐が使われ,$a_j=b_j$となる.これで十分性も示された.
                
                $\ell$が正整数であることから,最小の初期値は$\ell=1$のときの$2^{m+1}-1$である.単に仮定の下で式を求めるだけでなく,その式が仮定に矛盾しないことを逆向きに確かめる点が大切である.
                
                \smallskip
                \textbf{(3)~奇数区間の長さを決める.}
                $s=0$なら$a_0+1=t$は奇数であり,$a_0$は偶数なので答えは$0$である.
                $s\ge1$とする.$0\le j\le s$について
                \[
                  b_j=3^j2^{s-j}t-1
                \]
                とおく.$j<s$なら$2^{s-j}$は偶数だから,$b_j$は正の奇数である.先ほどと同じ計算で,これらの奇数の分岐から$b_{j+1}$が得られるので,$a_j=b_j$である.最後に
                \[
                  a_s=3^s t-1
                \]
                は偶数となる.よって,初項から続く奇数項は$a_0$から$a_{s-1}$までの$s$個である.
                
                例えば,$a_0=7$なら$a_0+1=2^3$なので,奇数項は3個続く.実際,
                \[
                  7,\ 11,\ 17,\ 26,\ldots
                \]
                となる.見た目には偶奇で複雑に分岐する数列でも,最初の奇数区間の長さは$a_0+1$に含まれる2の因子だけで決まる.
          % F3: 畳み込み和が一定になる数列
          \item \textbf{(1)~各項は階比型で求められる.}
                $n\ge1$に対し,係数を掛け合わせると
                \[
                  c_n=\prod_{j=0}^{n-1}\frac{2j+1}{2j+2}.
                \]
                分母は$2\cdot4\cdots2n=2^n n!$である.分子の奇数の積は,$(2n)!$から偶数の積を取り除いて
                \[
                  1\cdot3\cdots(2n-1)=\frac{(2n)!}{2^n n!}
                \]
                と表せる.よって
                \[
                  c_n=\frac{(2n)!}{4^n(n!)^2}.
                \]
                この式は$n=0$でも$c_0=1$を与える.
                
                \smallskip
                \textbf{(2)~求めたい和そのものを追う.}
                一般項は分かったが,これを直ちに$T_n$に代入すると,階乗を含む積の和になる.ここでは各項を展開するよりも,$T_n$と1つ前の$T_{n-1}$の関係を探す.
                まず,$n=k+(n-k)$を各項に掛けて
                \[
                \begin{aligned}
                  nT_n
                   &=\sum_{k=0}^{n}kc_kc_{n-k}\\
                   &\quad+\sum_{k=0}^{n}(n-k)c_kc_{n-k}.
                \end{aligned}
                \]
                第2の和は,添字を$\ell=n-k$に変えると
                \[
                  \sum_{\ell=0}^{n}\ell c_{n-\ell}c_\ell
                \]
                となり,第1の和と同じである.したがって
                \[
                  nT_n=2\sum_{k=1}^{n}kc_kc_{n-k}.
                \]
                ここで,$k=0$の項は$0$なので取り除いた.
                もとの漸化式を1つ前の添字で書くと
                \[
                  2kc_k=(2k-1)c_{k-1}\quad(k\ge1).
                \]
                これを使い,$j=k-1$に変えれば
                \[
                \begin{aligned}
                  nT_n
                   &=\sum_{k=1}^{n}(2k-1)c_{k-1}c_{n-k}\\
                   &=\sum_{j=0}^{n-1}(2j+1)c_jc_{n-1-j}
                \end{aligned}
                \]
                となり,設問の等式を得る.
                
                右辺を$U_n$とする.添字を$\ell=n-1-j$に入れ替えてから,再び$\ell$を$j$と書けば
                \[
                  U_n=\sum_{j=0}^{n-1}(2n-2j-1)c_jc_{n-1-j}.
                \]
                元の式とこの式を加えると,係数の和は
                \[
                  (2j+1)+(2n-2j-1)=2n
                \]
                となる.よって
                \[
                  2U_n=2n\sum_{j=0}^{n-1}c_jc_{n-1-j}
                        =2nT_{n-1}.
                \]
                したがって,$nT_n=U_n=nT_{n-1}$である.$n\ge1$なので$n$で割り,
                \[
                  T_n=T_{n-1}
                \]
                を得る.これは差が$0$の等差型と見ることができる.$T_0=c_0^2=1$より,すべての$n\ge0$で$T_n=1$である.
                
                一般項の代入より先に,求めたい量に漸化式を立てた.これは本編の総和型と同じく,和を1本の数列として扱う考え方である.この問題では,積の中の2つの添字の和が一定であることと,その添字を交換しても積が変わらないことが鍵になった.
                
                \smallskip
                \textbf{(3)~二項係数の恒等式へ戻す.}
                (1)より
                \[
                \begin{aligned}
                  c_kc_{n-k}
                    &=\frac1{4^k}\binom{2k}{k}
                      \frac1{4^{n-k}}\binom{2n-2k}{n-k}\\
                    &=\frac1{4^n}\binom{2k}{k}
                      \binom{2n-2k}{n-k}.
                \end{aligned}
                \]
                したがって,$T_n=1$の式を$4^n$倍すると,(3)の等式が得られる.
                例えば,$n=2$では
                \[
                  1\cdot6+2\cdot2+6\cdot1=16=4^2
                \]
                である.各$c_n$は複雑な階乗で表される一方,同じ添字の和をもつ積を集めると常に$1$になる.各項の姿からは見えにくい性質を,和の漸化式が明らかにしている.
          % F4: 内分を繰り返す平行四辺形
          \item 図形が変わるたびに四つの頂点を別々に追うと,計算が増える.
                平行四辺形を決める二本の辺と,図形の中心に注目する.
                \par\noindent
                \textbf{(1) 内分をベクトルの式にする.}
                位置ベクトルを$\boldsymbol a_n,\boldsymbol b_n,
                \boldsymbol c_n,\boldsymbol d_n$とする.
                内分の条件より,
                \[
                 \begin{aligned}
                 \boldsymbol a_{n+1}&=(1-t)\boldsymbol a_n+t\boldsymbol b_n,\\
                 \boldsymbol b_{n+1}&=(1-t)\boldsymbol b_n+t\boldsymbol c_n,\\
                 \boldsymbol c_{n+1}&=(1-t)\boldsymbol c_n+t\boldsymbol d_n,\\
                 \boldsymbol d_{n+1}&=(1-t)\boldsymbol d_n+t\boldsymbol a_n.
                 \end{aligned}
                \]
                いま,$A_nB_nC_nD_n$が平行四辺形で,中心の位置ベクトルが$\boldsymbol o$なら,
                \[
                 \boldsymbol a_n+\boldsymbol c_n
                 =\boldsymbol b_n+\boldsymbol d_n=2\boldsymbol o.
                \]
                したがって,
                \[
                 \begin{aligned}
                 \boldsymbol a_{n+1}+\boldsymbol c_{n+1}
                  &=(1-t)\,2\boldsymbol o+t\,2\boldsymbol o\\
                  &=2\boldsymbol o,
                 \end{aligned}
                \]
                同様に$\boldsymbol b_{n+1}+\boldsymbol d_{n+1}=2\boldsymbol o$である.
                二本の対角線は同じ点$O$で二等分されるので,次の図形も平行四辺形で,中心は変わらない.
                初めの図形は平行四辺形だから,この性質は数学的帰納法ですべての$n$について成り立つ.
                \par\noindent
                \textbf{(2) 面積を数列にする.}
                内側の四角形の周囲には四つの三角形が残る.
                たとえば頂点$A_n$にある三角形$A_nA_{n+1}D_{n+1}$では,
                二辺がそれぞれ$|A_nB_n|$の$t$倍と$|A_nD_n|$の$1-t$倍である.
                よって面積は$\frac12t(1-t)S_n$になる.
                ほかの三つも同じ面積なので,
                \[
                 \begin{aligned}
                 S_{n+1}
                 &=S_n-4\cdot\frac12t(1-t)S_n\\
                 &=\{(1-t)^2+t^2\}S_n.
                 \end{aligned}
                \]
                ここで初めて,図形の問いが本編の等比型になった.
                \,$\rho=(1-t)^2+t^2$と置けば,
                \[
                 S_n=S_0\rho^n.
                \]
                $\rho>0$なので面積は常に正であり,各平行四辺形は退化しない.
                \par\noindent
                \textbf{(3) 面積の収束だけで終えない.}
                面積が$0$に近づくだけでは,細長い図形になる可能性を除けない.
                そこで,二本の辺の長さも調べる.
                \[
                 \boldsymbol p_n=\boldsymbol b_n-\boldsymbol a_n,
                 \qquad
                 \boldsymbol q_n=\boldsymbol d_n-\boldsymbol a_n
                \]
                とすると,内分の式と平行四辺形の性質より,
                \[
                 \begin{cases}
                 \boldsymbol p_{n+1}=(1-t)\boldsymbol p_n+t\boldsymbol q_n,\\
                 \boldsymbol q_{n+1}=-t\boldsymbol p_n+(1-t)\boldsymbol q_n.
                 \end{cases}
                \]
                連立式から一つの数列を取り出すため,長さの二乗を足した
                \[
                 E_n=|\boldsymbol p_n|^2+|\boldsymbol q_n|^2
                \]
                を考える.内積の交差項が打ち消し合うので,
                \[
                 \begin{aligned}
                 E_{n+1}
                 &=|(1-t)\boldsymbol p_n+t\boldsymbol q_n|^2\\
                 &\quad+|-t\boldsymbol p_n+(1-t)\boldsymbol q_n|^2\\
                 &=\rho E_n.
                 \end{aligned}
                \]
                よって$E_n=\rho^nE_0$である.
                $0<t<1$より$0<\rho=1-2t(1-t)<1$である.
                二乗は非負だから,両方の辺の長さは$0$へ近づく.
                また,$O$は対角線の中点なので,
                \[
                 \overrightarrow{OA_n}=-\frac{\boldsymbol p_n+\boldsymbol q_n}{2},
                 \qquad
                 |OA_n|\leq\frac{|\boldsymbol p_n|+|\boldsymbol q_n|}{2}.
                \]
                右辺は$0$に近づく.ほかの頂点についても$\boldsymbol p_n,\boldsymbol q_n$の符号が変わるだけだから,
                同じ評価ができる.したがって四つの頂点はすべて$O$に近づく.
                \par\noindent
                \textbf{(4) 正方形では複素数の積になる.}
                正方形が反時計回りに並んでいるとき,$\boldsymbol q_n$は$\boldsymbol p_n$を反時計回りに$\pi/2$回転したものである.
                それぞれを表す複素数を$z_n,w_n$とすると,$w_n=iz_n$である.
                \,$t=1/2$を連立式に入れると,
                \[
                 z_{n+1}=\frac{1+i}{2}z_n,
                 \qquad
                 w_{n+1}=iz_{n+1}.
                \]
                したがって次も正方形であり,
                \[
                 z_n=\left(\frac{1+i}{2}\right)^n z_0.
                \]
                ここで,
                \[
                 \frac{1+i}{2}
                 =\frac1{\sqrt2}\left(\cos\frac\pi4+i\sin\frac\pi4\right).
                \]
                ド・モアブルの定理より,
                \[
                 z_n=2^{-n/2}
                 \left(\cos\frac{n\pi}{4}+i\sin\frac{n\pi}{4}\right)z_0.
                \]
                よって辺長は$\ell\,2^{-n/2}$,有向線分の回転角は反時計回りに$n\pi/4$である.
                \par
                正方形では内分という操作が,中心を固定した回転と縮小そのものになる.
                一般の平行四辺形については,二本の辺を混ぜる連立式と面積の等比性が成り立つが,
                図形全体が毎回同じ平面回転と一様な縮小を受けるとは限らない.
          % F5: 積分を数列にして円錐を切る
          \item ここでは数列の項が積分で定義されている.
                各積分を別々に計算するのでなく,べきの指数が異なる積分どうしの関係を作る.
                \par\noindent
                \textbf{(1) 帰着先となる積分を求める.}
                $I_1$では$u=\sin\theta$と置く.
                \,$0\leq\theta\leq\pi/3$では$\cos\theta>0$だから,
                \[
                 \begin{aligned}
                 I_1&=\int_0^{\sqrt3/2}\frac{du}{1-u^2}\\
                 &=\frac12\int_0^{\sqrt3/2}
                       \left(\frac1{1+u}+\frac1{1-u}\right)du\\
                 &=\frac12\left[\log\frac{1+u}{1-u}\right]_0^{\sqrt3/2}\\
                 &=\frac12\log\frac{2+\sqrt3}{2-\sqrt3}
                 =\log(2+\sqrt3).
                 \end{aligned}
                \]
                最後には$(2+\sqrt3)(2-\sqrt3)=1$を用いた.
                また,
                \[
                 I_2=[\tan\theta]_0^{\pi/3}=\sqrt3.
                \]
                \par\noindent
                \textbf{(2) 二つ前の積分へ帰着させる.}
                $(\tan\theta)'=1/\cos^2\theta$を利用して,
                \[
                 I_n=\int_0^{\pi/3}
                      \frac1{\cos^{n-2}\theta}\cdot\frac1{\cos^2\theta}\,d\theta
                \]
                を部分積分する.
                \,$n\geq3$のとき,
                \[
                 \begin{aligned}
                 I_n
                 &=\left[\frac{\sin\theta}{\cos^{n-1}\theta}\right]_0^{\pi/3}\\
                 &\quad -(n-2)\int_0^{\pi/3}
                           \frac{\sin^2\theta}{\cos^n\theta}\,d\theta\\
                 &=2^{n-2}\sqrt3-(n-2)(I_n-I_{n-2}).
                 \end{aligned}
                \]
                したがって,
                \[
                 (n-1)I_n=(n-2)I_{n-2}+2^{n-2}\sqrt3.
                \]
                偶数番目は$I_2$,奇数番目は$I_1$へ順に帰着できる.
                一般項を求めなくても,必要な積分を漸化式で計算できる点が重要である.
                \,$n=3$を入れると,
                \[
                 2I_3=I_1+2\sqrt3,
                 \qquad
                 I_3=\sqrt3+\frac12\log(2+\sqrt3).
                \]
                \par\noindent
                \textbf{(3) 空間図形から必要な積分を見つける.}
                高さ$z$で水平に切った円錐の断面は,半径$z$の円$x^2+y^2\leq z^2$である.
                \,$z<1$のとき$x\geq1$の部分はなく,$1\leq z\leq2$のときだけ弓形が現れる.
                弓形の円弧の中点の方向を$x$軸の正方向とし,弓形の半分の中心角を$\theta$とする.
                \[
                 z\cos\theta=1,\qquad 0\leq\theta\leq\frac\pi3.
                \]
                \begin{center}
                \begin{tikzpicture}[scale=0.78]
                 \fill[gray!20] (1,{-sqrt(3)}) arc[start angle=-60,end angle=60,radius=2] -- cycle;
                 \draw (0,0) circle[radius=2];
                 \draw[thick] (1,{-sqrt(3)}) -- (1,{sqrt(3)});
                 \draw[dashed] (0,0) -- (1,{sqrt(3)}) node[midway,left] {$z$};
                 \draw[dashed] (0,0) -- (1,{-sqrt(3)});
                 \draw[->] (-2.2,0) -- (2.5,0) node[right] {$x$};
                 \draw (0.45,0) arc[start angle=0,end angle=60,radius=0.45];
                 \node at (0.55,0.30) {$\theta$};
                 \node[below left] at (0,0) {$O$};
                 \node[below right] at (1,0) {$H$};
                 \node[above right] at (1,{sqrt(3)}) {$E$};
                 \node[below right] at (1,{-sqrt(3)}) {$F$};
                 \node[below] at (0.5,0) {$1$};
                 \draw (0.84,0) -- (0.84,0.16) -- (1,0.16);
                \end{tikzpicture}
                \end{center}
                図は水平断面で,網掛け部分が求める立体の断面である.
                \,$O$はこの断面の円の中心を表す.
                その面積を$A(z)$とすると,扇形から三角形を引いて,
                \[
                 \begin{aligned}
                 A(z)&=\frac12z^2\cdot2\theta
                      -\frac12\cdot2\sqrt{z^2-1}\cdot1\\
                     &=z^2\theta-\sqrt{z^2-1}.
                 \end{aligned}
                \]
                よって体積は$V=\int_1^2 A(z)\,dz$である.
                \,$z=1/\cos\theta$と置けば,
                \[
                 \sqrt{z^2-1}=\tan\theta,
                 \qquad dz=\frac{\sin\theta}{\cos^2\theta}\,d\theta.
                \]
                したがって,
                \[
                 \begin{aligned}
                 V&=\int_0^{\pi/3}\frac{\theta\sin\theta}{\cos^4\theta}\,d\theta\\
                  &\quad-\int_0^{\pi/3}\frac{\sin^2\theta}{\cos^3\theta}\,d\theta.
                 \end{aligned}
                \]
                第一項では$(\cos^{-3}\theta)'=3\sin\theta/\cos^4\theta$を使って部分積分する.
                第二項では$\sin^2\theta=1-\cos^2\theta$を使うと,
                \[
                 \begin{aligned}
                 V&=\left[\frac{\theta}{3\cos^3\theta}\right]_0^{\pi/3}
                      -\frac13I_3-(I_3-I_1)\\
                  &=\frac{8\pi}{9}-\frac43I_3+I_1\\
                  &=\frac{8\pi}{9}-\frac{4\sqrt3}{3}
                      +\frac13\log(2+\sqrt3).
                 \end{aligned}
                \]
                断面の半径と切断位置の関係を角度で表すと,最初に調べた$I_3$が自然に現れた.
                積分の漸化式は数列だけの技法ではなく,図形から生まれる積分を計算するためにも使える.
          % F6: 一般項を求めずに減衰の速さを決める
          \item 分母に$a_n$があるが,本編の分数型と違って分母が二乗になっている.
                まず逆数を取って,一般項が求められるかだけでなく,増え方を評価できるかを調べる.
                \par\noindent
                \textbf{(1) 逆数の階差を足し合わせる.}
                $a_1=1>0$であり,$a_n>0$なら$(1+2a_n)^2>0$なので$a_{n+1}>0$である.
                よって数学的帰納法からすべての項が正で,$b_n=1/a_n$も定義できる.
                逆数を取ると,
                \[
                 \begin{aligned}
                 b_{n+1}&=\frac{(1+2a_n)^2}{a_n}\\
                         &=b_n+4+4a_n.
                 \end{aligned}
                \]
                したがって$b_{n+1}-b_n=4+4a_n>4$である.
                \,$k=1,\ldots,n-1$について足すと,
                \[
                 b_n-b_1=4(n-1)+4\sum_{k=1}^{n-1}a_k.
                \]
                $b_1=1$より,
                \[
                 b_n=4n-3+4\sum_{k=1}^{n-1}a_k.
                \]
                $n\geq2$なら右辺の総和は正だから,$b_n>4n-3$となる.
                ここまでの変形は,逆数化したあとに階差を総和する本編の技法である.
                ただし総和に$a_k$が残っているため,この式だけでは一般項は得られていない.
                \par\noindent
                \textbf{(2) 項が小さくなることを平均へ移す.}
                正の数の逆数を取ると不等号が逆向きになるので,
                \[
                 0<a_n<\frac1{4n-3}\qquad(n\geq2).
                \]
                挟み撃ちの原理より$a_n\to0$である.
                しかし平均にも$n$が入っているので,項の極限だけで終えず,平均を直接評価する.
                任意の正数$\varepsilon$を取る.
                \,$a_n\to0$より,$k\geq N$なら$0<a_k<\varepsilon$となる正整数$N$が存在する.
                \,$C=\sum_{k=1}^{N-1}a_k$と置けば,$C$は$n$によらない定数である.
                \,$n\geq N$のとき,
                \[
                 \begin{aligned}
                 0<\frac1n\sum_{k=1}^{n}a_k
                 &=\frac Cn+\frac1n\sum_{k=N}^{n}a_k\\
                 &\leq\frac Cn+\frac{n-N+1}{n}\varepsilon\\
                 &\leq\frac Cn+\varepsilon.
                 \end{aligned}
                \]
                十分大きい$n$では$C/n<\varepsilon$でもあるため,平均は$2\varepsilon$より小さい.
                \,$\varepsilon$はいくらでも小さく取れるので,
                \[
                 \lim_{n\to\infty}\frac1n\sum_{k=1}^{n}a_k=0.
                \]
                この証明では,大きかった最初の有限個の項の影響は$n$で割ると消え,
                残りの項はすべて小さい,という二つの点を使っている.
                \par\noindent
                \textbf{(3) 逆数の増え方から元の減衰を読む.}
                (1)の等式を$n$で割ると,
                \[
                 \frac{b_n}{n}=4-\frac3n+
                            \frac4n\sum_{k=1}^{n-1}a_k.
                \]
                (2)より,
                \[
                 0\leq\frac1n\sum_{k=1}^{n-1}a_k
                 \leq\frac1n\sum_{k=1}^{n}a_k\longrightarrow0.
                \]
                したがって$b_n/n\to4$であり,
                \[
                 na_n=\frac1{b_n/n}\longrightarrow\frac14.
                \]
                \par
                最初に求めた上界は,$a_n\to0$を示すための中間結果である.
                それによって階差の追加項$4a_n$の平均が消え,$b_n$は最終的に$4n$と同じ割合で増えると分かった.
                結論$na_n\to1/4$は,$a_n$の減少の速さが$1/(4n)$のそれに近づくことを表す.
                このように,一般項を閉じた式で表さなくても,漸化式から数列の主要な性質を決められる.
          % F7: 1往復で整列できる札の配置
          \item \textbf{(1)~小さい例と, 左端に残る札.}
                $2$枚なら, 昇順の配置はそのまま, 降順の配置は最初の比較で昇順になる.
                よって$c_2=2$である.
                $3$枚では往路の2回の比較で$3$が右端に移る.
                復路では右端の比較は交換を起こさず, 残る2枚が最後の比較で昇順になる.
                したがってすべての配置が成功し,\,$c_3=3!=6$である.
                
                一般に, 最初の比較の直後, 左端の札は
                \[
                  b=\min(A_1,A_2)
                \]
                である. 往路の残りの比較は左端に触れず, 復路でも左端を含む比較は最後に1度だけ行う.
                したがって$b$は最終的に左から1番目か2番目に位置する.
                成功配置ならその数字は$1$または$2$でなければならない.
                これで$\min(A_1,A_2)\le2$が示された.
                ただし, この条件だけで成功が決まるわけではない. 残りの札の並びも調べる必要がある.
                
                \smallskip
                \textbf{(2)~取り除いた札を戻せる対応を作る.}
                比較・交換の結果は, 数字の大小関係だけで決まる.
                そこで札を1枚取り除いた際には, 残りの数字を小さい順に$1,2,\ldots,n-1$と読み替える.
                この読み替えは比較・交換の挙動を変えない.
                
                まず, 成功配置を次の4種類に分ける.
                \[
                \begin{array}{ll}
                \text{\textcircled{1}}&A_1=1,\\
                \text{\textcircled{2}}&A_2=1,\\
                \text{\textcircled{3}}&A_1=2,\ A_2\ge3,\\
                \text{\textcircled{4}}&A_2=2,\ A_1\ge3.
                \end{array}
                \]
                これらは互いに重ならず,\,(1)よりすべての成功配置を尽くす.
                
                \textbf{\textcircled{1},\textcircled{2}の数え方.}
                \textcircled{1}では左端の$1$を取り除く.
                \textcircled{2}では最初の比較によって$1$が左端へ移るので, その$1$を取り除く.
                どちらの場合も, その後の往路は残る$n-1$枚の往路と同じである.
                復路も残る$n-1$枚についての操作を行い, 最後の$1$との比較では交換しない.
                よって元の配置が成功することと, 縮約した配置が成功することは同値である.
                
                逆に,\,$n-1$枚の成功配置の数字をすべて$1$増やし,\,$1$を先頭に挿入すれば\textcircled{1}の配置が得られる.
                $1$を先頭から2番目に挿入すれば\textcircled{2}の配置が得られる.
                いずれも逆操作は一意であるから, 各種類は$c_{n-1}$通りである.
                
                \textbf{\textcircled{3}の数え方.}
                この場合は最初の2枚が$(2,A_2)$で,\,$A_2\ge3$だから最初の比較では交換しない.
                左端の$2$を取り除き, 数字$1$はそのまま,\,$3,4,\ldots,n$をそれぞれ$2,3,\ldots,n-1$に読み替える.
                すると縮約配置の先頭は$1$ではない.
                
                ここで, 元の配置と縮約配置の操作を対応させる.
                元の配置の往路が終わると, 配列は
                \[
                 (2,B_1,B_2,\ldots,B_{n-1})
                \]
                となる. 後ろの$n-1$枚は, 縮約配置の往路後の配列の数字を元に戻したものに等しい.
                復路では, まずこの後ろの$n-1$枚に対し, 縮約問題と同じ比較・交換が行われる.
                
                右から左へ進む1回の操作では, 最小の札$1$は必ずその部分の左端へ移る.
                したがって復路の最後の比較の直前には, 元の配列は
                \[
                 (2,1,D_2,\ldots,D_{n-1})
                \]
                となり, 最後に$(2,1)$が交換される.
                この最終配列が昇順になることは, 縮約問題の最終配列が昇順になることと同値である.
                
                逆に, 先頭が$1$でない$n-1$枚の成功配置で, 数字$2$以上を$1$ずつ増やし,
                先頭に$2$を挿入すれば, 必ず\textcircled{3}の成功配置となる.
                よって求める個数は,\,$n-1$枚の成功配置のうち先頭が$1$でないものの個数である.
                先頭が$1$である成功配置は, 上で示した\textcircled{1}の対応により$c_{n-2}$通りだから,
                \textcircled{3}は$c_{n-1}-c_{n-2}$通りである.
                
                \textbf{\textcircled{4}の数え方.}
                \textcircled{3}の最初の2枚を交換すると,\,$(A_1,2)$で$A_1\ge3$という\textcircled{4}の配置になる.
                最初の比較を済ませれば両者は同じ配列となるので, 成否も一致する.
                この交換は互いに逆だから,\textcircled{4}も$c_{n-1}-c_{n-2}$通りである.
                
                以上を足し合わせると,
                \begin{align*}
                 c_n
                 &=2c_{n-1}+2(c_{n-1}-c_{n-2})\\
                 &=4c_{n-1}-2c_{n-2}\quad(n\ge3).
                \end{align*}
                ここで引いた$c_{n-2}$は, 足し算で生じた曖昧な重複ではなく,
                \textcircled{3}の縮約先から除かなければならない「先頭が$1$の配置」の数である.
                
                \smallskip
                \textbf{(3)~数える問題から, 3項間漸化式へ.}
                特性方程式は
                \[
                 x^2-4x+2=0
                \]
                であり, 根は$2+\sqrt2,2-\sqrt2$である.
                したがって定数$P,Q$を用いて
                \[
                 c_n=P(2+\sqrt2)^{n-1}+Q(2-\sqrt2)^{n-1}
                \]
                と書ける. $c_1=1,c_2=2$より
                \[
                 P+Q=1,\qquad
                 (2+\sqrt2)P+(2-\sqrt2)Q=2.
                \]
                これを解くと$P=Q=1/2$なので,
                \[
                 c_n=\frac{(2+\sqrt2)^{n-1}+(2-\sqrt2)^{n-1}}2.
                \]
                例えば$c_4=20,c_5=68$となる.
                無理数を含む式で配置の個数という整数が表されるのは,
                2つのべきの展開に現れる$\sqrt2$の奇数次の項が相殺されるためである.
                この問題では, 本編の解法を使うまでの「何を数え, 何を取り除くか」の設計が中心である.
          % F8: 巨大な項の整除と, 添字の互除法
          \item \textbf{(1)~余りも漸化式に従う.}
                この数列の最初の項は
                \[
                 1,2,5,26,677,\ldots
                \]
                であり, 値を直接計算して調べ続ける方法はすぐに難しくなる.
                一方,\,$a_n$を$5$で割った余りが$r$なら,
                $a_{n+1}$の余りは$r^2+1$を$5$で割った余りである.
                そこで余りだけを調べると,
                \[
                 1\longmapsto2\longmapsto0\longmapsto1
                \]
                となる. 同じ余りに戻れば, それ以後の余りも同じ規則で繰り返される.
                したがって
                \[
                 a_n\equiv
                 \begin{cases}
                 1 & (n\equiv1\pmod3),\\
                 2 & (n\equiv2\pmod3),\\
                 0 & (n\equiv0\pmod3)
                 \end{cases}
                 \quad(\bmod5).
                \]
                よって$a_n$が$5$の倍数となるのは,\,$n$が$3$の正の倍数であるときに限る.
                
                \smallskip
                \textbf{(2)~法を数列の項にする.}
                補助的に$a_0=0$と定める.
                すると$a_1=a_0^2+1$も成り立つので, 漸化式は$n\ge0$で使える.
                任意の正整数$m$を固定し,
                \[
                 a_{m+j}\equiv a_j\pmod{a_m}
                 \quad(j\ge0)
                \]
                を示そう.
                $j=0$では両辺が$0$と合同である.
                $j$で成り立つと仮定すると,
                \begin{align*}
                 a_{m+j+1}
                 &=a_{m+j}^2+1\\
                 &\equiv a_j^2+1=a_{j+1}\pmod{a_m}.
                \end{align*}
                よって数学的帰納法により, すべての$j\ge0$で成り立つ.
                
                ここで$n=qm+r\ (0\le r<m)$と割り算すると,
                上の合同式を繰り返し使うことで
                \[
                 a_n\equiv a_r\pmod{a_m}
                \]
                が得られる.
                さらに$a_0=0,a_1=1$であり,\,$n\ge1$では
                \[
                 a_{n+1}-a_n=a_n(a_n-1)+1>0
                \]
                だから, 数列は$a_0$から厳密に増加する.
                したがって$0\le r<m$に対して
                \[
                 0\le a_r<a_m,
                \]
                しかも$a_r=0$となるのは$r=0$のときだけである.
                よって
                \begin{align*}
                 a_m\mid a_n
                 &\Longleftrightarrow a_r=0\\
                 &\Longleftrightarrow r=0\\
                 &\Longleftrightarrow m\mid n.
                \end{align*}
                この議論は$m=1$でも成り立つ. このとき余りの添字は必ず$r=0$で,
                $a_1=1$はすべての整数を割り切る.
                なお,\,$a_3=5$であるから,\,(1)はこの結果の$m=3$という特別な場合でもある.
                
                \smallskip
                \textbf{(3)~項の互除法を, 添字の互除法へ移す.}
                再び$n=qm+r\ (0\le r<m)$とおく.
                $a_n\equiv a_r\pmod{a_m}$であるから,\,$a_n$と$a_r$の差は$a_m$の整数倍である.
                よって$a_m$と$a_n$の共通の約数は,\,$a_m$と$a_r$の共通の約数に一致する.
                したがって
                \[
                 \gcd(a_m,a_n)=\gcd(a_m,a_r).
                \]
                これは, 添字に対して行うユークリッドの互除法に沿って,
                項の最大公約数を保ったまま添字を小さくできることを意味する.
                添字の互除法は,\,$d=\gcd(m,n)$と$0$で終わるので,
                \[
                 \gcd(a_m,a_n)=\gcd(a_d,a_0)=a_d.
                \]
                最後の等号は,\,$a_0=0$であり,\,$\gcd(x,0)=x$だからである.
                
                求める例では,
                \begin{align*}
                 3040&=2026+1014,\\
                 2026&=1014+1012,\\
                 1014&=1012+2,\\
                 1012&=506\cdot2
                \end{align*}
                だから,\,$\gcd(2026,3040)=2$である.
                従って
                \[
                 \gcd(a_{2026},a_{3040})=a_2=2.
                \]
                項がどれほど大きくなっても, 添字の小さな整数計算だけでこの最大公約数が決まる.
                漸化式は項の値を求めるだけでなく, 整除のような「項どうしの関係」を取り出す道具にもなる.
          % F9: 整数方程式の解を生み出す漸化式
          \item \textbf{(1)~1つの解から次の解へ.}
                $x=3$のとき, 方程式は
                \[
                 9+y^2+z^2=3yz
                \]
                となる. $z$を固定して$y$の2次方程式とみれば,
                もう1つの根は$3z-y$であることが, 根と係数の関係から予想できる.
                実際に確かめると,
                \begin{align*}
                 &9+z^2+(3z-y)^2-3z(3z-y)\\
                 &\qquad=9+y^2+z^2-3yz=0.
                \end{align*}
                よって$(3,z,3z-y)$も方程式を満たす.
                しかも$3z-y$は整数であり,\,$0<y\le z$より
                \[
                 3z-y\ge2z>z>0.
                \]
                したがって正の整数解である.
                この操作で新しくできた2つの数も$z<3z-y$を満たすため, 同じ操作を繰り返せる.
                
                \smallskip
                \textbf{(2)~解を作る操作を数列にする.}
                最初の解$(3,3,3)$から(1)の操作を繰り返すと,\,$y,z$は
                \[
                 (3,3),\ (3,6),\ (6,15),\ (15,39),\ldots
                \]
                と変化する.
                この隣接する2つの数を$u_n,u_{n+1}$とすれば, 次の数は$3u_{n+1}-u_n$となる.
                これが問題で与えた3項間漸化式の意味である.
                
                特性方程式は
                \[
                 t^2-3t+1=0
                \]
                であり, 2つの根を
                \[
                 \alpha=\frac{3+\sqrt5}2,\qquad
                 \beta=\frac{3-\sqrt5}2
                \]
                とおく.
                $u_n=P\alpha^n+Q\beta^n$と書くと,\,$u_0=u_1=3$より
                \[
                 P+Q=3,\qquad \alpha P+\beta Q=3.
                \]
                これを解けば
                \[
                 P=\frac{3(\sqrt5-1)}{2\sqrt5},\qquad
                 Q=\frac{3(\sqrt5+1)}{2\sqrt5}.
                \]
                したがって
                \begin{align*}
                 u_n
                 &=P\alpha^n+Q\beta^n\\
                 &=\frac32\left\{\alpha^n+\beta^n
                 -\frac{\alpha^n-\beta^n}{\sqrt5}\right\}.
                \end{align*}
                
                \smallskip
                \textbf{(3)~一般項とは別に, 整数性と方程式を保証する.}
                初期値は整数であり, 漸化式は整数の3倍と整数の差で次の項を作る.
                したがって数学的帰納法により, すべての$u_n$は整数である.
                また$u_0=u_1=3>0,u_2=6>u_1$である.
                $u_{n+1}>u_n>0$なら
                \[
                 u_{n+2}=3u_{n+1}-u_n>2u_{n+1}>u_{n+1}
                \]
                だから,\,$n\ge1$で$u_{n+1}>u_n>0$が成り立つ.
                
                さらに, 隣接する2項から作る量
                \[
                 I_n=u_n^2+u_{n+1}^2-3u_nu_{n+1}
                \]
                を考える.
                漸化式を代入すると,
                \begin{align*}
                 I_{n+1}
                 &=u_{n+1}^2+(3u_{n+1}-u_n)^2\\
                 &\quad{}-3u_{n+1}(3u_{n+1}-u_n)\\
                 &=u_n^2+u_{n+1}^2-3u_nu_{n+1}=I_n.
                \end{align*}
                従ってこの量は変わらず,
                \[
                 I_n=I_0=9+9-27=-9.
                \]
                よって
                \[
                 3^2+u_n^2+u_{n+1}^2=3u_nu_{n+1}.
                \]
                以上から$(3,u_n,u_{n+1})$は正の整数解である.
                $n\ge1$で$u_{n+1}$が厳密に増加するので, これらの三つ組は互いに異なり, 解は無限個ある.
                ここでは特定の操作で得られる解の系列を作っており, すべての解の分類までは要求していない.
                
                最後に比の極限を求める.
                $0<\beta<1<\alpha$であり,\,$P,Q>0$だから,
                \[
                 \frac{u_{n+1}}{u_n}
                 =\alpha\,
                 \frac{P+Q(\beta/\alpha)^{n+1}}
                      {P+Q(\beta/\alpha)^n}
                 \longrightarrow\alpha.
                \]
                一般項には無理数が現れるが, 漸化式が整数性を保証する.
                さらに整数どうしの比が無理数$\alpha$へ近づくことまで分かる.
                解法としての3項間漸化式が, 整数解を無限に生成する仕組みと, その増加の速さを同時に記述している.
          % F10: 元のペアをすべて解消する組み替え
          \item \textbf{(1) 小さい場合で数える対象を確認する.}
                $n=1$では組み替える相手がいないため,\,$S(1)=0$である.
                $n=2$では$(A_1,B_2),(A_2,B_1)$の1通りである.
                $n=3$では,\,$A_1,A_2,A_3$の相手の添字が
                $(2,3,1)$または$(3,1,2)$となる2通りなので,\,$S(3)=2$である.
                $S(0)=1$は, 要素がない場合の「何も作らない組合せ」を1通りと数える約束である.
                \par\smallskip
                \textbf{(2) 相手を1個決め, 残りを小さい問題に戻す.}
                $A_1$の相手$B_j$は,\,$j\ne1$の$n-1$通りから選ぶ.
                以後は$j$を固定する.
                \par
                まず$A_j$の相手が$B_1$なら,\,$A_1,A_j,B_1,B_j$を除いた残りの$n-2$組を,
                元のペアが残らないように作ればよい.
                これは$S(n-2)$通りである.
                \par
                次に$A_j$の相手が$B_1$でないとする.
                $B_1$の相手を$A_k$とすると,\,$k\ne1,j$である.
                2組のペア
                \[
                 (A_1,B_j),\quad(A_k,B_1)
                \]
                から$A_1,B_1$を除き,\,$(A_k,B_j)$の1組に置き換える.
                $k\ne j$なので, この新しいペアも元のペアではない.
                したがって, 添字$2,\ldots,n$の$n-1$組を完全に組み替えた形になる.
                \par
                逆に, そのような$n-1$組の組合せから$B_j$の相手$A_k$を見つけ,
                \[
                 (A_k,B_j)\ \longmapsto
                 (A_1,B_j),(A_k,B_1)
                \]
                と戻せば, 元の組合せが一意に復元できる.
                つまり, 対応は一対一であり, 後者は$S(n-1)$通りである.
                この「小さくした問題に戻せること」が, 個数の漸化式を作る根拠である.
                2つの場合は重ならないので,
                \[
                 S(n)=(n-1)\{S(n-1)+S(n-2)\}.
                \]
                \par\smallskip
                \textbf{(3) 変動する係数を, 差と階乗で整理する.}
                この漸化式は係数が$n$に依存するため, 定数係数の3項間漸化式としては解けない.
                そこで
                \[
                 T_n=S(n)-nS(n-1)\quad(n\ge1)
                \]
                とおく. $n\ge2$では(2)から
                \begin{align*}
                 T_n
                 &= (n-1)S(n-2)-S(n-1)\\
                 &= -T_{n-1}.
                \end{align*}
                $T_1=S(1)-S(0)=-1$より,\,$T_n=(-1)^n$である.
                よって
                \[
                 S(n)=nS(n-1)+(-1)^n.
                \]
                ここで$n$を階比として吸収するため,\,$P_n=S(n)/n!$とおく.
                \[
                 P_n-P_{n-1}=\frac{(-1)^n}{n!}.
                \]
                $P_0=S(0)/0!=1$より, 階差を足し合わせると
                \[
                 P_n=\sum_{k=0}^{n}\frac{(-1)^k}{k!}.
                \]
                したがって,
                \[
                 S(n)=n!\sum_{k=0}^{n}\frac{(-1)^k}{k!}.
                \]
                なお,\,$S(4)=9,\,S(5)=44$もこの式や(2)から求められる.
                本編の\textbf{等比型, 階比による正規化, 階差型}を順に使ったことになる.
                \par\smallskip
                \textbf{(4) 残る組数を指定して数える.}
                元のペアをちょうど$r$組残すには, まず残す$r$組を$\binom{n}{r}$通りから選ぶ.
                残りの$n-r$組では元のペアを1組も残せないので, その組み方は$S(n-r)$通りである.
                どの$r$組が残るかは完成形から一意に決まるため, 重複はない.
                よって答えは
                \[
                 \binom{n}{r}S(n-r)
                \]
                である. $r=n$でも$S(0)=1$によって1通りと数えられる.
                \par
                次に, 全組合せで残る元のペアの組数を合計する.
                固定したペア$(A_i,B_i)$が残る組合せは, 残りの相手を自由に割り当てる$(n-1)!$通りである.
                これを$i=1,\ldots,n$について足すと, 合計は$n(n-1)!$となる.
                元のペアが$r$組残る組合せはちょうど$r$回数えられるので, この合計でよい.
                全組合せは$n!$通りであるから, 平均値は$n(n-1)!/n!=1$である.
                \par\smallskip
                \noindent\textbf{【数学的な広がり：完全順列と確率】}\par\nopagebreak[4]
                元の位置に戻る要素のない並べ替えを\textbf{完全順列（錯置）}と呼ぶ.
                この問題はその数え上げである.
                また, 全$n!$通りを同じ確率で選ぶとき, 元のペアが1組も残らない確率は
                \[
                 P_n=\frac{S(n)}{n!}
                     =\sum_{k=0}^{n}\frac{(-1)^k}{k!}
                \]
                である. このように,\,$n!$で割る操作には, 漸化式を整理するだけでなく\textbf{個数を確率に直す意味}もある.
                \par
                さらに, この確率は$n$が大きくなると$1/e$に近づく.
                ここでは無限級数展開を仮定せず, 本編の二項定理と$e$の定義で確認しよう.
                \par
                まず$m\ge0$について,
                \begin{align*}
                 P_{2m+2}-P_{2m}
                 &= -\frac{2m+1}{(2m+2)!}<0,\\
                 P_{2m+3}-P_{2m+1}
                 &= \frac{2m+2}{(2m+3)!}>0.
                \end{align*}
                また,\,$P_{2m}-P_{2m+1}=1/(2m+1)!$である.
                したがって, 偶数番目の確率は減少し, 奇数番目の確率は増加し, 両者の差は0へ近づく.
                確率は$0$以上$1$以下なので, それぞれ収束し, 共通の極限$L$をもつ.
                \par
                その値を決めるため, 整数$M\ge2$に対して
                \[
                 b_{M,k}=\frac{\binom{M}{k}}{M^k}
                 \quad(0\le k\le M)
                \]
                とおく. すると
                \[
                 \frac{b_{M,k+1}}{b_{M,k}}
                 =\frac{M-k}{M(k+1)}\le1
                \]
                なので,\,$b_{M,k}$は$k$について非増加である.
                二項定理の有限和を,添字$k$が偶数または奇数の項で止める.
                残りを隣り合う2項ずつ組にすることで,
                $M\ge2\ell+1$に対して
                \begin{align*}
                 \sum_{k=0}^{2\ell+1}(-1)^k b_{M,k}
                 &\le\left(1-\frac1M\right)^M,\\
                 \left(1-\frac1M\right)^M
                 &\le\sum_{k=0}^{2\ell}(-1)^k b_{M,k}
                \end{align*}
                が分かる. 添字が奇数の項で止めた残りは「正の項$-$次の項」の和となり0以上で,
                添字が偶数の項で止めた残りはその符号が逆になる.
                \par
                ここで$\ell$を固定し,\,$M\to\infty$とする.
                固定した$k$について
                \[
                 b_{M,k}=\frac1{k!}
                 \prod_{j=0}^{k-1}\left(1-\frac jM\right)
                 \longrightarrow\frac1{k!}.
                \]
                $k=0$の積は$1$とする.
                一方,\,$t=M-1$とおくと,
                \[
                 \left(1-\frac1M\right)^M
                 =\left(1+\frac1t\right)^{-t-1}.
                \]
                本編の$e$の定義から,$(1+1/t)^t\to e$であり,
                $(1+1/t)\to1$なので,上の式は$1/e$へ近づく.
                よって
                \[
                 P_{2\ell+1}\le\frac1e\le P_{2\ell}.
                \]
                今度は$\ell\to\infty$とし, 挟み撃ちによって$L=1/e$を得る.
                \[
                 \boxed{\lim_{n\to\infty}\frac{S(n)}{n!}
                       =\frac1e\simeq0.368}.
                \]
                同じ偶奇の後の項は極限へ向かって厳密に進むため,
                $1/e$は隣り合う$P_n,P_{n+1}$の間に厳密に入る. よって$n\ge1$で
                \begin{align*}
                 \left|P_n-\frac1e\right|
                 &<\frac1{(n+1)!},\\
                 \left|S(n)-\frac{n!}{e}\right|
                 &<\frac1{n+1}\le\frac12.
                \end{align*}
                下段は上段を$n!$倍して得られる.
                つまり,\,$S(n)$は$n!/e$に最も近い整数である.
                場合の数を数える漸化式から, 階乗, 確率, 自然対数の底が一つにつながった.
        \end{enumerate}



\end{bookcolumns}

% chapter: 付録（問いの連鎖）

\chapter{付録}\label{章：付録}
\begin{bookcolumns}

  付録では,本編で選んだ量や変形を別の表現から見直す.
  \begin{itemize}
    \item 連立型や隣接3項間をまとめて更新したいなら,行列の基本から読む.
    \item 階差・総和と微積分の関係を見たいなら,\sectionref*{節：差分から始まる離散的な微積分}へ進む.微分・積分を使う部分は任意に読み進める.
    \item 数列全体の係数を扱いたいなら,\sectionref*{節：数列全体を一つの対象にまとめると,何が見えるか}へ進む.母関数の係数計算は,最初から無限級数の収束を仮定する計算とは区別する.
  \end{itemize}
  最初の一巡ではすべてを証明し直す必要はない.
  「本編のどの操作が現れたか」「新しい表現で何が見えたか」を一つずつ確かめてほしい.

% 付録：行列の基本

\section{行列の基本}\label{節：行列の基本}

本編では,複雑な\bookterm{recurrence}{漸化式}を等比型などの既知の形へ変えることを繰り返してきた.
付録では,その置き換えがなぜうまくいくのかを,\bookterm{matrix}{行列}や\bookterm{difference-operator}{差分},\bookterm{generating}{母関数}という言葉で捉え直す.
最初に,そのための道具となる\bookterm{matrix}{行列}の読み方と計算の仕方を説明しよう.
\bookterm{matrix}{行列}を初めて学ぶ読者も,数値例から順に確かめてほしい.
ここでは主に二行二列の\bookterm{matrix}{行列}を扱う.
後の節では,本編で求めた\bookterm{fixed}{不動点}や\bookterm{geometric}{等比数列}が,ここで学ぶ計算の中にそのまま現れる.

\subsection{行列とは何か}
  \BookFigMatrixMap


\bookterm{matrix}{行列}とは,数を長方形に並べ,全体を括弧で囲んだものである.
たとえば
\[
  A=\begin{pmatrix}2&1\\1&2\end{pmatrix}
\]
は,縦に二段,横に二列の数を並べた\bookterm{matrix}{行列}である.
横の並びを\textbf{行},縦の並びを\textbf{列}という.
上が第一行,下が第二行であり,左が第一列,右が第二列である.
並んでいる一つ一つの数を\textbf{成分}と呼ぶ.
この例の第一行・第二列の成分は$1$である.

また,二つの数を縦に並べた
\[
  \boldsymbol{v}=\begin{pmatrix}x\\y\end{pmatrix}
\]
を\textbf{\bookterm{column-vector}{列ベクトル}}と呼ぶ.
太字の$\boldsymbol{v}$は,二つの数をまとめて表すための記号である.
平面上の点$(x,y)$への矢印と考えることもできるが,以下ではまず「二つの値を入れた一組の記録」として読めばよい.

\bookterm{column-vector}{列ベクトル}の足し算と定数倍は,成分ごとに行う.
\[
  \begin{pmatrix}x\\y\end{pmatrix}
  +\begin{pmatrix}u\\v\end{pmatrix}
  =\begin{pmatrix}x+u\\y+v\end{pmatrix},
  \qquad
  c\begin{pmatrix}x\\y\end{pmatrix}
  =\begin{pmatrix}cx\\cy\end{pmatrix}.
\]
たとえば$\begin{pmatrix}3\\1\end{pmatrix}+\begin{pmatrix}1\\-1\end{pmatrix}
=\begin{pmatrix}4\\0\end{pmatrix}$である.
二つの\bookterm{column-vector}{列ベクトル}が等しいとは,上下それぞれの成分が等しいことをいう.
一般の二成分のベクトルも
\[
  \begin{pmatrix}x\\y\end{pmatrix}
  =\frac{x+y}{2}\begin{pmatrix}1\\1\end{pmatrix}
   +\frac{x-y}{2}\begin{pmatrix}1\\-1\end{pmatrix}
\]
と分けられる.上下の成分をそれぞれ足して確かめてほしい.
本編の\sectionref*{節：連立型}で和と差を求めた操作は,この二つの係数を求めることに対応する.
この対応は\bookproblemref{practice-basic}{7}{標準問7}で確認できる.
後で,この二つの方向が更新の定数倍とどう結び付くかを見る.

\subsubsection{行列を列ベクトルに掛ける}

先ほどの\bookterm{matrix}{行列}$A$を\bookterm{column-vector}{列ベクトル}に掛けるときは,次のように計算する.
\[
  \begin{pmatrix}2&1\\1&2\end{pmatrix}
  \begin{pmatrix}x\\y\end{pmatrix}
  =
  \begin{pmatrix}2x+y\\x+2y\end{pmatrix}.
\]
上の成分は,第一行の$2,\,1$を$x,y$にそれぞれ掛けて足したもの.
下の成分は,第二行の$1,\,2$を$x,y$にそれぞれ掛けて足したものである.
たとえば
\[
  A\begin{pmatrix}1\\0\end{pmatrix}
  =\begin{pmatrix}2\\1\end{pmatrix},
  \qquad
  A\begin{pmatrix}2\\1\end{pmatrix}
  =\begin{pmatrix}5\\4\end{pmatrix}.
\]
この\bookterm{matrix}{行列}は,二つの値$x,y$を受け取り,新しい二つの値$2x+y,x+2y$を作る規則を表している.

一般には
\[
  \begin{pmatrix}p&q\\r&s\end{pmatrix}
  \begin{pmatrix}x\\y\end{pmatrix}
  =\begin{pmatrix}px+qy\\rx+sy\end{pmatrix}
\]
と定める.
したがって,連立した二つの式
\[
  x'=px+qy,\qquad y'=rx+sy
\]
を,\bookterm{matrix}{行列}と\bookterm{column-vector}{列ベクトル}の積一つにまとめられる.
ここで$x',y'$は,更新後の値を表す記号である.

また,成分ごとに展開すれば
\[
  A(c\boldsymbol{u}+d\boldsymbol{v})
  =cA\boldsymbol{u}+dA\boldsymbol{v}
\]
が成り立つ.
これは,二つの\bookterm{column-vector}{列ベクトル}を足してから更新しても,別々に更新してから足しても結果が同じ,という性質である.
後で初期値を二つの部分に分けて考える際に使う.

\subsection{行列の積は,更新を続けて行うことを表す}

\bookterm{matrix}{行列}をもう一度掛けると,同じ規則で二回更新できる.
先ほどの例なら
\[
  \begin{pmatrix}1\\0\end{pmatrix}
  \ \longmapsto\
  \begin{pmatrix}2\\1\end{pmatrix}
  \ \longmapsto\
  \begin{pmatrix}5\\4\end{pmatrix}
\]
と進む.
この二回分の更新を一つにまとめた\bookterm{matrix}{行列}を$A^2$と書く.
実際に文字$x,y$で計算すると
\[
  \begin{aligned}
    A\left(A\begin{pmatrix}x\\y\end{pmatrix}\right)
    &=A\begin{pmatrix}2x+y\\x+2y\end{pmatrix}\\
    &=\begin{pmatrix}5x+4y\\4x+5y\end{pmatrix}.
  \end{aligned}
\]
したがって
\[
  A^2=\begin{pmatrix}5&4\\4&5\end{pmatrix}
\]
である.
各成分を単に二乗するのではなく,更新を二回行った結果に合わせて積を定めている.

異なる\bookterm{matrix}{行列}$A,B$でも,まず$B$,次に$A$を掛ける操作を$AB$と書く.
$B$の各列に$A$を掛けた結果を並べると,積$AB$が得られる.
具体的には
\[
  \begin{aligned}
    &\begin{pmatrix}p&q\\r&s\end{pmatrix}
     \begin{pmatrix}u&v\\w&z\end{pmatrix}\\
    &\qquad=
     \begin{pmatrix}
       pu+qw&pv+qz\\
       ru+sw&rv+sz
     \end{pmatrix}.
  \end{aligned}
\]
右の\bookterm{matrix}{行列}から先に作用するため,一般には$AB$と$BA$は一致しない.

何も変えない更新は
\[
  I=\begin{pmatrix}1&0\\0&1\end{pmatrix},
  \qquad
  I\begin{pmatrix}x\\y\end{pmatrix}
  =\begin{pmatrix}x\\y\end{pmatrix}
\]
で表される.
$I$を\textbf{\bookterm{identity}{単位行列}}という.
数の累乗で$a^0=1$とするのに対応して,\bookterm{matrix}{行列}では$A^0=I$と定める.

\subsection{固有ベクトルと固有値：更新が定数倍になる場合}
  \BookFigEigen


毎回\bookterm{matrix}{行列}の積を計算するのは手間がかかる.
そこで,\bookterm{matrix}{行列}を掛けても定数倍になるだけの\bookterm{column-vector}{列ベクトル}を探してみよう.
先ほどの$A$について
\[
  \begin{aligned}
    A\begin{pmatrix}1\\1\end{pmatrix}
      &=\begin{pmatrix}3\\3\end{pmatrix}
       =3\begin{pmatrix}1\\1\end{pmatrix},\\
    A\begin{pmatrix}1\\-1\end{pmatrix}
      &=\begin{pmatrix}1\\-1\end{pmatrix}
  \end{aligned}
\]
である.
第一の\bookterm{column-vector}{列ベクトル}は一回ごとに$3$倍され,第二の\bookterm{column-vector}{列ベクトル}は変わらない.
したがって$k$回更新した結果は
\[
  A^k\begin{pmatrix}1\\1\end{pmatrix}
    =3^k\begin{pmatrix}1\\1\end{pmatrix},
  \qquad
  A^k\begin{pmatrix}1\\-1\end{pmatrix}
    =\begin{pmatrix}1\\-1\end{pmatrix}
\]
とすぐに分かる.

一般に,零でない\bookterm{column-vector}{列ベクトル}$\boldsymbol{u}$が
\[
  A\boldsymbol{u}=\lambda\boldsymbol{u}
\]
を満たすとき,\,$\boldsymbol{u}$を$A$の\textbf{\bookterm{eigenvector}{固有ベクトル}},\,$\lambda$を対応する\textbf{\bookterm{eigenvalue}{固有値}}という.
\bookterm{eigenvalue}{固有値}は,そのベクトルに対して一回の更新が何倍になるかを表す数である.
零ベクトル$\boldsymbol{0}=\begin{pmatrix}0\\0\end{pmatrix}$は,どんな$\lambda$でもこの式を満たしてしまうため,\bookterm{eigenvector}{固有ベクトル}から除く.
負の\bookterm{eigenvalue}{固有値}では矢印の向きが反転し,\bookterm{eigenvalue}{固有値}$0$では零ベクトルになる場合もある.

\subsubsection{ベクトルを二つの部分に分ける}

\bookterm{eigenvector}{固有ベクトル}が見つかると,ほかのベクトルへの作用も計算しやすくなる.
たとえば
\[
  \begin{pmatrix}1\\0\end{pmatrix}
  =\frac12\begin{pmatrix}1\\1\end{pmatrix}
   +\frac12\begin{pmatrix}1\\-1\end{pmatrix}
\]
である.
右辺を実際に足せば,上が$1$,下が$0$になる.
\sectionref*{節：一次独立・基底・次元}では,独立な解の\bookterm{combination}{線形結合}で\bookterm{initial}{初期条件}を合わせた.
ここでは二つの独立な\bookterm{column-vector}{列ベクトル}を\bookterm{basis}{基底}に選び,元のベクトルを同じように\bookterm{combination}{線形結合}で表している.
\bookterm{sequence}{数列}全体を分けた見方を,今度は更新前の状態の分解に使っているのである.
二つの部分を別々に$k$回更新すれば
\[
  \begin{aligned}
    A^k\begin{pmatrix}1\\0\end{pmatrix}
      &=\frac{3^k}{2}\begin{pmatrix}1\\1\end{pmatrix}
        +\frac12\begin{pmatrix}1\\-1\end{pmatrix}\\
      &=\begin{pmatrix}\frac{3^k+1}{2}\\\frac{3^k-1}{2}\end{pmatrix}.
  \end{aligned}
\]
大きな累乗を一回ずつ計算する代わりに,定数倍になる二つの部分を追えばよい.
本編で等比型に帰着させると計算が簡単になったのと同じように,\bookterm{matrix}{行列}でも一定の倍率で変わる部分を見つけることが鍵になる.

一般に,二つの異なる\bookterm{eigenvalue}{固有値}$\lambda_1,\lambda_2$に対応する\bookterm{eigenvector}{固有ベクトル}$\boldsymbol{u}_1,\boldsymbol{u}_2$があれば,どんな二成分のベクトルも
\[
  \boldsymbol{v}=c\boldsymbol{u}_1+d\boldsymbol{u}_2
\]
と表せる.
異なる\bookterm{eigenvalue}{固有値}に対応する二つの\bookterm{eigenvector}{固有ベクトル}は,互いの定数倍にならないからである.
係数$c,d$は,上下の成分を比較した連立方程式で求められる.
その後は
\[
  A^k\boldsymbol{v}
  =c\lambda_1^k\boldsymbol{u}_1+d\lambda_2^k\boldsymbol{u}_2
\]
で計算できる.
$k=0$のときは元の分解式を読み,\,$k\ge1$で\bookterm{eigenvalue}{固有値}が$0$なら対応する部分は零になる.

\subsection{行ベクトルと左固有ベクトル}

数を横に並べた$L=(\alpha,\beta)$を\textbf{\bookterm{row-vector}{行ベクトル}}という.
\bookterm{column-vector}{列ベクトル}$\boldsymbol{\ell}$を横に並べ直す操作を\textbf{\bookterm{transpose}{転置}}といい,\,$\boldsymbol{\ell}^{\mathsf T}$と書く.
たとえば$\boldsymbol{\ell}=\begin{pmatrix}\alpha\\\beta\end{pmatrix}$なら,\,$\boldsymbol{\ell}^{\mathsf T}=(\alpha,\beta)$である.
\bookterm{row-vector}{行ベクトル}と\bookterm{column-vector}{列ベクトル}の積は,対応する成分を掛けて足す.
\[
  (\alpha,\beta)\begin{pmatrix}x\\y\end{pmatrix}
  =\alpha x+\beta y.
\]
\bookterm{row-vector}{行ベクトル}を\bookterm{matrix}{行列}に掛けるときは,\bookterm{matrix}{行列}の各列との積を横に並べる.
たとえば,先ほどの$A$について
\[
  \begin{aligned}
    (1,\,1)A&=(3,\,3)=3(1,\,1),\\
    (1,-1)A&=(1,-1)
  \end{aligned}
\]
である.
一般に,零でない\bookterm{row-vector}{行ベクトル}$L$が$LA=\gamma L$を満たすとき,\,$L$を\textbf{\bookterm{left-eigen}{左固有ベクトル}}という.
この場合,\bookterm{column-vector}{列ベクトル}$\boldsymbol{v}$から作った数$L\boldsymbol{v}$は,更新後には
\[
  L(A\boldsymbol{v})=(LA)\boldsymbol{v}
  =\gamma L\boldsymbol{v}
\]
となる.
たとえば$L=(1,\,1)$なら,\,$x+y$という和が更新によって$3$倍される.
\bookterm{eigenvector}{固有ベクトル}は定数倍になる\bookterm{column-vector}{列ベクトル}を表し,\bookterm{left-eigen}{左固有ベクトル}は定数倍になる量を取り出す係数を表す.

\subsection{固有値を求める条件：行列式}

\bookterm{eigenvector}{固有ベクトル}の成分を比較すれば,\bookterm{eigenvalue}{固有値}を求める条件が得られる.
その条件を一つの式にまとめるため,\bookterm{matrix}{行列}式という記号を導入しよう.
二行二列の場合には
\[
  \det\begin{pmatrix}a&b\\c&d\end{pmatrix}=ad-bc
\]
と定める.
\bookterm{matrix}{行列}そのものではなく,四つの成分から計算した一つの数である.

二つの式$ax+by=0,\ cx+dy=0$について,\,$ad-bc\ne0$なら解は$x=y=0$だけになる.
たとえば,第一の式を$d$倍したものから第二の式を$b$倍したものを引くと,\,$(ad-bc)x=0$となる.
同様に$(ad-bc)y=0$も得られる.
逆に$ad-bc=0$なら,零でない解が存在する.
実際,\,$(a,b)\ne(0,\,0)$なら$(x,y)=(b,-a)$が両方の式を満たす.
$(a,b)=(0,\,0)$の場合も,残る一つの式から零でない解を選べる.

\bookterm{eigenvalue}{固有値}の式$A\boldsymbol{u}=\lambda\boldsymbol{u}$は
\[
  (\lambda I-A)\boldsymbol{u}=\boldsymbol{0}
\]
と書き直せる.
\bookterm{matrix}{行列}の引き算は対応する成分ごとに行うので,\,$A=\begin{pmatrix}p&q\\r&s\end{pmatrix}$なら
\[
  \lambda I-A
  =\begin{pmatrix}\lambda-p&-q\\-r&\lambda-s\end{pmatrix}.
\]
零でない解が存在する条件は
\[
  \begin{aligned}
    \det(\lambda I-A)
      &=(\lambda-p)(\lambda-s)-qr\\
      &=\lambda^2-(p+s)\lambda+(ps-qr)=0.
  \end{aligned}
\]
これを\bookterm{matrix}{行列}$A$の\textbf{\bookterm{characteristic}{特性方程式}}という.
この方程式を解いて$\lambda$を求め,成分の連立方程式を解けば,対応する\bookterm{eigenvector}{固有ベクトル}も見つかる.
たとえば$A=\begin{pmatrix}2&1\\1&2\end{pmatrix}$では
\[
  \lambda^2-4\lambda+3
  =(\lambda-3)(\lambda-1)=0
\]
だから,\bookterm{eigenvalue}{固有値}は$3,\,1$である.
これは数値例で確かめた倍率と一致する.

\bookterm{recurrence}{漸化式}では,次の項を計算するために記録しておく値の組を\textbf{状態}と呼ぶ.
また,\bookterm{matrix}{行列}式が$0$でないときは,更新後の二つの値から更新前の二つの値を一意に求め直せる.
このように更新を逆向きに戻せる\bookterm{matrix}{行列}を\textbf{可逆な\bookterm{matrix}{行列}}という.

次節では,これらの計算を本編の\bookterm{recurrence}{漸化式}に当てはめていこう.
特に,本編で選んだ「何を引くか」「どのような比や和を作るか」が,\bookterm{matrix}{行列}では何を選ぶことに対応するのかを考える.

% 付録・第1部：行列の基本から漸化式へ

\section{行列の基本から,特性方程式を見直す}\label{節：特性方程式を行列で見直す}

本節では,まず行列の読み方と計算の仕方を説明する.
その後,本編で使った「等比数列になるように置き換える」という工夫を,行列の言葉で見直していこう.
行列を初めて学ぶ読者も,最初の数値例から順に確かめてほしい.
ここでは主に二行二列の行列を扱う.

\subsection{行列とは何か}

行列とは,数を長方形に並べ,全体を括弧で囲んだものである.
たとえば
\[
  A=\begin{pmatrix}2&1\\1&2\end{pmatrix}
\]
は,縦に二段,横に二列の数を並べた行列である.
横の並びを\textbf{行},縦の並びを\textbf{列}という.
上が第一行,下が第二行であり,左が第一列,右が第二列である.
並んでいる一つ一つの数を\textbf{成分}と呼ぶ.
この例の第一行・第二列の成分は$1$である.

また,二つの数を縦に並べた
\[
  \boldsymbol{v}=\begin{pmatrix}x\\y\end{pmatrix}
\]
を\textbf{列ベクトル}と呼ぶ.
太字の$\boldsymbol{v}$は,二つの数をまとめて表すための記号である.
平面上の点$(x,y)$への矢印と考えることもできるが,以下ではまず「二つの値を入れた一組の記録」として読めばよい.

列ベクトルの足し算と定数倍は,成分ごとに行う.
\[
  \begin{pmatrix}x\\y\end{pmatrix}
  +\begin{pmatrix}u\\v\end{pmatrix}
  =\begin{pmatrix}x+u\\y+v\end{pmatrix},
  \qquad
  c\begin{pmatrix}x\\y\end{pmatrix}
  =\begin{pmatrix}cx\\cy\end{pmatrix}.
\]
たとえば$\begin{pmatrix}3\\1\end{pmatrix}+\begin{pmatrix}1\\-1\end{pmatrix}
=\begin{pmatrix}4\\0\end{pmatrix}$である.
二つの列ベクトルが等しいとは,上下それぞれの成分が等しいことをいう.

\subsubsection{行列を列ベクトルに掛ける}

先ほどの行列$A$を列ベクトルに掛けるときは,次のように計算する.
\[
  \begin{pmatrix}2&1\\1&2\end{pmatrix}
  \begin{pmatrix}x\\y\end{pmatrix}
  =
  \begin{pmatrix}2x+y\\x+2y\end{pmatrix}.
\]
上の成分は,第一行の$2,\,1$を$x,y$にそれぞれ掛けて足したもの.
下の成分は,第二行の$1,\,2$を$x,y$にそれぞれ掛けて足したものである.
たとえば
\[
  A\begin{pmatrix}1\\0\end{pmatrix}
  =\begin{pmatrix}2\\1\end{pmatrix},
  \qquad
  A\begin{pmatrix}2\\1\end{pmatrix}
  =\begin{pmatrix}5\\4\end{pmatrix}.
\]
この行列は,二つの値$x,y$を受け取り,新しい二つの値$2x+y,x+2y$を作る規則を表している.

一般には
\[
  \begin{pmatrix}p&q\\r&s\end{pmatrix}
  \begin{pmatrix}x\\y\end{pmatrix}
  =\begin{pmatrix}px+qy\\rx+sy\end{pmatrix}
\]
と定める.
したがって,連立した二つの式
\[
  x'=px+qy,\qquad y'=rx+sy
\]
を,行列と列ベクトルの積一つにまとめられる.
ここで$x',y'$は,更新後の値を表す記号である.

また,成分ごとに展開すれば
\[
  A(c\boldsymbol{u}+d\boldsymbol{v})
  =cA\boldsymbol{u}+dA\boldsymbol{v}
\]
が成り立つ.
これは,二つの列ベクトルを足してから更新しても,別々に更新してから足しても結果が同じ,という性質である.
後で初期値を二つの部分に分けて考える際に使う.

\subsection{行列の積は,更新を続けて行うことを表す}

行列をもう一度掛けると,同じ規則で二回更新できる.
先ほどの例なら
\[
  \begin{pmatrix}1\\0\end{pmatrix}
  \ \longmapsto\
  \begin{pmatrix}2\\1\end{pmatrix}
  \ \longmapsto\
  \begin{pmatrix}5\\4\end{pmatrix}
\]
と進む.
この二回分の更新を一つにまとめた行列を$A^2$と書く.
実際に文字$x,y$で計算すると
\[
  \begin{aligned}
    A\left(A\begin{pmatrix}x\\y\end{pmatrix}\right)
    &=A\begin{pmatrix}2x+y\\x+2y\end{pmatrix}\\
    &=\begin{pmatrix}5x+4y\\4x+5y\end{pmatrix}.
  \end{aligned}
\]
したがって
\[
  A^2=\begin{pmatrix}5&4\\4&5\end{pmatrix}
\]
である.
各成分を単に二乗するのではなく,更新を二回行った結果に合わせて積を定めている.

異なる行列$A,B$でも,まず$B$,次に$A$を掛ける操作を$AB$と書く.
$B$の各列に$A$を掛けた結果を並べると,積$AB$が得られる.
具体的には
\[
  \begin{aligned}
    &\begin{pmatrix}p&q\\r&s\end{pmatrix}
     \begin{pmatrix}u&v\\w&z\end{pmatrix}\\
    &\qquad=
     \begin{pmatrix}
       pu+qw&pv+qz\\
       ru+sw&rv+sz
     \end{pmatrix}.
  \end{aligned}
\]
右の行列から先に作用するため,一般には$AB$と$BA$は一致しない.

何も変えない更新は
\[
  I=\begin{pmatrix}1&0\\0&1\end{pmatrix},
  \qquad
  I\begin{pmatrix}x\\y\end{pmatrix}
  =\begin{pmatrix}x\\y\end{pmatrix}
\]
で表される.
$I$を\textbf{単位行列}という.
数の累乗で$a^0=1$とするのに対応して,行列では$A^0=I$と定める.

ここで,二つの数列が
\[
  x_{n+1}=2x_n+y_n,\qquad
  y_{n+1}=x_n+2y_n
\]
を満たす場合を考えよう.
$n$番目の二つの値を
\[
  \boldsymbol{v}_n=\begin{pmatrix}x_n\\y_n\end{pmatrix}
\]
とまとめれば,\,$\boldsymbol{v}_{n+1}=A\boldsymbol{v}_n$である.
このように,次の値を計算するために持っておく値の組を\textbf{状態}と呼ぶ.
初期状態から一回,二回と進めると
\[
  \boldsymbol{v}_2=A\boldsymbol{v}_1,\qquad
  \boldsymbol{v}_3=A^2\boldsymbol{v}_1
\]
なので,一般に
\[
  \boldsymbol{v}_n=A^{n-1}\boldsymbol{v}_1
\]
となる.
初項から第$n$項までの更新は$n-1$回である.
行列の累乗が分かれば,連立漸化式の一般項も分かることになる.

\subsection{固有ベクトルと固有値：更新が定数倍になる場合}

毎回行列の積を計算するのは手間がかかる.
そこで,行列を掛けても定数倍になるだけの列ベクトルを探してみよう.
先ほどの$A$について
\[
  \begin{aligned}
    A\begin{pmatrix}1\\1\end{pmatrix}
      &=\begin{pmatrix}3\\3\end{pmatrix}
       =3\begin{pmatrix}1\\1\end{pmatrix},\\
    A\begin{pmatrix}1\\-1\end{pmatrix}
      &=\begin{pmatrix}1\\-1\end{pmatrix}
  \end{aligned}
\]
である.
第一の列ベクトルは一回ごとに$3$倍され,第二の列ベクトルは変わらない.
したがって$k$回更新した結果は
\[
  A^k\begin{pmatrix}1\\1\end{pmatrix}
    =3^k\begin{pmatrix}1\\1\end{pmatrix},
  \qquad
  A^k\begin{pmatrix}1\\-1\end{pmatrix}
    =\begin{pmatrix}1\\-1\end{pmatrix}
\]
とすぐに分かる.

一般に,零でない列ベクトル$\boldsymbol{u}$が
\[
  A\boldsymbol{u}=\lambda\boldsymbol{u}
\]
を満たすとき,\,$\boldsymbol{u}$を$A$の\textbf{固有ベクトル},\,$\lambda$を対応する\textbf{固有値}という.
固有値は,そのベクトルに対して一回の更新が何倍になるかを表す数である.
零ベクトル$\boldsymbol{0}=\begin{pmatrix}0\\0\end{pmatrix}$は,どんな$\lambda$でもこの式を満たしてしまうため,固有ベクトルから除く.
負の固有値では矢印の向きが反転し,固有値$0$では零ベクトルになる場合もある.

\subsubsection{初期状態を二つに分けると,一般項が出る}

初期値を$(x_1,y_1)=(1,\,0)$として,先ほどの連立漸化式を解いてみよう.
初期状態は
\[
  \begin{pmatrix}1\\0\end{pmatrix}
  =\frac12\begin{pmatrix}1\\1\end{pmatrix}
   +\frac12\begin{pmatrix}1\\-1\end{pmatrix}
\]
と分けられる.
右辺を実際に足せば,上が$1$,下が$0$になる.
二つの部分を別々に$n-1$回更新すると
\[
  \begin{pmatrix}x_n\\y_n\end{pmatrix}
  =\frac{3^{n-1}}2\begin{pmatrix}1\\1\end{pmatrix}
   +\frac12\begin{pmatrix}1\\-1\end{pmatrix}.
\]
上下の成分を読み取って
\[
  x_n=\frac{3^{n-1}+1}{2},\qquad
  y_n=\frac{3^{n-1}-1}{2}
\]
を得る.
初期状態を,更新の簡単な二つの部分に分けたことが解法の要点である.

本編の連立型では,和と差を
\[
  u_n=x_n+y_n,\qquad v_n=x_n-y_n
\]
と置いた.
もとの二つの式を足す,引くという計算から
\[
  u_{n+1}=3u_n,\qquad v_{n+1}=v_n
\]
となる.
そして
\[
  x_n=\frac{u_n+v_n}{2},\qquad
  y_n=\frac{u_n-v_n}{2}
\]
で元に戻せる.
ここに現れた倍率$3,\,1$は,先ほどの固有値と同じである.

和と差を作る係数を横に並べると,\,$(1,\,1)$と$(1,-1)$になる.
このように数を横に並べたものを\textbf{行ベクトル}という.
行ベクトルと列ベクトルの積は,対応する成分を掛けて足す.
\[
  (\alpha,\beta)\begin{pmatrix}x\\y\end{pmatrix}
  =\alpha x+\beta y.
\]
行ベクトルを行列に掛けるときは,行列の各列との積を横に並べる.
たとえば$(1,\,1)A=(3,\,3)=3(1,\,1)$である.
一般の行列$A$でも,零でない行ベクトル$L$が
\[
  LA=\gamma L
\]
を満たせば,\,$c_n=L\boldsymbol{v}_n$は
\[
  c_{n+1}=LA\boldsymbol{v}_n=\gamma c_n
\]
を満たす.
この$L$を\textbf{左固有ベクトル}と呼ぶ.
本編で等比数列になるように$\alpha x_n+\beta y_n$の係数を選ぶ方法は,この行ベクトルを探す計算として理解できる.
固有ベクトルで状態を分ける方法と,等比数列になる量を取り出す方法は,同じ更新を異なる側から見ている.

\subsection{本編の特性方程式型を,行列で書き直す}\label{節：一つの値の更新を行列で書く}

次に,本編で扱った
\[
  a_{n+1}=pa_n+q\qquad(p\ne1)
\]
に戻ろう.
ここでは一つの値だけを更新しているが,定数$q$の足し算がある.
この足し算を行列の掛け算の中に含めるため,\,$a_n$と一緒に$1$も並べておく.
すると
\[
  \begin{pmatrix}p&q\\0&1\end{pmatrix}
  \begin{pmatrix}a_n\\1\end{pmatrix}
  =\begin{pmatrix}pa_n+q\\1\end{pmatrix}
  =\begin{pmatrix}a_{n+1}\\1\end{pmatrix}.
\]
上の成分では$q$に$1$を掛けることで,必要な足し算を作っている.
下の成分は毎回$1$のままなので,次の更新でも同じ計算ができる.

この行列を$H$とし,本編と同じく
\[
  \alpha=p\alpha+q,\qquad
  \alpha=\frac{q}{1-p}
\]
を満たす$\alpha$を取る.
この値から始めると次の項も変わらないので,\,$\alpha$は不動点である.
行列で計算すると
\[
  H\begin{pmatrix}\alpha\\1\end{pmatrix}
  =\begin{pmatrix}\alpha\\1\end{pmatrix},
  \qquad
  H\begin{pmatrix}1\\0\end{pmatrix}
  =p\begin{pmatrix}1\\0\end{pmatrix}.
\]
これで,更新が簡単な二つの列ベクトルが見つかった.

初期状態を
\[
  \begin{pmatrix}a_1\\1\end{pmatrix}
  =\begin{pmatrix}\alpha\\1\end{pmatrix}
   +(a_1-\alpha)\begin{pmatrix}1\\0\end{pmatrix}
\]
と分ける.
第一の部分は変わらず,第二の部分だけが一回ごとに$p$倍されるから
\[
  \begin{pmatrix}a_n\\1\end{pmatrix}
  =\begin{pmatrix}\alpha\\1\end{pmatrix}
   +(a_1-\alpha)p^{n-1}\begin{pmatrix}1\\0\end{pmatrix}.
\]
上の成分を読めば
\[
  a_n=\alpha+(a_1-\alpha)p^{n-1}
\]
となる.
$p=0$では$a_2$以降が$\alpha=q$になるので直接読めばよい.

たとえば$a_{n+1}=2a_n+1,\ a_1=0$なら$\alpha=-1$である.
初期状態は
\[
  \begin{pmatrix}0\\1\end{pmatrix}
  =\begin{pmatrix}-1\\1\end{pmatrix}
   +\begin{pmatrix}1\\0\end{pmatrix}
\]
と分かれ,後半だけが$2$倍ずつになる.
したがって$a_n=-1+2^{n-1}$である.
本編の$b_n=a_n-\alpha$という置き換えは,この「不動点からのずれ」だけを取り出していたのである.

ここで用いた$\alpha=p\alpha+q$は,不動点を求める方程式である.
一方,行列の固有値を求める方程式も線形代数では「特性方程式」と呼ばれる.
二つは異なる式であり,不動点$\alpha$そのものを固有値と呼ぶわけではない.
上の行列$H$の固有値は$1$と$p$である.

$p=1$なら,もとの式は$a_{n+1}=a_n+q$である.
更新を$k$回続けると
\[
  H^k=\begin{pmatrix}1&kq\\0&1\end{pmatrix}
\]
となり,\,$a_n=a_1+(n-1)q$を得る.
$q\ne0$では不動点がなく,加えた$q$が更新回数だけ積み重なる.
これは後で見る,特性根が重なると一般項に$n$が現れる例につながっている.

\subsection{二つの項を記録すれば,三項間漸化式も行列になる}\label{節：状態を二つ持つと,二階の漸化式になる}

三項間漸化式では,次の項を計算するために直前の二つの項が必要である.
たとえばフィボナッチ数列
\[
  F_0=0,\quad F_1=1,\quad
  F_{n+2}=F_{n+1}+F_n\quad(n\ge0)
\]
を考える.
状態を$\begin{pmatrix}F_{n+1}\\F_n\end{pmatrix}$とすれば
\[
  \begin{pmatrix}1&1\\1&0\end{pmatrix}
  \begin{pmatrix}F_{n+1}\\F_n\end{pmatrix}
  =\begin{pmatrix}F_{n+1}+F_n\\F_{n+1}\end{pmatrix}
  =\begin{pmatrix}F_{n+2}\\F_{n+1}\end{pmatrix}.
\]
上の成分で新しい項を作り,下の成分に直前の項を残している.
実際,
\[
  \begin{pmatrix}1\\0\end{pmatrix}
  \longmapsto\begin{pmatrix}1\\1\end{pmatrix}
  \longmapsto\begin{pmatrix}2\\1\end{pmatrix}
  \longmapsto\begin{pmatrix}3\\2\end{pmatrix}
\]
と進む.

一般に
\[
  a_{n+2}=s a_{n+1}+t a_n\qquad(n\ge0)
\]
は
\[
  \boldsymbol{w}_n=\begin{pmatrix}a_{n+1}\\a_n\end{pmatrix},
  \qquad
  B=\begin{pmatrix}s&t\\1&0\end{pmatrix}
\]
と置けば,\,$\boldsymbol{w}_{n+1}=B\boldsymbol{w}_n$になる.
直前の二項を必要とするため,この形を\textbf{二階の漸化式}とも呼ぶ.
本編の$a_{n+2}+pa_{n+1}+qa_n=0$では,\,$s=-p,\ t=-q$とすればよい.

\subsubsection{なぜ本編と同じ特性方程式が出るのか}

$B$を掛けると$\lambda$倍になる,零でない列ベクトルを探そう.
成分を$x,y$と書けば
\[
  B\begin{pmatrix}x\\y\end{pmatrix}
  =\lambda\begin{pmatrix}x\\y\end{pmatrix}
  \quad\Longleftrightarrow\quad
  \begin{cases}
    sx+ty=\lambda x,\\
    x=\lambda y.
  \end{cases}
\]
二つ目の式から,\,$y=0$なら$x=0$になってしまう.
そこで$y\ne0$と分かるので,ベクトル全体を$y$で割り,\,$y=1,\ x=\lambda$としてよい.
一つ目の式は
\[
  s\lambda+t=\lambda^2,
  \qquad
  \lambda^2-s\lambda-t=0
\]
となる.
これが固有値を求める方程式である.

本編では,等比数列の形$a_n=\lambda^n$を漸化式に当てはめて,同じ方程式を得た.
行列では「二項をまとめた記録が一回ごとに$\lambda$倍になる条件」を調べている.
数列が等比的に進むことと,状態ベクトルが定数倍になることは,同じ条件を別の形で表している.
$\lambda=0$が根になる場合も,行列の方程式はそのまま使える.

二つの異なる根$\lambda_1,\lambda_2$があれば
\[
  \begin{pmatrix}\lambda_1\\1\end{pmatrix},
  \qquad
  \begin{pmatrix}\lambda_2\\1\end{pmatrix}
\]
が固有ベクトルになる.
初期状態を
\[
  \begin{pmatrix}a_1\\a_0\end{pmatrix}
  =c\begin{pmatrix}\lambda_1\\1\end{pmatrix}
   +d\begin{pmatrix}\lambda_2\\1\end{pmatrix}
\]
と分けるには,\,$c+d=a_0,\ c\lambda_1+d\lambda_2=a_1$を解けばよい.
根が異なるので,\,$c,d$は一意に決まる.
$n$回更新した後は,下の成分から
\[
  a_n=c\lambda_1^n+d\lambda_2^n
\]
を得る.
根が$0$の場合,\,$n=0$での係数は初期状態の式から読み,\,$n\ge1$ではその成分は$0$になる.

フィボナッチ数列なら特性方程式は$\lambda^2-\lambda-1=0$で,根は
\[
  \varphi=\frac{1+\sqrt5}{2},\qquad
  \psi=\frac{1-\sqrt5}{2}
\]
である.
初期状態$\begin{pmatrix}1\\0\end{pmatrix}$に合わせる条件は
\[
  c+d=0,\qquad c\varphi+d\psi=1.
\]
$\varphi-\psi=\sqrt5$より$c=1/\sqrt5,\ d=-1/\sqrt5$と決まり,
\[
  F_n=\frac{\varphi^n-\psi^n}{\sqrt5}
\]
を得る.
$|\psi|<1<\varphi$なので,\,$\psi^n$の成分は小さくなり,\,$\varphi^n$の成分が大きくなる.
固有値は一般項の計算に加えて,更新を重ねたときの増え方も教えている.
異なる二つの根が複素数の場合も,成分や係数を複素数まで広げれば同じ計算ができる.
実数の初期値に合わせると,共役な二つの根の寄与が組になり,数列の項は実数になる.

\subsubsection{行列式を使うと,固有値の条件をまとめられる}

ここまでのように成分を比較すれば,固有値は求められる.
その条件を一つの式にまとめる記号が\textbf{行列式}である.
二行二列の場合には
\[
  \det\begin{pmatrix}a&b\\c&d\end{pmatrix}=ad-bc
\]
と定める.
行列そのものではなく,四つの成分から計算した一つの数である.

二つの式$ax+by=0,\ cx+dy=0$について,\,$ad-bc\ne0$なら解は$x=y=0$だけになる.
たとえば,第一の式を$d$倍したものから第二の式を$b$倍したものを引くと,\,$(ad-bc)x=0$となる.
同様に$(ad-bc)y=0$も得られる.
逆に$ad-bc=0$なら,零でない解が存在する.
実際,\,$(a,b)\ne(0,\,0)$なら$(x,y)=(b,-a)$が両方の式を満たす.
$(a,b)=(0,\,0)$の場合も,残る一つの式から零でない解を選べる.

固有値の式$B\boldsymbol{u}=\lambda\boldsymbol{u}$は,
\[
  (\lambda I-B)\boldsymbol{u}=\boldsymbol{0}
\]
と書き直せる.
行列の引き算は対応する成分ごとに行うので
\[
  \lambda I-B
  =\begin{pmatrix}\lambda-s&-t\\-1&\lambda\end{pmatrix}.
\]
零でない解が存在する条件は
\[
  \begin{aligned}
    \det(\lambda I-B)
      &=(\lambda-s)\lambda-(-t)(-1)\\
      &=\lambda^2-s\lambda-t=0.
  \end{aligned}
\]
こうして,先ほど成分から求めた特性方程式が再び得られる.
一般の行列$A=\begin{pmatrix}p&q\\r&s\end{pmatrix}$でも,同じ計算で
\[
  \det(\lambda I-A)
  =\lambda^2-(p+s)\lambda+(ps-qr)=0
\]
が固有値を求める特性方程式になる.

\subsubsection{特性根が重なるとき}

特性方程式が重解$r\ne0$を持つときには,異なる二つの固有ベクトルに分ける先ほどの方法が使えない.
一般項に$n$が現れる理由は,もとの漸化式を直接変形すると分かりやすい.
重解のときは
\[
  a_{n+2}=2r a_{n+1}-r^2a_n
\]
なので,\,$a_n=r^n b_n$と置いて$r^{n+2}$で割ると
\[
  b_{n+2}-2b_{n+1}+b_n=0.
\]
これは
\[
  b_{n+2}-b_{n+1}=b_{n+1}-b_n
\]
ということだから,\,$b_n$は等差数列になる.
したがって
\[
  b_n=C+Dn,\qquad a_n=(C+Dn)r^n
\]
である.
等比的な倍率$r^n$を取り除くと,残りに等差的な増え方が現れるのである.
$r=0$が重解なら$a_{n+2}=0$なので,\,$a_2$以降はすべて$0$として直接扱う.

\subsection{一次分数型は,二つの成分の比として読む}

ここまで,列ベクトルの成分そのものを数列の項として読んだ.
今度は成分の\textbf{比}を読むと,本編の一次分数型につながることを見てみよう.
\[
  a_{n+1}=\frac{pa_n+q}{ra_n+s}
\]
に対し,\,$r\ne0,\ ps-qr\ne0$とし,すべての項で$ra_n+s\ne0$と仮定する.
初めを$(x_1,y_1)=(a_1,\,1)$として
\[
  \begin{pmatrix}x_{n+1}\\y_{n+1}\end{pmatrix}
  =M\begin{pmatrix}x_n\\y_n\end{pmatrix},
  \qquad M=\begin{pmatrix}p&q\\r&s\end{pmatrix}
\]
と更新する.
$a_n=x_n/y_n$であれば
\[
  \begin{aligned}
    \frac{x_{n+1}}{y_{n+1}}
    &=\frac{px_n+qy_n}{rx_n+sy_n}\\
    &=\frac{p(x_n/y_n)+q}{r(x_n/y_n)+s}
     =\frac{pa_n+q}{ra_n+s}.
  \end{aligned}
\]
また,\,$y_{n+1}=(ra_n+s)y_n$なので,分母が零にならないという仮定のもとで$y_n\ne0$が保たれる.
したがって,各回の成分比はもとの漸化式の項になる.
ベクトル全体を零でない数で何倍しても比は同じである.
このように成分比に注目して行列の作用を見る方法を,射影変換の見方という.

不動点$\alpha$については
\[
  p\alpha+q=\alpha(r\alpha+s)
\]
なので
\[
  M\begin{pmatrix}\alpha\\1\end{pmatrix}
  =(r\alpha+s)\begin{pmatrix}\alpha\\1\end{pmatrix}.
\]
つまり,不動点に対応する列ベクトルは固有ベクトルであり,固有値は$r\alpha+s$である.
$\alpha$自身が固有値なのではなく,ベクトルの二成分の比になっている.

異なる二つの不動点$\alpha,\beta$がある場合を考え,
\[
  \lambda_\alpha=r\alpha+s,\qquad
  \lambda_\beta=r\beta+s
\]
と置く.
各状態を
\[
  \begin{pmatrix}x_n\\y_n\end{pmatrix}
  =c_n\begin{pmatrix}\alpha\\1\end{pmatrix}
   +d_n\begin{pmatrix}\beta\\1\end{pmatrix}
\]
と分ければ
\[
  c_{n+1}=\lambda_\alpha c_n,\qquad
  d_{n+1}=\lambda_\beta d_n
\]
となる.
二つの係数が,別々の等比数列として進むのである.

本編の置き換えとの関係を,成分から確かめよう.
$x_n=\alpha c_n+\beta d_n,\ y_n=c_n+d_n$だから
\[
  x_n-\alpha y_n=(\beta-\alpha)d_n,
  \qquad
  x_n-\beta y_n=(\alpha-\beta)c_n.
\]
したがって
\[
  \frac{a_n-\alpha}{a_n-\beta}
  =\frac{x_n-\alpha y_n}{x_n-\beta y_n}
  =-\frac{d_n}{c_n}.
\]
以下では$a_1\ne\beta$とする.
このとき$c_1\ne0$であり,\,$ps-qr\ne0$より$\lambda_\alpha\ne0$なので,\,$c_n\ne0$が保たれる.
ここで$\lambda_\alpha=0$なら$M\begin{pmatrix}\alpha\\1\end{pmatrix}=\boldsymbol{0}$となり,行列式が零でないことに反する.
$a_1=\beta$の場合は,\,$a_n=\beta$という定数列として扱えばよい.

係数の更新規則より
\[
  \frac{a_{n+1}-\alpha}{a_{n+1}-\beta}
  =\frac{\lambda_\beta}{\lambda_\alpha}
   \frac{a_n-\alpha}{a_n-\beta}.
\]
本編で$b_n=(a_n-\alpha)/(a_n-\beta)$を等比数列にできたのは,この二つの係数の比を取り出していたからである.
不動点の方程式
\[
  r z^2+(s-p)z-q=0
\]
の解と係数の関係から$r(\alpha+\beta)=p-s$なので,
\[
  p-r\alpha=r\beta+s,\qquad
  p-r\beta=r\alpha+s
\]
となる.
よって公比は本編と同じ
\[
  \frac{\lambda_\beta}{\lambda_\alpha}
  =\frac{p-r\alpha}{p-r\beta}
\]
である.

たとえば$a_{n+1}=2a_n/(1-a_n)$では
\[
  M=\begin{pmatrix}2&0\\-1&1\end{pmatrix}
\]
で,不動点$0,-1$に対応する固有値はそれぞれ$1,\,2$になる.
$\alpha=0,\ \beta=-1$とすると
\[
  b_n=\frac{a_n}{a_n+1},\qquad b_{n+1}=2b_n
\]
である.
この$2$は,二つの固有値の比$2/1$から現れていた.

\subsection{三角関数型漸化式からチェビシェフ多項式へ}\label{節：三角関数型漸化式からチェビシェフ多項式へ}

特性根が複素数になると,数列には振動が現れる.
実係数の漸化式では複素根は共役な組で現れ,一般には$\rho>0,\ 0<\theta<\pi$の条件下で
\[
  \rho(\cos\theta+i\sin\theta),
  \qquad
  \rho(\cos\theta-i\sin\theta)
\]
と書ける.
特性方程式の係数を比べると,対応する漸化式は
\[
  a_{n+2}=2\rho\cos\theta\,a_{n+1}-\rho^2a_n
\]
である.
ド・モアブルの公式から,特性根の $n$ 乗は振幅 $\rho^n$ と角度 $n\theta$ を持つ.
そのため実数解は
\[
  a_n=\rho^n\{C\cos(n\theta)+D\sin(n\theta)\}
\]
という形になる.
$\rho>1$ なら振幅は大きくなり,$\,0<\rho<1$ なら小さくなる.

特に振幅が変わらない $\rho=1$ の場合,特性根は $\cos\theta\pm i\sin\theta$ であり,漸化式は
\[
  a_{n+2}=2\cos\theta\,a_{n+1}-a_n
\]
となる.
余弦と正弦の加法定理を使うと,$\,\cos(n\theta)$ も $\,\sin(n\theta)$ もこの同じ漸化式を満たす.
したがって,二つの初期値に合わせて定数 $C,D$ を選べば
\[
  a_n=C\cos(n\theta)+D\sin(n\theta)
\]
と表せる.
複素数の根は,ただ計算を複雑にするものではない.
実数の数列にひそむ振動と,その振幅の増減を記録している.

ここで $x=\cos\theta$ と書き,角度の整数倍を $x$ の多項式で表せないか考えてみよう.
余弦の加法定理
\[
  \cos((n+1)\theta)+\cos((n-1)\theta)
  =2\cos\theta\cos(n\theta)
\]
を変形すると
\[
  \cos((n+1)\theta)
  =2x\cos(n\theta)-\cos((n-1)\theta)
\]
となる.
右辺は,一つ前と二つ前の余弦から次の余弦を作っている.
そこで同じ規則を多項式の列に移し,
\begin{align*}
  &T_0(x)=1,\qquad T_1(x)=x,\\
  &T_{n+1}(x)=2xT_n(x)-T_{n-1}(x)
  \qquad(n\ge1)
\end{align*}

と定める.
最初の二つが多項式で,次も足し算と掛け算だけで作られるので,すべての $T_n(x)$ は多項式になる.
この列をチェビシェフ多項式という.

この定義から
\[
  T_2(x)=2x^2-1,\qquad
  T_3(x)=4x^3-3x
\]
が得られる.
これらは二倍角・三倍角の公式の右辺と同じ形をしている.
実際,
\[
  T_n(\cos\theta)=\cos(n\theta)
\]
が成り立つ.
導出を確かめよう.
$n=0,1$ では
\[
  T_0(\cos\theta)=1=\cos0,\qquad
  T_1(\cos\theta)=\cos\theta
\]
で一致する.
また余弦の加法定理から
\[
  \cos((n+1)\theta)
  =2\cos\theta\cos(n\theta)-\cos((n-1)\theta)
\]
である.
これは $T_{n+1}(\cos\theta)$ を作る漸化式と同じである.
したがって,二つの初めの値が一致し,次の値を作る規則も一致しているので,帰納法によりすべての $n$ で
\[
  T_n(\cos\theta)=\cos(n\theta)
\]
となる.
チェビシェフ多項式は,余弦の整数倍の角を,角度を使わずに $x=\cos\theta$ の多項式として記録している.

この関係は,本編で扱った三角関数型漸化式
\[
  a_{n+1}=2a_n^2-1
\]
にも戻ってくる.
初項が $-1\le a_1\le1$ なら,ある角 $\theta$ を選んで $a_1=\cos\theta$ と書ける.
二倍角の公式と $T_2(x)=2x^2-1$ から
\[
  a_{n+1}=T_2(a_n)
\]
であり,$\,a_n=\cos(2^{n-1}\theta)$ と予想できる.
この予想は初項で正しく,成り立つと仮定すれば
\[
  a_{n+1}
  =2\cos^2(2^{n-1}\theta)-1
  =\cos(2^n\theta)
\]
となるので,帰納法で確かめられる.
さらに $a_1=\cos\theta$ と $T_m(\cos\theta)=\cos(m\theta)$ を使えば
\[
  a_n=\cos(2^{n-1}\theta)
  =T_{2^{n-1}}(a_1)
\]
である.
たとえば本編の $a_1=\sqrt3/2=\cos(\pi/6)$ では
\[
  a_n=T_{2^{n-1}}\left(\frac{\sqrt3}{2}\right)
  =\cos\left(2^{n-1}\frac{\pi}{6}\right).
\]
つまり,本編の二次式による更新を何度も繰り返すことは,角度を二倍ずつ進めることと同じである.
その更新を多項式の言葉で書き直したものがチェビシェフ多項式だった.

より一般に,$\,T_m(T_k(x))=T_{mk}(x)$ という合成の関係も,同じ角度の式から分かる.
$x=\cos\theta$ とすれば
\[
  T_m(T_k(x))
  =T_m(\cos(k\theta))
  =\cos(mk\theta)
  =T_{mk}(x).
\]
この等式は $[-1,1]$ のすべての $x$ で成り立ち,両辺が多項式なので,多項式の恒等式として実数全体で成り立つ.
したがって $T_2$ を繰り返し合成する操作は,添字を二倍ずつ進める $T_2,T_4,T_8,\ldots$ にまとまる.

\subsection{非線形な更新も,不動点の近くでは線形に見える}\label{節：非線形な更新も,不動点の近くでは線形に見える}

ここまでは,漸化式を行列で表し,線形な更新の方向ごとの倍率を見てきた.
最後に,一次分数型のような非線形の更新規則でも,不動点の近くでは同じ見方ができることを確かめよう.
一項前から次が決まる漸化式
\[
  a_{n+1}=f(a_n)
\]
を考え,$\,\alpha$ を不動点 $f(\alpha)=\alpha$ とする.
そこからのずれを $e_n=a_n-\alpha$ と置くと
\[
  e_{n+1}=f(\alpha+e_n)-f(\alpha)
\]
である.
$f$ が $\alpha$ で微分可能なら,微分係数の定義から
\[
  \frac{f(\alpha+e)-f(\alpha)}{e}\longrightarrow f'(\alpha)
  \qquad(e\longrightarrow0)
\]
である.
つまり,ずれ $e$ が十分小さい範囲では
\[
  f(\alpha+e)-f(\alpha)\approx f'(\alpha)e
\]
となる.
曲線を不動点のすぐ近くまで拡大すると接線に見える,という一次近似を使ったのである.

この近似は,収束についても手がかりを与える.
もし $|f'(\alpha)|<1$ なら,$\,|f'(\alpha)|$ より大きく $1$ より小さい数 $c$ を選べる.
微分係数の定義により,$\,a_n$ が $\alpha$ に十分近ければ
\[
  |e_{n+1}|
  =|f(\alpha+e_n)-f(\alpha)|
  \le c|e_n|
\]
となる.
誤差は一歩ごとに少なくとも一定の割合で縮み,いったん十分近くに入ればその近くにとどまって $\alpha$ へ近づく.
したがって $|f'(\alpha)|<1$ は,不動点が近くの初期値を引き寄せることを示す.
先ほどの行列では固有値が各方向の倍率を表していた.
ここでは微分係数 $f'(\alpha)$ が,不動点近くでのずれの倍率を表している.
線形代数と微分は別々の話に見えるが,どちらも「一歩ごとの変化を倍率で測る」道具としてつながっている.

% 付録：階差型・総和型から微積分と近似計算へ

\section{一歩の変化を細かくすると,微積分になるか}\label{節：差分から始まる離散的な微積分}
  \BookFigDifference


前節では,本編で等比数列を作った置き換えを,行列の作用として捉え直した.
ここでは,\sectionref{節：階差型}の「差から元の数列を求める」という考え方を出発点にしよう.
その逆向きの操作は,\sectionref{節：総和型}で使った「和から差を取って一般項を取り出す」という操作である.
この二つを一緒に見ると,差分と和が微分と積分にどう対応するかが分かる.

以下では,最初の項を$a_0$と書くことがある.
数列$a_0,a_1,a_2,\ldots$の隣り合う項の差
\[
  \Delta a_n=a_{n+1}-a_n
\]
を一階差分という.
本編で使った$a_{n+1}-a_n$に,\,$\Delta$という記号と「差分」という名前を付けたのである.
たとえば$a_n=n^2$なら
\[
  \Delta a_n=2n+1,\qquad \Delta^2a_n=2.
\]
二次式から差を取ると一次式になり,もう一回差を取ると定数になる.
ここで$\Delta^2a_n$は,\,$\Delta a_{n+1}-\Delta a_n$の意味である.

本編の特性方程式型の別解でも,添字を一つずらして引くことで定数項$q$を消し,
\[
  \Delta a_{n+1}=p\,\Delta a_n
\]
という等比型を作った.
その計算も,差分によって扱いやすい量を取り出した例だった.

逆に,差分を足し合わせれば元の数列を求め直せる.
\[
  \begin{aligned}
    a_n-a_0
      &=(a_1-a_0)+\cdots+(a_n-a_{n-1})\\
      &=\sum_{k=0}^{n-1}\Delta a_k.
  \end{aligned}
\]
途中の項が打ち消され,最初と最後だけが残る.
本編の階差型で和を取った操作は,離散的な積分と見直せる.
総和型で$S_{n+1}-S_n=a_{n+1}$とした操作は,その逆である.

この見方は,\sectionref{節：階差等比複合型}で引く多項式$g(n)$を探したことにもつながる.
特に$p=1$なら,\,$g(n+1)-g(n)=f(n)$を満たす式を探す問題になる.
たとえば$S_n=\sum_{k=1}^n k^2$については
\[
  S_n-S_{n-1}=n^2,\qquad S_0=0.
\]
二次式の差を逆向きに戻すので,三次式
$S_n=An^3+Bn^2+Cn+D$を候補にする.
その差は
\[
  S_n-S_{n-1}
  =3An^2+(-3A+2B)n+(A-B+C).
\]
係数比較から$\displaystyle A=\frac{1}{3}$, $\displaystyle B=\frac{1}{2}$, $\displaystyle C=\frac{1}{6}$となり,\,$S_0=0$から$D=0$である.
この候補が同じ初期値と差の関係を満たすことにより
\[
  S_n=\frac{n(n+1)(2n+1)}6
\]
が確かめられる.
和の公式を覚えて使うだけでなく,「差を取ると目的の式に戻るものは何か」と考えて作ることができた.

ここからは微分と一次近似を使い,差分と連続的な変化の関係を見よう.
時刻$t_n=nh$の値を$a_n$とすると,一単位時間あたりの変化は
\[
  \frac{a_{n+1}-a_n}{h}
\]
である.
ある微分可能な関数$y$の値$a_n=y(t_n)$を記録したものなら,刻み幅$h$を小さくすると,この差分商は微分係数$y'(t_n)$に近づく.

逆に,微分方程式$y'=F(t,y)$を数値で解くには,今の点での傾きから
\[
  a_{n+1}=a_n+hF(t_n,a_n)
\]
と次の値を予測できる.
これをオイラー法という.
正確な連続解の値と一致するとは限らないが,接線の方向へ一歩ずつ進む漸化式で,曲線を近似している.

本編の等比型と比較しやすい例として,\,$y'=-y$, $y(0)=1$を考えよう.
連続解は$y(t)=e^{-t}$で,時間が進むと$0$に近づく.
オイラー法では
\[
  a_{n+1}=(1-h)a_n,\qquad a_0=1
\]
だから,\,$a_n=(1-h)^n$である.
本編の等比型で調べたように,\,$0$への収束条件は$|1-h|<1$, すなわち$0<h<2$になる.
$0<h<1$では正のまま減少するが,\,$1<h<2$では符号を交互に変えながら近づく.
さらに$h>2$では絶対値が増大する.
本当の解が減衰していても,刻み幅を大きくすると,近似計算の列はまったく違う動きになりうる.
本編で等比数列の公比を調べた知識は,数値計算が安定するかを判断する道具にもなる.

また,正の時刻$t$を$N$等分して$\displaystyle h=\frac{t}{N}$とすれば,時刻$t$における近似値は
\[
  a_N=\left(1-\frac{t}{N}\right)^N
  \longrightarrow e^{-t}.
\]
刻み幅を細かくすることと,近似値が連続解へ近づくことが,この例では基本的な指数関数の極限で確認できる.
一般の微分方程式でも同様の収束を期待するが,その保証には関数や刻み幅の条件を調べる必要がある.

次に,\sectionref{節：ニュートン法型}で扱った近似計算へ戻ろう.
本編では,接線と横軸の交点を次の値に選ぶことで
\[
  a_{n+1}=\frac12\left(a_n+\frac2{a_n}\right)
\]
を作り,\,$\sqrt2$への誤差がほぼ二乗されることを見た.
ここではその具体例から,なぜ一般のニュートン法でも誤差が二乗程度になるのかへ進む.

方程式$g(x)=0$に対するニュートン法の更新は
\[
  x_{n+1}=x_n-\frac{g(x_n)}{g'(x_n)}
\]
である.
$g(\alpha)=0$, $g'(\alpha)\ne0$とし,根$\alpha$の近くで$g$は二回連続微分可能とする.
本編のずれの置き換えと同じく,\,$e_n=x_n-\alpha$とおくと
\[
  e_{n+1}
  =\frac{g'(x_n)e_n-g(x_n)}{g'(x_n)}.
\]
分子は,接線による変化と実際の変化との差である.
二階微分$g''$の絶対値がこの近くで$M$以下なら
\[
  \begin{aligned}
    g'(x_n)e_n-g(x_n)
      &=\int_\alpha^{x_n}\{g'(x_n)-g'(t)\}\,dt,\\
    |g'(x_n)e_n-g(x_n)|
      &\le \frac{M}{2}|e_n|^2.
  \end{aligned}
\]
二つ目の式は,平均値の定理から
$|g'(x_n)-g'(t)|\le M|x_n-t|$として積分すれば得られる.
積分の向きが逆になる$x_n<\alpha$でも,絶対値を取れば同じ評価になる.
また,\,$g'(\alpha)\ne0$と連続性により,十分近くでは$|g'(x_n)|\ge m$となる正の定数$m$を選べる.
したがって
\[
  |e_{n+1}|\le\frac{M}{2m}|e_n|^2.
\]
誤差が十分小さければ,次の誤差はその二乗の定数倍以下になる.
いったん十分近くから始めれば近くにとどまり,急速に根へ近づくこともこの評価から分かる.
本編の$\sqrt2$の例で得た正確な式
\[
  e_{n+1}=\frac{e_n^2}{2a_n}
\]
は,この一般的な性質が特に見えやすい形で現れたものだった.

前節の不動点の見方でも確かめられる.
更新関数$\displaystyle N(x)=x-\frac{g(x)}{g'(x)}$について
\[
  N'(\alpha)
  =\left.\frac{g(x)g''(x)}{g'(x)^2}\right|_{x=\alpha}=0.
\]
つまり,根のところでは一次のずれの倍率が零になり,二次のずれが残る.
本編で等比型へ帰着させたときは一定倍率でずれを追ったが,ニュートン法ではずれが小さいほど,次に残る割合も小さくなるのである.

接線の傾きを毎回微分で求めるのが難しいなら,過去二点を結ぶ割線の傾きで代用できる.
その更新は
\[
  x_{n+1}
  =x_n-g(x_n)\frac{x_n-x_{n-1}}{g(x_n)-g(x_{n-1})}
\]
となり,割線法と呼ばれる.
分母が零でない範囲で用いる.
直前の二項を記録して次を作るという点は,前節の二成分の状態と共通している.
本編で与えられた漸化式を解いた経験は,このように近似計算の規則を作り,その収束を調べることにもつながる.

% 付録再編草案・第3部：数列全体と母関数

\section{数列全体を一つの対象にまとめると,何が見えるか}

ここまでは,次の項をどう作るか,隣り合う項の差から何が読めるかを見てきた.
しかし,漸化式が決めているのは一つの項だけではない.
初項から先の全項をつなぐ,数列全体の骨組みである.
その骨組みを一度に持ち運べる形へまとめたら,何が見えるだろう.

数列
\[
  a_0,a_1,a_2,a_3,\ldots
\]
に対して,項の番号を印にして
\[
  A(x)=a_0+a_1x+a_2x^2+a_3x^3+\cdots
  =\sum_{n=0}^{\infty}a_nx^n
\]
と並べる.
これを母関数という.
$x$ を一つ掛けるごとに項の番号が一つ増えるので,\,$x$ は「何番目の項か」を記録する目盛りのような役割を持つ.
最初は,これを収束する関数としてではなく,数列を係数として並べた形式的な記号として扱おう.
母関数の力は,無限個の項を一つに圧縮することよりも,数列の操作を代数の操作に言い換えられることにある.

たとえば $A(x)$ に $x$ を掛けると,
\[
  xA(x)=a_0x+a_1x^2+a_2x^3+\cdots
\]
となる.
数列で言えば,初めに $0$ を置いて,その後の項を一つ右へずらした形である.
また,二つの数列の項を
\[
  c_n=\sum_{k=0}^{n}a_kb_{n-k}
\]
のように組み合わせる操作は畳み込みと呼ばれ,母関数を掛けたときの $x^n$ の係数になる.
一つの対象を作る方法を二つに分ける,あるいは二つの部品を組み合わせるとき,この畳み込みが自然に現れる.
つまり母関数では,「一つ先へずらす」は $x$ 倍に,「二つの選択を組み合わせる」は積に翻訳される.

この翻訳をフィボナッチ数列で追ってみよう.
\[
  F_0=0,\quad F_1=1,\qquad F_{n+2}=F_{n+1}+F_n
\]
とし,\,$F(x)=\sum_{n=0}^{\infty}F_nx^n$ と置く.
左辺の初めの項を取り出すと,
\[
  F(x)-x=\sum_{n=0}^{\infty}F_{n+2}x^{n+2}
\]
である.
漸化式 $F_{n+2}=F_{n+1}+F_n$ を代入すると,右辺の二つの和はそれぞれ $xF(x)$ と $x^2F(x)$ になる.
したがって
\[
  F(x)-x=xF(x)+x^2F(x)
\]
であり,
\[
  F(x)=\frac{x}{1-x-x^2}
\]
を得る.
項を一つずつ計算する代わりに,数列全体が満たす一つの代数方程式を解いた.
漸化式は母関数の分母に姿を変えたのである.

ここで分母を特性方程式と比べてみる.
\[
  1-x-x^2=0
\]
に対して $r=1/x$ と置くと,\,$r^2-r-1=0$ となる.
つまり母関数の分母の零点は,特性方程式の根の逆数である.
これは偶然ではない.
特性方程式は「一歩ごとの倍率」を調べ,母関数は「全項を並べたときの係数関係」を調べている.
行列の固有値,特性根,母関数の分母は,同じ更新規則を違う入口から見たものだった.

フィボナッチ数列をもう少し遠くから眺めると,長期的な増え方も見えてくる.
特性根は
\[
  \varphi=\frac{1+\sqrt5}{2},
  \qquad
  \psi=\frac{1-\sqrt5}{2}
\]
であり,\,$|\psi|<1<\varphi$ である.
一般項
\[
  F_n=\frac{\varphi^n-\psi^n}{\sqrt5}
\]
では,\,$\psi^n$ の影響は次第に小さくなる.
したがって
\[
  F_n\sim\frac{\varphi^n}{\sqrt5}
\]
となる.
ここで $\sim$ は,二つの量の比が $1$ に近づくという意味である.
大きな $n$ でどの根が勝つかを見れば,正確な一般項を毎回計算しなくても成長の景色が分かる.
母関数の分母からも同じ特性根が取り出されるため,固有値・一般項・母関数・漸近的な増大は一つの話につながる.

母関数が本領を発揮するのは,項どうしの組合せが漸化式に現れるときである.
カタラン数は,たとえば括弧の付け方や二分木の形を数える数列で,
\[
  C_0=1,\qquad C_{n+1}=\sum_{k=0}^{n}C_kC_{n-k}
\]
を満たす.
全体を左側の大きさ $k$ と右側の大きさ $n-k$ に分けると,左の形と右の形を一つずつ選ぶので積 $C_kC_{n-k}$ が現れる.
すべての分け方を足すため,右辺は畳み込みになっている.
母関数 $C(x)=\sum_{n=0}^{\infty}C_nx^n$ を作ると,畳み込みは積に変わり,
\[
  C(x)=1+xC(x)^2
\]
となる.
項ごとには複雑に見えた再帰が,母関数についての二次方程式へ変わった.
これを解くと
\[
  C(x)=\frac{1-\sqrt{1-4x}}{2x}
\]
である.
もう一方の符号では $x=0$ の近くで値が発散してしまい,定数項 $C_0=1$ を持つ母関数にならないため,こちらの枝を選ぶ.
母関数を使うと,非線形な漸化式を力ずくで代数方程式へ移し,そこから項の情報を引き出せる.

この発想は,チェビシェフ多項式にも戻ってくる.
その漸化式に母関数を作ると,
\[
  \sum_{n=0}^{\infty}T_n(x)t^n
  =\frac{1-xt}{1-2xt+t^2}
\]
となる.
これは,漸化式に対応する項を一つずつ並べ,同じ次数の係数を整理すれば得られる.
三角関数で定義した多項式列を,今度は一つの有理関数にまとめている.
漸化式,三角関数,母関数は別々の定義ではなく,同じ列を引き出すための三つの窓になっている.
このように,一つの対象に複数の表現を持たせると,見えにくかった性質が別の窓から見えてくる.

さらに遠くには,「数列を満たす関数を作る」話も広がっている.
階乗は整数 $n$ で $(n+1)!=(n+1)n!$ を満たすが,階乗を整数以外でも使いたい場面がある.
その拡張の一つがガンマ関数で,正の $x$ に対して
\[
  \Gamma(x)=\int_0^\infty t^{x-1}e^{-t}\,dt
\]
と定められる.
部分積分をすると
\[
  \Gamma(x+1)=x\Gamma(x)
\]
が現れ,整数では $\Gamma(n+1)=n!$ となる.
数列の外へ関数を作ると,階乗が積分や確率に現れる新しい景色へつながる.
本節ではその理論を追わないが,チェビシェフ多項式と同じく,漸化式が関数を生み出す一つの例として見ておこう.

確率に目を向けると, 状態の移り変わりも漸化式で記録できる.
晴れ・雨の順に確率を並べた列ベクトルを $\boldsymbol p_n$ とし, 晴れの翌日に晴れる確率を $0.8$, 雨の翌日に晴れる確率を $0.3$ とすると,
\[
  \boldsymbol p_{n+1}
  =
  \begin{pmatrix}0.8&0.3\\0.2&0.7\end{pmatrix}
  \boldsymbol p_n
  =P\boldsymbol p_n
\]
と書ける.
各列の和は $1$ なので, どの状態から出発しても次の確率の合計は $1$ に保たれる.
長く更新したときの定常分布は $P\boldsymbol p=\boldsymbol p$ を満たし, どの初期状態からそこへ近づくかは固有値で調べられる.
線形代数で見た状態の更新が, 確率の予測にもそのまま現れている.

数論では, ユークリッドの互除法が余りを次々に作る更新として表せる.
$a_{n-1}$ を $a_n$ で割った商を $q_n$, 余りを $a_{n+1}$ とすれば,
\[
  a_{n-1}=q_na_n+a_{n+1},
  \qquad 0\le a_{n+1}<a_n
\]
となる.
正の余りは毎回小さくなるため, この計算は有限回で止まり, 最大公約数が求まる.
ここでは一般項を求めず, 更新のたびに小さくなるという性質からアルゴリズムの停止を示した.

同じ数論の見方として, 整数係数が一定の二階漸化式を $m$ で割った余りで考えてみよう.
状態 $(a_n,a_{n+1})$ の各成分は $0,1,\ldots,m-1$ のいずれかなので, 状態は高々 $m^2$ 通りしかない.
次の状態が今の状態だけで決まるなら, いつか同じ状態が現れ, その後は同じ並びを繰り返す.
したがって, 余りの列は周期的になる.
たとえばフィボナッチ数列を $2$ で割った余りは
\[
  0,1,1,0,1,1,0,1,1,\ldots
\]
と周期 $3$ で繰り返す.
実数としての増大や収束を調べるときと, 余りとしての周期を調べるときでは, 同じ更新規則から見えてくる性質が変わる.

一般項を簡単な式で書きにくい数列にも, 調べられることは残っている.
たとえば
\[
  a_0=0,\qquad a_{n+1}=a_n^2+1
\]
では $0,1,2,5,26,\ldots$ と増える.
$n\ge2$ では $a_n\ge2$ であり, $a_{n+1}\ge a_n^2$ だから, 帰納的に
\[
  a_n\ge 2^{\,2^{n-2}}
\]
となる.
一般項を求めなくても, 増加が非常に速いことは分かる.
また, 数列が単調で有界と分かれば収束し, 更新が $a_{n+1}=f(a_n)$ で $f$ が連続なら, 極限 $L$ は $L=f(L)$ を満たす.
単調性, 有界性, 不動点, 増加速度は, 一般項とは別の問いに答えてくれる.

本編では, 漸化式から一般項を取り出す方法を整理した.
付録では, 同じ更新規則を行列で表し, 一歩の差を細かく見て微積分へつなぎ, 全項を母関数にまとめた.
確率の定常分布, 整数の余りの周期, 一般項を使わない増加の評価も, 見たい性質に応じて更新規則を読み替えた例である.
漸化式を理解する道具は, 一般項を得るためだけでなく, その数列がどのように動き, 何を数え, どのような性質を保つかを知るために広がっている.


\section{演習の操作を別の表現で見直す}\label{節：演習の操作を別の表現で見直す}
  新しい道具を学んだ後に,すでに解いた問題へ戻ろう.
  ここでは\bookproblemref{practice-entrance}{3}{発展問3},
  \bookproblemref{practice-entrance}{4}{問4},
  \bookproblemref{practice-entrance}{9}{問9}で行った操作を,
  \bookterm{generating}{母関数}・回転・\bookterm{invariant}{不変量}という見方で確かめる.

  \subsection{畳み込み和が一定になる理由}
  発展問3の$c_0=1$, $c_{n+1}=\frac{2n+1}{2n+2}c_n$から,
  \[
    c_n=\frac1{4^n}\binom{2n}{n}
  \]
  を得た.その\bookterm{generating}{母関数}を$C(x)=\sum_{n\ge0}c_nx^n$とする.
  まず,級数の各項を微分する形式的な操作を使う.
  収束する関数の微分を仮定するのではなく,
  $x^n$を$nx^{n-1}$へ移して係数を並べ直す操作として読む.
  元の\bookterm{recurrence}{漸化式}を$(n+1)c_{n+1}=(n+\frac12)c_n$と書けば,
  \[
    (1-x)C'(x)=\frac12C(x)
  \]
  を得る.積の係数は有限和なので,係数を比較することで通常の積の微分法則も使える.
  $D(x)=C(x)^2=\sum_{n\ge0}d_nx^n$と置くと,
  \[
    (1-x)D'(x)=2C(x)(1-x)C'(x)=D(x).
  \]
  $x^n$の係数を比較すれば$(n+1)d_{n+1}-nd_n=d_n$,すなわち$d_{n+1}=d_n$である.
  $d_0=c_0^2=1$なので,
  \[
    C(x)^2=\frac1{1-x}=\sum_{n\ge0}x^n
  \]
  の$x^n$の係数から
  \[
    \sum_{k=0}^n c_kc_{n-k}=1
  \]
  を得る.問題の解法は有限和を変形し,この解法は\bookterm{sequence}{数列}全体の積へ移した.
  どちらも同じ\bookterm{convolution}{畳み込み}を扱っている.

  \subsection{中点の更新を,回転と面積から見直す}
  発展問4では,平行四辺形の各辺上で同じ比率$t$ $(0<t<1)$の分点を取り,
  新しい平行四辺形を作った.連続する辺ベクトルを$u,v$とする.
  頂点の更新を引き算すれば,
  \[
    u'=(1-t)u+tv,\qquad v'=-tu+(1-t)v
  \]
  である.従って辺の組を混ぜる\bookterm{matrix}{行列}は
  \[
    M=\begin{pmatrix}1-t&t\\-t&1-t\end{pmatrix}
  \]
  となる.\bookterm{matrix}{行列}式は$\rho=(1-t)^2+t^2$だから,
  平行四辺形の面積は一回ごとに$\rho$倍となる.
  \bookterm{eigenvalue}{固有値}は$1-t\pm it$であり,その絶対値の二乗も$\rho$である.

  正方形で,向きを反時計回りに取った辺を複素数として$v=iu$と書けるときは,
  \[
    u'=((1-t)+it)u
  \]
  である.$t=\frac12$なら倍率は$\frac1{\sqrt2}$,回転角は$\frac\pi4$になる.
  これが問題の中点更新の辺長と向きに一致する.
  \begin{bookpoint}
    \textbf{\bookterm{matrix}{行列}が混ぜる対象を確認する}\par
    一般の平行四辺形では$u,v$は直交する等長の辺ではない.
    辺の組を混ぜる\bookterm{matrix}{行列}を,そのまま平面上の相似変換と読むことはできない.
    面積の倍率の説明は一般の場合に使え,上の回転の説明には正方形という条件を使う.
  \end{bookpoint}

  \subsection{整数方程式を保つ更新を,不変量から見る}
  発展問9では,$u_0=3$, $u_1=3$, $u_{n+2}=3u_{n+1}-u_n$から,
  $(3,u_n,u_{n+1})$が$x^2+y^2+z^2=xyz$の解になることを示した.
  初期値では
  \[
    u_n^2+u_{n+1}^2+9=3u_nu_{n+1}
  \]
  が成り立つ.一般の$x,y$に対して
  \[
    y^2+(3y-x)^2-3y(3y-x)=x^2+y^2-3xy
  \]
  だから,左辺の差は更新の前後で変わらず$-9$に保たれる.
  その結果,$u_nu_{n+2}=u_{n+1}^2+9$となり,
  \[
    u_{n+2}=\frac{u_{n+1}^2+9}{u_n}
  \]
  とも書ける.正値性は$u_0=u_1=3$と,
  $u_2=6$以降の増加から確認できる.
  \sectionref*{節：不変量から線形の関係}の定数$1$を$9$へ変えた更新だった.
  元の問題では\bookterm{linear}{線形}の更新から保存を示し,本編の新しい例では保存から\bookterm{linear}{線形}の更新を取り出した.
  同じ関係を二つの向きから使っている.

  \begin{bookpoint}
    \textbf{別解から何を持ち帰るか}\par
    計算が短くなるだけでなく,足す意味・面積の意味・保存される関係が見える.
    新しい表現を使った後は,元の問題の条件と答えへ戻って対応を確かめよう.
  \end{bookpoint}


\end{bookcolumns}

% chapter: 参考文献
\chapter{参考文献}\label{章：参考文献}

本書に関係する書籍と,本文に対応する個別ウェブ資料をまとめた.
ウェブ資料は説明・公式の照合や補足学習の参照先である.
本書と資料で型の名前・初項の番号・\bookterm{normalize}{正規化}が異なる場合があるため,
対応箇所を手掛かりに,条件と記法を照合して読んでほしい.

\begingroup
\raggedright
\section{書籍}
\begin{enumerate}[label={[\arabic*]}]
  \item \begin{minipage}[t]{\linewidth}\raggedright
  松坂和夫,\,『新装版 数学読本1』,\,岩波書店,\,2019年.
  ISBN 978-4-00-029877-3.
  \par\noindent\url{https://www.hanmoto.com/bd/isbn/9784000298773}\par
  \end{minipage}\par\smallskip

  \item \begin{minipage}[t]{\linewidth}\raggedright
  松坂和夫,\,『新装版 数学読本2』,\,岩波書店,\,2019年.
  ISBN 978-4-00-029878-0.
  \par\noindent\url{https://www.books.or.jp/book-details/9784000298780}\par
  \end{minipage}\par\smallskip

  \item \begin{minipage}[t]{\linewidth}\raggedright
  松坂和夫,\,『新装版 数学読本3』,\,岩波書店,\,2019年.
  ISBN 978-4-00-029879-7.
  \par\noindent\url{https://www.hanmoto.com/bd/isbn/9784000298797}\par
  \end{minipage}\par\smallskip

  \item \begin{minipage}[t]{\linewidth}\raggedright
  松坂和夫,\,『新装版 数学読本4』,\,岩波書店,\,2019年.
  ISBN 978-4-00-029880-3.
  \par\noindent\url{https://www.hanmoto.com/bd/isbn/9784000298803}\par
  \end{minipage}\par\smallskip

  \item \begin{minipage}[t]{\linewidth}\raggedright
  松坂和夫,\,『新装版 数学読本5』,\,岩波書店,\,2019年.
  ISBN 978-4-00-029881-0.
  \par\noindent\url{https://www.hanmoto.com/bd/isbn/9784000298810}\par
  \end{minipage}\par\smallskip

  \item \begin{minipage}[t]{\linewidth}\raggedright
  松坂和夫,\,『新装版 数学読本6』,\,岩波書店,\,2019年.
  ISBN 978-4-00-029882-7.
  \par\noindent\url{https://www.hanmoto.com/bd/isbn/9784000298827}\par
  \end{minipage}\par\smallskip

  \item \begin{minipage}[t]{\linewidth}\raggedright
  宮腰忠,\,『高校数学+α：基礎と論理の物語』,\,共立出版,\,2004年.
  ISBN 978-4-320-01768-9.
  \par\noindent\url{https://www.kyoritsu-pub.co.jp/book/b10010328.html}\par
  \end{minipage}\par\smallskip

\end{enumerate}

\section{本編の漸化式：おいしい数学}
\begin{enumerate}[resume,label={[\arabic*]}]
  \item \begin{minipage}[t]{\linewidth}\raggedright
  おいしい数学,\,「\bookterm{recurrence}{漸化式}の全パターン紹介」.
  \par{\small 対応：第3章の全体像.\par}
  \par\noindent\url{https://hiraocafe.com/note/zenkashiki.html}\par
  \noindent（2026年10月9日閲覧）.
  \end{minipage}\par\smallskip

  \item \begin{minipage}[t]{\linewidth}\raggedright
  おいしい数学,\,「2-1型（等差型）の\bookterm{recurrence}{漸化式}」.
  \par{\small 対応：基本4パターン：等差型.\par}
  \par\noindent\url{https://hiraocafe.com/note/zenkashiki2-1.html}\par
  \noindent（2026年10月9日閲覧）.
  \end{minipage}\par\smallskip

  \item \begin{minipage}[t]{\linewidth}\raggedright
  おいしい数学,\,「2-2型（等比型）の\bookterm{recurrence}{漸化式}」.
  \par{\small 対応：基本4パターン：等比型.\par}
  \par\noindent\url{https://hiraocafe.com/note/zenkashiki2-2.html}\par
  \noindent（2026年10月9日閲覧）.
  \end{minipage}\par\smallskip

  \item \begin{minipage}[t]{\linewidth}\raggedright
  おいしい数学,\,「2-3型（\bookterm{difference}{階差}型）の\bookterm{recurrence}{漸化式}」.
  \par{\small 対応：基本4パターン：\bookterm{difference}{階差}型.\par}
  \par\noindent\url{https://hiraocafe.com/note/zenkashiki2-3.html}\par
  \noindent（2026年10月9日閲覧）.
  \end{minipage}\par\smallskip

  \item \begin{minipage}[t]{\linewidth}\raggedright
  おいしい数学,\,「2-8型（\bookterm{ratio}{階比}型）の\bookterm{recurrence}{漸化式}」.
  \par{\small 対応：基本4パターン：\bookterm{ratio}{階比}型／\bookterm{auxiliary}{補助数列}の\bookterm{normalize}{正規化}.\par}
  \par\noindent\url{https://hiraocafe.com/note/zenkashiki2-8.html}\par
  \noindent（2026年10月9日閲覧）.
  \end{minipage}\par\smallskip

  \item \begin{minipage}[t]{\linewidth}\raggedright
  おいしい数学,\,「2-4型（\bookterm{characteristic}{特性方程式}型）の\bookterm{recurrence}{漸化式}」.
  \par{\small 対応：\bookterm{characteristic}{特性方程式}型／\bookterm{auxiliary}{補助数列}の設計.\par}
  \par\noindent\url{https://hiraocafe.com/note/zenkashiki2-4.html}\par
  \noindent（2026年10月9日閲覧）.
  \end{minipage}\par\smallskip

  \item \begin{minipage}[t]{\linewidth}\raggedright
  おいしい数学,\,「2-5型（指数型）の\bookterm{recurrence}{漸化式}」.
  \par{\small 対応：\bookterm{difference}{階差}等比複合型：指数的な付加項.\par}
  \par\noindent\url{https://hiraocafe.com/note/zenkashiki2-5.html}\par
  \noindent（2026年10月9日閲覧）.
  \end{minipage}\par\smallskip

  \item \begin{minipage}[t]{\linewidth}\raggedright
  おいしい数学,\,「2-7型（関数スライド型）の\bookterm{recurrence}{漸化式}」.
  \par{\small 対応：\bookterm{difference}{階差}等比複合型：多項式の付加項／\bookterm{auxiliary}{補助数列}の設計.\par}
  \par\noindent\url{https://hiraocafe.com/note/zenkashiki2-7.html}\par
  \noindent（2026年10月9日閲覧）.
  \end{minipage}\par\smallskip

  \item \begin{minipage}[t]{\linewidth}\raggedright
  おいしい数学,\,「2-6型（逆数型）の\bookterm{recurrence}{漸化式}」.
  \par{\small 対応：分数型：逆数型.\par}
  \par\noindent\url{https://hiraocafe.com/note/zenkashiki2-6.html}\par
  \noindent（2026年10月9日閲覧）.
  \end{minipage}\par\smallskip

  \item \begin{minipage}[t]{\linewidth}\raggedright
  おいしい数学,\,「2-9型（\bookterm{logarithm}{対数}型）の\bookterm{recurrence}{漸化式}」.
  \par{\small 対応：\bookterm{logarithm}{対数}型.\par}
  \par\noindent\url{https://hiraocafe.com/note/zenkashiki2-9.html}\par
  \noindent（2026年10月9日閲覧）.
  \end{minipage}\par\smallskip

  \item \begin{minipage}[t]{\linewidth}\raggedright
  おいしい数学,\,「2-10型（1次分数型）の\bookterm{recurrence}{漸化式}」.
  \par{\small 対応：分数型：1次分数型.\par}
  \par\noindent\url{https://hiraocafe.com/note/zenkashiki2-10.html}\par
  \noindent（2026年10月9日閲覧）.
  \end{minipage}\par\smallskip

  \item \begin{minipage}[t]{\linewidth}\raggedright
  おいしい数学,\,「3-1型の\bookterm{recurrence}{漸化式}（隣接3項間\bookterm{recurrence}{漸化式}の基本）」.
  \par{\small 対応：隣接3項間\bookterm{recurrence}{漸化式}型.\par}
  \par\noindent\url{https://hiraocafe.com/note/zenkashiki3-1.html}\par
  \noindent（2026年10月9日閲覧）.
  \end{minipage}\par\smallskip

  \item \begin{minipage}[t]{\linewidth}\raggedright
  おいしい数学,\,「3-2型の\bookterm{recurrence}{漸化式}（隣接3項間\bookterm{recurrence}{漸化式}の応用）」.
  \par{\small 対応：隣接3項間\bookterm{recurrence}{漸化式}型の拡張.\par}
  \par\noindent\url{https://hiraocafe.com/note/zenkashiki3-2.html}\par
  \noindent（2026年10月9日閲覧）.
  \end{minipage}\par\smallskip

  \item \begin{minipage}[t]{\linewidth}\raggedright
  おいしい数学,\,「和を含んだ\bookterm{recurrence}{漸化式}」.
  \par{\small 対応：総和型.\par}
  \par\noindent\url{https://hiraocafe.com/note/zenkashikisn.html}\par
  \noindent（2026年10月9日閲覧）.
  \end{minipage}\par\smallskip

  \item \begin{minipage}[t]{\linewidth}\raggedright
  おいしい数学,\,「連立\bookterm{recurrence}{漸化式}」.
  \par{\small 対応：連立型.\par}
  \par\noindent\url{https://hiraocafe.com/note/zenkashikirenritsu.html}\par
  \noindent（2026年10月9日閲覧）.
  \end{minipage}\par\smallskip

\end{enumerate}

\section{本編と付録：高校数学の美しい物語}
\begin{enumerate}[resume,label={[\arabic*]}]
  \item \begin{minipage}[t]{\linewidth}\raggedright
  学びTimes「高校数学の美しい物語」,\,「\bookterm{recurrence}{漸化式}の\bookterm{characteristic}{特性方程式}の意味とうまくいく理由」.
  \par{\small 対応：\bookterm{characteristic}{特性方程式}型／\bookterm{auxiliary}{補助数列}の設計／隣接多項間.\par}
  \par\noindent\url{https://manabitimes.jp/math/1299}\par
  \noindent（2026年10月9日閲覧）.
  \end{minipage}\par\smallskip

  \item \begin{minipage}[t]{\linewidth}\raggedright
  学びTimes「高校数学の美しい物語」,\,「三項間\bookterm{recurrence}{漸化式}の3通りの解き方」.
  \par{\small 対応：隣接3項間\bookterm{recurrence}{漸化式}型／付録の\bookterm{matrix}{行列}による見直し.\par}
  \par\noindent\url{https://manabitimes.jp/math/697}\par
  \noindent（2026年10月9日閲覧）.
  \end{minipage}\par\smallskip

  \item \begin{minipage}[t]{\linewidth}\raggedright
  学びTimes「高校数学の美しい物語」,\,「一次分数型の\bookterm{recurrence}{漸化式}の解法と例題」.
  \par{\small 対応：分数型：1次分数型.\par}
  \par\noindent\url{https://manabitimes.jp/math/1331}\par
  \noindent（2026年10月9日閲覧）.
  \end{minipage}\par\smallskip

  \item \begin{minipage}[t]{\linewidth}\raggedright
  学びTimes「高校数学の美しい物語」,\,「\bookterm{matrix}{行列}の\bookterm{eigenvalue}{固有値}・\bookterm{eigenvector}{固有ベクトル}の定義と具体的な計算方法」.
  \par{\small 対応：付録：\bookterm{matrix}{行列}の基本、図5.2.\par}
  \par\noindent\url{https://manabitimes.jp/math/1008}\par
  \noindent（2026年10月9日閲覧）.
  \end{minipage}\par\smallskip

  \item \begin{minipage}[t]{\linewidth}\raggedright
  学びTimes「高校数学の美しい物語」,\,「\bookterm{matrix}{行列}の対角化の意味と具体的な計算方法」.
  \par{\small 対応：付録：\bookterm{characteristic}{特性方程式}を\bookterm{matrix}{行列}で見直す.\par}
  \par\noindent\url{https://manabitimes.jp/math/1133}\par
  \noindent（2026年10月9日閲覧）.
  \end{minipage}\par\smallskip

  \item \begin{minipage}[t]{\linewidth}\raggedright
  学びTimes「高校数学の美しい物語」,\,「\bookterm{chebyshev}{チェビシェフ多項式}」.
  \par{\small 対応：虚数解型隣接3項間／付録：\bookterm{chebyshev}{チェビシェフ多項式}.\par}
  \par\noindent\url{https://manabitimes.jp/math/729}\par
  \noindent（2026年10月9日閲覧）.
  \end{minipage}\par\smallskip

  \item \begin{minipage}[t]{\linewidth}\raggedright
  学びTimes「高校数学の美しい物語」,\,「\bookterm{sequence}{数列}の\bookterm{generating}{母関数}の意味とその応用例」.
  \par{\small 対応：付録：\bookterm{sequence}{数列}全体と\bookterm{generating}{母関数}.\par}
  \par\noindent\url{https://manabitimes.jp/math/658}\par
  \noindent（2026年10月9日閲覧）.
  \end{minipage}\par\smallskip

  \item \begin{minipage}[t]{\linewidth}\raggedright
  学びTimes「高校数学の美しい物語」,\,「カタラン数の意味と\bookterm{recurrence}{漸化式}」.
  \par{\small 対応：付録：\bookterm{convolution}{畳み込み}とカタラン数.\par}
  \par\noindent\url{https://manabitimes.jp/math/657}\par
  \noindent（2026年10月9日閲覧）.
  \end{minipage}\par\smallskip

  \item \begin{minipage}[t]{\linewidth}\raggedright
  学びTimes「高校数学の美しい物語」,\,「オイラー法をわかりやすく解説」.
  \par{\small 対応：付録：一歩の変化と微積分.\par}
  \par\noindent\url{https://manabitimes.jp/math/2536}\par
  \noindent（2026年10月9日閲覧）.
  \end{minipage}\par\smallskip

  \item \begin{minipage}[t]{\linewidth}\raggedright
  学びTimes「高校数学の美しい物語」,\,「ガンマ関数（階乗の一般化）の定義と性質」.
  \par{\small 対応：付録：階乗からガンマ関数へ.\par}
  \par\noindent\url{https://manabitimes.jp/math/960}\par
  \noindent（2026年10月9日閲覧）.
  \end{minipage}\par\smallskip

\end{enumerate}

\section{公式と条件：NIST DLMF}
\begin{enumerate}[resume,label={[\arabic*]}]
  \item \begin{minipage}[t]{\linewidth}\raggedright
  NIST Digital Library of Mathematical Functions,\,「§3.8 Nonlinear Equations」.
  \par{\small 対応：\bookterm{newton}{ニュートン法}型／付録の近似計算.\par}
  \par\noindent\url{https://dlmf.nist.gov/3.8}\par
  \noindent（2026年10月9日閲覧）.
  \end{minipage}\par\smallskip

  \item \begin{minipage}[t]{\linewidth}\raggedright
  NIST Digital Library of Mathematical Functions,\,「§3.3 Interpolation」.
  \par{\small 対応：付録の\bookterm{interpolation}{ラグランジュ補間}コラム.\par}
  \par\noindent\url{https://dlmf.nist.gov/3.3}\par
  \noindent（2026年10月9日閲覧）.
  \end{minipage}\par\smallskip

  \item \begin{minipage}[t]{\linewidth}\raggedright
  NIST Digital Library of Mathematical Functions,\,「§18.9 Recurrence Relations and Derivatives」.
  \par{\small 対応：付録のチェビシェフ・\bookterm{legendre}{ルジャンドル多項式}.\par}
  \par\noindent\url{https://dlmf.nist.gov/18.9}\par
  \noindent（2026年10月9日閲覧）.
  \end{minipage}\par\smallskip

  \item \begin{minipage}[t]{\linewidth}\raggedright
  NIST Digital Library of Mathematical Functions,\,「§18.3 Definitions」.
  \par{\small 対応：付録の\bookterm{legendre}{ルジャンドル多項式}の直交性.\par}
  \par\noindent\url{https://dlmf.nist.gov/18.3}\par
  \noindent（2026年10月9日閲覧）.
  \end{minipage}\par\smallskip

  \item \begin{minipage}[t]{\linewidth}\raggedright
  NIST Digital Library of Mathematical Functions,\,「§18.12 Generating Functions」.
  \par{\small 対応：付録の\bookterm{chebyshev}{チェビシェフ多項式}の\bookterm{generating}{母関数}.\par}
  \par\noindent\url{https://dlmf.nist.gov/18.12}\par
  \noindent（2026年10月9日閲覧）.
  \end{minipage}\par\smallskip

  \item \begin{minipage}[t]{\linewidth}\raggedright
  NIST Digital Library of Mathematical Functions,\,「§5.2 Definitions」.
  \par{\small 対応：付録のガンマ関数の定義.\par}
  \par\noindent\url{https://dlmf.nist.gov/5.2}\par
  \noindent（2026年10月9日閲覧）.
  \end{minipage}\par\smallskip

  \item \begin{minipage}[t]{\linewidth}\raggedright
  NIST Digital Library of Mathematical Functions,\,「§5.5 Functional Relations」.
  \par{\small 対応：付録のガンマ関数の更新.\par}
  \par\noindent\url{https://dlmf.nist.gov/5.5}\par
  \noindent（2026年10月9日閲覧）.
  \end{minipage}\par\smallskip

\end{enumerate}

\section{整数の更新と数え上げ}
\begin{enumerate}[resume,label={[\arabic*]}]
  \item \begin{minipage}[t]{\linewidth}\raggedright
  Keith Conrad,\,\textit{Pell's Equation I}.
  \par{\small 対応：\bookterm{pell}{ペル方程式}の具体例・整数解を作る更新.\par}
  \par\noindent\url{https://kconrad.math.uconn.edu/blurbs/ugradnumthy/pelleqn1.pdf}\par
  \noindent（2026年10月10日閲覧）.
  \end{minipage}\par\smallskip

  \item \begin{minipage}[t]{\linewidth}\raggedright
  Keith Conrad,\,\textit{Pell's Equation II}.
  \par{\small 対応：\bookterm{pell}{ペル方程式}と\bookterm{continued}{連分数}の近似分数の関係.\par}
  \par\noindent\url{https://kconrad.math.uconn.edu/blurbs/ugradnumthy/pelleqn2.pdf}\par
  \noindent（2026年10月10日閲覧）.
  \end{minipage}\par\smallskip

  \item \begin{minipage}[t]{\linewidth}\raggedright
  NIST Digital Library of Mathematical Functions,
  \,$\S$26.3 \textit{Lattice Paths: Binomial Coefficients}.
  \par{\small 対応：\bookterm{double}{二重数列}の格子経路,二項係数とその加法公式.\par}
  \par\noindent\url{https://dlmf.nist.gov/26.3}\par
  \noindent（2026年10月10日閲覧）.
  \end{minipage}\par\smallskip

\end{enumerate}
\endgroup


\end{document}
